% Generated by build_chapter.py from the GY1987 files. Do not edit by hand; edit the files and rebuild.
\chapter{The academic year 1987--88}\label{ch:1987}
\setlength{\emergencystretch}{3em}
\hyphenation{Par-tha-sa-ra-thy}
\noindent\textbf{Scope and editorial date: 11 November 2026.}
This chapter covers all 38 catalogue records, 714--751, in notebook
order, from October 1987 through April 1988. The archive windows overlap;
the Russian text and scans determine the narrower page attributions.
Record 719 precedes 718, and Gelfand's February exercises (733) precede
Beilinson's programme (732). The year heading and programme remain
records without invented mathematics. P001460 is a reading list.
P001478 explicitly begins 12 September 1988; P001479 continues that
next academic year and is excluded, with the boundary documented in
the corrections. The source catalogue is \citet{mccme2026gelfand}.

\paragraph{The mathematical year.}
Beilinson's course moves from sheaf cohomology, resolutions and derived
categories to microlocal and perverse sheaves, then to Morse and gauge
theory. Gelfand's interventions connect symmetric functions, Schubert
polynomials, matroids and Hecke algebras with explicit calculations.
Other sustained strands are dynamical zeta functions, joint stochastic
spectra, toric vector bundles, conformal field theory, principal angles,
quantum integrable systems, exceptional bundles, graph colouring and
quantum-group link invariants. Their coexistence is a feature of the
seminar; it does not imply that every juxtaposition is a precise theorem.

\begin{center}
\begin{tabular}{@{}lrrrrrrr@{}}
\toprule
Class & A & B & C & D & E & F & Total\\
\midrule
Items & 11 & 16 & 3 & 0 & 235 & 22 & 287\\
\bottomrule
\end{tabular}
\end{center}
The occurrence index has 287 items; the problem index has 27
definite questions or conjectures. A--F describe historical fate,
not mathematical correctness or originality. The separate assessment
fields preserve that distinction. An established theorem can have a
false literal notebook formula, and an unresolved wording can concern
a well-developed subject. No unsuccessful search is promoted into an
orphan-idea claim.

\paragraph{What became of the talks.}
Exact theorem matching is more informative than subject matching.
Klyachko's filtration construction led to the 1989 Russian paper
(English translation 1990), submitted in February 1988, including
cohomology and characteristic classes
\citep[Sections~2--4]{klyachko1989equivariant}.
Conversely, some material that looks like a subsequent discovery was
already circulating: the Schubert-character construction had a 1986
preprint and 1987 announcement (item 1987-720.21). Reshetikhin's reports
carry 1987 preprint identifiers, while the scanned editions have 1988
imprints; these are distinguished rather than treated as one date.
Publication chronology is therefore checked
against first circulation wherever it could change a class.

\paragraph{Open questions and the selected development.}
The joint stochastic-spectrum question, GP-1987-21, concerns a whole
multiset, not the Karpelevich region for one eigenvalue. The development
gives complete proofs in dimensions two and three, structural necessary
conditions, and a five-dimensional counterexample to sufficiency of
all power-sum inequalities. Recent work still addresses the general
realization problem \citep{johnsonpaparella2025perron}.
The microlocal reconstruction question, GP-1987-46, retains the
codimension-two essential-image conditions explicitly singled out in
\citet[Problem~9.3]{kashiwaravilonen2014codimension}; a fully faithful
embedding alone is insufficient. The named Riemann hypothesis offers
no bounded seminar-specific route. These latter questions receive
precise status entries without manufactured partial solutions.
Three additional annotations repair a bracket that fails Jacobi,
settle the finite-dimensional tensor-angle range, and disprove the
notebook's equivalence between integral
homology-sphere topology and vanishing twisted first cohomology.
All supplied proofs are marked, independently rechecked, and make no
novelty claim; computational output is included.

\paragraph{Continuity and limits.}
The Schubert material continues 1985-661.1--8, matroid representability
continues 1985-658.12--16, the Virasoro discussion returns to
GP-1982-03 and 1985-657, and determinant lines connect with
1985-655 and 1985-666. Earlier occurrence identifiers are retained
without pretending that a thematic resemblance is the same problem.
Primary papers, author manuscripts, Math-Net records and scans carry
the argument; metadata-only checks and unresolved leads remain
explicit. The 85 corrections distinguish mistranscription,
translation, ambiguity, faithful mathematical error and provenance.
Searches establish the bounded conclusions recorded here, not the
absence of all unpublished work.

\GStalk{714}{1987-10-05}{A. A. Beilinson}{P001385-P001388}{Sheaf cohomology and the preparation for Verdier duality}
Following the previous notebook's discussion of sheaves, Beilinson develops sheafification, exactness, flabby resolutions, and cohomology. A finite-poset exercise connects the general formalism with combinatorial sheaves; the last page begins the compact-support dual construction needed for Verdier duality. These are foundational results, with hypotheses suppressed in several classroom formulations.

\begin{GSitem}{1987-714.1}{E}{The abelian category of sheaves}
\GSfield{Kind; source}{theorem; \GSpage{P001385-P001386}}
\GSfield{The notes}{Kernels are defined sectionwise and cokernels by associating a sheaf to the sectionwise cokernel; Exercise 2 identifies kernel of the cokernel with cokernel of the kernel.}
\GSfield{Statement}{For a topological space \(X\), sheaves of complex vector spaces form an abelian category. For \(f:\mathcal F\to\mathcal G\), \((\ker f)(U)=\ker(\mathcal F(U)\to\mathcal G(U))\), while \(\operatorname{coker}f\) is the sheafification of the sectionwise cokernel; \(\operatorname{coim}f\simeq\operatorname{im}f\).}
\GSfield{Fate}{This is pre-seminar sheaf homological algebra; see \citet[\S 3]{grothendieck1957homological} and the modern sheafification construction in \citet[Tag 007X]{stacks2026foundations}.}
\GSfield{Present status}{As of 11 November 2026, this remains the standard abelian-category structure; the sectionwise cokernel need not itself be a sheaf.}
\GSfield{Assessment}{literal statement: correct; publication match: exact; verification: read the result; development: none; confidence: high}
\end{GSitem}
\begin{GSitem}{1987-714.2}{E}{Sheafification and its universal property}
\GSfield{Kind; source}{construction; \GSpage{P001385-P001386}}
\GSfield{The notes}{The notes glue compatible local representatives, denote the associated sheaf by \(C^+\), and require unique factorization of every map from \(C\) to a sheaf.}
\GSfield{Statement}{For every presheaf \(C\) of complex vector spaces there is a sheaf \(aC\) and a map \(C\to aC\) such that \(\operatorname{Hom}(aC,\mathcal F)\cong\operatorname{Hom}(C,\mathcal F)\) naturally for every sheaf \(\mathcal F\); \(C_x\to(aC)_x\) is an isomorphism for every \(x\).}
\GSfield{Reconstruction}{The stalk formulation disambiguates the abbreviated gluing description; a single compatible-family plus construction requires a separated input, or a second application.}
\GSfield{Fate}{\citet[Lemmas 6.17.1--6.17.2 and the following universal property, Tag 007X]{stacks2026foundations} gives the construction. The notation \(C^+\) here names the completed associated sheaf, not necessarily one application of the usual plus construction.}
\GSfield{Present status}{As of 11 November 2026, the universal property uniquely determines sheafification up to unique isomorphism.}
\GSfield{Assessment}{literal statement: correct with missing hypotheses; publication match: exact; verification: read the result; development: none; confidence: high}
\end{GSitem}
\begin{GSitem}{1987-714.3}{E}{Exactness is tested on stalks}
\GSfield{Kind; source}{theorem; \GSpage{P001386}}
\GSfield{The notes}{Exercise 3 says a sequence of sheaves is exact if and only if all its stalk sequences are exact.}
\GSfield{Statement}{A sequence \(\mathcal F_1\to\mathcal F_2\to\mathcal F_3\) of sheaves of vector spaces on \(X\) is exact at \(\mathcal F_2\) if and only if \((\mathcal F_1)_x\to(\mathcal F_2)_x\to(\mathcal F_3)_x\) is exact for every \(x\in X\).}
\GSfield{Fate}{Sheafification preserves stalks, and filtered colimits of vector spaces are exact; this is part of the classical sheaf formalism of \citet[\S 3]{grothendieck1957homological}.}
\GSfield{Present status}{As of 11 November 2026, the statement holds without separation or local-compactness assumptions on the space.}
\GSfield{Assessment}{literal statement: correct; publication match: exact; verification: read the result; development: none; confidence: high}
\end{GSitem}
\begin{GSitem}{1987-714.4}{E}{Flabby sheaves and the local-system caveat}
\GSfield{Kind; source}{definition; \GSpage{P001386}}
\GSfield{The notes}{A sheaf is called flabby when every section on an open subset extends to all of \(X\). The following sentence says that locally constant sheaves are not flabby.}
\GSfield{Statement}{\(\mathcal F\) is flabby if \(\mathcal F(X)\to\mathcal F(U)\) is surjective for every open \(U\subset X\), equivalently all open restriction maps are surjective. A nonzero constant sheaf on a connected locally connected space with a disconnected open subset is not flabby. The zero sheaf and every sheaf on a discrete space give exceptions to a blanket claim about local systems.}
\GSfield{Reconstruction}{The extra topological and nonzero hypotheses state the situation illustrated by the notebook drawing.}
\GSfield{Fate}{The definition and acyclicity are classical; \citet[Definition 20.12.1 and Lemma 20.12.3, Tag 09SV]{stacks2026foundations}. The exceptions follow directly from the definition.}
\GSfield{Present status}{As of 11 November 2026, the definition is correct, but the unqualified assertion about locally constant sheaves is false.}
\GSfield{Assessment}{literal statement: false as written; publication match: exact; verification: read the result; development: none; confidence: high}
\end{GSitem}
\begin{GSitem}{1987-714.5}{E}{Sheaves and flabby sheaves on a finite poset}
\GSfield{Kind; source}{problem; \GSpage{P001386-P001387}}
\GSfield{The notes}{Problem 4 equips a finite poset \(P\) with \(\overline{\{x\}}=\{y:y\le x\}\), asks for sheaves as representations, and asks which are flabby.}
\GSfield{Statement}{With upper sets open, a sheaf on finite \(P\) is equivalently a covariant functor \(F:P\to\operatorname{Vect}_{\mathbb C}\). Its sections on an upper set \(U\) are \(\varprojlim_{x\in U}F(x)\). It is flabby exactly when \(F(x)\to\varprojlim_{y>x}F(y)\) is surjective for every \(x\in P\), with the empty limit zero.}
\GSfield{Reconstruction}{The inverse-limit criterion and its two-direction argument are our explicit solution of the recorded exercise. \ours{} The stated short argument is supplied in this companion.}
\GSfield{Fate}{The diagram description is elementary sheaf theory; the stated flabbiness test is reconstructed directly: extend a compatible section over an upper set by adding its remaining points in reverse order. Necessity follows by restricting from \(\{y:y\ge x\}\) to \(\{y:y>x\}\). Compare the definition in \citet[Tag 09SV]{stacks2026foundations}.}
\GSfield{Present status}{As of 11 November 2026, the exercise has this finite and effective answer; surjectivity on individual comparable pairs alone is insufficient at branching points.}
\GSfield{Assessment}{literal statement: correct; publication match: partial; verification: read the result; development: annotated; confidence: high}
\end{GSitem}
\begin{GSitem}{1987-714.6}{E}{Combinatorial sheaves and the poset of strata}
\GSfield{Kind; source}{claim; \GSpage{P001387}}
\GSfield{The notes}{The map sending a point to its stratum is said to induce an equivalence from sheaves on the poset of strata to combinatorial sheaves on \(X\).}
\GSfield{Statement}{For a finite regular cell complex with its face poset, cellular sheaves can be encoded by vector spaces on cells and compatible maps along face inclusions. For an arbitrary stratification the stated equivalence requires additional hypotheses: its exact truth depends on what is meant here by a combinatorial sheaf and on incidence and monodromy data.}
\GSfield{Reconstruction}{The regular-cell formulation is a limited modern interpretation, not a claim that the notebook imposed that condition.}
\GSfield{Fate}{The finite-poset assertion is precise, but the notes refer back to definitions in a previous notebook. The present pages do not provide the assumptions needed to identify the claimed equivalence for all stratifications; the general sheaf formalism in \citet[\S 3]{grothendieck1957homological} does not supply them automatically.}
\GSfield{Present status}{As of 11 November 2026, a standard cellular interpretation is available, while the unrestricted stratified-space claim cannot be certified from this record alone.}
\GSfield{Assessment}{literal statement: not determinable; publication match: partial; verification: indirect; development: none; confidence: medium}
\end{GSitem}
\begin{GSitem}{1987-714.7}{E}{Exactness of sections when the kernel is flabby}
\GSfield{Kind; source}{theorem; \GSpage{P001387}}
\GSfield{The notes}{The left-exact global-section functor becomes exact on a short exact sequence when its first term is flabby.}
\GSfield{Statement}{If \(0\to\mathcal F\to\mathcal G\to\mathcal H\to0\) is exact and \(\mathcal F\) is flabby, then \(0\to\Gamma(X,\mathcal F)\to\Gamma(X,\mathcal G)\to\Gamma(X,\mathcal H)\to0\) is exact.}
\GSfield{Fate}{\citet[Lemma 20.12.3 and its proof, Tag 09SV]{stacks2026foundations} proves precisely the needed lifting property. It belongs to the older flabby-resolution method.}
\GSfield{Present status}{As of 11 November 2026, the assertion is valid on every topological space.}
\GSfield{Assessment}{literal statement: correct; publication match: exact; verification: read the result; development: none; confidence: high}
\end{GSitem}
\begin{GSitem}{1987-714.8}{E}{Flabby resolutions and sheaf cohomology}
\GSfield{Kind; source}{construction; \GSpage{P001387-P001388}}
\GSfield{The notes}{Every sheaf embeds in a flabby one; iterating on cokernels gives a resolution, and its global-section cohomology is independent of the resolution.}
\GSfield{Statement}{For every sheaf \(\mathcal F\), the map into the sheaf \(U\mapsto\prod_{x\in U}\mathcal F_x\) is injective and the target is flabby. Iteration gives \(0\to\mathcal F\to I^0\to I^1\to\cdots\). The groups \(H^i(\Gamma(X,I^\bullet))\) are independent of a chosen flabby resolution and equal \(R^i\Gamma(X,\mathcal F)\).}
\GSfield{Reconstruction}{The product-of-stalks construction makes the asserted embedding explicit; independence is via the acyclic-resolution comparison theorem.}
\GSfield{Fate}{\citet[\S 3]{grothendieck1957homological} supplies the derived-functor framework; flabby acyclicity is checked in \citet[Lemma 20.12.3, Tag 09SV]{stacks2026foundations}. This is exposition, not a new 1987 construction.}
\GSfield{Present status}{As of 11 November 2026, the resolution computes sheaf cohomology in all nonnegative degrees.}
\GSfield{Assessment}{literal statement: correct; publication match: exact; verification: read the result; development: none; confidence: high}
\end{GSitem}
\begin{GSitem}{1987-714.9}{E}{Cohomology with compact supports}
\GSfield{Kind; source}{construction; \GSpage{P001388}}
\GSfield{The notes}{The same resolution construction is proposed with compactly supported sections on a locally compact space.}
\GSfield{Statement}{For a locally compact Hausdorff space \(X\), \(\Gamma_c(X,\mathcal F)\) consists of sections with compact support, and \(H_c^i(X,\mathcal F)=R^i\Gamma_c(X,\mathcal F)\). A resolution by \(c\)-soft sheaves, in particular an appropriate flabby resolution, computes these groups.}
\GSfield{Reconstruction}{Hausdorffness and the acyclicity condition are supplied explicitly.}
\GSfield{Fate}{This is the standard compact-support construction underlying the duality formalism already presented in \citet[Chapters II--III]{hartshorne1966residues}; only the publisher chapter metadata was reopened for that older volume.}
\GSfield{Present status}{As of 11 November 2026, the locally compact Hausdorff setting gives the usual compact-support theory; unrestricted support claims require a specified family of supports.}
\GSfield{Assessment}{literal statement: correct with missing hypotheses; publication match: partial; verification: publisher metadata checked; development: none; confidence: high}
\end{GSitem}
\begin{GSitem}{1987-714.10}{E}{The dual of compactly supported sections}
\GSfield{Kind; source}{construction; \GSpage{P001388}}
\GSfield{The notes}{The last construction sets \(\mathcal F^*(V)=\Gamma_c(V,\mathcal F)^*\), with restrictions dual to extension by zero, and relates it to flabbiness.}
\GSfield{Statement}{For sheaves of complex vector spaces on locally compact Hausdorff \(X\), set \(D_c\mathcal F(U)=\operatorname{Hom}_{\mathbb C}(\Gamma_c(U,\mathcal F),\mathbb C)\). If \(\mathcal F\) is \(c\)-soft, compact-support Mayer--Vietoris makes this presheaf a sheaf. Its restriction maps are surjective because linear functionals extend from vector subspaces, so it is flabby.}
\GSfield{Reconstruction}{The sufficient c-soft hypothesis and the Mayer--Vietoris explanation specify the missing sheaf condition.}
\GSfield{Fate}{This is the elementary compact-support dual used in derived duality; compare the 1966 duality framework \citep[Chapters II--III]{hartshorne1966residues}. It is not an assertion that naive pointwise dualization of an arbitrary sheaf is Verdier duality.}
\GSfield{Present status}{As of 11 November 2026, the construction is valid under the stated acyclicity condition; derived duality requires applying it to a suitable complex.}
\GSfield{Assessment}{literal statement: correct with missing hypotheses; publication match: partial; verification: indirect; development: none; confidence: high}
\end{GSitem}

\GStalk{715}{1987-10-12}{I. M. Gelfand}{P001389-P001390}{Dyson's constant term and Good's proof}
Gelfand presents the Dyson constant-term identity and the short interpolation proof attributed in the notebook to a name written as \ru{Вуда}. The exact recurrence and bibliographic year identify I. J. Good's 1970 proof. The final induction ends at the top of P001390, before Veselov's heading.

\begin{GSitem}{1987-715.1}{E}{Dyson's constant-term identity}
\GSfield{Kind; source}{theorem; \GSpage{P001389-P001390}}
\GSfield{The notes}{For nonnegative integral \(a_j\), the constant term of \(F(x;a)=\prod_{i\ne j}(1-x_j/x_i)^{a_j}\) is asserted to be \((\sum a_j)!/\prod a_j!\).}
\GSfield{Statement}{For every \(n\ge1\) and \(a\in\mathbb Z_{\ge0}^n\), \[\operatorname{CT}_{x_1,\ldots,x_n}\prod_{i\ne j}(1-x_j/x_i)^{a_j}=\frac{(a_1+\cdots+a_n)!}{a_1!\cdots a_n!}.\]}
\GSfield{Fate}{\citet[p.~1884]{good1970short} states this exact identity and records the 1962 proofs by Gunson and Wilson; Good's paper was received 26 December 1969 and published June 1970.}
\GSfield{Present status}{As of 11 November 2026, the identity is an established theorem, and the seminar occurrence is expository.}
\GSfield{Assessment}{literal statement: correct; publication match: exact; verification: read the result; development: none; confidence: high}
\end{GSitem}
\begin{GSitem}{1987-715.2}{E}{Good's interpolation recurrence}
\GSfield{Kind; source}{method; \GSpage{P001389-P001390}}
\GSfield{The notes}{Interpolation is applied to a function printed as identically zero, but the ensuing sum is one and is used to reduce the exponents one at a time.}
\GSfield{Statement}{For distinct nonzero \(x_i\), interpolation of \(1\) gives \(\sum_j\prod_{i\ne j}(1-x_j/x_i)^{-1}=1\). Thus the constant terms satisfy \(G(a)=\sum_jG(a-e_j)\) if all \(a_j>0\), the deletion rule when \(a_j=0\), and \(G(0)=1\). These rules uniquely determine the multinomial coefficient.}
\GSfield{Fate}{The complete proof is \citet[p.~1884, equations (1)--(3)]{good1970short}. The source explicitly interpolates the constant function one.}
\GSfield{Present status}{As of 11 November 2026, the proof works after correcting that constant; no limiting interpretation of a zero exponent is needed.}
\GSfield{Assessment}{literal statement: false as written; publication match: exact; verification: read the result; development: annotated; confidence: high}
\end{GSitem}

\GStalk{716}{1987-10-12}{A. P. Veselov}{P001390-P001391}{Quadrics, the Neumann system, and two Poisson brackets}
Veselov compares geodesics on a quadric with the Neumann oscillator on a sphere using Kn\"orrer's time change and Gauss map. He then presents a spectral-operator interpretation and two Hamiltonian descriptions in dimension three. Gelfand asks an inverse problem for a rank-one matrix perturbation; its uniqueness needs qualifications absent from the shorthand question.

\begin{GSitem}{1987-716.1}{E}{Geodesics and the constrained Neumann equation}
\GSfield{Kind; source}{theorem; \GSpage{P001390}}
\GSfield{The notes}{The two columns write \(\ddot q=\lambda Aq\) on \((Aq,q)=1\), and \(x^{\prime\prime}=-Ax+ux\) on \(|x|^2=1\).}
\GSfield{Statement}{Let \(A\) be positive definite symmetric. A geodesic \(q(s)\) on \(q^TAq=1\), in arclength, satisfies \(q_{ss}=\lambda Aq\), where \(\lambda=-q_s^TAq_s/\|Aq\|^2\). Motion on the unit sphere in potential \(x^TAx/2\) satisfies \(x^{\prime\prime}=-Ax+ux\), with \(u=x^TAx-\|x^{\prime}\|^2\).}
\GSfield{Reconstruction}{Positive definiteness and arclength fix the ellipsoid case and the signs of the multipliers.}
\GSfield{Fate}{These constrained equations are classical. The derivation and normalization are reproduced by \citet[\S\S 2--4]{veselovwu2020geodesic}, whose later work explicitly distinguishes Kn\"orrer's 1982 theorem from its own hyperboloid scattering results.}
\GSfield{Present status}{As of 11 November 2026, both equations follow directly by differentiating their constraints.}
\GSfield{Assessment}{literal statement: correct with missing hypotheses; publication match: exact; verification: read the result; development: none; confidence: high}
\end{GSitem}
\begin{GSitem}{1987-716.2}{E}{Knörrer's Gauss-map correspondence}
\GSfield{Kind; source}{theorem; \GSpage{P001390}}
\GSfield{The notes}{The notes credit H. Kn\"orrer, 1982: after changing the parameter, \(x=Aq/|Aq|\) takes geodesics into Neumann trajectories.}
\GSfield{Statement}{For a nonconstant ellipsoid geodesic as above, choose \(d\tau/ds=\sqrt{-\lambda(s)}\). Then \(x(\tau)=Aq(\tau)/\|Aq(\tau)\|\) satisfies the Neumann equation with matrix \(A\). The image is a fixed level of an additional Neumann integral; it is not all Neumann trajectories for the one fixed quadric.}
\GSfield{Reconstruction}{The explicit time change and restriction on the converse are supplied from the primary later calculation.}
\GSfield{Fate}{\citet[Theorem 4.1, equations (28)--(33)]{veselovwu2020geodesic} gives the calculation and explicitly attributes it to Knörrer 1982. The original DOI was located \citep{knorrer1982geodesics}, but the original full paper could not be opened.}
\GSfield{Present status}{As of 11 November 2026, the forward correspondence is established; the extra integral matters for a converse.}
\GSfield{Assessment}{literal statement: correct with missing hypotheses; publication match: exact; verification: read the result; development: none; confidence: high}
\end{GSitem}
\begin{GSitem}{1987-716.3}{E}{The scalar spectral operator and its gauge transform}
\GSfield{Kind; source}{construction; \GSpage{P001390-P001391}}
\GSfield{The notes}{Each coordinate is an eigenfunction of \(L=-d^2/dt^2+u\); a first-derivative operator \(\widetilde L=-d^2/dt^2+v\,d/dt\) is called conjugate to it by the Gauss or Miura substitution.}
\GSfield{Statement}{If \(x_i^{\prime\prime}=-a_ix_i+ux_i\), then \(Lx_i=a_ix_i\). For \(r>0\) with \(v=2r^{\prime}/r\), \[r^{-1}(-D^2+vD)r=-D^2+\frac{v^2}{4}-\frac{v^{\prime}}2.\] In the Gauss substitution take \(r=\|Aq\|\), so constant multiplication by \(A\) commutes with these scalar operators.}
\GSfield{Reconstruction}{The gauge multiplier, potential, interval, and absence of an imposed operator domain are made explicit. \oursComputed{} The independent checks are recorded in the early-record verification log.}
\GSfield{Fate}{The first identity is immediate from the differential equation. The gauge calculation is independently expanded here; it agrees with the time-change calculation in \citet[\S 4]{veselovwu2020geodesic}.}
\GSfield{Present status}{As of 11 November 2026, this is an operator conjugacy on any interval where the multiplier is nonzero, not a global spectral assertion without boundary conditions.}
\GSfield{Assessment}{literal statement: correct with missing hypotheses; publication match: partial; verification: read the result; development: annotated; confidence: high}
\end{GSitem}
\begin{GSitem}{1987-716.4}{E}{Recovering a rank-one perturbation from its spectrum}
\GSfield{Kind; source}{problem; \GSpage{P001391}; I. M. Gelfand}
\GSfield{The notes}{Gelfand fixes a symmetric matrix \(A\), gives the eigenvalues of \(A+K\) with \(\operatorname{rank}K=1\), and asks how to recover \(K\).}
\GSfield{Statement}{For \(A=\operatorname{diag}(\alpha_1,\ldots,\alpha_n)\) with distinct real entries and symmetric positive rank-one \(K=vv^T\), a prescribed characteristic polynomial \(p_\beta(z)=\prod_j(z-\beta_j)\) requires \[v_i^2=-\frac{p_\beta(\alpha_i)}{\prod_{k\ne i}(\alpha_i-\alpha_k)}.\] Nonnegative right sides, with at least one strictly positive, are necessary and sufficient for rank exactly one. The signs of nonzero \(v_i\), modulo a common sign, remain undetermined. For arbitrary \(K=uv^T\), only the products \(u_iv_i\) are recovered.}
\GSfield{Reconstruction}{Simple spectrum, positivity, and the distinction between a matrix and its diagonal-sign-conjugacy class are additional assumptions; the residue argument is our verification. \oursComputed{} The independent checks are recorded in the early-record verification log.}
\GSfield{Fate}{The matrix determinant lemma yields \(p_\beta(z)/p_\alpha(z)=1-\sum_i v_i^2/(z-\alpha_i)\); taking residues proves the formula. This is a classical rank-one secular equation, related to the spectral viewpoint in \citet[\S 2]{veselovwu2020geodesic}, but the notebook's bare reconstruction question is not uniquely solvable.}
\GSfield{Present status}{As of 11 November 2026, the normalized symmetric inverse problem is explicit; uniqueness of the matrix in a fixed basis is false without extra data.}
\GSfield{Assessment}{literal statement: correct with missing hypotheses; publication match: partial; verification: indirect; development: annotated; confidence: high}
\end{GSitem}
\begin{GSitem}{1987-716.5}{E}{The Euclidean Lie--Poisson realization of Neumann motion}
\GSfield{Kind; source}{claim; \GSpage{P001391}}
\GSfield{The notes}{The notes set \(M=[\dot x,x]\), \(p=x\), use the standard \(e(3)\) brackets, and take \(H=(M^2+(Ap,p))/2\).}
\GSfield{Statement}{With \(\varepsilon_{123}=1\), brackets \(\{M_i,M_j\}=\varepsilon_{ijk}M_k\), \(\{M_i,p_j\}=\varepsilon_{ijk}p_k\), \(\{p_i,p_j\}=0\), and convention \(\dot f=\{f,H\}\), the Hamiltonian gives \(\dot p=M\times p\) and \(\dot M=(Ap)\times p\). On \(p^2=1\), \(M\cdot p=0\), one has \(M=p\times\dot p\) and \(\ddot p=-Ap+(p^TAp-\dot p^2)p\).}
\GSfield{Reconstruction}{Orientation and Hamiltonian-vector-field conventions are specified to make the sign comparison meaningful. \oursComputed{} The independent checks are recorded in the early-record verification log.}
\GSfield{Fate}{The equations are verified by contracting the displayed brackets; the Neumann normalization agrees with \citet[\S 3]{veselovwu2020geodesic}. The handwritten \([\dot x,x]\) has the opposite sign under the usual cross-product convention.}
\GSfield{Present status}{As of 11 November 2026, the Hamiltonian realization is standard, with the momentum sign corrected or an opposite cross-product convention declared.}
\GSfield{Assessment}{literal statement: false as written; publication match: partial; verification: read the result; development: annotated; confidence: high}
\end{GSitem}
\begin{GSitem}{1987-716.6}{E}{The second bracket fails Jacobi as written}
\GSfield{Kind; source}{claim; \GSpage{P001391}}
\GSfield{The notes}{The second table reads \(\{M_i,M_j\}_2=\varepsilon_{ijk}a_kM_k\), \(\{M_i,p_j\}_2=\varepsilon_{ijk}a_kp_k\), \(\{p_i,p_j\}_2=\varepsilon_{ijk}M_k\), with \(F=p^2/2\). It is said to give another Hamiltonian description.}
\GSfield{Statement}{The displayed table is not a Poisson bracket for a general diagonal \(A\): its Jacobiator on \((M_1,p_1,p_2)\) is \((a_3-a_2)M_2\). Its formal vector field for \(F=p^2/2\) is also the negative of the first recorded flow. One algebraically valid repair keeps \(\{M_i,M_j\}_2=\varepsilon_{ijk}a_kM_k\), changes \(\{M_i,p_j\}_2\) to \(\varepsilon_{ijk}a_jp_k\), and changes \(\{p_i,p_j\}_2\) to \(-\varepsilon_{ijk}M_k\). This repaired bracket is Poisson and gives \(\dot M=Ap\times p,\ \dot p=M\times p\) with the same \(F\).}
\GSfield{Reconstruction}{The mixed coefficient and the momentum-momentum sign are changed explicitly; simply reversing time does not cure the failed Jacobi identity. \oursComputed{} The independent checks are recorded in the early-record verification log.}
\GSfield{Fate}{The scan confirms the literal indices and sign. The Jacobi counterexample, all coordinate Jacobi identities for the repaired table, and the two Hamiltonian vector fields are independently calculated in the accompanying script. The underlying Neumann system is classical \citep[\S 3]{veselovwu2020geodesic}; that citation is not an exact source for this chosen repair. The independent associative-matrix construction in Section~\ref{dev:poisson-repair} proves Jacobi and compatibility for the repair.}
\GSfield{Present status}{As of 11 November 2026, the literal second Hamiltonian description is false; the displayed repair is a verified reconstruction, not a uniquely established reading of the intended historical formula.}
\GSfield{Assessment}{literal statement: false as written; publication match: partial; verification: indirect; development: annotated; confidence: high}
\end{GSitem}

\GStalk{717}{1987-10-12}{A. A. Beilinson}{P001392-P001394}{Complexes, homotopy, and localization}
Beilinson moves from poset representations to complexes over an abelian category, homotopies, and quasi-isomorphisms. Roof diagrams and the Ore construction introduce the derived category. The same localization is revisited on 23 November; these occurrences are kept separate, with a cross-link rather than treated as new research.

\begin{GSitem}{1987-717.1}{E}{Combinatorial sheaves as representations, recalled}
\GSfield{Kind; source}{definition; \GSpage{P001392}}
\GSfield{The notes}{The opening reminder assigns objects to points of a poset and compatible morphisms to ordered pairs.}
\GSfield{Statement}{A representation of a poset \(P\) in an abelian category \(\mathcal A\) is a family \(F_x\) with maps \(f_{xy}:F_x\to F_y\) for \(x\le y\), satisfying \(f_{xx}=1\) and \(f_{yz}f_{xy}=f_{xz}\). For finite \(P\) with upper-set topology this is the corresponding cellular sheaf diagram.}
\GSfield{Fate}{This repeats the finite-poset construction in occurrence 1987-714.5; its elementary sheaf interpretation uses the same restriction and gluing axioms \citep[Tag 007X]{stacks2026foundations}.}
\GSfield{Present status}{As of 11 November 2026, this is an established categorical encoding, with the direction fixed by the notebook's topology.}
\GSfield{Assessment}{literal statement: correct; publication match: exact; verification: read the result; development: none; confidence: high}
\end{GSitem}
\begin{GSitem}{1987-717.2}{E}{The abelian category of complexes and cohomology}
\GSfield{Kind; source}{definition; \GSpage{P001392-P001393}}
\GSfield{The notes}{The notes introduce \(C(\mathcal A)\), bounded variants, differentials, morphisms, and \(H^i\).}
\GSfield{Statement}{For an abelian category \(\mathcal A\), \(C(\mathcal A)\) has cochain complexes \((X^i,d^i)\), \(d^{i+1}d^i=0\), and degree-preserving chain maps. It is abelian with kernels and cokernels degreewise, and \(H^i(X)=\ker d^i/\operatorname{im}d^{i-1}\). The subcategories \(C^b,C^+,C^-\) impose bounded, bounded-below, or bounded-above terms.}
\GSfield{Fate}{This predates the seminar; the derived-category account was already published in \citet[Chapter I]{hartshorne1966residues}. Modern definitions and cohomology are in \citet[Chapters 12--13]{stacks2026foundations}.}
\GSfield{Present status}{As of 11 November 2026, these are standard definitions; cohomology need not be exact as a functor on complexes.}
\GSfield{Assessment}{literal statement: correct; publication match: exact; verification: read the result; development: none; confidence: high}
\end{GSitem}
\begin{GSitem}{1987-717.3}{E}{Homotopy and the quotient category}
\GSfield{Kind; source}{construction; \GSpage{P001393}}
\GSfield{The notes}{A null homotopy is written \(f=d\psi+\psi d\); quotienting by such maps defines \(K(\mathcal A)\).}
\GSfield{Statement}{Chain maps \(f,g:X\to Y\) are homotopic if \(f^i-g^i=d_Y^{i-1}h^i+h^{i+1}d_X^i\) for maps \(h^i:X^i\to Y^{i-1}\). Null-homotopic maps form an additive two-sided ideal. The quotient \(K(\mathcal A)\) has the same objects as \(C(\mathcal A)\), and every \(H^i\) factors through it.}
\GSfield{Fate}{\citet[\S 13.8 and \S 13.11, Tag 05RR]{stacks2026foundations} gives precisely this construction; its pre-1987 setting is \citet[Chapter I]{hartshorne1966residues}.}
\GSfield{Present status}{As of 11 November 2026, homotopy is weaker than isomorphism of complexes and stronger than the derived equivalence generated by quasi-isomorphisms.}
\GSfield{Assessment}{literal statement: correct; publication match: exact; verification: read the result; development: none; confidence: high}
\end{GSitem}
\begin{GSitem}{1987-717.4}{E}{Quasi-isomorphisms}
\GSfield{Kind; source}{definition; \GSpage{P001393-P001394}}
\GSfield{The notes}{A map inducing isomorphisms in all cohomological degrees is called a quasi-isomorphism.}
\GSfield{Statement}{A chain map \(f:X\to Y\) is a quasi-isomorphism if \(H^i(f)\) is an isomorphism for every integer \(i\). Such maps need not be homotopy equivalences in an arbitrary abelian category.}
\GSfield{Fate}{\citet[\S 13.11, Lemma 13.11.2]{stacks2026foundations} identifies precisely this multiplicative system, in the classical derived-category construction.}
\GSfield{Present status}{As of 11 November 2026, the distinction from homotopy equivalence is essential and remains unchanged.}
\GSfield{Assessment}{literal statement: correct; publication match: exact; verification: read the result; development: none; confidence: high}
\end{GSitem}
\begin{GSitem}{1987-717.5}{E}{The Ore construction for a roof}
\GSfield{Kind; source}{theorem; \GSpage{P001394}}
\GSfield{The notes}{A roof with a quasi-isomorphism on one leg can be replaced by a common outgoing diagram; the construction is used to compose fractions.}
\GSfield{Statement}{If \(s:Z\to X\) is a quasi-isomorphism and \(f:Z\to Y\) is a chain map, the homotopy pushout \(P=\operatorname{Cone}((s,-f):Z\to X\oplus Y)\) admits maps \(X\to P\), \(Y\to P\) commuting in \(K(\mathcal A)\), and \(Y\to P\) is a quasi-isomorphism. The dual statement uses a homotopy pullback.}
\GSfield{Reconstruction}{The explicit cone and the location of the outgoing quasi-isomorphism make the diagrammatic lemma precise.}
\GSfield{Fate}{This is the calculus of quasi-isomorphisms underlying \citet[Lemma 13.11.2 and \S 13.5]{stacks2026foundations}; the cone construction is classical \citep[Chapter I]{hartshorne1966residues}.}
\GSfield{Present status}{As of 11 November 2026, the outgoing quasi-isomorphism is on the opposite leg; commutativity is in the homotopy category.}
\GSfield{Assessment}{literal statement: correct; publication match: exact; verification: read the result; development: none; confidence: high}
\end{GSitem}
\begin{GSitem}{1987-717.6}{E}{The derived category as roofs}
\GSfield{Kind; source}{construction; \GSpage{P001394}}
\GSfield{The notes}{Morphisms are roofs \(X\leftarrow Z\to Y\), the left map a quasi-isomorphism, modulo common refinement.}
\GSfield{Statement}{For an essentially small abelian category, \(D(\mathcal A)=K(\mathcal A)[\mathrm{Qis}^{-1}]\). A morphism is represented by \(X\xleftarrow{s}Z\xrightarrow{f}Y\) with \(s\) a quasi-isomorphism. Equivalent roofs have common refinements in \(K(\mathcal A)\); composition is obtained by an Ore square. The localization is universal among functors inverting quasi-isomorphisms.}
\GSfield{Reconstruction}{The size convention is stated to avoid assuming that arbitrary localization morphisms form sets.}
\GSfield{Fate}{\citet[Definition 13.11.3, Tag 05RR]{stacks2026foundations}; the foundational account appears in \citet[Chapter I]{hartshorne1966residues}. Earlier companion occurrences include 1971-266.3 and 1977-438.15.}
\GSfield{Present status}{As of 11 November 2026, this is the established derived-category construction; large categories require the usual universe or existence conventions.}
\GSfield{Assessment}{literal statement: correct with missing hypotheses; publication match: exact; verification: read the result; development: none; confidence: high}
\end{GSitem}

\GStalk{719}{1987-10-19}{I. M. Gelfand}{P001395-P001396}{Summation formulas and a parabolic double-coset algebra}
Gelfand derives a triangular system relating derivative sums to endpoint differences, leading to the Euler--Maclaurin formula. The same short contribution recalls the Coxeter presentation of a symmetric group and introduces the subalgebra invariant on both sides under a Young subgroup. The second topic ends at the top of P001396, before Beilinson resumes.

\begin{GSitem}{1987-719.1}{E}{Taylor telescoping and universal summation coefficients}
\GSfield{Kind; source}{method; \GSpage{P001395}}
\GSfield{The notes}{The sums \(S_r=\sum_{j=1}^n f^{(r)}(j)\) are expressed by a triangular system involving endpoint differences, with a determinant solution.}
\GSfield{Statement}{For a polynomial \(f\) and an antiderivative \(F\), \[F(n+1)-F(1)=\sum_{r\ge0}\frac{S_r}{(r+1)!},\qquad f^{(p)}(n+1)-f^{(p)}(1)=\sum_{r\ge1}\frac{S_{p+r}}{r!}.\] All sums terminate, and the unit triangular system can be inverted.}
\GSfield{Reconstruction}{Polynomial input supplies a rigorous domain for the notebook's formal manipulation.}
\GSfield{Fate}{This is a direct polynomial version of the classical Euler--Maclaurin method. \citet[\S 2.10(i), equation 2.10.1]{nist2026dlmf} states the analytic formula with its remainder.}
\GSfield{Present status}{As of 11 November 2026, the polynomial identity is exact; applying an infinite Taylor expansion to an arbitrary function is unjustified.}
\GSfield{Assessment}{literal statement: correct with missing hypotheses; publication match: exact; verification: read the result; development: none; confidence: high}
\end{GSitem}
\begin{GSitem}{1987-719.2}{E}{Bernoulli coefficients and Euler--Maclaurin summation}
\GSfield{Kind; source}{theorem; \GSpage{P001395}}
\GSfield{The notes}{The universal coefficients are recovered using exponentials; the notes propose the answer for any function.}
\GSfield{Statement}{Writing \(t/(e^t-1)=\sum_{r\ge0}B_rt^r/r!\), a polynomial \(f=F^{\prime}\) satisfies \[\sum_{j=1}^nf(j)=F(n+1)-F(1)+\sum_{r\ge1}\frac{B_r}{r!}\bigl(f^{(r-1)}(n+1)-f^{(r-1)}(1)\bigr).\] For nonpolynomial functions one uses a finite truncation with the Euler--Maclaurin remainder and sufficient differentiability.}
\GSfield{Reconstruction}{The exact polynomial case is separated from the analytic finite-remainder theorem.}
\GSfield{Fate}{\citet[equations 24.2.1 and 2.10.1]{nist2026dlmf} verifies the generating function, sign \(B_1=-1/2\), and remainder. This was classical long before 1987.}
\GSfield{Present status}{As of 11 November 2026, the coefficient sequence is established; the unrestricted no-remainder statement is false.}
\GSfield{Assessment}{literal statement: false as written; publication match: exact; verification: read the result; development: annotated; confidence: high}
\end{GSitem}
\begin{GSitem}{1987-719.3}{E}{Symmetric-group generators and the Young-subgroup algebra}
\GSfield{Kind; source}{definition; \GSpage{P001395-P001396}}
\GSfield{The notes}{The simple transpositions obey Coxeter relations; for \(H=S_k\times S_{n-k}\), the algebra consists of \(x\in\mathbb C[S_n]\) satisfying \(hx=xh=x\) for all \(h\in H\).}
\GSfield{Statement}{The group \(S_n\) is generated by \(s_i=(i,i+1)\), with \(s_i^2=1\), \(s_is_j=s_js_i\) for \(|i-j|>1\), and \(s_is_{i+1}s_i=s_{i+1}s_is_{i+1}\). For \(e_H=|H|^{-1}\sum_{h\in H}h\), the displayed invariant algebra is \(e_H\mathbb C[S_n]e_H\), with unit \(e_H\); it has a basis indexed by \(H\backslash S_n/H\).}
\GSfield{Reconstruction}{The averaging idempotent identifies the algebra and clarifies that its unit need not be the ambient group-algebra unit.}
\GSfield{Fate}{This is the ordinary finite-group double-coset construction, related to the symmetric-group representation framework of \citet[Chapter I, \S 7]{macdonald1995symmetric}. No further theorem or unanswered question about this algebra is recorded here.}
\GSfield{Present status}{As of 11 November 2026, the definition is complete after making its unit explicit.}
\GSfield{Assessment}{literal statement: correct; publication match: partial; verification: indirect; development: none; confidence: high}
\end{GSitem}

\GStalk{718}{1987-10-19}{A. A. Beilinson}{P001396-P001398}{Truncation, mapping cones, and distinguished triangles}
Beilinson continues the derived-category course after Gelfand's intervention. He explains how cohomology in one degree identifies a complex in the derived category, contrasts this with splitting in several degrees, and introduces mapping cones and exact triangles. Several signs and the ordering of cone summands require explicit conventions.

\begin{GSitem}{1987-718.1}{E}{Truncation and a complex with one cohomology object}
\GSfield{Kind; source}{construction; \GSpage{P001396}}
\GSfield{The notes}{The complex is cut off using a kernel, and, when only its zeroth cohomology survives, it is compared with that object in degree zero.}
\GSfield{Statement}{If \(H^i(X)=0\) for \(i\ne0\), the canonical truncation \(\tau_{\le0}X\), with \((\tau_{\le0}X)^0=\ker d^0\), has quasi-isomorphisms \(X\leftarrow\tau_{\le0}X\to H^0(X)[0]\). Hence \(X\simeq H^0(X)[0]\) in \(D(\mathcal A)\).}
\GSfield{Fate}{This is standard truncation in the derived category, already part of \citet[Chapter I]{hartshorne1966residues}; compare \citet[Chapter 13]{stacks2026foundations}.}
\GSfield{Present status}{As of 11 November 2026, the conclusion is canonical in the derived category and does not imply a chain-level isomorphism.}
\GSfield{Assessment}{literal statement: correct; publication match: exact; verification: read the result; development: none; confidence: high}
\end{GSitem}
\begin{GSitem}{1987-718.2}{E}{Cohomology objects do not generally split a complex}
\GSfield{Kind; source}{claim; \GSpage{P001396-P001397}}
\GSfield{The notes}{The notes warn that several cohomology groups do not determine a complex as their direct sum, and that a splitting, when available, need not be canonical.}
\GSfield{Statement}{In a general abelian category a complex need not be isomorphic in \(D(\mathcal A)\) to its graded cohomology with zero differential: the truncation triangles can have nonzero connecting extension classes. If every short exact sequence in \(\mathcal A\) splits, every complex is homotopy equivalent to its graded cohomology, although choices enter a chain-level splitting.}
\GSfield{Reconstruction}{The extension-class explanation specifies the warning rather than imputing a particular example to the speaker.}
\GSfield{Fate}{The distinction is built into the triangle formalism of \citet[Chapter I]{hartshorne1966residues} and \citet[\S\S 13.9--13.11]{stacks2026foundations}.}
\GSfield{Present status}{As of 11 November 2026, nonformality is a real phenomenon; the semisimple hypothesis is sufficient for the stated splitting.}
\GSfield{Assessment}{literal statement: correct; publication match: partial; verification: read the result; development: none; confidence: high}
\end{GSitem}
\begin{GSitem}{1987-718.3}{E}{When is the homotopy category abelian?}
\GSfield{Kind; source}{problem; \GSpage{P001397}}
\GSfield{The notes}{The exercise says \(K(\mathcal A)\) is abelian exactly when \(\mathcal A\) is semisimple; the handwritten explanation speaks of sums of irreducibles.}
\GSfield{Statement}{For an abelian category \(\mathcal A\), \(K(\mathcal A)\) is abelian if and only if every short exact sequence in \(\mathcal A\) splits. In that case it is equivalent to the category of graded objects of \(\mathcal A\). Here semisimple means split exact; a decomposition into simple objects requires the usual additional finiteness or category hypotheses.}
\GSfield{Reconstruction}{The split-exact formulation avoids silently imposing finite length. \ours{} The stated short argument is supplied in this companion.}
\GSfield{Fate}{The implication from split exactness follows by splitting cycles and boundaries. Conversely an abelian triangulated category is split exact, so every morphism has a generalized inverse. For a monomorphism \(i:U\to V\) in \(\mathcal A\), apply this to its degree-zero map to obtain \(g\) with \(igi=i\); cancellation in \(\mathcal A\) gives \(gi=1_U\), forcing all exact sequences in \(\mathcal A\) to split. The triangulated structure is checked in \citet[Proposition 13.10.3]{stacks2026foundations}. Earlier occurrence 1971-266.2 discusses the general failure of abelianness.}
\GSfield{Present status}{As of 11 November 2026, the exercise is a settled structural criterion, with the meaning of semisimple specified.}
\GSfield{Assessment}{literal statement: correct with missing hypotheses; publication match: partial; verification: indirect; development: annotated; confidence: high}
\end{GSitem}
\begin{GSitem}{1987-718.4}{E}{Mapping cone and its elementary homotopies}
\GSfield{Kind; source}{construction; \GSpage{P001397-P001398}}
\GSfield{The notes}{The cone is \(C_f^i=Y^i\oplus X^{i+1}\), but the differential is written with the reversed coordinate ordering. The cone of the identity is contractible; the cone of zero splits.}
\GSfield{Statement}{For a chain map \(f:X\to Y\), put \(C_f^i=Y^i\oplus X^{i+1}\) and \(d(y,x)=(d_Yy+f(x),-d_Xx)\). Then \(C_{1_X}\) is contractible and \(C_0=Y\oplus X[1]\). If \(i:Y\to C_f\) is the inclusion, \(C_i\) is homotopy equivalent to \(X[1]\).}
\GSfield{Reconstruction}{The differential has been reordered to match the declared direct sum.}
\GSfield{Fate}{\citet[\S 13.9, definition and elementary cone lemmas]{stacks2026foundations} gives the standard construction. These statements were established well before the seminar \citep[Chapter I]{hartshorne1966residues}.}
\GSfield{Present status}{As of 11 November 2026, all three cone properties are valid with consistent shift and summand conventions.}
\GSfield{Assessment}{literal statement: false as written; publication match: exact; verification: read the result; development: none; confidence: high}
\end{GSitem}
\begin{GSitem}{1987-718.5}{E}{Distinguished triangles and the sign of rotation}
\GSfield{Kind; source}{definition; \GSpage{P001397-P001398}}
\GSfield{The notes}{Exact triangles are those isomorphic to cone triangles, followed by a continuing sequence of shifts.}
\GSfield{Statement}{The triangle \(X\xrightarrow fY\to C_f\to X[1]\), with one consistent cone convention, is distinguished; all triangles isomorphic to such ones are distinguished. If \(X\xrightarrow fY\xrightarrow gZ\xrightarrow hX[1]\) is distinguished, its rotation is \(Y\xrightarrow gZ\xrightarrow hX[1]\xrightarrow{-f[1]}Y[1]\). Its cohomology sequence is exact.}
\GSfield{Reconstruction}{The negative sign on the rotated final arrow is made explicit; the notebook's continuing sequence alone does not state a sign convention.}
\GSfield{Fate}{\citet[Proposition 13.10.3, proof of TR2, and equation 13.11.1.1]{stacks2026foundations} checks the rotation sign and long exact sequence. The 1966 published account predates this exposition \citep[Chapter I]{hartshorne1966residues}.}
\GSfield{Present status}{As of 11 November 2026, the rotation sign is indispensable even though signs can be moved by coherent changes of cone convention.}
\GSfield{Assessment}{literal statement: correct with missing hypotheses; publication match: exact; verification: read the result; development: none; confidence: high}
\end{GSitem}

\GStalk{720}{1987-11-09}{I. M. Gelfand}{P001399-P001405}{Symmetric functions and Schubert polynomials after Lascoux}
Gelfand develops symmetric functions in alphabet notation, determinant identities, and Schubert polynomials via divided differences. The finite coinvariant ring, its pairing, and the symmetric-group action connect these formulas with flag cohomology. The final request for a character interpretation had already been answered in the appropriate Borel-module category by the 1986 preprint and 1987 announcement of Kra\'skiewicz and Pragacz; the notebook's word irreducible needs qualification.

\begin{GSitem}{1987-720.1}{E}{Complete functions of a formal alphabet}
\GSfield{Kind; source}{definition; \GSpage{P001399}}
\GSfield{The notes}{For a formal sum \(P\) the coefficients \(S^r(P)\) are defined by a product generating function, with examples allowing multiplicities and differences of alphabets.}
\GSfield{Statement}{For a finite alphabet with multiplicities \(P=\sum_jm_ju_j\), \(m_j\in\mathbb Z\), define \[\sum_{r\ge0}h_r[P]t^r=\prod_j(1-tu_j)^{-m_j},\qquad h_0[P]=1,\quad h_r[P]=0\ (r<0).\] Positive multiplicities give ordinary complete symmetric functions; negative ones define virtual alphabets formally.}
\GSfield{Fate}{This is classical symmetric-function notation; \citet[Chapter I, \S 2]{macdonald1995symmetric}. The first edition appeared in 1979.}
\GSfield{Present status}{As of 11 November 2026, the construction is a formal power-series identity; no analytic convergence is required.}
\GSfield{Assessment}{literal statement: correct; publication match: exact; verification: read the result; development: none; confidence: high}
\end{GSitem}
\begin{GSitem}{1987-720.2}{E}{The alphabet homomorphism to triangular Toeplitz matrices}
\GSfield{Kind; source}{theorem; \GSpage{P001399-P001400}}
\GSfield{The notes}{The matrix \(M(P)\) has consecutive complete functions along diagonals and is called Hankel; the notes assert \(M(P+Q)=M(P)M(Q)\).}
\GSfield{Statement}{Let \(M(P)_{ij}=h_{j-i}[P]\) for \(i,j\ge0\). Then \(M(P)\) is an upper triangular Toeplitz matrix with diagonal one and \(M(P+Q)=M(P)M(Q)\), since \(h_r[P+Q]=\sum_{a+b=r}h_a[P]h_b[Q]\). Every matrix entry involves a finite sum.}
\GSfield{Fate}{The identity follows by multiplying the generating series in \citet[Chapter I, \S 2]{macdonald1995symmetric}. The matrix is Toeplitz, since its entries depend on the difference of the indices.}
\GSfield{Present status}{As of 11 November 2026, the homomorphism is correct; the notebook's matrix-name label is incorrect.}
\GSfield{Assessment}{literal statement: false as written; publication match: exact; verification: read the result; development: annotated; confidence: high}
\end{GSitem}
\begin{GSitem}{1987-720.3}{E}{Cauchy--Binet through exterior powers}
\GSfield{Kind; source}{theorem; \GSpage{P001400}}
\GSfield{The notes}{The notes use minors of a matrix and the multiplicativity of its exterior powers.}
\GSfield{Statement}{For compatible finite matrices \(A,B\), and ordered \(k\)-element index sets \(I,J\), \[\det(AB)_{I,J}=\sum_{|L|=k}\det A_{I,L}\det B_{L,J}.\] Equivalently \(\bigwedge^k(AB)=(\bigwedge^kA)(\bigwedge^kB)\); the same identity applies to the triangular matrices here because relevant sums are finite.}
\GSfield{Fate}{The identity is classical multilinear algebra and underlies the determinant manipulations in \citet[Chapter I, \S\S 3--5]{macdonald1995symmetric}.}
\GSfield{Present status}{As of 11 November 2026, the equality is valid over every commutative coefficient ring.}
\GSfield{Assessment}{literal statement: correct; publication match: partial; verification: read the result; development: none; confidence: high}
\end{GSitem}
\begin{GSitem}{1987-720.4}{E}{Skew Schur functions as minors}
\GSfield{Kind; source}{definition; \GSpage{P001400}}
\GSfield{The notes}{Minors of \(M(P)\), indexed by two ordered integer lists, are denoted \(S_{I\mid J}(P)\).}
\GSfield{Statement}{For partitions \(\mu\subseteq\lambda\), padded to a common length \(r\), the skew Schur function is \[s_{\lambda/\mu}[P]=\det\bigl(h_{\lambda_i-\mu_j-i+j}[P]\bigr)_{1\le i,j\le r}.\] Reversing the notebook's increasing index order gives this usual partition notation.}
\GSfield{Reconstruction}{The partition convention is modern; it does not replace the notebook's increasing-code convention without an index change.}
\GSfield{Fate}{\citet[Chapter I, \S 5, skew Jacobi--Trudi identities]{macdonald1995symmetric}.}
\GSfield{Present status}{As of 11 November 2026, these minors give the standard skew Schur functions; an explicit index reversal is necessary when comparing formulas.}
\GSfield{Assessment}{literal statement: correct; publication match: exact; verification: read the result; development: none; confidence: high}
\end{GSitem}
\begin{GSitem}{1987-720.5}{E}{Alternants and the Jacobi--Trudi identity}
\GSfield{Kind; source}{theorem; \GSpage{P001400-P001401}}
\GSfield{The notes}{The notes define Schur polynomials by a quotient of alternants and ask for the determinant identity relating this quotient to complete functions.}
\GSfield{Statement}{For a partition \(\lambda\) of length at most \(n\), \[s_\lambda(x_1,\ldots,x_n)=\frac{\det(x_i^{\lambda_j+n-j})}{\det(x_i^{n-j})}=\det(h_{\lambda_i-i+j}(x_1,\ldots,x_n)).\] The quotient is a symmetric polynomial with integral coefficients, since the alternating numerator is divisible by the Vandermonde.}
\GSfield{Fate}{\citet[Chapter I, \S 3, equation (3.4) and historical notes]{macdonald1995symmetric} attributes the determinant identity to Jacobi.}
\GSfield{Present status}{As of 11 November 2026, the requested proof concerns a classical exact identity.}
\GSfield{Assessment}{literal statement: correct; publication match: exact; verification: read the result; development: none; confidence: high}
\end{GSitem}
\begin{GSitem}{1987-720.6}{E}{Schur coefficients and their representation meaning}
\GSfield{Kind; source}{question; \GSpage{P001400-P001401}}
\GSfield{The notes}{The text asks for explicit coefficients and a representation-theoretic interpretation; an abbreviated reference to a Gelfand--Zelevinsky 1986 paper is attached.}
\GSfield{Statement}{In characteristic zero, \(s_\lambda\) is the character of the Schur functor \(S_\lambda V\) for \(\operatorname{GL}(V)\). In \(s_\lambda=\sum_\mu K_{\lambda\mu}m_\mu\), \(K_{\lambda\mu}\) counts semistandard tableaux of shape \(\lambda\) and content \(\mu\). These give a precise coefficient interpretation, though not a single factorial expression for all coefficients.}
\GSfield{Reconstruction}{The coefficient basis and the polynomial representation category are supplied explicitly.}
\GSfield{Fate}{The character and tableau descriptions predate 1987; \citet[Chapter I, \S 5 and Appendix A]{macdonald1995symmetric}. The notebook's terse Gelfand--Zelevinsky bibliographic pointer could not be matched securely to a particular theorem and is not used as evidence of priority.}
\GSfield{Present status}{As of 11 November 2026, the ordinary Schur question has an established answer; the exact intended 1986 citation remains unresolved.}
\GSfield{Assessment}{literal statement: correct with missing hypotheses; publication match: partial; verification: read the result; development: none; confidence: high}
\end{GSitem}
\begin{GSitem}{1987-720.7}{F}{The K-functor and Sato comparison}
\GSfield{Kind; source}{remark; \GSpage{P001400}}
\GSfield{The notes}{A parenthetical passage connects the determinant formalism with a K-functor and Sato; no operator, category, or assertion is defined there.}
\GSfield{Statement}{No uniquely specified mathematical statement can be recovered from this aside: neither the K-functor nor the proposed Sato correspondence is defined in the supplied pages.}
\GSfield{Fate}{Targeted searches for the alphabet determinant formalism and Sato theory located related symmetric-function literature but no exact result identifiable from the fragment. The surrounding determinant identities remain independently identified.}
\GSfield{Present status}{As of 11 November 2026, this is a provenance and interpretation gap, not an established unsolved theorem.}
\GSfield{Assessment}{literal statement: not determinable; publication match: none found; verification: unresolved lead; development: none; confidence: low}
\end{GSitem}
\begin{GSitem}{1987-720.8}{E}{Plücker relations as the determinant background}
\GSfield{Kind; source}{claim; \GSpage{P001401}}
\GSfield{The notes}{The note connects the identities among minors with Grassmannians and Plücker relations.}
\GSfield{Statement}{For a rank-\(r\) matrix, its ordered maximal minors \(p_I\) satisfy the quadratic Pl\"ucker relations \[\sum_{a=1}^{r+1}(-1)^a p_{i_1\cdots i_{r-1}j_a}\,p_{j_1\cdots\widehat{j_a}\cdots j_{r+1}}=0,\] with each index list interpreted alternately. These are the defining equations of the Pl\"ucker embedded Grassmannian.}
\GSfield{Reconstruction}{The index convention and ambient Grassmannian have been made explicit.}
\GSfield{Fate}{This is the classical determinant geometry also used in \citet[\S 1]{gelfandserganova1987combinatorial}; the notebook returns to the exchange identity in record 725. The full quadratic-generation statement is explicitly recalled as classical in \citet[\S 3, opening paragraph]{speyersturmfels2003tropical}.}
\GSfield{Present status}{As of 11 November 2026, the relation is established; no new higher Plücker constraint is specified here.}
\GSfield{Assessment}{literal statement: correct; publication match: exact; verification: read the result; development: none; confidence: high}
\end{GSitem}
\begin{GSitem}{1987-720.9}{E}{Addition and multiplication of alphabets}
\GSfield{Kind; source}{theorem; \GSpage{P001401}}
\GSfield{The notes}{Besides the addition formula for complete functions, the notes expand a complete function of the product alphabet in products of Schur functions.}
\GSfield{Statement}{For finite alphabets \(X=(x_i)\), \(Y=(y_j)\), \[h_n[X+Y]=\sum_{a+b=n}h_a[X]h_b[Y],\qquad h_n[XY]=\sum_{\lambda\vdash n}s_\lambda[X]s_\lambda[Y].\] Here \(XY\) is the alphabet \((x_iy_j)\); Schur functions with too many rows specialize to zero.}
\GSfield{Fate}{\citet[Chapter I, \S 4, Cauchy identity]{macdonald1995symmetric}.}
\GSfield{Present status}{As of 11 November 2026, both are exact polynomial identities, valid for virtual alphabets by formal specialization.}
\GSfield{Assessment}{literal statement: correct; publication match: exact; verification: read the result; development: none; confidence: high}
\end{GSitem}
\begin{GSitem}{1987-720.10}{E}{Plethysm and its Schur expansion}
\GSfield{Kind; source}{construction; \GSpage{P001401}}
\GSfield{The notes}{The notes form complete functions of an alphabet whose entries are themselves monomials and write an expansion in Schur functions.}
\GSfield{Statement}{For positive integers \(m,n\), \(h_n[h_m]=\sum_{\lambda\vdash mn}c_\lambda s_\lambda\), where \(c_\lambda\in\mathbb Z_{\ge0}\) is the multiplicity of \(S_\lambda V\) in \(\operatorname{Sym}^n(\operatorname{Sym}^mV)\) in characteristic zero and sufficiently large dimension. Plethysm is defined by substitution of the eigenvalue alphabet.}
\GSfield{Reconstruction}{The characteristic-zero representation interpretation explains the coefficients without claiming an explicit rule absent from the notes.}
\GSfield{Fate}{\citet[Chapter I, \S 8, especially (8.2) and (8.10)]{macdonald1995symmetric}. The notes state the operation and expansion, not a quantified positivity conjecture or a uniform combinatorial rule.}
\GSfield{Present status}{As of 11 November 2026, the existence and nonnegativity of this expansion are classical; the fragment does not justify entering a modern plethysm open problem.}
\GSfield{Assessment}{literal statement: correct; publication match: exact; verification: read the result; development: none; confidence: high}
\end{GSitem}
\begin{GSitem}{1987-720.11}{E}{Flag cohomology and Schubert classes}
\GSfield{Kind; source}{theorem; \GSpage{P001401-P001402}}
\GSfield{The notes}{Schubert cells are compared with differential forms and a polynomial presentation of flag cohomology.}
\GSfield{Statement}{For the complete flag manifold \(\operatorname{Fl}(\mathbb C^N)\), \[H^*(\operatorname{Fl}(\mathbb C^N),\mathbb Z)\simeq\mathbb Z[x_1,\ldots,x_N]/(e_1,\ldots,e_N),\qquad\deg x_i=2.\] Schubert classes form an integral basis; their indices are permutations and their complex codimensions are inversion lengths.}
\GSfield{Fate}{\citet[\S\S 2--3]{bernsteingelfandgelfand1973schubert} treats the polynomial and Schubert descriptions. Its 1973 publication precedes this seminar. See earlier occurrences 1971-239.1--2 and 1985-661.3. Earlier occurrences include 1971-239.1--2 and 1985-661.1--2. The opened original note \citep[\S 4, p.~449]{lascouxschutzenberger1982schubert} explicitly recalls the integral flag quotient and Schubert basis, and the product-duality assertion following Theorem 4.3.}
\GSfield{Present status}{As of 11 November 2026, the cohomology is torsion free and concentrated in even real degrees.}
\GSfield{Assessment}{literal statement: correct; publication match: exact; verification: read the result; development: none; confidence: high}
\end{GSitem}
\begin{GSitem}{1987-720.12}{E}{Divided differences and nil-Coxeter relations}
\GSfield{Kind; source}{definition; \GSpage{P001402}}
\GSfield{The notes}{The operator \(\partial_i=(1-s_i)/(x_i-x_{i+1})\) is introduced, with nilpotence, commuting, and braid relations.}
\GSfield{Statement}{On \(\mathbb Z[x_1,x_2,\ldots]\), define \(\partial_if=(f-s_if)/(x_i-x_{i+1})\). Then \(\partial_i^2=0\), \(\partial_i\partial_j=\partial_j\partial_i\) for \(|i-j|>1\), and \(\partial_i\partial_{i+1}\partial_i=\partial_{i+1}\partial_i\partial_{i+1}\). If \(g\) is \(s_i\)-invariant then \(\partial_i(gf)=g\partial_if\).}
\GSfield{Fate}{The divided-difference formalism appears in \citet[\S 3]{bernsteingelfandgelfand1973schubert}; the later primary Schubert-module paper recalls it in \citet[\S 1]{kraskiewiczpragacz2004schubert}. The divided-difference sequence is also recorded in 1985-661.4--8.}
\GSfield{Present status}{As of 11 November 2026, the operators are integral because the numerator is divisible by the displayed linear factor.}
\GSfield{Assessment}{literal statement: correct; publication match: exact; verification: read the result; development: none; confidence: high}
\end{GSitem}
\begin{GSitem}{1987-720.13}{E}{The Hecke deformation of the quadratic relation}
\GSfield{Kind; source}{remark; \GSpage{P001402}}
\GSfield{The notes}{Alongside the nilpotent relation, the notes mention the quadratic relation \((\sigma_i+1)(\sigma_i-q)=0\) as its q-analogue.}
\GSfield{Statement}{The type-A Hecke algebra is generated by \(T_i\) with braid and distant-commutation relations and \((T_i-q)(T_i+1)=0\). This is a deformation of the symmetric-group presentation at \(q=1\). The divided-difference operators themselves satisfy nilpotence and cannot satisfy this Hecke quadratic relation for generic \(q\).}
\GSfield{Reconstruction}{The distinction between the nil-Coxeter operator and a Hecke generator prevents interpreting both quadratic relations simultaneously.}
\GSfield{Fate}{The exact Hecke presentation is revisited in record 726 and occurs in \citet[\S 1]{kerov1987realizations}; the notebook is juxtaposing related algebras, not identifying their unchanged generators.}
\GSfield{Present status}{As of 11 November 2026, the analogy is valid after specifying different operators or generators.}
\GSfield{Assessment}{literal statement: correct with missing hypotheses; publication match: partial; verification: read the result; development: none; confidence: high}
\end{GSitem}
\begin{GSitem}{1987-720.14}{E}{Schubert polynomials from dominant codes}
\GSfield{Kind; source}{construction; \GSpage{P001402-P001403}}
\GSfield{The notes}{A weakly decreasing code \(I\) is assigned the monomial \(x^I\); divided differences change adjacent entries by \((i_p,i_{p+1})\mapsto(i_{p+1},i_p-1)\) when \(i_p>i_{p+1}\).}
\GSfield{Statement}{For a permutation \(w\), let \(c_i(w)=\#\{j>i:w(j)<w(i)\}\). There is a unique stable family \(\mathfrak S_w\) with \(\mathfrak S_{w_0}=x_1^{N-1}x_2^{N-2}\cdots x_{N-1}\) in \(S_N\), and \(\partial_i\mathfrak S_w=\mathfrak S_{ws_i}\) when \(\ell(ws_i)<\ell(w)\), zero otherwise. If the code is weakly decreasing, \(\mathfrak S_w=x^{c(w)}\).}
\GSfield{Reconstruction}{The permutation code fixes the notebook's indexing convention.}
\GSfield{Fate}{Lascoux--Sch\"utzenberger's 1982 announcement is the pre-seminar source \citep{lascouxschutzenberger1982schubert}; the definition and code convention are read in \citet[\S 1]{kraskiewiczpragacz2004schubert}. Earlier occurrence: 1985-661.7. Compare the generation from the longest permutation in 1985-661.7.}
\GSfield{Present status}{As of 11 November 2026, the construction is independent of a reduced word by the braid relations.}
\GSfield{Assessment}{literal statement: correct; publication match: exact; verification: read the result; development: none; confidence: high}
\end{GSitem}
\begin{GSitem}{1987-720.15}{E}{The worked chain from code 730 to code 025}
\GSfield{Kind; source}{construction; \GSpage{P001403}}
\GSfield{The notes}{The notes compute \(Y_{730}\to Y_{702}\to Y_{062}\to Y_{025}\) using \(\partial_2,\partial_1,\partial_2\) and complete-function determinants.}
\GSfield{Statement}{The divided-difference rules give \[\partial_2\partial_1\partial_2(x_1^7x_2^3)=Y_{025}=s_{(5,2)}(x_1,x_2,x_3)=h_5h_2-h_6h_1,\] where operator composition acts from right to left and each \(h_i\) uses the three-variable alphabet in the final equality.}
\GSfield{Reconstruction}{ \oursComputed{} The independent checks are recorded in the early-record verification log.}
\GSfield{Fate}{The general rules are those of \citet[\S 1]{kraskiewiczpragacz2004schubert} and Jacobi--Trudi \citep[Chapter I, (3.4)]{macdonald1995symmetric}. The concrete equality is checked by independent symbolic polynomial expansion in the accompanying script.}
\GSfield{Present status}{As of 11 November 2026, the displayed example is correct after maintaining the changing alphabets in the intermediate determinants.}
\GSfield{Assessment}{literal statement: correct; publication match: exact; verification: read the result; development: annotated; confidence: high}
\end{GSitem}
\begin{GSitem}{1987-720.16}{E}{Schur functions as special Schubert polynomials}
\GSfield{Kind; source}{theorem; \GSpage{P001404}}
\GSfield{The notes}{For an increasing code \(0\le i_1\le\cdots\le i_n\) followed by zeroes, the notes identify \(Y_I\) with the Schur function previously defined.}
\GSfield{Statement}{If a permutation is Grassmannian with its only possible descent at \(n\), then \(\mathfrak S_w=s_\lambda(x_1,\ldots,x_n)\), where \(\lambda=(c_n(w),\ldots,c_1(w))\). Equivalently its first \(n\) code entries are weakly increasing and the later entries zero.}
\GSfield{Fate}{This is a classical specialization of the Schubert construction, recalled in \citet[\S 1]{kraskiewiczpragacz2004schubert} and originating in the Lascoux--Sch\"utzenberger theory \citep{lascouxschutzenberger1982schubert}.}
\GSfield{Present status}{As of 11 November 2026, the identity explains why the previous worked example ends at an ordinary Schur function.}
\GSfield{Assessment}{literal statement: correct; publication match: exact; verification: read the result; development: none; confidence: high}
\end{GSitem}
\begin{GSitem}{1987-720.17}{E}{The finite Schubert basis and its cardinality}
\GSfield{Kind; source}{theorem; \GSpage{P001404}}
\GSfield{The notes}{For \(n+1\) variables, codes satisfy \(0\le i_1\le n,\ldots,0\le i_{n+1}\le0\); the notes then give the dimension as \(n!\), despite listing six elements for \(n=2\).}
\GSfield{Statement}{For \(R_N=\mathbb Z[x_1,\ldots,x_N]/(e_1,\ldots,e_N)\), the classes \(\mathfrak S_w\), \(w\in S_N\), form a free integral basis, hence \(\operatorname{rank}_{\mathbb Z}R_N=N!\). In the notebook's \(N=n+1\) convention the count is \((n+1)!\). Stable Schubert polynomials indexed by permutations outside \(S_N\), specialized to these variables, vanish in this quotient.}
\GSfield{Fate}{The flag-cohomology basis is classical \citep[\S\S 2--3]{bernsteingelfandgelfand1973schubert}. The count also follows independently by multiplying the \((n+1),n,\ldots,1\) choices for the code. The opened original note \citep[\S 4, p.~449]{lascouxschutzenberger1982schubert} explicitly recalls the integral flag quotient and Schubert basis, and the product-duality assertion following Theorem 4.3.}
\GSfield{Present status}{As of 11 November 2026, the basis assertion is correct; the printed factorial loses one in its argument.}
\GSfield{Assessment}{literal statement: false as written; publication match: exact; verification: read the result; development: annotated; confidence: high}
\end{GSitem}
\begin{GSitem}{1987-720.18}{E}{The symmetric-group action in the Schubert basis}
\GSfield{Kind; source}{question; \GSpage{P001404}; I. M. Gelfand}
\GSfield{The notes}{Gelfand asks how variable permutation by \(S_n\) is written in Schubert polynomials.}
\GSfield{Statement}{The action on the coinvariant algebra is determined by \[s_i[P]=[P]-(x_i-x_{i+1})[\partial_iP].\] Thus apply \(\partial_i\mathfrak S_w\), multiply by the indicated linear form, and reduce in the finite Schubert basis. Over \(\mathbb Q\), the ungraded coinvariant representation is the regular representation of \(S_N\).}
\GSfield{Reconstruction}{The quotient, the rational coefficient field for the regular representation, and the algorithmic meaning of an explicit formula are specified.}
\GSfield{Fate}{The Weyl-group action on flag cohomology is already a subject of \citet[\S 3]{bernsteingelfandgelfand1973schubert}. The displayed operator formula follows immediately from the definition, so the recorded request is computational exposition rather than an unresolved existence problem.}
\GSfield{Present status}{As of 11 November 2026, the action is explicit; the graded pieces are not each regular representations.}
\GSfield{Assessment}{literal statement: correct with missing hypotheses; publication match: exact; verification: indirect; development: none; confidence: high}
\end{GSitem}
\begin{GSitem}{1987-720.19}{E}{Poincaré pairing and adjoint divided differences}
\GSfield{Kind; source}{theorem; \GSpage{P001404-P001405}}
\GSfield{The notes}{The pairing is obtained from \(\partial_1\partial_2\partial_1(PQ)\) in three variables; divided differences are self-adjoint, and a code-complement rule gives the dual basis.}
\GSfield{Statement}{On the coinvariant algebra for \(S_N\), define \((P,Q)=\epsilon(\partial_{w_0}(PQ))\), where \(\epsilon\) sets all variables to zero. Then \((\partial_iP,Q)=(P,\partial_iQ)\). With the standard Schubert indexing, \((\mathfrak S_u,\mathfrak S_v)=\delta_{v,w_0u}\); equivalently \(c_i(v)=N-i-c_i(u)\).}
\GSfield{Reconstruction}{The augmentation and indexing convention clarify the notebook's phrase that the answer is a number in the quotient. \oursComputed{} The independent checks are recorded in the early-record verification log.}
\GSfield{Fate}{The pairing is integration on the flag manifold and the operators are its push-pull operators \citep[\S 3]{bernsteingelfandgelfand1973schubert}. A separate finite-rank polynomial check tests the code complement and the adjoint identity. The opened original note \citep[\S 4, p.~449]{lascouxschutzenberger1982schubert} explicitly recalls the integral flag quotient and Schubert basis, and the product-duality assertion following Theorem 4.3.}
\GSfield{Present status}{As of 11 November 2026, the pairing is perfect over the integers, and evaluation after the longest divided difference is necessary for arbitrary degrees.}
\GSfield{Assessment}{literal statement: correct with missing hypotheses; publication match: exact; verification: read the result; development: annotated; confidence: high}
\end{GSitem}
\begin{GSitem}{1987-720.20}{E}{Freeness over symmetric polynomials}
\GSfield{Kind; source}{theorem; \GSpage{P001405}}
\GSfield{The notes}{The polynomial ring in three variables is said to be free over its symmetric subring with Schubert-polynomial basis.}
\GSfield{Statement}{For \(N\ge1\), \(\mathbb Z[x_1,\ldots,x_N]\) is a free module over \(\mathbb Z[e_1,\ldots,e_N]\), with basis \(\{\mathfrak S_w:w\in S_N\}\). The basis has \(N!\) elements; it is not the entire stable family indexed by all finite permutations.}
\GSfield{Fate}{This is the polynomial form of the coinvariant-basis theorem \citep[\S 3]{bernsteingelfandgelfand1973schubert}; the modern Schubert convention is fixed in \citet[\S 1]{kraskiewiczpragacz2004schubert}.}
\GSfield{Present status}{As of 11 November 2026, the integral freeness statement holds with the finite indexing set made explicit.}
\GSfield{Assessment}{literal statement: correct with missing hypotheses; publication match: exact; verification: indirect; development: none; confidence: high}
\end{GSitem}
\begin{GSitem}{1987-720.21}{E}{Schubert polynomials as characters}
\GSfield{Kind; source}{problem; \GSpage{P001405}; I. M. Gelfand}
\GSfield{The notes}{Gelfand requests a generalization of character theory whose characters are Schubert functions; the line also mentions irreducible representations without naming a group or category.}
\GSfield{Statement}{For each finite permutation \(w\), one can construct a finite-dimensional cyclic polynomial module \(S_w(E_\bullet)\) for the triangular group preserving a complete flag, whose diagonal-torus character is \(\mathfrak S_w\). It is generally not irreducible. A nonsymmetric Schubert polynomial cannot be the character of a polynomial \(\operatorname{GL}_N\)-representation, whose torus character is symmetric.}
\GSfield{Reconstruction}{The representation category is supplied by the matched theorem; irreducibility is explicitly not imported into it.}
\GSfield{Fate}{\citet[Theorem 4.1]{kraskiewiczpragacz2004schubert} proves the exact Borel-module realization. Its introduction dates the preprint to 1986 and its announcement to 1987; the full article appeared in 2004. Thus this is class E at the seminar date, not a 2004 solution of a new 1987 problem. Related occurrence 1985-661.3 concerns Schubert polynomial representatives, not this character-realization question.}
\GSfield{Present status}{As of 11 November 2026, the character realization is established in the Borel category; the notebook does not identify a category supporting its extra irreducibility wording.}
\GSfield{Assessment}{literal statement: correct with missing hypotheses; publication match: partial; verification: read the result; development: none; confidence: high; problem \hyperref[prob:GP-1987-01]{GP-1987-01}}
\end{GSitem}
\begin{GSitem}{1987-720.22}{F}{A theory organizing determinant formulas}
\GSfield{Kind; source}{remark; \GSpage{P001404}}
\GSfield{The notes}{In discussing Muir's determinant books, Gelfand proposes building a theory behind the formulas, by analogy with integrable systems.}
\GSfield{Statement}{The remark specifies neither a class of determinant identities, an invariant, nor a proposed equivalence. It supports a research programme, but no uniquely quantified conjecture or theorem.}
\GSfield{Fate}{The adjacent Schur and Schubert identities have exact classical matches. Searches for the Muir and integrable-system connection do not turn this unspecific sentence into a definite outstanding problem.}
\GSfield{Present status}{As of 11 November 2026, the programme remains too broad for a mathematical fate classification beyond F.}
\GSfield{Assessment}{literal statement: not determinable; publication match: none found; verification: unresolved lead; development: none; confidence: low}
\end{GSitem}

\GStalk{721}{1987-11-09}{M. M. Kapranov}{P001405-P001406}{Double categories, two-categories, and higher fundamental groupoids}
Kapranov introduces squares with horizontal and vertical composition, the structure now conventionally called a double category, and its interchange axiom. Kontsevich contrasts this with two-morphisms between parallel arrows. A tensor-contraction example, Gelfand's boundary-value suggestion, and the homotopy squares of a filtered space give three distinct motivations.

\begin{GSitem}{1987-721.1}{E}{Double categories and the interchange law}
\GSfield{Kind; source}{definition; \GSpage{P001405-P001406}}
\GSfield{The notes}{The heading says 2-categories, while the data comprise objects, horizontal arrows, vertical arrows, and squares with two compositions.}
\GSfield{Statement}{A strict double category is a category object in categories: its squares have horizontal and vertical boundary arrows and associative unital compositions \(\circ_h,\circ_v\) satisfying \[(A\circ_vB)\circ_h(C\circ_vD)=(A\circ_hC)\circ_v(B\circ_hD)\] whenever the four squares have matching boundaries.}
\GSfield{Reconstruction}{The modern name separates this structure from the strict 2-category in the later intervention.}
\GSfield{Fate}{\citet[\S 1, equations (1.1)--(1.5)]{brownspencer1976double} provides these exact four kinds of data and the interchange law; the paper appeared in 1976.}
\GSfield{Present status}{As of 11 November 2026, this is established double-category terminology; the notebook uses a broader informal label.}
\GSfield{Assessment}{literal statement: correct; publication match: exact; verification: read the result; development: none; confidence: high}
\end{GSitem}
\begin{GSitem}{1987-721.2}{E}{Why naive corner-index contraction fails interchange}
\GSfield{Kind; source}{claim; \GSpage{P001406}}
\GSfield{The notes}{Four-index tensors are multiplied by contracting the two indices on a common edge; the note says this example does not satisfy interchange.}
\GSfield{Statement}{Write a square tensor as \(A(k,l;i,j)\). Define horizontal contraction by \((A\circ_hC)(k,r;i,s)=\sum_{j,l}A(k,l;i,j)C(l,r;j,s)\), and vertical contraction by \((A\circ_vB)(k,l;u,v)=\sum_{i,j}A(k,l;i,j)B(i,j;u,v)\). With indices \(0,1\), take the only nonzero entries of \(A,C,B,D\) to be, respectively, \(A(0,0;0,0)=C(0,0;0,0)=B(0,1;0,0)=D(1,0;0,0)=1\). The outer all-zero entry of \((A\circ_hC)\circ_v(B\circ_hD)\) is 1, while that of \((A\circ_vB)\circ_h(C\circ_vD)\) is 0.}
\GSfield{Reconstruction}{The full vertical contraction convention and sparse counterexample are ours; the notebook gives only a drawing of vertical composition. \oursComputed{} The independent checks are recorded in the early-record verification log.}
\GSfield{Fate}{The failure follows by two direct finite contractions, independently checked in the accompanying script. The relevant required axiom is \citet[equation (1.5)]{brownspencer1976double}.}
\GSfield{Present status}{As of 11 November 2026, the proposed naive contraction does fail; suitable categorical tensor constructions must keep different boundary data.}
\GSfield{Assessment}{literal statement: correct; publication match: partial; verification: read the result; development: annotated; confidence: high}
\end{GSitem}
\begin{GSitem}{1987-721.3}{E}{Strict two-categories and Cat}
\GSfield{Kind; source}{definition; \GSpage{P001406}; M. L. Kontsevich}
\GSfield{The notes}{Kontsevich proposes morphisms between morphisms with identical endpoints; Cat is the example. The note calls this a special case obtained by collapsing vertical arrows.}
\GSfield{Statement}{A strict 2-category has objects, 1-morphisms, and 2-morphisms between parallel 1-morphisms, with strictly associative unital horizontal and vertical compositions satisfying interchange. In \(\mathrm{Cat}\), these are categories, functors, and natural transformations. A strict 2-category embeds as a double category whose vertical arrows are identities.}
\GSfield{Fate}{This is pre-1987 higher-category formalism; \citet[\S 1 and references to B\'enabou, Gray, and Kelly--Street]{brownspencer1976double} distinguishes the double-category structure and its specializations.}
\GSfield{Present status}{As of 11 November 2026, the two descriptions are related, but their general definitions are not interchangeable.}
\GSfield{Assessment}{literal statement: correct; publication match: partial; verification: read the result; development: none; confidence: high}
\end{GSitem}
\begin{GSitem}{1987-721.4}{F}{Harmonic boundary data as a categorical suggestion}
\GSfield{Kind; source}{remark; \GSpage{P001406}; I. M. Gelfand}
\GSfield{The notes}{Gelfand draws a square with \(\Delta u=0\), records \(u\) and \(\partial u/\partial n\) on its boundary, and asks Keldysh, Chentsov, and Gelfand to consider the idea.}
\GSfield{Statement}{The data suggest composing boundary-value relations by gluing domains. No spaces of admissible Cauchy data, corner conditions, morphisms, or composition law are defined, so an exact double-category theorem cannot be extracted.}
\GSfield{Fate}{The boundary-value and categorical keywords were checked against the surrounding record and the later higher-category thread. No precise publication can be matched to this fragment without choosing new analytic definitions.}
\GSfield{Present status}{As of 11 November 2026, this is a mathematically motivated programme but an indeterminate statement, not an identified unsolved analytic problem.}
\GSfield{Assessment}{literal statement: not determinable; publication match: none found; verification: unresolved lead; development: none; confidence: low}
\end{GSitem}
\begin{GSitem}{1987-721.5}{E}{Homotopy squares of a filtered space}
\GSfield{Kind; source}{construction; \GSpage{P001406}}
\GSfield{The notes}{For \(X_0\subset X_1\subset X\), vertices lie in \(X_0\), edges are paths in \(X_1\), and squares map \((I^2,\partial I^2)\) to \((X,X_1)\), modulo homotopy.}
\GSfield{Statement}{For a filtered space \(X_0\subset X_1\subset X\), take paths in \(X_1\) with endpoints in \(X_0\), and squares in \(X\) whose edges lie in \(X_1\) and vertices in \(X_0\). Homotopies must preserve these filtrations and fix vertices; the resulting homotopy double-groupoid construction has horizontal and vertical gluing. At a basepoint, its relative structure is encoded by the crossed module \(\pi_2(X,X_1)\to\pi_1(X_1)\).}
\GSfield{Reconstruction}{The filtered and vertex-fixed homotopy conventions are essential additions to the shorthand.}
\GSfield{Fate}{\citet[\S 1, Example 1; introduction; Theorem A]{brownspencer1976double} explains squares, their double compositions, and the crossed-module relation. Theorem A circulated in a 1972 preliminary report, before its 1976 publication.}
\GSfield{Present status}{As of 11 November 2026, the construction is classical once the homotopies and connections are specified; arbitrary unrelative homotopy would not preserve the necessary boundary data.}
\GSfield{Assessment}{literal statement: correct with missing hypotheses; publication match: partial; verification: read the result; development: none; confidence: high}
\end{GSitem}

\GStalk{722}{1987-11-23}{A. A. Beilinson}{P001407}{Derived categories: the cone lemma and roofs, revisited}
This one-page continuation recalls the categories of complexes, homotopy classes, and derived morphisms. It states the cone proof of the Ore lemma and the equivalence relation on roofs. The outgoing quasi-isomorphism is written on the wrong leg of the opening diagram, an error visible in the original scan.

\begin{GSitem}{1987-722.1}{E}{The cone lemma with its outgoing quasi-isomorphism}
\GSfield{Kind; source}{theorem; \GSpage{P001407}}
\GSfield{The notes}{The initial roof has \(Z\to X\) a quasi-isomorphism; the outgoing diagram labels \(X\to Z_1\), rather than \(Y\to Z_1\), as a quasi-isomorphism.}
\GSfield{Statement}{For \(Z\xrightarrow{s}X\), \(Z\xrightarrow fY\), with \(s\) a quasi-isomorphism, there is a homotopy-commutative square with \(X\to P\) and a quasi-isomorphism \(Y\to P\). One may take \(P=\operatorname{Cone}((s,-f):Z\to X\oplus Y)\).}
\GSfield{Fate}{This repeats occurrence 1987-717.5 and the established multiplicative-system construction \citep[Lemma 13.11.2 and \S 13.5]{stacks2026foundations}. The wrong-leg version cannot be true for \(s=1_Z\) and an arbitrary map \(f\).}
\GSfield{Present status}{As of 11 November 2026, the correct lemma holds in the homotopy category; the original diagram requires an emendation.}
\GSfield{Assessment}{literal statement: false as written; publication match: exact; verification: read the result; development: annotated; confidence: high}
\end{GSitem}
\begin{GSitem}{1987-722.2}{E}{Homotopies, localization, and refinement of roofs}
\GSfield{Kind; source}{construction; \GSpage{P001407}}
\GSfield{The notes}{The page repeats \(f^i=dg^i+g^{i+1}d\), inverts quasi-isomorphisms in \(K(\mathcal A)\), and generates equivalence of roofs by precomposing with a quasi-isomorphism.}
\GSfield{Statement}{With complexes as objects, \(C(\mathcal A)\to K(\mathcal A)\to D(\mathcal A)\) first identifies homotopic chain maps and then inverts every quasi-isomorphism. A roof \((s,f)\) is equivalent to \((s\chi,f\chi)\) for a quasi-isomorphism \(\chi:Z_1\to Z\); the equivalence relation generated by these refinements is the derived-morphism relation.}
\GSfield{Fate}{\citet[Definition 13.11.3 and localization construction]{stacks2026foundations}; compare occurrences 1987-717.3--6 and earlier 1971-266.3.}
\GSfield{Present status}{As of 11 November 2026, this is a repeated foundational construction rather than a new problem or later development.}
\GSfield{Assessment}{literal statement: correct; publication match: exact; verification: read the result; development: none; confidence: high}
\end{GSitem}

\GStalk{723}{1987-11-23}{I. S. Zakharevich}{P001408}{Bernoulli coding, the baker map, and hyperbolicity}
The short talk illustrates sensitivity to finite observational precision using the two-sided binary shift and the baker transformation. A closing comparison invokes geodesic flow on a negatively curved manifold. Measure-theoretic coding, topological conjugacy, and deterministic loss of exact information must be distinguished.

\begin{GSitem}{1987-723.1}{E}{The two-sided Bernoulli shift and baker map}
\GSfield{Kind; source}{construction; \GSpage{P001408}}
\GSfield{The notes}{The notes depict a bi-infinite sequence of zeroes and ones and encode the two halves as binary coordinates of a square.}
\GSfield{Statement}{Let \(\Sigma=\{0,1\}^{\mathbb Z}\) with product measure and \((\sigma a)_j=a_{j+1}\). Set \(x=.a_0a_1\cdots\), \(y=.a_{-1}a_{-2}\cdots\). Outside the null set of ambiguous dyadic expansions, this gives a measure-preserving conjugacy with \[B(x,y)=(2x-\lfloor2x\rfloor,(y+\lfloor2x\rfloor)/2).\]}
\GSfield{Reconstruction}{The formula, probability measure, and dyadic exception are supplied explicitly.}
\GSfield{Fate}{The coding is checked directly by moving the binary point. The shift and hyperbolic examples belong to the pre-1987 theory surveyed in \citet[hyperbolic examples]{smale1967dynamical}; only that survey's metadata was reopened in this run.}
\GSfield{Present status}{As of 11 November 2026, the measure-theoretic conjugacy is exact almost everywhere; the square and the Cantor product are not globally homeomorphic.}
\GSfield{Assessment}{literal statement: correct with missing hypotheses; publication match: partial; verification: indirect; development: none; confidence: high}
\end{GSitem}
\begin{GSitem}{1987-723.2}{E}{Finite observation and loss of predictive information}
\GSfield{Kind; source}{claim; \GSpage{P001408}}
\GSfield{The notes}{Iterating the shift is described as gradually losing information about a point when only finitely many digits are known.}
\GSfield{Statement}{For independent fair binary coordinates, if only \(a_j\) with \(|j|\le m\) are observed, then any unobserved coordinate \(a_k\), \(|k|>m\), still has conditional probabilities \(1/2,1/2\). Under \(\sigma^n\), the observed window shifts, and future digits outside the initial finite window are not determined by it. The complete two-sided shift remains invertible.}
\GSfield{Reconstruction}{Conditional probability is our precise interpretation of the informal information claim.}
\GSfield{Fate}{This precise finite-window statement is a direct consequence of product independence, not a new theorem about destruction of exact state information. It clarifies the standard symbolic-dynamics example associated with \citet{smale1967dynamical}.}
\GSfield{Present status}{As of 11 November 2026, the observation-level assertion is correct; exact invertible dynamics does not erase its full initial data.}
\GSfield{Assessment}{literal statement: correct with missing hypotheses; publication match: partial; verification: indirect; development: none; confidence: high}
\end{GSitem}
\begin{GSitem}{1987-723.3}{E}{Negative curvature as a geometric hyperbolicity example}
\GSfield{Kind; source}{claim; \GSpage{P001408}}
\GSfield{The notes}{The closing line compares the binary example with geodesics on a negatively curved manifold.}
\GSfield{Statement}{If \(M\) is a closed smooth Riemannian manifold with strictly negative sectional curvature, its geodesic flow on the unit tangent bundle is Anosov: the tangent bundle splits into flow, stable, and unstable subbundles with uniform exponential contraction in forward or backward time.}
\GSfield{Reconstruction}{The closed-manifold hypothesis and the exact hyperbolicity assertion supply a limited rigorous interpretation.}
\GSfield{Fate}{\citet{anosov1967geodesic} is the original 1967 monograph; its Math-Net publication record was opened. The notebook does not state a literal conjugacy of this flow to a binary shift.}
\GSfield{Present status}{As of 11 November 2026, uniform hyperbolicity holds under compactness and negative curvature; arbitrary noncompact negatively curved manifolds need additional uniformity conditions.}
\GSfield{Assessment}{literal statement: correct with missing hypotheses; publication match: partial; verification: publisher metadata checked; development: none; confidence: high}
\end{GSitem}

\GStalk{724}{1987-11-23}{G. A. Margulis}{P001409-P001410}{Rational iteration, Sullivan's theorem, and fractal questions}
Margulis moves from affine maps and squaring to the Fatou and Julia sets of a rational map, then sketches Sullivan's quasiconformal proof excluding wandering domains. Gelfand's interventions include circle linearization, elementary measure and dimension exercises, and a programme for discrete-group representations on fractal sets. The final disk-packing sketch omits the iterative rule needed for either residual-set question.

\begin{GSitem}{1987-724.1}{E}{Affine iteration and invariant circles}
\GSfield{Kind; source}{claim; \GSpage{P001409}}
\GSfield{The notes}{For \(z\mapsto az+b\), the notes say that \(|a|=1\) gives invariant circles carrying rotations.}
\GSfield{Statement}{If \(a\ne1\), translating by \(z_*=b/(1-a)\) conjugates the map to \(w\mapsto aw\). For \(|a|=1\), \(a\ne1\), the circles \(|z-z_*|=r\) are invariant. If \(a=1,b\ne0\), the map is a translation and has no fixed point or invariant Euclidean circle; if \(a=1,b=0\), it is the identity.}
\GSfield{Fate}{The correction follows by direct substitution and supplies the missing exceptional case in the recorded elementary example. The ensuing nonlinear setting is the rational dynamics of \citet[Introduction]{sullivan1985quasiconformal}.}
\GSfield{Present status}{As of 11 November 2026, the statement is valid after excluding the nonzero translation case.}
\GSfield{Assessment}{literal statement: false as written; publication match: partial; verification: read the result; development: annotated; confidence: high}
\end{GSitem}
\begin{GSitem}{1987-724.2}{E}{Smooth conjugacy of circle diffeomorphisms to rotations}
\GSfield{Kind; source}{question; \GSpage{P001409}; I. M. Gelfand}
\GSfield{The notes}{Gelfand says the circle-conjugacy question is nontrivial and names Denjoy, Arnold, and Herman.}
\GSfield{Statement}{An orientation-preserving \(C^2\) circle diffeomorphism with irrational rotation number is topologically conjugate to the corresponding rotation. If it is \(C^\infty\) and its rotation number satisfies a Diophantine condition \(|\alpha-p/q|\ge Cq^{-2-\nu}\) for some \(C,\nu>0\) and all rational \(p/q\), the conjugacy can be chosen \(C^\infty\). Smoothness alone does not ensure a smooth conjugacy for every irrational rotation number.}
\GSfield{Reconstruction}{The modern formulation separates topological and smooth conjugacy, which the notes do not spell out.}
\GSfield{Fate}{\citet[Introduction and main theorem]{yoccoz1984conjugaison} proves the smooth Diophantine statement, generalizing Herman, in 1984. It predates this intervention.}
\GSfield{Present status}{As of 11 November 2026, the regularity and arithmetic qualifications are essential; the remark is about an established theory, not an explicitly new conjecture.}
\GSfield{Assessment}{literal statement: correct; publication match: partial; verification: read the result; development: none; confidence: high}
\end{GSitem}
\begin{GSitem}{1987-724.3}{F}{Representations on sets of intermediate Hausdorff dimension}
\GSfield{Kind; source}{conjecture; \GSpage{P001409}; I. M. Gelfand}
\GSfield{The notes}{Gelfand asks for a representation theory of infinite discrete groups on functions on sets of intermediate Hausdorff dimension, suggesting that some classical theory might reappear.}
\GSfield{Statement}{The fragment specifies no class of groups, fractal spaces, measures, function spaces, actions, or meaning of recovering the classical theory. Consequently there is no definite universal existence or classification statement to test.}
\GSfield{Fate}{Targeted searches for Gelfand, representations of discrete groups, and Hausdorff dimension found related boundary and fractal representation constructions, but no publication with the quantifiers of an identifiable question here. Those thematic developments are not treated as a solution.}
\GSfield{Present status}{As of 11 November 2026, the programme remains indeterminate from this record; F records the missing formulation and does not assert that no such representations exist.}
\GSfield{Assessment}{literal statement: not determinable; publication match: none found; verification: unresolved lead; development: none; confidence: low}
\end{GSitem}
\begin{GSitem}{1987-724.4}{E}{A point outside short intervals around the rationals}
\GSfield{Kind; source}{problem; \GSpage{P001409}; I. M. Gelfand}
\GSfield{The notes}{The school-level question covers the nth rational by an interval of length \(\varepsilon/2^n\) and asks for a number left uncovered.}
\GSfield{Statement}{For any enumeration \((q_n)_{n\ge1}\) of \(\mathbb Q\) and open intervals \(I_n\ni q_n\) with \(|I_n|=\varepsilon2^{-n}\), \(m(\bigcup_nI_n)\le\varepsilon\). Hence for \(0<\varepsilon<1\), \([0,1]\setminus\bigcup_nI_n\) has measure at least \(1-\varepsilon\), is nonempty, and consists of irrationals. A particular uncovered number depends on the chosen enumeration and intervals.}
\GSfield{Reconstruction}{The unit interval and epsilon bound turn the informal task into a precise existence assertion.}
\GSfield{Fate}{This is countable subadditivity of Lebesgue measure, explicitly reconstructed here. The basic measure setting is described in \citet[\S 2.5]{hutchinson1981fractals}; no fixed numerical example can be inferred without the missing interval data.}
\GSfield{Present status}{As of 11 November 2026, the existence claim is elementary; the source does not specify enough data to name one universal uncovered real number.}
\GSfield{Assessment}{literal statement: correct with missing hypotheses; publication match: partial; verification: indirect; development: none; confidence: high}
\end{GSitem}
\begin{GSitem}{1987-724.5}{E}{Hausdorff dimension of the Cantor set}
\GSfield{Kind; source}{problem; \GSpage{P001409}; I. M. Gelfand}
\GSfield{The notes}{Gelfand asks for the Hausdorff dimension of the perfect Cantor set, without spelling out its construction.}
\GSfield{Statement}{For the middle-thirds Cantor set \(C\), defined by \(C=\tfrac13C\cup(\tfrac23+\tfrac13C)\), \(\dim_HC=\log2/\log3\), and its Hausdorff measure in that dimension is finite and positive. An arbitrary perfect nowhere-dense subset of the line need not have this dimension.}
\GSfield{Reconstruction}{Middle thirds is the conventional interpretation, explicitly marked because the page only says Cantor perfect set.}
\GSfield{Fate}{\citet[\S 5.3(1), with the examples in \S 3.3]{hutchinson1981fractals} gives the self-similar dimension theorem. The paper appeared as a 1979 ANU report before its 1981 publication; this particular Cantor dimension is older still.}
\GSfield{Present status}{As of 11 November 2026, the standard middle-thirds exercise is settled; the construction must be specified to make the answer unique.}
\GSfield{Assessment}{literal statement: correct with missing hypotheses; publication match: exact; verification: read the result; development: none; confidence: high}
\end{GSitem}
\begin{GSitem}{1987-724.6}{E}{Squaring and one-sided binary dynamics}
\GSfield{Kind; source}{construction; \GSpage{P001409}}
\GSfield{The notes}{For \(z\mapsto z^2\), the unit circle is invariant, the inside tends to zero, the outside to infinity, and the circle map resembles the one-sided baker map.}
\GSfield{Statement}{If \(|z|<1\), \(z^{2^n}\to0\); if \(|z|>1\), it tends to infinity on the Riemann sphere. On \(z=e^{2\pi i\theta}\), iteration is \(\theta\mapsto2\theta\pmod1\). Binary coding gives the one-sided shift apart from dyadic ambiguity; this map is two-to-one, unlike the invertible two-sided baker transformation.}
\GSfield{Fate}{The formulas follow directly from exponentiation. This same base example introduces \citet[p.~401]{sullivan1985quasiconformal}.}
\GSfield{Present status}{As of 11 November 2026, the dynamics is completely explicit and the one-sided qualifier is mathematically significant.}
\GSfield{Assessment}{literal statement: correct with missing hypotheses; publication match: exact; verification: read the result; development: none; confidence: high}
\end{GSitem}
\begin{GSitem}{1987-724.7}{E}{The Fatou set and the Julia set}
\GSfield{Kind; source}{definition; \GSpage{P001410}}
\GSfield{The notes}{The Fatou set is the open set of local normality of the iterates, and its closed complement is the Julia set; the family is called compact.}
\GSfield{Statement}{For a rational map \(R:\widehat{\mathbb C}\to\widehat{\mathbb C}\) of degree at least two, \(F(R)\) consists of points with a neighborhood on which \(\{R^n:n\ge0\}\) is normal for locally uniform convergence in the spherical metric. \(J(R)=\widehat{\mathbb C}\setminus F(R)\) is closed. Normal means relative sequential compactness in this topology, not that the set of iterates is itself closed and compact.}
\GSfield{Fate}{\citet[p.~403]{sullivan1985quasiconformal} states this stable-set definition and its relation to Fatou's normal-family formulation.}
\GSfield{Present status}{As of 11 November 2026, these are the standard definitions with the spherical topology specified.}
\GSfield{Assessment}{literal statement: correct with missing hypotheses; publication match: exact; verification: read the result; development: none; confidence: high}
\end{GSitem}
\begin{GSitem}{1987-724.8}{E}{An empty Fatou set requires the rational-map setting}
\GSfield{Kind; source}{claim; \GSpage{P001410}}
\GSfield{The notes}{The notes say the normality set may be empty, after beginning with polynomials and allowing rational functions more generally.}
\GSfield{Statement}{A rational map of degree at least two can have \(F(R)=\varnothing\), for example a Latt\`es map induced by an expanding endomorphism of an elliptic curve. A polynomial of degree at least two always has a nonempty basin of infinity and therefore a nonempty Fatou set.}
\GSfield{Reconstruction}{The rational-versus-polynomial qualification is supplied explicitly; the example identifies the standard scope.}
\GSfield{Fate}{The distinction is between the two ambient classes named in the record. The rational setting of \citet[Introduction]{sullivan1985quasiconformal} permits empty stable set; the polynomial exception follows from growth outside a sufficiently large disk.}
\GSfield{Present status}{As of 11 November 2026, the assertion is valid for rational maps, not for all maps in the narrower polynomial subcase.}
\GSfield{Assessment}{literal statement: correct with missing hypotheses; publication match: partial; verification: indirect; development: none; confidence: high}
\end{GSitem}
\begin{GSitem}{1987-724.9}{E}{Sullivan: no wandering Fatou components}
\GSfield{Kind; source}{theorem; \GSpage{P001410}}
\GSfield{The notes}{Sullivan is credited with showing that every component orbit stabilizes and that there are finitely many fixed components.}
\GSfield{Statement}{For every rational map \(R\) of degree at least two and every Fatou component \(U\), there exist integers \(m\ge0,p\ge1\) such that \(R^{m+p}(U)=R^m(U)\). There are finitely many periodic component cycles. Eventual periodicity does not mean that every component eventually becomes fixed under \(R\).}
\GSfield{Fate}{\citet[Theorem 1, p.~404]{sullivan1985quasiconformal} states eventual periodicity exactly; the following discussion treats the finite periodic types. The article appeared in September 1985, before the seminar.}
\GSfield{Present status}{As of 11 November 2026, the no-wandering-domain theorem is established; replacing periodic cycles by eventual fixed points overstates it.}
\GSfield{Assessment}{literal statement: false as written; publication match: exact; verification: read the result; development: annotated; confidence: high}
\end{GSitem}
\begin{GSitem}{1987-724.10}{E}{The quasiconformal deformation argument}
\GSfield{Kind; source}{method; \GSpage{P001410}}
\GSfield{The notes}{The proof sketch supposes an infinite orbit of components and changes the conformal structures while keeping the map conformal.}
\GSfield{Statement}{The strategy is to assign compatible Beltrami data along a hypothetical wandering orbit, integrate it by the measurable Riemann mapping theorem, and obtain too many independent quasiconformal deformations for the finite-dimensional space of rational maps of the fixed degree. Producing the independent deformations and controlling conjugacy is part of the proof, not a consequence of parameter counting alone.}
\GSfield{Fate}{\citet[\S\S 7--10]{sullivan1985quasiconformal} implements the method. The introduction dates earlier lectures to 1978, 1980--81, and 1981; the final article is 1985.}
\GSfield{Present status}{As of 11 November 2026, the notebook accurately signals the proof mechanism but does not record a complete proof.}
\GSfield{Assessment}{literal statement: correct with missing hypotheses; publication match: partial; verification: read the result; development: none; confidence: high}
\end{GSitem}
\begin{GSitem}{1987-724.11}{F}{The area of a disk-packing residual set}
\GSfield{Kind; source}{problem; \GSpage{P001410}}
\GSfield{The notes}{A square filled with successively smaller tangent disks is drawn, with the instruction to inscribe them and prove that the remainder has measure zero.}
\GSfield{Statement}{Given a specified countable family of disjoint open disks \(D_j\subset Q\), the proposed conclusion would be \(m_2(Q\setminus\bigcup_jD_j)=0\). The notebook does not specify how the disks are chosen, which gaps must be filled, or a maximality or uniformity condition. Disjointness alone allows a positive-area remainder.}
\GSfield{Fate}{The scan shows only a tangency sketch and the word \ru{вписываем}. Searches for residual disk packings distinguish Apollonian and arbitrary packings; their theorems have different hypotheses. None can be identified as this exact unstated algorithm.}
\GSfield{Present status}{As of 11 November 2026, the intended packing process is indeterminate; no general zero-area theorem is asserted on the available data.}
\GSfield{Assessment}{literal statement: not determinable; publication match: none found; verification: unresolved lead; development: none; confidence: low}
\end{GSitem}
\begin{GSitem}{1987-724.12}{F}{The dimension of the same packing remainder}
\GSfield{Kind; source}{problem; \GSpage{P001410}}
\GSfield{The notes}{The second question asks for Hausdorff dimension strictly below two.}
\GSfield{Statement}{For \(K=Q\setminus\bigcup_jD_j\), the requested claim is \(\dim_HK<2\). Even the additional condition \(m_2(K)=0\) would not imply it: zero measure and strict dimension deficit are different properties. The missing disk-selection rule prevents a uniquely specified conjecture.}
\GSfield{Fate}{No exact packing algorithm or quantitative scaling assumption was found in the supplied page or matched by targeted literature searches. Self-similar dimension results such as \citet[\S 5.3]{hutchinson1981fractals} require contraction data absent here.}
\GSfield{Present status}{As of 11 November 2026, the strict dimension bound cannot be classified as a definite open or solved problem without the omitted construction.}
\GSfield{Assessment}{literal statement: not determinable; publication match: thematic only; verification: unresolved lead; development: none; confidence: low}
\end{GSitem}

\GStalk{725}{1987-11-23}{I. M. Gelfand}{P001411-P001412}{Matroids, bimatroids, and determinant relations}
Gelfand revisits representability through the Fano configuration, basis exchange, and the nonzero minors of a matrix. Several questions are already classical, while the comparison with an uncertain named author and the request to discretize an entire syzygy chain lack a determinate target. The literal determinant identity contains several errors; a corrected exchange identity explains the intended argument.

\begin{GSitem}{1987-725.1}{E}{Vector matroids and symmetric basis exchange}
\GSfield{Kind; source}{definition; \GSpage{P001411}}
\GSfield{The notes}{The independent subfamilies of a finite vector configuration define its matroid. A displayed axiom permits simultaneous exchange in two bases.}
\GSfield{Statement}{Let \(E\) be finite. A nonempty family \(\mathcal B\) of equal-cardinality subsets is a matroid basis family if, for \(B_1,B_2\in\mathcal B\) and \(y\in B_2\setminus B_1\), there exists \(x\in B_1\setminus B_2\) such that both \(B_1-x+y\) and \(B_2-y+x\) belong to \(\mathcal B\). This symmetric exchange formulation is equivalent to the usual basis axiom. For a vector configuration, the bases are its maximal independent subfamilies.}
\GSfield{Reconstruction}{Nonemptiness, equal cardinality, and exclusion of the exchanged element from the other basis are explicit.}
\GSfield{Fate}{This is classical matroid theory; the representation and independence definitions are stated in \citet[Theorem 1.6 and \S 1]{oxleybriefly}, while the basis-based bimatroid reformulation is recalled in \citet[Proposition 2.2]{rohrleulirsch2026bimatroids}. It repeats the subject of occurrence 1985-658.12.}
\GSfield{Present status}{As of 11 November 2026, the finite nonempty-family and equal-cardinality conventions make the displayed axiom precise.}
\GSfield{Assessment}{literal statement: correct with missing hypotheses; publication match: exact; verification: read the result; development: none; confidence: high}
\end{GSitem}
\begin{GSitem}{1987-725.2}{E}{The Fano configuration and characteristic two}
\GSfield{Kind; source}{problem; \GSpage{P001411}}
\GSfield{The notes}{The seven-point diagram is to be realized over \(\mathbb F_2\), but not over \(\mathbb R\) or \(\mathbb C\).}
\GSfield{Statement}{The rank-three Fano matroid \(F_7\), whose lines are the dependent triples in \(\mathbb F_2^3\setminus\{0\}\), is representable over a field \(k\) if and only if \(\operatorname{char}k=2\). In particular its seven indicated incidence relations cannot all arise from a rank-three real or complex vector configuration with that matroid.}
\GSfield{Fate}{\citet[Theorem 4.1]{oxleybriefly} states the exact field criterion. The seven nonzero binary vectors give the recorded realization. This is the same exercise as 1985-658.14.}
\GSfield{Present status}{As of 11 November 2026, the problem is settled and is an expository obstruction to representability over the specified fields.}
\GSfield{Assessment}{literal statement: correct; publication match: exact; verification: read the result; development: none; confidence: high}
\end{GSitem}
\begin{GSitem}{1987-725.3}{E}{Which finite matroids are real or complex representable?}
\GSfield{Kind; source}{problem; \GSpage{P001411}}
\GSfield{The notes}{The notebook calls the general question of representability over \(\mathbb R\) or \(\mathbb C\) unsolved, without specifying a type of classification.}
\GSfield{Statement}{For a finite rank-\(r\) matroid on \([n]\) with basis list \(\mathcal B\), representability over \(k=\mathbb R\) or \(\mathbb C\) is equivalent to the existence of \(A\in k^{r\times n}\) satisfying \(\det A_B\ne0\) for \(B\in\mathcal B\), and \(\det A_I=0\) for every other \(r\)-subset \(I\). These are finitely many polynomial conditions, and their solvability is decidable over both fields.}
\GSfield{Reconstruction}{The finite-input decision interpretation is made explicit; it is not claimed to supply an efficient or illuminating combinatorial classification.}
\GSfield{Fate}{\citet[\S 2, pp.~15--42, and \S 3, pp.~43--46]{tarski1951decision} explicitly states quantifier elimination over the reals and its complex extension. The nonvanishing conditions can equivalently be encoded by an auxiliary equation \(u\prod_{B\in\mathcal B}\det A_B=1\). Thus the finite decision problem was settled before 1987. Occurrences 1985-658.13 and 1985-658.16 already connect representability with polynomial systems.}
\GSfield{Present status}{As of 11 November 2026, a finite decision procedure exists; the source specifies no stronger structural classification or complexity target that could responsibly be declared open.}
\GSfield{Assessment}{literal statement: correct with missing hypotheses; publication match: partial; verification: read the result; development: annotated; confidence: high; problem \hyperref[prob:GP-1987-02]{GP-1987-02}}
\end{GSitem}
\begin{GSitem}{1987-725.4}{E}{Matroids representable over every field}
\GSfield{Kind; source}{theorem; \GSpage{P001411}}
\GSfield{The notes}{The notes say that representability over every field has been characterized.}
\GSfield{Statement}{For a finite matroid \(M\), the following are equivalent: \(M\) is representable over every field; \(M\) has a totally unimodular real matrix representation; and \(M\) has no minor isomorphic to \(U_{2,4}\), \(F_7\), or \(F_7^*\). Such matroids are called regular.}
\GSfield{Fate}{\citet[Lemma 4.4 and Theorem 4.5]{oxleybriefly} gives the exact equivalences and attributes them to Tutte 1958. The original Tutte paper was located but could not be opened in this run; the checked modern author exposition is therefore the immediate evidence.}
\GSfield{Present status}{As of 11 November 2026, the statement is an established characterization, with no new 1987 question recorded.}
\GSfield{Assessment}{literal statement: correct; publication match: exact; verification: indirect; development: none; confidence: high}
\end{GSitem}
\begin{GSitem}{1987-725.5}{E}{The bimatroid of nonzero minors}
\GSfield{Kind; source}{definition; \GSpage{P001411-P001412}}
\GSfield{The notes}{For a matrix with row set \(E\) and column set \(F\), the notes retain all pairs of equally sized subsets indexing a nonzero minor.}
\GSfield{Statement}{For a matrix \(A\in k^{E\times F}\), define \(\mathcal R(A)=\{(I,J):I\subseteq E,\ J\subseteq F,\ |I|=|J|,\ \det A_{I,J}\ne0\}\), including \((\varnothing,\varnothing)\) with determinant one. This records joint row-column dependence at every size, rather than just the row and column matroids.}
\GSfield{Fate}{\citet[Example 2.3]{rohrleulirsch2026bimatroids} gives this exact definition and recalls its origin in \citet{kung1978bimatroids}, published in 1978. The latter original remained inaccessible, so its attribution is indirect; the former result was read.}
\GSfield{Present status}{As of 11 November 2026, the construction was established before this seminar and has an exact distinguished-basis matroid encoding.}
\GSfield{Assessment}{literal statement: correct; publication match: exact; verification: read the result; development: none; confidence: high}
\end{GSitem}
\begin{GSitem}{1987-725.6}{E}{An axiom system for bimatroids}
\GSfield{Kind; source}{problem; \GSpage{P001411-P001412}}
\GSfield{The notes}{The explicit question asks for axioms on the collection of nonzero square minors.}
\GSfield{Statement}{For finite disjoint \(E,F\), a collection \(\mathcal R\) of pairs \((I,J)\) with \(|I|=|J|\) is an abstract bimatroid exactly when \(\{(E\setminus I)\sqcup J:(I,J)\in\mathcal R\}\) is a matroid basis family on \(E\sqcup F\) containing \(E\). Thus nonemptiness, the pair \((\varnothing,\varnothing)\), and the ordinary basis-exchange axiom after this substitution are a complete axiom system. In the matrix case the extended matroid is represented by \([I_E\mid A]\).}
\GSfield{Fate}{\citet[Proposition 2.2 and Example 2.3]{rohrleulirsch2026bimatroids} proves this equivalence in the exposition of Kung's 1978 theory. This is a later checked source for an older answer, not a 2026 solution date.}
\GSfield{Present status}{As of 11 November 2026, the recorded axiomatization problem is solved by an explicit exchange system; abstract bimatroids need not be realizable matrices.}
\GSfield{Assessment}{literal statement: correct; publication match: exact; verification: read the result; development: annotated; confidence: high; problem \hyperref[prob:GP-1987-03]{GP-1987-03}}
\end{GSitem}
\begin{GSitem}{1987-725.7}{F}{The comparison with the uncertain named bimatroid theory}
\GSfield{Kind; source}{question; \GSpage{P001412}}
\GSfield{The notes}{The page asks how a bimatroid notion in Gelfand--Serganova, UMN, relates to a name transcribed tentatively as \ru{Кунтом}.}
\GSfield{Statement}{An exact comparison would require two specified definitions and a map between their objects. The named passage in Gelfand--Serganova and the second surname are not identified unambiguously by the record. Kung is a plausible bibliographic lead, but this cannot determine the missing first definition.}
\GSfield{Fate}{The full English text of \citet{gelfandserganova1987combinatorial} was opened and searched; it develops ordinary and Coxeter matroids, but no occurrence of bimatroid identified the intended definition. \citet{kung1978bimatroids} is a relevant 1978 lead, with its modern formulation checked in \citet{rohrleulirsch2026bimatroids}. Neither fact proves the requested equivalence.}
\GSfield{Present status}{As of 11 November 2026, the comparison is indeterminate from the supplied citation and name; this is not a claim that the two theories differ or that an equivalence is unknown.}
\GSfield{Assessment}{literal statement: not determinable; publication match: thematic only; verification: unresolved lead; development: none; confidence: low}
\end{GSitem}
\begin{GSitem}{1987-725.8}{E}{A determinant exchange identity and basis exchange}
\GSfield{Kind; source}{problem; \GSpage{P001412}}
\GSfield{The notes}{Equation (6) has a sign \((-1)^{d-i}\), repeats \(y_1\) in the first bracket, and inserts \(x_i\) into the second; the next sentence uses a nonzero summand to obtain exchange.}
\GSfield{Statement}{For \(x_1,\ldots,x_d,y_1,\ldots,y_d\in k^d\), with brackets denoting ordered determinants, \[[x_1,\ldots,x_d][y_1,\ldots,y_d]=\sum_{i=1}^d[y_i,x_2,\ldots,x_d]\,[y_1,\ldots,y_{i-1},x_1,y_{i+1},\ldots,y_d].\] If both left factors are nonzero, some summand is nonzero, proving simultaneous exchange of \(x_1\) with a suitable \(y_i\).}
\GSfield{Reconstruction}{The insertion positions are fixed exactly, so no additional alternating sign is needed. \oursComputed{} The independent checks are recorded in the early-record verification log.}
\GSfield{Fate}{The corrected identity follows by expanding \(x_1\) in the \(y\)-basis and applying multilinearity; polynomial identity then removes the assumption that the \(y\)-matrix is invertible. It is part of the classical determinant relations recalled in \citet[\S 3, opening paragraph]{speyersturmfels2003tropical}. Direct symbolic checks give a counterexample to the literal notebook expression.}
\GSfield{Present status}{As of 11 November 2026, the intended exchange argument is valid after correcting the indices and the redundant sign.}
\GSfield{Assessment}{literal statement: false as written; publication match: partial; verification: indirect; development: annotated; confidence: high}
\end{GSitem}
\begin{GSitem}{1987-725.9}{E}{All relations among maximal minors}
\GSfield{Kind; source}{claim; \GSpage{P001412}}
\GSfield{The notes}{The notes claim that equation (6) supplies all relations among the determinant generators.}
\GSfield{Statement}{Let \(S=k[p_I: I\subseteq[n],\ |I|=r]\), and let \(S\to k[a_{ij}]\) send \(p_I\) to the corresponding maximal minor of a generic \(r\times n\) matrix. Its kernel is the Pl\"ucker ideal, generated by the full family of quadratic Pl\"ucker relations. One erroneous identity or one chosen instance does not by itself generate this ideal.}
\GSfield{Reconstruction}{The generic matrix, coefficient ring, homomorphism, and full relation family are specified.}
\GSfield{Fate}{\citet[\S 3, opening paragraph]{speyersturmfels2003tropical} explicitly states the integral version and quadratic generation as a classical fact. The notebook's assertion is correct only when (6) denotes the correctly indexed family of determinant exchanges, with the generic maximal-minor interpretation fixed.}
\GSfield{Present status}{As of 11 November 2026, the defining relations of the Grassmannian are established; the source formula requires the preceding emendation.}
\GSfield{Assessment}{literal statement: correct with missing hypotheses; publication match: exact; verification: read the result; development: none; confidence: high}
\end{GSitem}
\begin{GSitem}{1987-725.10}{E}{Writing the syzygy chain for a fixed Plücker ring}
\GSfield{Kind; source}{problem; \GSpage{P001412}}
\GSfield{The notes}{The requested Hilbert syzygy chain passes from determinant generators to relations and then to relations among relations; part (a) asks to write the whole chain.}
\GSfield{Statement}{For specified integers \(0<r<n\) and a field \(k\), the natural algebraic interpretation is a graded free resolution of \(S/I_{r,n}\) over \(S=k[p_I:|I|=r]\). A finite such resolution exists, with length at most \(\binom nr\), by Hilbert's syzygy theorem. Finite polynomial presentations permit effective syzygy computation. The record does not specify minimality, characteristic, equivariance, or whether a uniform closed formula is required.}
\GSfield{Reconstruction}{A fixed generic Pluecker ring is a limited interpretation of the preceding determinant discussion, not a supplied source hypothesis.}
\GSfield{Fate}{The finite-existence bound follows from \citet[Proposition 10.114.2, Tag 00OQ]{stacks2026foundations}. \citet[Chapter 7, \S 7.3]{weyman2003syzygies} specifically treats syzygies of Pl\"ucker ideals; the chapter description was checked, not a full resolution formula. Consequently neither that metadata nor the general theorem is presented as verification of a particular uniform chain intended by the speaker.}
\GSfield{Present status}{As of 11 November 2026, the finite existence and effective fixed-input task are settled; the stronger intended presentation is not specified well enough to assign it an open-problem status.}
\GSfield{Assessment}{literal statement: correct with missing hypotheses; publication match: partial; verification: indirect; development: none; confidence: medium}
\end{GSitem}
\begin{GSitem}{1987-725.11}{F}{Discretizing the whole syzygy chain}
\GSfield{Kind; source}{problem; \GSpage{P001412}}
\GSfield{The notes}{Part (b) asks to discretize the whole chain, with the gloss to write combinatorial relations.}
\GSfield{Statement}{No combinatorial objects, equivalence relation, coefficient theory, or criterion for recovering the algebraic syzygies is given. Thus the requested discretization is not a uniquely defined mathematical construction; the existence of a graded free resolution alone does not supply it.}
\GSfield{Fate}{The Pl\"ucker-syzygy literature located through \citet[\S 7.3]{weyman2003syzygies} addresses algebraic resolutions. Its metadata does not identify the extra combinatorial target in this sentence, and no exact publication match was found.}
\GSfield{Present status}{As of 11 November 2026, the omitted meaning of discretization prevents a definite solved or open classification.}
\GSfield{Assessment}{literal statement: not determinable; publication match: thematic only; verification: unresolved lead; development: none; confidence: low}
\end{GSitem}
\begin{GSitem}{1987-725.12}{E}{What information remains beyond the row and column matroids?}
\GSfield{Kind; source}{question; \GSpage{P001412}}
\GSfield{The notes}{The notes suggest constructing a bimatroid from two matroids, then warn that a bimatroid contains additional information and ask what it is.}
\GSfield{Statement}{The pair of row and column matroids does not determine the bimatroid. Over \(\mathbb R\), \(I_2\) and \(\left(\begin{smallmatrix}1&1\\1&2\end{smallmatrix}\right)\) both have free row and column matroids of rank two, but the first has two zero \(1\times1\) minors and the second has none. The extra data are the compatible joint ranks \(\rho(I,J)=\operatorname{rank}A_{I,J}\), equivalently all nonzero square-minor pairs.}
\GSfield{Reconstruction}{The two matrices and the joint-rank answer supply an explicit solution; the source does not contain this example. \oursComputed{} The independent checks are recorded in the early-record verification log.}
\GSfield{Fate}{The counterexample is computed directly. \citet[Proposition 2.6]{rohrleulirsch2026bimatroids} gives the relative-rank characterization and its equivalence with the extended matroid. It also explains why separate row and column restrictions lose joint incidence data. The final cohomological motivation in the notes specifies no separate invariant or theorem.}
\GSfield{Present status}{As of 11 November 2026, the insufficiency of two matroids has an elementary witness and a precise joint-rank description.}
\GSfield{Assessment}{literal statement: correct; publication match: partial; verification: read the result; development: annotated; confidence: high; problem \hyperref[prob:GP-1987-04]{GP-1987-04}}
\end{GSitem}

\GStalk{726}{1987-11-30}{I. M. Gelfand}{P001413-P001416}{Kerov’s Hecke representations and Frobenius characters}
Gelfand explicitly expounds Kerov's work: Hecke generators, Young-tableau adjacency, and rational weights give an orthogonal-form construction. The seminar then asks for Hecke characters and recalls the ordinary Frobenius formula. The latter ends above the divider on P001416; the following Ruelle-zeta discussion belongs to record 727.

\begin{GSitem}{1987-726.1}{E}{The two presentations of the type-A Hecke algebra}
\GSfield{Kind; source}{definition; \GSpage{P001413}}
\GSfield{The notes}{Idempotents \(e_i\) satisfy the modified braid relation; the exercise changes to \(\sigma_i=(q+1)e_i-1\).}
\GSfield{Statement}{For \(q\in\mathbb C^\times\), \(q\ne-1\), put \(\tau=q/(q+1)^2\). The relations \(e_i^2=e_i\), \(e_ie_j=e_je_i\) for \(|i-j|>1\), and \(e_ie_{i+1}e_i-\tau e_i=e_{i+1}e_ie_{i+1}-\tau e_{i+1}\) become \((\sigma_i-q)(\sigma_i+1)=0\), the ordinary braid relation, and distant commutativity. The \(\sigma\)-presentation itself also makes sense at \(q=-1\).}
\GSfield{Reconstruction}{Nonzero q and exclusion of minus one state the domain used by the cited paper; the standard presentation extends further.}
\GSfield{Fate}{\citet[\S 1, equations (1a)--(1c)]{kerov1987realizations} uses these exact generators. The change of generators is invertible only for \(q\ne-1\), a restriction hidden by the notebook's parameter formula.}
\GSfield{Present status}{As of 11 November 2026, the two presentations agree on their stated common domain.}
\GSfield{Assessment}{literal statement: correct with missing hypotheses; publication match: exact; verification: read the result; development: none; confidence: high}
\end{GSitem}
\begin{GSitem}{1987-726.2}{E}{Dimension and symmetric-group specialization}
\GSfield{Kind; source}{problem; \GSpage{P001413}}
\GSfield{The notes}{At \(q=1\) the algebra is \(\mathbb C[S_N]\); the next exercise asks for \(\dim H_N(q)=N!\).}
\GSfield{Statement}{In the standard presentation, \(T_w=T_{i_1}\cdots T_{i_\ell}\) for a reduced expression \(w=s_{i_1}\cdots s_{i_\ell}\) is well defined, and \(\{T_w:w\in S_N\}\) is a basis of \(H_N(q)\). Consequently \(\dim H_N(q)=N!\). At \(q=1\), \(T_i^2=1\) and the algebra identifies with \(\mathbb C[S_N]\).}
\GSfield{Fate}{\citet[\S 1, p.~156]{kerov1987realizations} records the dimension and specialization in the already established theory. The same Hecke-basis topic occurs in 1979-483.2.}
\GSfield{Present status}{As of 11 November 2026, the dimension is independent of the parameter; semisimplicity is a separate issue.}
\GSfield{Assessment}{literal statement: correct; publication match: exact; verification: read the result; development: none; confidence: high}
\end{GSitem}
\begin{GSitem}{1987-726.3}{E}{The modular analogy and exceptional parameters}
\GSfield{Kind; source}{claim; \GSpage{P001413}}
\GSfield{The notes}{Dipper--James is recommended under the handwritten name Tipper; the theory is called analogous to modular symmetric-group representations, even over \(\mathbb C\), followed by the non-root-of-unity assumption.}
\GSfield{Statement}{For \(q\in\mathbb C^\times\), define \(e(q)\) as the least positive integer \(e\) with \(1+q+\cdots+q^{e-1}=0\), or infinity if none exists. Then \(H_N(q)\) is semisimple if and only if \(e(q)>N\). In the semisimple case its simple modules are indexed by partitions of \(N\); at exceptional roots of unity the modular analogy becomes substantive, and simples arise as nonzero simple quotients of Specht modules.}
\GSfield{Reconstruction}{The quantum characteristic makes the informal modular analogy precise without claiming that all generic cases are modular.}
\GSfield{Fate}{\citet[\S 1, p.~156]{kerov1987realizations} states this distinction and attributes it to \citet{dipperjames1986representations}. The publisher confirms January 1986 and pages 20--52; its full result was not read. Kerov's full statement was read.}
\GSfield{Present status}{As of 11 November 2026, the generic hypothesis removes the nonsemisimple cases from the subsequent construction; complex coefficients alone do not ensure semisimplicity.}
\GSfield{Assessment}{literal statement: correct with missing hypotheses; publication match: exact; verification: read the result; development: none; confidence: high}
\end{GSitem}
\begin{GSitem}{1987-726.4}{E}{Young diagrams, tableaux, and axial adjacency}
\GSfield{Kind; source}{definition; \GSpage{P001413-P001414}}
\GSfield{The notes}{Diagrams are partitions, written with column lengths. A tableau is an increasing chain of diagrams, and its axial distance compares the cells containing \(k\) and \(k+1\).}
\GSfield{Statement}{A standard tableau of shape \(\lambda\vdash N\) is a chain \(\varnothing=\lambda^{(0)}\subset\lambda^{(1)}\subset\cdots\subset\lambda^{(N)}=\lambda\), with one cell added each time. If cell \(k\) has row-column coordinates \((i_k,j_k)\), set \(r_k(t)=(i_{k+1}-j_{k+1})-(i_k-j_k)\), \(1\le k<N\). Say \(s\sim_k t\) when their chains agree except possibly at stage \(k\). For \(N\ge2\), the recorded axial parameter is \(h(\lambda)=1+\max_{t,k}|r_k(t)|\).}
\GSfield{Fate}{\citet[\S 2, p.~157]{kerov1987realizations} gives the same definitions with \(k<N\). The notebook's upper limit \(k\le N\) refers to a nonexistent cell \(N+1\). Partitions also classify conjugacy classes of \(S_N\) by cycle lengths, but that labeling is distinct from the tableau construction.}
\GSfield{Present status}{As of 11 November 2026, the chain and axial-distance definitions are standard once the index range and row-column convention are fixed.}
\GSfield{Assessment}{literal statement: false as written; publication match: exact; verification: read the result; development: none; confidence: high}
\end{GSitem}
\begin{GSitem}{1987-726.5}{E}{The one- and two-tableau adjacency blocks}
\GSfield{Kind; source}{claim; \GSpage{P001414}}
\GSfield{The notes}{The two assertions distinguish axial distance \(\pm1\) from larger absolute value.}
\GSfield{Statement}{If \(|r_k(t)|=1\), then the only standard tableau \(s\sim_k t\) is \(t\). If \(|r_k(t)|>1\), there is exactly one other such tableau, obtained by interchanging \(k\) and \(k+1\), and it satisfies \(r_k(s)=-r_k(t)\). Thus each adjacency class has one or two elements.}
\GSfield{Fate}{\citet[\S 2, before Definition 1, p.~157]{kerov1987realizations} states this exact combinatorial fact. It follows directly from the increasing row and column conditions for a standard tableau.}
\GSfield{Present status}{As of 11 November 2026, the adjacency blocks are the elementary combinatorics behind the representation matrices.}
\GSfield{Assessment}{literal statement: correct; publication match: exact; verification: read the result; development: none; confidence: high}
\end{GSitem}
\begin{GSitem}{1987-726.6}{E}{Kerov’s tableau matrices}
\GSfield{Kind; source}{construction; \GSpage{P001414-P001415}}
\GSfield{The notes}{The matrix entry is the square root of the product of two weights when \(s\sim_k t\), and zero otherwise; the theorem says these matrices represent \(H_N(q)\).}
\GSfield{Statement}{For \(q>0\), \(q\ne1\), and \(\lambda\vdash N\), define \(S^{(k)}_{tt}=W_{r_k(t)}\), \(S^{(k)}_{st}=\sqrt{W_{r_k(s)}W_{r_k(t)}}\) for the distinct adjacent tableau, and zero otherwise, taking the nonnegative square root. With the weights of the next item and \(\tau=q/(q+1)^2\), the assignment \(e_k\mapsto S^{(k)}\) satisfies the Hecke relations and realizes the corresponding Specht module. The limit \(q=1\) recovers Young's orthogonal form. Generic complex parameters also work with coherent square-root choices.}
\GSfield{Reconstruction}{The positive-real case supplies an unambiguous square-root convention; the broader generic-complex statement requires the convention in the original construction.}
\GSfield{Fate}{\citet[\S 2, Definition 1, Theorem 1 and Remarks 1--2, pp.~157--158]{kerov1987realizations} proves the representation. The notebook expressly introduces this as a report on Kerov's existing work. Its journal date is 1987; an exact earlier circulation day was not established.}
\GSfield{Present status}{As of 11 November 2026, this is an established construction already being expounded; arbitrary independent complex square-root choices should not be substituted for the coherent convention.}
\GSfield{Assessment}{literal statement: correct with missing hypotheses; publication match: exact; verification: read the result; development: none; confidence: high}
\end{GSitem}
\begin{GSitem}{1987-726.7}{E}{The rational weights and their corrected identities}
\GSfield{Kind; source}{construction; \GSpage{P001415}}
\GSfield{The notes}{The recurrence starts at \(r\ge0\), despite \(W_0\) being undefined, and equation (19) says \(W_rW_{1-r}=1\). Explicit rational formulas follow.}
\GSfield{Statement}{For \(q\ne0,-1\) with the displayed denominators nonzero, and \(r\ge1\), \[W_r=\frac{q(q^{r-1}-1)}{(q+1)(q^r-1)},\qquad W_{-r}=\frac{q^{r+1}-1}{(q+1)(q^r-1)}.\] Then \(W_{-1}=1\), \(W_1=0\), \(W_r+W_{-r}=1\), and \(W_{-(r+1)}=1-\tau/W_{-r}\) for \(r\ge1\). For \(r\ge2\), \(W_rW_{1-r}=\tau\), not one. At \(q=1\) the limits are \(W_{\pm r}=(r\mp1)/(2r)\).}
\GSfield{Reconstruction}{Initial values, valid integer ranges, and the removable q=1 singularity are stated explicitly. \oursComputed{} The independent checks are recorded in the early-record verification log.}
\GSfield{Fate}{\citet[Appendix, equations (P1)--(P5), pp.~166--167]{kerov1987realizations} supplies the recurrence and product identity. Independent rational-function simplification confirms the factor \(\tau\) and the \(q=1\) limits.}
\GSfield{Present status}{As of 11 November 2026, the formulas are consistent after correcting the product constant and excluding the undefined zero index; poles require separate limiting representation theory.}
\GSfield{Assessment}{literal statement: false as written; publication match: exact; verification: read the result; development: annotated; confidence: high}
\end{GSitem}
\begin{GSitem}{1987-726.8}{B}{A Frobenius formula for Hecke characters}
\GSfield{Kind; source}{problem; \GSpage{P001415}}
\GSfield{The notes}{After constructing representations, Gelfand asks to find the characters of the Hecke algebra, then turns to the ordinary symmetric-group formula.}
\GSfield{Statement}{In the ring of symmetric functions in \(X\), for generic semisimple \(H_N(q)\), normalized by \((T_i-q)(T_i+1)=0\), define \[Q_r(X;q)=\frac{1}{q-1}[t^r]\prod_j\frac{1-x_jt}{1-qx_jt},\qquad Q_\mu=\prod_aQ_{\mu_a}.\] If \(w_\mu\) is the product of consecutive disjoint cycles of lengths \(\mu_a\), then \[Q_\mu(X;q)=\sum_{\lambda\vdash N}\chi_q^\lambda(T_{w_\mu})s_\lambda(X).\] This gives generic irreducible character values on the standard cycle representatives; at \(q=1\), \(Q_r\to p_r\).}
\GSfield{Reconstruction}{The generic parameter, generator normalization, and cycle representatives make the broad character request precise. \oursComputed{} The independent checks are recorded in the early-record verification log.}
\GSfield{Fate}{\citet{ram1991frobenius} published a precise Frobenius formula in 1991. The checked full formula is \citet[equation (4.10), p.~266, and Theorem 4.11]{halversonram1999bitraces}, after replacing their Hall--Littlewood \(q_r(X;q^{-1})\) by \(q^rq_r/(q-1)\). That paper explicitly credits several independent discoveries, including Vershik--Kerov; this classification makes no priority claim.}
\GSfield{Present status}{As of 11 November 2026, the generic Hecke-character problem has explicit symmetric-function formulas. Nonsemisimple root-of-unity simple characters are a different problem, not solved by this generic identity.}
\GSfield{Assessment}{literal statement: correct with missing hypotheses; publication match: partial; verification: read the result; development: none; confidence: high; problem \hyperref[prob:GP-1987-05]{GP-1987-05}}
\end{GSitem}
\begin{GSitem}{1987-726.9}{E}{The ordinary Frobenius character formula}
\GSfield{Kind; source}{theorem; \GSpage{P001413-P001416}}
\GSfield{The notes}{The opening task asks for symmetric-group characters; the closing alternant expands in power-sum products, with the coefficients equal to character values up to a simple constant.}
\GSfield{Statement}{For \(\lambda\vdash N\), \[s_\lambda=\sum_{\mu\vdash N}\frac{\chi^\lambda(\mu)}{z_\mu}p_\mu,\qquad z_\mu=\prod_{r\ge1}r^{m_r(\mu)}m_r(\mu)!,\qquad p_\mu=\prod_a p_{\mu_a}.\] Thus the omitted constant is the centralizer order \(z_\mu\). The alternant quotient is \(s_\lambda\) after reversing the notebook's increasing exponent convention to \(\lambda_j+n-j\).}
\GSfield{Reconstruction}{The total-degree condition and the exact centralizer factor replace the abbreviated coefficient description.}
\GSfield{Fate}{\citet[Chapter I, \S 7, especially equation (7.3) and the characteristic map]{macdonald1995symmetric} gives the precise Frobenius normalization. This is the classical result already cited in the notebook, not a 1987 invention.}
\GSfield{Present status}{As of 11 November 2026, the coefficients determine all irreducible symmetric-group characters, with partitions of total size N and the stated normalization.}
\GSfield{Assessment}{literal statement: correct with missing hypotheses; publication match: exact; verification: read the result; development: none; confidence: high}
\end{GSitem}

\GStalk{727}{1987-11-30}{Unidentified reporter}{P001416}{Ruelle zeta functions and transfer determinants}
The short report names a 1987 paper of D. H. Mayer and records a weighted periodic-orbit zeta function and a quotient of transfer-operator determinants. The date comes from the catalogue heading; the reporter is not identified. The preceding Hecke-character material on the same scan belongs to record 726. Analytic hypotheses needed for Fredholm determinants are absent.

\begin{GSitem}{1987-727.1}{E}{Weighted periodic-orbit zeta function}
\GSfield{Kind; source}{definition; \GSpage{P001416}}
\GSfield{The notes}{The notes define \(\zeta_T(z,\varphi)=\exp\sum_{n\ge1}\frac{z^n}{n}\sum_{\xi\in\operatorname{Fix}T^n}\prod_{k=0}^{n-1}\varphi(T^kx)\).}
\GSfield{Statement}{If every fixed-point set is finite, the same formula with \(x=\xi\) defines a formal series. It defines a holomorphic germ when its weighted fixed-point sums grow at most exponentially. The product contains exactly \(n\) factors.}
\GSfield{Reconstruction}{The bound for convergence and the correction of the dummy variable are explicit additions. No global continuation is claimed from the definition alone.}
\GSfield{Fate}{The bibliographical pointer is \citet[abstract, pp.~105--115]{mayer1987location}, published in August 1987 before this report; it places the definition in an already established thermodynamic-formalism setting. The reporter is not identified as Mayer.}
\GSfield{Present status}{As of 11 November 2026: The weighted periodic-orbit definition is established; analytic continuation depends on the dynamical and function-space hypotheses, and does not follow merely from the displayed formal series.}
\GSfield{Assessment}{literal statement: correct with missing hypotheses; publication match: exact; verification: publisher metadata checked; development: none; confidence: high}
\end{GSitem}
\begin{GSitem}{1987-727.2}{E}{A quotient of transfer-operator Fredholm determinants}
\GSfield{Kind; source}{theorem; \GSpage{P001416}}
\GSfield{The notes}{For a full-branch piecewise linear interval map the notes write \(\zeta_T=\det(1-zL_1)/\det(1-zL_2)\), with \(L_1f=\varphi\sum_j\rho_j^{-1}f\circ\psi_j\) and \(L_2f=\varphi\sum_jf\circ\psi_j\).}
\GSfield{Statement}{One precise version assumes finitely many holomorphic contracting inverse branches \(\psi_j\) on a complex domain, holomorphic weights, and a holomorphic function space on which the induced operators are nuclear and satisfy the fixed-point trace formula. Put \(L_1f(w)=\sum_j\psi_j\prime(w)\varphi(w)f(\psi_jw)\), \(L_2f(w)=\sum_j\varphi(w)f(\psi_jw)\). For affine branches, subtraction of the two fixed-point traces gives the weighted periodic-orbit sum, hence the displayed determinant quotient near zero and by continuation. Periodic endpoint multiplicities must follow the chosen coding.}
\GSfield{Reconstruction}{The analytic setting, signed derivative \(\psi_j\prime\), and boundary convention are supplied. Keeping \(\varphi(w)\) outside the sum is compatible with the cyclic product along periodic orbits. A pole-location theorem itself is not stated in the notebook.}
\GSfield{Fate}{\citet[abstract]{mayer1987location} explicitly uses nuclear operators on holomorphic Banach spaces. Only that abstract and metadata were read; the precise model stated here is a reconstruction using the trace identity, not a claim that every hypothesis was read in Mayer's text.}
\GSfield{Present status}{As of 11 November 2026: Established in analytic transfer-operator settings. The unrestricted interval-map statement is incomplete, and the title about pole location is not enough evidence for a theorem that all poles are real.}
\GSfield{Assessment}{literal statement: correct with missing hypotheses; publication match: partial; verification: publisher metadata checked; development: none; confidence: medium}
\end{GSitem}

\GStalk{728}{1987-12-07}{Unidentified reporter; intervention by I. M. Gelfand}{P001417-P001418}{Finite-state dynamical zeta functions and stochastic spectra}
The speaker proves the rationality of the zeta function of a shift of finite type, then recodes a finite-memory weight as a finite matrix. Gelfand interrupts to distinguish the set of individual stochastic eigenvalues from the set of complete spectra. The latter, explicitly marked as a problem, is the substantive open question developed in this chapter.

\begin{GSitem}{1987-728.1}{E}{Finite-type shift and periodic-point counts}
\GSfield{Kind; source}{construction; \GSpage{P001417}}
\GSfield{The notes}{The alphabet \(F\), allowed neighbour pairs \(M\subset F\times F\), bi-infinite admissible sequences and the shift \(T\) are defined. The adjacency operator is then introduced.}
\GSfield{Statement}{For the zero-one adjacency matrix \(A\) of a finite directed graph, the two-sided edge-admissible shift has \(N_n=\#\operatorname{Fix}(T^n)=\operatorname{tr}(A^n)\). The trace counts closed walks with a distinguished starting position, exactly as fixed sequences are counted.}
\GSfield{Fate}{This is the finite-matrix part of the dynamical-zeta theory already underlying the references of \citet{mayer1987location}. Its proof is the elementary closed-walk count written in the notes.}
\GSfield{Present status}{As of 11 November 2026: The count is exact for finite alphabets; no mixing, irreducibility, or aperiodicity is needed.}
\GSfield{Assessment}{literal statement: correct; publication match: exact; verification: publisher metadata checked; development: none; confidence: high}
\end{GSitem}
\begin{GSitem}{1987-728.2}{E}{Rationality of the zeta function of a finite-state shift}
\GSfield{Kind; source}{theorem; \GSpage{P001417}}
\GSfield{The notes}{The speaker derives \(\zeta_T(s)=\det(1-sA)^{-1}\) from \(N_n=\operatorname{tr}A^n\), initially for \(|s|<1/\#F\).}
\GSfield{Statement}{For every finite square matrix \(A\), \(\exp(\sum_{n\ge1}\operatorname{tr}(A^n)s^n/n)=\det(I-sA)^{-1}\) as a formal series, hence as analytic germs. For a zero-one \(m\)-state shift, \(|s|<1/m\) is a sufficient convergence disk; the rational expression continues meromorphically to the whole complex plane.}
\GSfield{Fate}{The source itself supplies the trace-log argument; the finite-state case predates the analytic-transfer work cited in \citet{mayer1987location}.}
\GSfield{Present status}{As of 11 November 2026: Completely settled; the convergence bound in the notebook need not be optimal.}
\GSfield{Assessment}{literal statement: correct; publication match: exact; verification: publisher metadata checked; development: none; confidence: high}
\end{GSitem}
\begin{GSitem}{1987-728.3}{E}{Single-eigenvalue constraints for a stochastic matrix}
\GSfield{Kind; source}{remark; \GSpage{P001418}; I. M. Gelfand}
\GSfield{The notes}{Gelfand recalls that a stochastic matrix has eigenvalue \(1\) and spectrum in the unit disk, and mentions work of Dynkin and Dmitriev on the possible position of an individual eigenvalue.}
\GSfield{Statement}{For a real row-stochastic \(n\times n\) matrix \(P\), \(P\mathbf1=\mathbf1\) and \(\rho(P)\le\|P\|_\infty=1\); nonreal eigenvalues occur in conjugate pairs. Determining the union of possible individual eigenvalues projects away all correlations between eigenvalues and is a different question from realizing a whole multiset.}
\GSfield{Fate}{The elementary bound follows directly from the stated stochastic hypotheses. The exact single-eigenvalue region was determined by \citet{karpelevich1951roots}; its opened record confirms publication in1951, received19May1950. The source attribution to Dynkin and Dmitriev concerns earlier work and is retained. The joint problem is assessed separately in item1987-728.4.}
\GSfield{Present status}{As of 11 November 2026: These elementary necessary conditions are known and are far from a sufficient joint-spectrum criterion.}
\GSfield{Assessment}{literal statement: correct; publication match: partial; verification: publisher metadata checked; development: none; confidence: high}
\end{GSitem}
\begin{GSitem}{1987-728.4}{C}{All eigenvalues of a stochastic matrix simultaneously}
\GSfield{Kind; source}{problem; \GSpage{P001418}; I. M. Gelfand}
\GSfield{The notes}{Gelfand expressly asks for the region of complete eigenvalue lists in \(\mathbb C^n\), emphasizing that the known plane region describes only its projection. The Russian marks this \ru{Задача!!}.}
\GSfield{Statement}{For each fixed \(n\), characterize the multisets \(\Lambda=(1,\lambda_2,\ldots,\lambda_n)\) for which a real matrix \(P\ge0\), \(P\mathbf1=\mathbf1\), has characteristic polynomial \(\prod_{j=1}^n(t-\lambda_j)\). Multiplicity is algebraic; reducible and non-diagonalizable matrices are allowed. No doubly stochastic, symmetry or prescribed-zero-pattern condition is imposed.}
\GSfield{Fate}{The order-three characterization predates the seminar: \citet{loewylondon1978inverse} appeared in 1978, with its precise alternatives recorded in \citet[p.~506]{ranteng2024patterns}; the latter treats prescribed order-three zero patterns in 2024. \citet[Theorems~3.9--3.12; Sections~4--5, 8]{johnsonpaparella2025perron}, first circulated in September 2024 and published in 2025, develops geometry and realizing families for joint stochastic spectra. \citet[Section~2]{croninlaffey2018diagonalizable}, preprint January 2017 and publication July 2018, records the five-dimensional obstruction used below. Section~\ref{dev:stochastic} gives complete low-dimensional proofs and an explicit account of what these sources do not solve.}
\GSfield{Present status}{As of 11 November 2026: low-dimensional realizations and many structural tests are known, but the checked literature supplies no explicit general characterization in orders at least five. Fixed-order membership for effectively encoded algebraic input is decidable by real quantifier elimination \citep[Chapters~2, 13--14]{basupollackroy2006algorithms}; this does not provide the requested structural description. The chapter proves the two- and three-state cases and exact counterexamples, without claiming novelty or a solution for arbitrary order.}
\GSfield{Assessment}{literal statement: correct; publication match: partial; verification: read the result; development: developed in this chapter; confidence: high; problem \hyperref[prob:GP-1987-21]{GP-1987-21}}
\GSfield{Worth developing}{A bounded development is worthwhile because exact small-dimensional realizations and explicit moment obstructions sharpen the distinction the notes make between a spectral projection and a full spectrum.}
\end{GSitem}
\begin{GSitem}{1987-728.5}{E}{Finite-memory weights reduce to a finite matrix}
\GSfield{Kind; source}{theorem; \GSpage{P001418}}
\GSfield{The notes}{The speaker replaces the alphabet by admissible blocks of length \(k\), gives the overlap transition rule, and obtains \(\zeta_T(s,\varphi)=\det(I-sL)^{-1}\). The orbit product in this page is printed with upper limit \(n\).}
\GSfield{Statement}{If \(\varphi\) depends on \(k\) consecutive symbols, let \(L_{\alpha\beta}=\varphi(\alpha)\) when the admissible blocks overlap and concatenate, and zero otherwise. Then \(\operatorname{tr}L^n\) is the sum of \(\prod_{j=0}^{n-1}\varphi(T^jx)\) over period-n fixed sequences, and the weighted zeta function is \(\det(I-sL)^{-1}\).}
\GSfield{Reconstruction}{The upper bound is corrected from n to n-1, as required by the stated transition matrix; no arbitrary long-range potential is covered.}
\GSfield{Fate}{This is the same finite-dimensional trace-log identity as item 1987-728.2, applied to the higher-block presentation; it is expository material in the transfer-operator setting of \citet{mayer1987location}.}
\GSfield{Present status}{As of 11 November 2026: The corrected finite-memory theorem is exact. The literal product of n+1 weights is false: for a one-symbol shift of constant weight c it gives \((1-sc)^{-c}\), rather than \((1-sc)^{-1}\).}
\GSfield{Assessment}{literal statement: false as written; publication match: exact; verification: publisher metadata checked; development: annotated; confidence: high}
\end{GSitem}

\GStalk{729}{1987-12-07}{A. A. Beilinson}{P001419-P001422}{Adelic string amplitudes, Tamagawa volumes, and epsilon factors}
Beilinson explains the product formula, the functional equation of the Riemann zeta function, adelic amplitudes attributed to Freund--Witten, and Tamagawa volume calculations. He then moves to Fourier transforms, determinant lines and Laumon's local epsilon factors. The suggested adelic analogue of moduli-space volumes remains insufficiently specified. The determinant formula must be read with care: its literal sign convention already fails on a constant local system.

\begin{GSitem}{1987-729.1}{E}{The rational product formula and p-adic arithmetic}
\GSfield{Kind; source}{theorem; \GSpage{P001419}}
\GSfield{The notes}{After Gelfand urges the audience to learn p-adic arithmetic and root extraction, Beilinson writes \(\prod_v|a|_v=1\), with \(a\in\mathbb Q\), and checks \(a=1/2\).}
\GSfield{Statement}{For \(a\in\mathbb Q^\times\), normalize \(|p|_p=p^{-1}\) and take the usual real absolute value. Then \(|a|_\infty\prod_p|a|_p=1\), with only finitely many nonunit factors. Learning root extraction means solving equations when solutions exist; it does not assert that \(\mathbb Q_p\) is algebraically closed.}
\GSfield{Reconstruction}{The nonzero hypothesis and normalization are supplied.}
\GSfield{Fate}{The product formula and the adelic measure framework are classical; \citet[Chapters~I--II]{serre1981adeles} predates the talk. The displayed example follows immediately from the valuation of 2.}
\GSfield{Present status}{As of 11 November 2026: The theorem is exact after excluding a=0. No new assertion that arbitrary p-adic elements admit every root is attributed to Gelfand.}
\GSfield{Assessment}{literal statement: correct with missing hypotheses; publication match: exact; verification: publisher metadata checked; development: none; confidence: high}
\end{GSitem}
\begin{GSitem}{1987-729.2}{E}{Completed zeta functional equation and a wrong power of pi}
\GSfield{Kind; source}{theorem; \GSpage{P001419}}
\GSfield{The notes}{Equation (2) is the completed Riemann zeta functional equation. Equation (3) rewrites it with a factor \(\pi^{1-s}\) and an Euler product, explicitly saying analytic continuation is understood.}
\GSfield{Statement}{With \(\Lambda(s)=\pi^{-s/2}\Gamma(s/2)\zeta(s)\), one has \(\Lambda(s)=\Lambda(1-s)\) meromorphically. Consequently \(\zeta(1-s)/\zeta(s)=\pi^{1/2-s}\Gamma(s/2)/\Gamma((1-s)/2)\), wherever the quotient is defined and then by meromorphic continuation. The formal product \(\prod_p(1-p^{-s})/(1-p^{s-1})\) is not a pair of ordinarily convergent Euler products on any common half-plane.}
\GSfield{Reconstruction}{The gamma ratio is corrected by direct division of equation (2).}
\GSfield{Fate}{\citet[25.4.3--25.4.4]{nist2026dlmf} gives the established symmetric equation. The notes' exponent follows neither from it nor from a change of gamma convention.}
\GSfield{Present status}{As of 11 November 2026: The equation is classical; the exponent 1-s is a faithful notebook error and must be 1/2-s. Analytic continuation, explicitly mentioned in the source, remains essential.}
\GSfield{Assessment}{literal statement: false as written; publication match: exact; verification: read the result; development: annotated; confidence: high}
\end{GSitem}
\begin{GSitem}{1987-729.3}{B}{Adelic factorization of string amplitudes}
\GSfield{Kind; source}{claim; \GSpage{P001419}}
\GSfield{The notes}{The gamma ratio is described as approximately an open-string Veneziano--Witten amplitude, with a product over primes. Three recent papers are named: Volovich, Manin, and Witten with an unnamed coauthor.}
\GSfield{Statement}{The identifiable reference is Freund--Witten's factorization of appropriate four-particle open- and closed-string amplitudes into p-adic factors, after analytic continuation and normalization. The single-variable gamma ratio displayed here is an illustrative specialization; it does not specify the full three Mandelstam variables, on-shell constraint, or normalization of a complete amplitude.}
\GSfield{Reconstruction}{The amplitude identification comes from the named topic and opened abstract; no missing multivariable formula is attributed to the speaker.}
\GSfield{Fate}{\citet[abstract]{freundwitten1987adelic} identifies the missing coauthor and states the Veneziano and Virasoro--Shapiro product relations. Journal publication was 17 December 1987, ten days after the seminar; the report itself is evidence of earlier circulation. Only metadata and abstract were read.}
\GSfield{Present status}{As of 11 November 2026: The particular paper is identified with high confidence. The notebook's schematic one-variable expression is only a partial match to its complete amplitude formulas.}
\GSfield{Assessment}{literal statement: correct with missing hypotheses; publication match: partial; verification: publisher metadata checked; development: none; confidence: high}
\end{GSitem}
\begin{GSitem}{1987-729.4}{E}{Tamagawa volume of SL(n) and the real lattice covolume}
\GSfield{Kind; source}{theorem; \GSpage{P001420}}
\GSfield{The notes}{The shorthand \ru{Мера Тамагавы = 1} is expanded into
\[
\begin{aligned}
 \operatorname{vol}\bigl(\mathrm{SL}_n(\mathbb A)/\mathrm{SL}_n(\mathbb Q)\bigr)&=1,\\
 \operatorname{vol}\bigl(\mathrm{SL}_n(\mathbb R)/\mathrm{SL}_n(\mathbb Z)\bigr)&=\prod_{j=2}^n\zeta(j).
\end{aligned}
\]}
\GSfield{Statement}{For \(n\ge2\), use the restricted adelic product and Tamagawa Haar measures from a rational invariant top form. Then \(\tau(SL_n)=1\). With compatible real and p-adic normalization and \(\operatorname{vol}(SL_n(\mathbb Z_p))=\prod_{j=2}^n(1-p^{-j})\), strong approximation gives the stated real covolume. Replacing Haar measure by an arbitrary scalar multiple changes that covolume.}
\GSfield{Reconstruction}{The restricted product, quotient, and common normalization are made explicit.}
\GSfield{Fate}{These are classical volume calculations; \citet[Chapters~II--III, especially the contents at pp.~122 and 180]{serre1981adeles} records the relevant pre-1987 framework. Full scan access exceeded the reader limit, so verification is metadata, not a claim to have read its proof.}
\GSfield{Present status}{As of 11 November 2026: The normalized theorem is settled. The statement concerns the volume of an arithmetic quotient, not the total measure of the group or an unnormalized Haar measure.}
\GSfield{Assessment}{literal statement: correct with missing hypotheses; publication match: exact; verification: publisher metadata checked; development: none; confidence: high}
\end{GSitem}
\begin{GSitem}{1987-729.5}{E}{The SL(2) volume picture on the hyperbolic plane}
\GSfield{Kind; source}{construction; \GSpage{P001420}}
\GSfield{The notes}{The notes identify \(SL_2(\mathbb R)\) with disk automorphisms and use hyperbolic area times an angular measure on \(S^*H^2\).}
\GSfield{Statement}{Orientation-preserving hyperbolic-plane isometries form \(PSL_2(\mathbb R)=SL_2(\mathbb R)/\{\pm I\}\); they act simply transitively on oriented unit tangent frames. Hyperbolic area times fibre-angle measure is Haar measure on this frame group. Lifting it to \(SL_2\) requires accounting for the double cover before comparing with the Tamagawa normalization.}
\GSfield{Reconstruction}{The central quotient and double-cover normalization are supplied.}
\GSfield{Fate}{The construction is classical and illustrates the arithmetic-volume calculation of item 1987-729.4; see the real-measure and arithmetic-group sections listed in \citet{serre1981adeles}.}
\GSfield{Present status}{As of 11 November 2026: The geometric model is valid after distinguishing SL from PSL and fixing the angular period. It alone does not select the normalization giving zeta(2).}
\GSfield{Assessment}{literal statement: correct with missing hypotheses; publication match: partial; verification: publisher metadata checked; development: none; confidence: high}
\end{GSitem}
\begin{GSitem}{1987-729.6}{E}{Counting GL(n) and SL(n) over a finite field}
\GSfield{Kind; source}{theorem; \GSpage{P001420-P001421}}
\GSfield{The notes}{The audience is asked to count \(SL_n(\mathbb F_p)\); the next page counts successively independent columns of \(GL_n\) and divides by \(p-1\).}
\GSfield{Statement}{For a prime power \(q\), \(\#GL_n(\mathbb F_q)=\prod_{i=0}^{n-1}(q^n-q^i)\), because each new column must avoid the span of its predecessors. The surjective determinant map has equally sized fibres, so \(\#SL_n=\#GL_n/(q-1)=q^{n^2-1}\prod_{j=2}^n(1-q^{-j})\).}
\GSfield{Fate}{The argument is fully present in the source and belongs to the classical-group point counts listed at p.~20 of \citet{serre1981adeles}.}
\GSfield{Present status}{As of 11 November 2026: The formula is exact, including n=1 with the empty product. It supplies the local factors used in the previous volume formula.}
\GSfield{Assessment}{literal statement: correct; publication match: exact; verification: publisher metadata checked; development: none; confidence: high}
\end{GSitem}
\begin{GSitem}{1987-729.7}{F}{An adelic volume analogy for moduli of curves}
\GSfield{Kind; source}{question; \GSpage{P001421}}
\GSfield{The notes}{Manin is credited with an analogous formula for a moduli-space volume; Gelfand asks what p-adic volume and an adelic Teichmuller space would mean.}
\GSfield{Statement}{No precise moduli functor, genus, marked-point data, differential form or measure normalization, convergence condition, or equation is given. The strongest supported content is a request to formulate an adelic analogue of the preceding volume comparison.}
\GSfield{Fate}{Searches for Manin's arithmetical-physics reflections identify a contemporaneous programme, but neither the notebook nor the accessible bibliographic records supplies an exact conjecture to match. The explicit amplitude relation of \citet{freundwitten1987adelic} concerns a different, specified object.}
\GSfield{Present status}{As of 11 November 2026: No definite statement can be assigned a solved or open status from these notes. Later p-adic moduli and string constructions do not by themselves identify this particular question.}
\GSfield{Assessment}{literal statement: not determinable; publication match: thematic only; verification: unresolved lead; development: none; confidence: low}
\end{GSitem}
\begin{GSitem}{1987-729.8}{E}{Fourier transformation moves local information to infinity}
\GSfield{Kind; source}{method; \GSpage{P001421}}
\GSfield{The notes}{A compactly supported generalized function on the line is Fourier transformed to a smooth function; the notes say its singularities become asymptotic behaviour at infinity, then mention Laumon's use of Witten's Morse-theoretic ideas.}
\GSfield{Statement}{For \(u\in\mathcal E\prime(\mathbb R)\), \(\widehat u(z)=\langle u(x),e^{-izx}\rangle\) is entire of exponential type and of polynomial growth on the real axis. Smoothness and singular support are reflected in directional decay of localized Fourier transforms. A single unspecified asymptotic expansion need not determine every singularity; the full Fourier transform determines u.}
\GSfield{Reconstruction}{The entire extension and localized decay formulation clarify the informal word singularities; they are not asserted to be verbatim in the notes.}
\GSfield{Fate}{\citet[Introduction, pp.~131--134]{laumon1987fourier} explicitly develops the finite-field Fourier analogy and acknowledges Witten's Morse-theory inspiration. The real-distribution statement is standard background, not a theorem first proved there.}
\GSfield{Present status}{As of 11 November 2026: The real Fourier statement is established. Laumon's arithmetic stationary-phase construction is a precise analogue, not an identification of real distributions with l-adic sheaves.}
\GSfield{Assessment}{literal statement: correct with missing hypotheses; publication match: partial; verification: read the result; development: none; confidence: high}
\end{GSitem}
\begin{GSitem}{1987-729.9}{E}{Determinant line of a finite complex}
\GSfield{Kind; source}{construction; \GSpage{P001421-P001422}}
\GSfield{The notes}{The notes set \(\det H^\bullet=(\det H^0\otimes\det H^2)\otimes(\det H^1)^{-1}\) and assert equality with the analogous determinant of a complex.}
\GSfield{Statement}{For a bounded complex \(C^\bullet\) of finite-dimensional vector spaces, define \(\det C^\bullet=\bigotimes_i(\bigwedge^{\dim C^i}C^i)^{(-1)^i}\). The short exact sequences of boundaries, cycles and cohomology give a natural isomorphism \(\det C^\bullet\simeq\det H^\bullet(C)\), with coherent graded signs. Finiteness or a perfect-complex hypothesis is necessary.}
\GSfield{Reconstruction}{The boundedness, finite-dimensionality and graded-sign conditions are supplied.}
\GSfield{Fate}{The construction was established before the talk in \citet[Section~I]{knudsenmumford1976det}; the opened author-hosted scan confirms the publication, while the exact finite-dimensional isomorphism is also obtained directly from the two short exact sequences.}
\GSfield{Present status}{As of 11 November 2026: Standard determinant-functor theory. This valid isomorphism must be separated from the following divisor formula, whose displayed sign is inconsistent.}
\GSfield{Assessment}{literal statement: correct with missing hypotheses; publication match: exact; verification: publisher metadata checked; development: none; confidence: high}
\end{GSitem}
\begin{GSitem}{1987-729.10}{E}{The determinant-divisor formula has the wrong variance}
\GSfield{Kind; source}{claim; \GSpage{P001422}}
\GSfield{The notes}{For \(D=\sum a_ix_i\), the notes define \(\det F_D=\bigotimes_i(\det F_{x_i})^{a_i}\), then write \(\det H^\bullet(X,F)=\det F_K\) for the canonical divisor K.}
\GSfield{Statement}{The divisor construction is meaningful for a local system F and a specified divisor. The claimed natural identity with K is false under the immediately preceding alternating determinant convention. For a constant local system V on a compact genus-g curve, automorphisms of V act on \(\det R\Gamma(X,V)\) with determinant weight \(2-2g\), while \(\det V_K\) has weight \(2g-2\). Dualizing the right side removes this elementary sign obstruction, but a canonical general product formula requires additional local epsilon data and conventions.}
\GSfield{Reconstruction}{The comparison of determinant weights is our consistency check. We do not silently replace the source formula by an asserted canonical isomorphism with \(F_{-K}\).}
\GSfield{Fate}{\citet[Section~I]{knudsenmumford1976det} fixes the determinant convention; \citet[Chapter~III, Theorem~3.2.1.1]{laumon1987fourier} proves a precise arithmetic epsilon-product statement, not this unqualified identity of lines.}
\GSfield{Present status}{As of 11 November 2026: The literal formula is false functorially. The sign diagnosis is decisive even on a constant sheaf on the projective line; it does not establish a replacement theorem for arbitrary constructible sheaves.}
\GSfield{Assessment}{literal statement: false as written; publication match: partial; verification: read the result; development: annotated; confidence: high}
\end{GSitem}
\begin{GSitem}{1987-729.11}{E}{Laumon: local factors for global epsilon constants}
\GSfield{Kind; source}{theorem; \GSpage{P001422}}
\GSfield{The notes}{The discussion of determinants is presented as a report on Laumon's work in number theory and on a local description of the functional-equation constant.}
\GSfield{Statement}{For a smooth proper curve over a finite field of characteristic p and a constructible l-adic complex, \(l\ne p\), the global epsilon constant of its L-function is expressed as a product of locally defined epsilon factors, with the normalization and genus/rank factor specified by Laumon. A nonzero rational differential supplies the local additive characters; the resulting global expression is independent of that choice. Local Fourier transforms and stationary phase supply the proof.}
\GSfield{Reconstruction}{The arithmetic objects and hypotheses come from the identified paper, while the notebook gives only the motivating outline.}
\GSfield{Fate}{\citet[Theorem~3.2.1.1; Remark~3.2.1.3]{laumon1987fourier} contains the exact product theorem. It was already published in 1987 when Beilinson reported it in December.}
\GSfield{Present status}{As of 11 November 2026: Established. The field, coefficient and normalization hypotheses matter: this result cannot be read as a blanket complex-analytic determinant-line formula.}
\GSfield{Assessment}{literal statement: correct with missing hypotheses; publication match: partial; verification: read the result; development: none; confidence: high}
\end{GSitem}

\GStalk{730}{1987-12-07}{A. A. Klyachko}{P001422-P001424}{Equivariant toric bundles and compatible filtrations}
Klyachko associates decreasing filtrations to invariant divisors of a toric variety and states a compatibility criterion, then asks for cohomology, characteristic classes and classification. Simplicial-complex examples lead to the triangle and tetrahedral boundary. His 1989 paper, submitted in February 1988, supplies the precise smooth and singular compatibility conditions and geometric calculations. The tetrahedral decomposition requires the short-filtration restriction that the notebook omits.

\begin{GSitem}{1987-730.1}{E}{Toric varieties, fans, and equivariant vector bundles}
\GSfield{Kind; source}{definition; \GSpage{P001422-P001423}}
\GSfield{The notes}{The speaker defines a toric variety by an algebraic-torus action with an orbit isomorphic to the torus, associates a fan, and defines a compatible linear action on a vector bundle.}
\GSfield{Statement}{A normal toric variety over an algebraically closed field contains \(T\simeq(\mathbb G_m)^d\) as a dense open orbit, with multiplication extending to the variety. It corresponds to a rational polyhedral fan in \(N_\mathbb R\), where \(N=\operatorname{Hom}(\mathbb G_m,T)\) is the cocharacter lattice. A toric vector bundle is an algebraic vector bundle equipped with a T-action linear on fibres and covering the base action.}
\GSfield{Reconstruction}{Density, normality and the integral cocharacter structure are supplied.}
\GSfield{Fate}{This is the standard background recalled by \citet[Section~1]{klyachko1989equivariant}. The existence of an orbit merely isomorphic to T is not enough without density and normality.}
\GSfield{Present status}{As of 11 November 2026: Established definitions. The combinatorial lattice belongs to one-parameter subgroups, not an arbitrary vector-space identification with the Lie algebra.}
\GSfield{Assessment}{literal statement: correct with missing hypotheses; publication match: exact; verification: read the result; development: none; confidence: high}
\end{GSitem}
\begin{GSitem}{1987-730.2}{A}{Klyachko classification by compatible filtrations}
\GSfield{Kind; source}{theorem; \GSpage{P001423}}
\GSfield{The notes}{An equivariant bundle yields a vector space E and decreasing filtrations \(E^\alpha(i)\), one for each ray. The notes require the subspaces on each cone to generate a distributive lattice.}
\GSfield{Statement}{On a smooth toric variety \(X_\Sigma\), toric vector bundles correspond to finite-dimensional E with exhaustive separated decreasing integer filtrations \(E^\rho(i)\). On every cone \(\sigma\) there must be a decomposition \(E=\bigoplus_{u\in M}E_u\) with \(E^\rho(i)=\sum_{\langle u,v_\rho\rangle\ge i}E_u\) for all rays of that cone. Morphisms preserve every filtration. On a smooth cone simultaneous splitting, equivalently distributivity, supplies such integral characters. For singular cones the integral-character compatibility must be imposed explicitly.}
\GSfield{Reconstruction}{The smooth/singular distinction, finite filtration conditions, character lattice, and morphisms are supplied from the identified theorem.}
\GSfield{Fate}{\citet[Theorems~0.1.1 and 2.2.1]{klyachko1989equivariant} gives the precise equivalence. The paper was received on 3 February 1988, published in Russian in 1989 and translated in 1990; the seminar documents circulation in December 1987.}
\GSfield{Present status}{As of 11 November 2026: The theorem is established, including a precise singular version. The raw distributivity criterion without smoothness is incomplete.}
\GSfield{Assessment}{literal statement: correct with missing hypotheses; publication match: exact; verification: read the result; development: none; confidence: high}
\end{GSitem}
\begin{GSitem}{1987-730.3}{A}{Recover cohomology and characteristic classes from filtrations}
\GSfield{Kind; source}{problem; \GSpage{P001423}}
\GSfield{The notes}{The explicit tasks are to compute cohomology, Chern classes and related geometric quantities from the filtration data.}
\GSfield{Statement}{Given a smooth complete toric variety with fan \(\Sigma\) and a compatible filtered vector space representing \(\mathcal E\), determine the weight spaces of \(H^q(X_\Sigma,\mathcal E)\) and the Chern classes from that finite combinatorial data. Completeness ensures finite-dimensional global cohomology; for a noncomplete variety the characterwise construction still needs the corresponding finiteness interpretation.}
\GSfield{Reconstruction}{Completeness and the precise output invariants are added to make the question finite and testable.}
\GSfield{Fate}{\citet[Theorems~3.2.1, 4.1.1 and 4.2.1]{klyachko1989equivariant} gives Chern-class formulas, a complex computing cohomology and character formulas. Submission February 1988 followed this December 1987 report; publication was 1989/1990.}
\GSfield{Present status}{As of 11 November 2026: The stated recovery problem is solved by Klyachko's constructions. This is not an unsolved request for a new cohomology theory.}
\GSfield{Assessment}{literal statement: correct with missing hypotheses; publication match: exact; verification: read the result; development: none; confidence: high; problem \hyperref[prob:GP-1987-22]{GP-1987-22}}
\end{GSitem}
\begin{GSitem}{1987-730.4}{B}{Classification and moduli of compatible filtration data}
\GSfield{Kind; source}{problem; \GSpage{P001423}}
\GSfield{The notes}{The second broad task is classification of the bundles; the speaker then abstracts the compatibility rule to a finite simplicial complex with a filtration at every vertex.}
\GSfield{Statement}{For a fixed complex K, a fixed rank r and prescribed dimensions of the filtration steps, classify the configurations of flags in \(\mathbb C^r\) whose subspaces are simultaneously split on each simplex, modulo simultaneous \(GL_r\). For a toric fan, use the stronger character compatibility when needed. Constructing a parameter space and understanding its orbit equivalence is a precise interpretation of classification; it need not produce a finite list.}
\GSfield{Reconstruction}{Rank, flag dimensions and the quotient problem are supplied; the paper cited solves the toric fixed-data moduli problem, not every possible meaning of classification.}
\GSfield{Fate}{For toric bundles, \citet[Theorem~1.1 and Section~4]{payne2008moduli}, first circulated in May 2007 and published in 2008, constructs framed moduli with fixed equivariant Chern data and proves Murphy's law already in rank three. This develops the geometric moduli question, while not claiming a simple orbit list for every abstract simplicial complex.}
\GSfield{Present status}{As of 11 November 2026: Established moduli constructions and severe singularity phenomena explain the limits of a uniform elementary classification. The arbitrary-complex reformulation remains a representation problem with complexity depending on K, not one unqualified open conjecture.}
\GSfield{Assessment}{literal statement: correct with missing hypotheses; publication match: partial; verification: read the result; development: none; confidence: high; problem \hyperref[prob:GP-1987-23]{GP-1987-23}}
\end{GSitem}
\begin{GSitem}{1987-730.5}{E}{A filled simplex has only one-dimensional split summands}
\GSfield{Kind; source}{theorem; \GSpage{P001423}}
\GSfield{The notes}{For the triangle, all three filtrations are compatible and the speaker decomposes the representation into one-dimensional ones.}
\GSfield{Statement}{If all filtrations attached to the vertices of a filled simplex are simultaneously split, choose a common basis of E adapted to every filtration. The one-dimensional coordinate lines, endowed with their integer jump positions, give a direct-sum decomposition of the filtered object. This uses a filled simplex; its boundary has weaker compatibility.}
\GSfield{Reconstruction}{The distinction between a filled triangle and its boundary is made explicit.}
\GSfield{Fate}{This is the elementary simultaneous-basis content of the distributivity condition in \citet[Theorem~0.1.1]{klyachko1989equivariant}; no later classification is needed.}
\GSfield{Present status}{As of 11 November 2026: The splitting is exact for the stated common-basis condition. It must not be transferred to complexes where only proper faces are constrained.}
\GSfield{Assessment}{literal statement: correct; publication match: exact; verification: read the result; development: none; confidence: high}
\end{GSitem}
\begin{GSitem}{1987-730.6}{A}{Tetrahedral configurations and the short-filtration restriction}
\GSfield{Kind; source}{claim; \GSpage{P001423-P001424}}
\GSfield{The notes}{For the boundary of a tetrahedron the notes give four planes in general position in a three-dimensional space, their exterior powers, and assert that all representations split into those and one-dimensional objects.}
\GSfield{Statement}{The matching theorem concerns toric bundles on projective space whose ray filtrations are short: each has at most one proper nonzero subspace, with allowed integer shifts. For \(\mathbb P^3\), such bundles decompose into line bundles and the appropriate twists of exterior powers of its cotangent bundle. Their filtration configurations include the four general-position planes described in the notes. The same conclusion for arbitrary long filtrations does not follow.}
\GSfield{Reconstruction}{The short-filtration hypothesis is taken from the proved theorem. The source does not state it explicitly, so the literal unrestricted assertion is not endorsed.}
\GSfield{Fate}{\citet[Theorem~6.4.1]{klyachko1989equivariant} states the required short-filtration theorem; published 1989/1990 after its December 1987 exposition. The missing restriction is substantive.}
\GSfield{Present status}{As of 11 November 2026: The restricted theorem is established. The unqualified notebook sentence is not a classification of all vector bundles on projective three-space or of all arbitrary-length tetrahedral flag configurations.}
\GSfield{Assessment}{literal statement: correct with missing hypotheses; publication match: partial; verification: read the result; development: none; confidence: high}
\end{GSitem}
\begin{GSitem}{1987-730.7}{E}{Isolating tractable pieces of a representation problem}
\GSfield{Kind; source}{remark; \GSpage{P001424}; I. M. Gelfand}
\GSfield{The notes}{Gelfand urges finding the solvable portions of any classification problem and refers to the ideology of Gelfand--Ponomarev.}
\GSfield{Statement}{The supported content is a methodological recommendation: split direct-sum objects into indecomposables and identify representation-finite or otherwise tractable subclasses before seeking a uniform description. It is not a claim that every classification problem has a computable finite list.}
\GSfield{Fate}{The concrete antecedent is \citet{gelfandponomarev1971quadruples}, already discussed in the earlier index at 1970-198 and 1972-277; the present sheet later explicitly returns to the four-subspace problem.}
\GSfield{Present status}{As of 11 November 2026: A documented methodological link, with no new unsolved mathematical assertion.}
\GSfield{Assessment}{literal statement: correct; publication match: thematic only; verification: publisher metadata checked; development: none; confidence: high}
\end{GSitem}
\begin{GSitem}{1987-730.8}{E}{Kac: indecomposable dimension vectors are positive roots}
\GSfield{Kind; source}{theorem; \GSpage{P001424}}
\GSfield{The notes}{The notes connect dimensions at graph vertices to coordinates of roots of a Kac--Moody algebra, mention numbers of parameters, and cite Lecture Notes996 and an earlier Inventiones paper.}
\GSfield{Statement}{For a finite quiver without loops over an algebraically closed field, the dimension vectors of indecomposable representations are exactly the positive roots of the associated symmetric Kac--Moody root system. A positive real root gives a unique indecomposable up to isomorphism. For imaginary roots there can be parameter families; a bare dimension vector does not specify a unique object. The Tits quadratic form controls the parameter count in the precise representation-variety formulation.}
\GSfield{Reconstruction}{The field, quiver and real-versus-imaginary hypotheses are supplied, and no unverified universal moduli-dimension formula is substituted for the shorthand.}
\GSfield{Fate}{The exact volume cited is \citet[pp.~74--108]{kac1983roots}, published in1983 and listed on Kac's own bibliography. The linked scan was too large to read here; verification is limited to identification and metadata.}
\GSfield{Present status}{As of 11 November 2026: Kac's root theorem is established and predates the talk. It classifies dimension vectors and real-root objects, not every imaginary-root isomorphism class by a single root.}
\GSfield{Assessment}{literal statement: correct with missing hypotheses; publication match: exact; verification: publisher metadata checked; development: none; confidence: high}
\end{GSitem}

\GStalk{731}{}{-}{P001425}{Year 1988 heading}
A year heading only. It has no independent mathematical assertion and no occurrence-index items. The source then places Gelfand's exercises before Beilinson's programme, despite the numerical ordering of catalogue records.


\GStalk{733}{1988-02-08}{I. M. Gelfand}{P001425}{Exercises on determinants, real roots, limits and conformal maps}
Five elementary questions ask for Cauchy's determinant, real roots of Rodrigues'polynomial, discontinuities of pointwise limits, and conformal maps of a slit disk and a dumbbell domain. They are exercises in established mathematics. The narrow corridor is specified only pictorially, so a numerical map is not determined by the notes.

\begin{GSitem}{1987-733.1}{E}{Evaluate Cauchy's determinant}
\GSfield{Kind; source}{problem; \GSpage{P001425}}
\GSfield{The notes}{Problem1 asks for \(\det(1/(x_i-y_j))\).}
\GSfield{Statement}{If \(x_i\ne y_j\) for every i,j, then \[\det_{1\le i,j\le n}\frac1{x_i-y_j}=\frac{\prod_{i<j}(x_j-x_i)(y_i-y_j)}{\prod_{i,j}(x_i-y_j)}.\] Repeated x or y values give zero, consistently with the formula.}
\GSfield{Fate}{This is the classical double alternant, checked in \citet[Section~2.1, equation~(2.7)]{krattenthaler1999determinant}; his1999 exposition is not the first publication. It recurs from item1981-541.3 in the earlier index.}
\GSfield{Present status}{As of 11 November 2026: Solved before the seminar. The difference-denominator sign is included explicitly.}
\GSfield{Assessment}{literal statement: correct; publication match: exact; verification: read the result; development: none; confidence: high; problem \hyperref[prob:GP-1987-24]{GP-1987-24}}
\end{GSitem}
\begin{GSitem}{1987-733.2}{E}{Real roots of the Rodrigues polynomial}
\GSfield{Kind; source}{problem; \GSpage{P001425}}
\GSfield{The notes}{Problem2 asks to prove that the nth derivative of \((x^2-1)^n\) has only real roots.}
\GSfield{Statement}{For every integer \(n\ge1\), \(\frac{d^n}{dx^n}(x^2-1)^n=2^nn!P_n(x)\), where \(P_n\) is the Legendre polynomial. It has exactly n simple roots, all in \((-1,1)\). Repeated Rolle arguments applied to the endpoint roots of multiplicity n already prove the required reality.}
\GSfield{Fate}{The Rodrigues formula and orthogonal-polynomial zero theorem are recorded in \citet[18.5(ii),18.2(vi)]{nist2026dlmf}; item1968-160.15 records the earlier seminar connection.}
\GSfield{Present status}{As of 11 November 2026: Solved classical exercise. n=0 gives a nonzero constant and no roots.}
\GSfield{Assessment}{literal statement: correct; publication match: exact; verification: read the result; development: none; confidence: high; problem \hyperref[prob:GP-1987-25]{GP-1987-25}}
\end{GSitem}
\begin{GSitem}{1987-733.3}{E}{Discontinuities of a pointwise limit of continuous functions}
\GSfield{Kind; source}{problem; \GSpage{P001425}}
\GSfield{The notes}{Problem3 asks whether a pointwise limit on an interval can be discontinuous, have a dense set of discontinuities, or be discontinuous everywhere.}
\GSfield{Statement}{A real-valued pointwise limit of continuous functions on a nondegenerate real interval is Baire class1 and has a dense G-delta set of continuity points. It may be discontinuous (for example \(x^n\) on \([0,1]\)), and discontinuities may be dense (Thomae's function, with value \(1/q\) at a reduced rational \(p/q\) and0 at irrationals). It cannot be discontinuous everywhere.}
\GSfield{Fate}{\citet[p.~173, theorem and reference to the1899 thesis]{baire1900nouvelle} provides the classical criterion for pointwise continuous approximation. The category conclusion applies to interval domains, which are Baire spaces.}
\GSfield{Present status}{As of 11 November 2026: All three questions are settled: yes, yes, no. Pointwise convergence must not be replaced by uniform convergence.}
\GSfield{Assessment}{literal statement: correct; publication match: exact; verification: read the result; development: none; confidence: high; problem \hyperref[prob:GP-1987-26]{GP-1987-26}}
\end{GSitem}
\begin{GSitem}{1987-733.4}{E}{Conformally map the radially slit disk}
\GSfield{Kind; source}{problem; \GSpage{P001425}}
\GSfield{The notes}{Problem4(a) draws a disk cut along a radius and asks for a conformal map to a disk.}
\GSfield{Statement}{For \(\Omega=\{z:|z|<1\}\setminus[0,1)\), choose \(w=\sqrt z\) with \(0<\arg z<2\pi\). Then w maps onto the upper half-disk; \(h=(1+w)/(1-w)\) maps that onto the first quadrant. Therefore \(f(z)=(h^2-i)/(h^2+i)\) is a biholomorphism from \(\Omega\) to the unit disk. A rotated slit is treated by an initial rotation.}
\GSfield{Reconstruction}{The unit radius and positive-real slit choose coordinates for the picture; the construction is supplied here.}
\GSfield{Fate}{The existence is a classical case of the mapping theorem originating in \citet[pp.~37--39 of the checked transcription]{riemann1851grundlagen}. The elementary displayed composition is an annotation of the exercise, not a formula found in the notebook.}
\GSfield{Present status}{As of 11 November 2026: Solved. The slit endpoint0 is removed, so no normalization f(0)=0 is imposed.}
\GSfield{Assessment}{literal statement: correct; publication match: exact; verification: read the result; development: annotated; confidence: high; problem \hyperref[prob:GP-1987-27]{GP-1987-27}}
\end{GSitem}
\begin{GSitem}{1987-733.5}{E}{Conformally map the dumbbell domain with a thin corridor}
\GSfield{Kind; source}{problem; \GSpage{P001425}}
\GSfield{The notes}{Problem4(b) draws two unit disks joined by a corridor of width \(10^{-6}\), marks A at a disk centre, and asks for a disk map sending A to0.}
\GSfield{Statement}{If the drawn union is an open connected simply connected proper plane domain \(\Omega\) and \(A\in\Omega\), the Riemann mapping theorem gives a biholomorphism \(f:\Omega\to\mathbb D\) with \(f(A)=0\); fixing \(f\prime(A)>0\) makes it unique. Corridor width alone does not determine \(\Omega\): its length, boundary shape and junctions are also needed for a particular numerical formula.}
\GSfield{Reconstruction}{The topological hypotheses and derivative normalization are stated explicitly; no unseen corridor length is invented.}
\GSfield{Fate}{This is the same established mapping theorem as in \citet{riemann1851grundlagen}. Thin corridors can make numerical conformal parametrization poorly conditioned, but the notes give no quantitative estimate to verify.}
\GSfield{Present status}{As of 11 November 2026: Existence and normalized uniqueness are settled. A unique explicit numerical map cannot be extracted from the incomplete geometric dimensions in the sketch.}
\GSfield{Assessment}{literal statement: correct with missing hypotheses; publication match: exact; verification: read the result; development: none; confidence: high; problem \hyperref[prob:GP-1987-27]{GP-1987-27}}
\end{GSitem}

\GStalk{732}{1988-02-08}{A. A. Beilinson}{P001425}{Programme for learning string theory}
This is an agenda, not a theorem-bearing lecture: two-dimensional conformal theory; operator products; representations of infinite-dimensional Lie algebras; parafermions and \(W\)-algebras; vertex operators; reduction from affine \(\mathfrak{sl}_2\) to Virasoro; moduli; partition and correlation functions; and branching of characters. It gets a talk-list entry but no artificial mathematical items.


\GStalk{734}{1988-02-08}{D. V. Alekseevsky}{P001426-P001428}{G-structures, first jets, connections and torsion}
The lecture develops reductions of the frame bundle, their tensor examples and solder forms. First jets of sections become horizontal complements, giving connections and a torsion map. These are standard constructions, not a new existence theorem for canonical connections. The source contains two useful corrections: the principal action is transitive on each fibre, and the vertical tangent quotient belongs to \(P\), not to \(M\).

\begin{GSitem}{1987-734.1}{E}{A G-structure is a principal reduction}
\GSfield{Kind; source}{definition; \GSpage{P001426}}
\GSfield{The notes}{The first definition takes a principal subbundle \(P\subset F(M)\) for \(G\subset GL(V)\), then says that G acts transitively on P.}
\GSfield{Statement}{For a smooth n-manifold M and closed Lie subgroup \(G\subset GL(n,\mathbb R)\), a G-structure is a principal G-subbundle of its frame bundle. The right G-action is free and transitive on each fibre \(P_x\), not on all of P when the base has more than one point.}
\GSfield{Reconstruction}{The word each fibre repairs the source statement.}
\GSfield{Fate}{This is the standard definition recalled in \citet[Section~5.1]{alekseevskymichor1995cartan}; that later paper explicitly reviews older theory. The earlier index has the same reduction in1980-493.1.}
\GSfield{Present status}{As of 11 November 2026: Established; the global transitivity sentence is a faithful mathematical error in the notebook.}
\GSfield{Assessment}{literal statement: false as written; publication match: exact; verification: read the result; development: none; confidence: high}
\end{GSitem}
\begin{GSitem}{1987-734.2}{E}{Tensor fields of constant algebraic type define reductions}
\GSfield{Kind; source}{construction; \GSpage{P001426}}
\GSfield{The notes}{The examples include a metric, a tensor in one canonical form, an endomorphism with fixed Jordan form, an almost complex structure and a nondegenerate2-form.}
\GSfield{Statement}{If a smooth tensor field takes values in one \(GL(n)\)-orbit and its adapted frames form a smooth reduction, its stabilizer is the structure group. A positive metric yields \(O(n)\); \(J^2=-I\) in dimension2m yields \(GL(m,\mathbb C)\); a nondegenerate alternating2-form yields \(Sp(2m,\mathbb R)\). An almost symplectic form becomes symplectic only with the additional differential condition \(d\omega=0\). An endomorphism requires locally constant algebraic type for its adapted frames.}
\GSfield{Reconstruction}{The orbit/reduction smoothness and even-dimensional requirements are made explicit.}
\GSfield{Fate}{These are the elementary G-structure examples in the framework of \citet[Section~5.1]{alekseevskymichor1995cartan}; the note itself distinguishes almost symplectic from symplectic.}
\GSfield{Present status}{As of 11 November 2026: Established. A pointwise tensor classification does not imply differential integrability of the resulting structure.}
\GSfield{Assessment}{literal statement: correct with missing hypotheses; publication match: exact; verification: read the result; development: none; confidence: high}
\end{GSitem}
\begin{GSitem}{1987-734.3}{E}{Equivalent description by a solder form}
\GSfield{Kind; source}{theorem; \GSpage{P001426-P001427}}
\GSfield{The notes}{The second definition gives a principal G-bundle with a V-valued1-form \(\theta\), equivariance, and kernel equal to the vertical tangent space. The equivalence is explained using coframes.}
\GSfield{Statement}{For a faithful inclusion \(G\subset GL(V)\), a principal G-bundle \(P\to M\) with \(\dim M=\dim V\) and a horizontal equivariant form \(\theta\) satisfying \(\ker\theta_p=\ker d\pi_p\) determines a coframe \(T_xM\to V\) at each p over x. Equivariance \(R_g^*\theta=g^{-1}\theta\) makes these coframes a G-reduction, and conversely the tautological solder form of a reduction has these properties.}
\GSfield{Reconstruction}{The quotient's base space and the right-action convention are corrected.}
\GSfield{Fate}{\citet[Section~5.1]{alekseevskymichor1995cartan} recalls precisely this solder-form formulation. It was established G-structure theory before1988.}
\GSfield{Present status}{As of 11 November 2026: The equivalence is standard. The quotient used in the notebook must be \(T_pP/T_p^vP\), which is canonically \(T_xM\).}
\GSfield{Assessment}{literal statement: false as written; publication match: exact; verification: read the result; development: none; confidence: high}
\end{GSitem}
\begin{GSitem}{1987-734.4}{E}{First jets of sections as horizontal complements}
\GSfield{Kind; source}{construction; \GSpage{P001427}}
\GSfield{The notes}{The first jet \(j_x^1s\) is identified with a right inverse of the projection derivative and with a horizontal subspace.}
\GSfield{Statement}{For \(s(x)=p\), its first jet is determined by \(d s_x:T_xM\to T_pP\) with \(d\pi_p\circ d s_x=I\). Its image H complements the vertical space. Conversely every such linear splitting is realized by a local section to first order. At fixed p these splittings form an affine space modeled on \(\operatorname{Hom}(T_xM,\mathfrak g)\), identified with \(\operatorname{Hom}(V,\mathfrak g)\) using p.}
\GSfield{Fate}{\citet[Section~6.1]{alekseevskymichor1995cartan} gives exactly the jet/horizontal-space identification, as recalled standard theory.}
\GSfield{Present status}{As of 11 November 2026: Established. This is the unrestricted first-jet bundle, before any torsion normalization defining a smaller prolongation.}
\GSfield{Assessment}{literal statement: correct; publication match: exact; verification: read the result; development: none; confidence: high}
\end{GSitem}
\begin{GSitem}{1987-734.5}{E}{A principal connection as an equivariant horizontal distribution}
\GSfield{Kind; source}{definition; \GSpage{P001427-P001428}}
\GSfield{The notes}{The space of connections is written as the G-equivariant sections of \(P\prime\to P\); horizontal lifts give parallel transport along curves.}
\GSfield{Statement}{A principal connection is a smooth G-equivariant assignment \(p\mapsto H_p\) with \(T_pP=H_p\oplus V_pP\). Equivalently it is a \(\mathfrak g\)-valued1-form \(\omega\) reproducing fundamental vertical fields and satisfying \(R_g^*\omega=\operatorname{Ad}(g^{-1})\omega\). The horizontal-lift ODE defines parallel transport along piecewise smooth paths in its existence interval; for principal connections it extends along compact parameter intervals by local trivializations.}
\GSfield{Fate}{The principal-connection framework is recalled by \citet[Sections~2.6 and3]{alekseevskymichor1995cartan}. It is background, not a new canonical connection selected by every G-structure.}
\GSfield{Present status}{As of 11 November 2026: Established. The lecture defines the space of connections; it does not prove that torsion can always be set to zero.}
\GSfield{Assessment}{literal statement: correct; publication match: exact; verification: read the result; development: none; confidence: high}
\end{GSitem}
\begin{GSitem}{1987-734.6}{E}{The torsion function on the space of horizontal complements}
\GSfield{Kind; source}{definition; \GSpage{P001428}}
\GSfield{The notes}{Using the vertical parallelism and \(\theta|_H:H\to V\), the notes define \(t(p\prime)(\xi,\eta)=d\theta((\theta|_H)^{-1}\xi,(\theta|_H)^{-1}\eta)\).}
\GSfield{Statement}{For every horizontal complement H at p, this formula defines \(t(H)\in\operatorname{Hom}(\Lambda^2V,V)\). For a selected connection, it is the coefficient function of the torsion2-form \(d\theta+\omega\wedge\theta\) evaluated on horizontal vectors. The alternating2-form, rather than an unrestricted bilinear map, is the natural target.}
\GSfield{Fate}{\citet[Section~6.1]{alekseevskymichor1995cartan} records the same torsion function and its role in prolongation.}
\GSfield{Present status}{As of 11 November 2026: Standard construction. No invariant zero-torsion complement is guaranteed without an algebraic condition on the intrinsic torsion.}
\GSfield{Assessment}{literal statement: correct; publication match: exact; verification: read the result; development: none; confidence: high}
\end{GSitem}
\begin{GSitem}{1987-734.7}{E}{Change of horizontal complement and the Spencer alternation}
\GSfield{Kind; source}{claim; \GSpage{P001428}}
\GSfield{The notes}{The final sentence asserts \(G\)-equivariance and \(G^\prime=\operatorname{Hom}(V,\mathfrak g)\)-equivariance of torsion, without displaying the action on its target.}
\GSfield{Statement}{If \(H_A\) is obtained from H by adding the fundamental vertical vector corresponding to \(A(v)\), then with right-action convention \(t(H_A)=t(H)-\delta A\), where \(\delta A(v,w)=A(v)w-A(w)v\). Thus the class in \(\operatorname{Hom}(\Lambda^2V,V)/\delta\operatorname{Hom}(V,\mathfrak g)\) is independent of H. G acts linearly; the full space of changes of complement acts affinely by this alternation. The kernel of delta is the first algebraic prolongation.}
\GSfield{Reconstruction}{The omitted action and its sign are supplied by the right-principal-bundle convention.}
\GSfield{Fate}{\citet[Sections~6.1--6.2]{alekseevskymichor1995cartan} explicitly uses the Spencer map to normalize torsion and define prolongation.}
\GSfield{Present status}{As of 11 November 2026: Established after specifying the affine action and sign convention. The phrase equivariant alone does not mean torsion is invariant under all changes of connection.}
\GSfield{Assessment}{literal statement: correct with missing hypotheses; publication match: exact; verification: read the result; development: none; confidence: high}
\end{GSitem}

\GStalk{735}{1988-02-15}{I. M. Gelfand}{P001429}{Subspace angles and a thermodynamics reading suggestion}
Gelfand refers to his 1959 exposition of quasilinear equations and asks about two planes in four dimensions and complete invariants of two Hilbert subspaces. The former are principal angles and the latter are described by the theory of two projections. The remarks anticipate the separate problem sheet of 22 February.

\begin{GSitem}{1987-735.1}{E}{A reading suggestion on thermodynamics and quasilinear equations}
\GSfield{Kind; source}{remark; \GSpage{P001429}}
\GSfield{The notes}{Gelfand recommends his 1959 paper in Uspekhi, volume 14, issue 2, saying that three pages contain all thermodynamics.}
\GSfield{Statement}{The identifiable reference is the 1959 exposition Some problems in the theory of quasi-linear equations, pp.~87--158. The remark recommends its thermodynamic discussion; it does not state a new theorem or identify a three-page passage precisely enough to extract one.}
\GSfield{Fate}{\citet[publication record]{gelfand1959quasilinear} confirms the exact title, volume, pages and receipt on 17 November 1958. The full thermodynamic section was not read in this research.}
\GSfield{Present status}{As of 11 November 2026: The bibliographic identification is secure. The phrase all thermodynamics is an informal assessment, not an assertion of mathematical completeness.}
\GSfield{Assessment}{literal statement: correct; publication match: thematic only; verification: publisher metadata checked; development: none; confidence: high}
\end{GSitem}
\begin{GSitem}{1987-735.2}{E}{Complete invariants of two planes in four dimensions}
\GSfield{Kind; source}{question; \GSpage{P001429}}
\GSfield{The notes}{Gelfand asks what should be called the angle between two two-dimensional planes in four-dimensional Euclidean space.}
\GSfield{Statement}{For two linear two-planes \(U,W\) in a four-dimensional real inner-product space, choose orthonormal frame matrices \(Q_U,Q_W\). The singular values of \(Q_U^TQ_W\) are \(\cos\theta_1,\cos\theta_2\), where \(0\le\theta_1\le\theta_2\le\pi/2\). This ordered pair is a complete invariant under the full orthogonal group, and every pair in that range occurs.}
\GSfield{Reconstruction}{The word planes is interpreted as linear subspaces and the acting group as the full orthogonal group, consistently with the following problem sheet.}
\GSfield{Fate}{The singular-value characterization was already explicit in \citet[Section~3, Theorem~1]{bjorckgolub1973angles}, published in 1973. The earlier index records related angle questions in 1964-46.2. The same question is restated as problem 3 of record 737.}
\GSfield{Present status}{As of 11 November 2026: Solved classical finite-dimensional classification. One scalar angle loses information in this dimension.}
\GSfield{Assessment}{literal statement: correct; publication match: exact; verification: read the result; development: none; confidence: high; problem \hyperref[prob:GP-1987-28]{GP-1987-28}}
\end{GSitem}
\begin{GSitem}{1987-735.3}{E}{Complete invariants of two closed Hilbert subspaces}
\GSfield{Kind; source}{question; \GSpage{P001429}}
\GSfield{The notes}{The second question asks for all invariants of a pair of subspaces in a Hilbert space.}
\GSfield{Statement}{For an ordered pair of closed subspaces M,N of a real or complex Hilbert space, the four intersections \(M\cap N\), \(M\cap N^\perp\), \(M^\perp\cap N\), \(M^\perp\cap N^\perp\) contribute their Hilbert dimensions. On the remaining generic part \(K\oplus K\), the projections have forms \(P=\left(\begin{smallmatrix}I&0\\0&0\end{smallmatrix}\right)\) and \(Q=\left(\begin{smallmatrix}C^2&CS\\CS&S^2\end{smallmatrix}\right)\), where \(0<C<I\) means that C and \(S=(I-C^2)^{1/2}\) are injective positive contractions. Their endpoints may still belong to the continuous spectrum. The unitary or orthogonal equivalence class of C, together with the four dimensions, is complete.}
\GSfield{Reconstruction}{Closedness, the full isometry group and the endpoint-kernel convention are supplied; these are essential for a projection-based answer.}
\GSfield{Fate}{\citet[Theorems~1--3]{halmos1969subspaces} gives the two-projection canonical form and its completeness; the paper was received in April 1969. The statement is repeated and elaborated in record 737.}
\GSfield{Present status}{As of 11 November 2026: Solved for closed Hilbert subspaces. A mere set of spectral values is insufficient: spectral measure class and multiplicity are part of the invariant.}
\GSfield{Assessment}{literal statement: correct with missing hypotheses; publication match: exact; verification: read the result; development: none; confidence: high; problem \hyperref[prob:GP-1987-32]{GP-1987-32}}
\end{GSitem}

\GStalk{736}{1988-02-15}{A. B. Zamolodchikov}{P001429-P001434}{Conformal field theory and the Virasoro algebra}
The lecture starts from formal Euclidean correlation functions, operator products and the stress tensor; conservation and tracelessness lead to chiral fields, the stress-tensor OPE and Virasoro modes. Reflection positivity and primary fields are discussed. Several conditions implicit in unitary conformal theory must be restored, and the printed contravariant complex indices reverse the asserted chirality. The broad physical remarks about a smallest length and polyhedra do not determine mathematical theorems.

\begin{GSitem}{1987-736.1}{E}{Formal Euclidean correlations from a local action}
\GSfield{Kind; source}{construction; \GSpage{P001429}}
\GSfield{The notes}{The speaker writes correlation functions as \(Z^{-1}\int A_1(x_1)\cdots A_n(x_n)e^{-S[\varphi]}\mathcal D\varphi\), with \(S=\int\sigma(\varphi,\partial\varphi)\,d^2x\).}
\GSfield{Statement}{This is a formal prescription for multilinear Euclidean correlation distributions. To make it a construction one must specify a measure or regularized limit, composite fields and coincident-point renormalization, and ensure \(0<Z<\infty\) when an ordinary normalization is used. The displayed continuum symbol \(\mathcal D\varphi\) alone does not provide these data.}
\GSfield{Reconstruction}{Distributional interpretation and analytic existence conditions are supplied. The 1975 reconstruction source was checked at metadata/abstract level only.}
\GSfield{Fate}{The established operator/correlation formulation is described by \citet[Introduction and Section~2]{belavinpolyakovzamolodchikov1984infinite}; a rigorous Euclidean reconstruction framework predates it in \citet[abstract]{osterwalderschrader1975euclidean}.}
\GSfield{Present status}{As of 11 November 2026: Established formal language, not a general existence theorem for every local action. Specific rigorously defined models require additional analytic input.}
\GSfield{Assessment}{literal statement: correct with missing hypotheses; publication match: partial; verification: read the result; development: none; confidence: high}
\end{GSitem}
\begin{GSitem}{1987-736.2}{E}{Operator products and recursive determination of correlations}
\GSfield{Kind; source}{method; \GSpage{P001429-P001430}}
\GSfield{The notes}{The fields form an algebra with \(A_i(x)A_j(y)=\sum_k c_{ij}^{\,k}(x,y)A_k(y)\). The notes say these coefficients and the one-point functions determine all correlations by repeated expansion; translation invariance makes the coefficients functions of \(x-y\).}
\GSfield{Statement}{If the chosen field system has convergent OPE expansions in correlation functions on nested domains, and these expansions satisfy associativity, the OPE coefficients together with vacuum one-point functions determine higher correlation functions there by iteration and elsewhere by any justified analytic continuation. Formal expansions alone do not imply convergence or independence of the sequence of pairings.}
\GSfield{Reconstruction}{The convergence domains and compatibility of repeated expansions are restored.}
\GSfield{Fate}{\citet[Introduction, equation~(1.6)]{belavinpolyakovzamolodchikov1984infinite} explicitly formulates the OPE with convergence in a finite domain and associativity. This is background from 1984, not a result originating in the 1988 lecture.}
\GSfield{Present status}{As of 11 November 2026: Established under the stated OPE hypotheses. The notebook suppresses precisely the analytic and associativity conditions needed for the proposed recursion.}
\GSfield{Assessment}{literal statement: correct with missing hypotheses; publication match: exact; verification: read the result; development: none; confidence: high}
\end{GSitem}
\begin{GSitem}{1987-736.3}{E}{Stress tensor and the translation Ward identity}
\GSfield{Kind; source}{definition; \GSpage{P001430}}
\GSfield{The notes}{A conserved symmetric field \(T^{\mu\nu}\) is required to generate the derivative of a field by a surrounding contour integral. The repeated formula contracts \(\varepsilon_{\mu\nu}\,dx^\lambda\).}
\GSfield{Statement}{A translation-generating stress tensor satisfies \(\partial_\mu T^{\mu\nu}=0\) away from other insertions, with contact terms at insertions. With the normalization of the charge fixed, \(\oint T^{\mu\nu}(x)A(y)\varepsilon_{\mu\lambda}\,dx^\lambda=\partial_{y^\nu}A(y)\) expresses its action on a local field. Conservation justifies contour deformations avoiding other insertions.}
\GSfield{Reconstruction}{The lambda contraction and the contact-term qualification repair the displayed identity.}
\GSfield{Fate}{\citet[Section~2, equations~(2.1)--(2.10)]{belavinpolyakovzamolodchikov1984infinite} develops the stress-tensor Ward identity. The notation in the first notebook display has the correct free index; its repetition does not.}
\GSfield{Present status}{As of 11 November 2026: Standard Ward-identity construction. The literal repeated contraction is ill-formed, and conservation cannot be imposed without contact terms at coincident insertions.}
\GSfield{Assessment}{literal statement: false as written; publication match: exact; verification: read the result; development: none; confidence: high}
\end{GSitem}
\begin{GSitem}{1987-736.4}{E}{Reflection positivity requires a full quadratic form}
\GSfield{Kind; source}{claim; \GSpage{P001430}}
\GSfield{The notes}{For the reflection \(\vartheta(x^1,x^2)=(-x^1,x^2)\), the matrix \(\langle A_i^+(\vartheta x)A_j(x)\rangle\) is required to be positive definite for every x.}
\GSfield{Statement}{The Euclidean positivity axiom is positive semidefiniteness of \(\langle(\Theta F)^*F\rangle\) for every finite linear combination F of smeared field products supported in one open time half-space. The reflected form can have null vectors, which are quotiented out in Hilbert-space reconstruction. Positivity of the displayed same-point matrix is only a restricted necessary test and does not by itself imply the full axiom.}
\GSfield{Reconstruction}{The half-space, smearing, arbitrary finite sums and semidefinite quotient are missing in the notes.}
\GSfield{Fate}{The opened record for \citet[abstract]{osterwalderschrader1975euclidean} confirms the corrected Euclidean reconstruction theorem; full-text retrieval was blocked, so no claim is made to have read that proof. The full quadratic-form condition is the explicit modern reconstruction used here.}
\GSfield{Present status}{As of 11 November 2026: Established framework. Strict pointwise positivity is too strong in the presence of null fields and too weak as a substitute for reflection positivity on all field products.}
\GSfield{Assessment}{literal statement: correct with missing hypotheses; publication match: partial; verification: publisher metadata checked; development: none; confidence: high}
\end{GSitem}
\begin{GSitem}{1987-736.5}{E}{Construct substantive models of the field axioms}
\GSfield{Kind; source}{question; \GSpage{P001430}}
\GSfield{The notes}{After the algebra, symmetry and positivity requirements the speaker asks for meaningful examples of all this.}
\GSfield{Statement}{A definite conformal version asks for a local two-dimensional theory with associative OPE, a stress tensor and a positive state space. The unitary Ising conformal theory supplies such a standard example, with central charge \(c=1/2\) and primary weights \((0,0),(1/2,1/2),(1/16,1/16)\) in its diagonal local sector. This is an example answering the request, not a classification of all models.}
\GSfield{Reconstruction}{The Ising model is supplied as one precise response to the broad request. Section 6 of the cited paper gives its conformal weights and closed operator products.}
\GSfield{Fate}{\citet[Section~6, equations~(6.11)--(6.18)]{belavinpolyakovzamolodchikov1984infinite} already develops degenerate theories and their correlation functions in 1984; the lecture is expository. The exact positivity completion of a continuum path integral is a separate analytic question.}
\GSfield{Present status}{As of 11 November 2026: The existence request had established nontrivial examples before the talk. The notebook does not state a definite universal classification problem.}
\GSfield{Assessment}{literal statement: correct with missing hypotheses; publication match: partial; verification: read the result; development: none; confidence: high}
\end{GSitem}
\begin{GSitem}{1987-736.6}{E}{Conservation and tracelessness give chiral stress components}
\GSfield{Kind; source}{claim; \GSpage{P001431}}
\GSfield{The notes}{With \(z=x^1+ix^2\), the notebook defines \(T=T^{zz}\), \(\bar T=T^{\bar z\bar z}\), sets \(T^{z\bar z}=0\), then asserts \(\partial_{\bar z}T=\partial_z\bar T=0\).}
\GSfield{Statement}{For a genuine contravariant symmetric tensor, \(T^{zz}=T^{11}-T^{22}+2iT^{12}\) and \(T^{z\bar z}=T^{11}+T^{22}\). Conservation and zero trace imply \(\partial_zT^{zz}=0\), so the displayed contravariant T has the opposite chirality. The conventional holomorphic component is the covariant \(T_{zz}=(T^{11}-T^{22}-2iT^{12})/4\), satisfying \(\partial_{\bar z}T_{zz}=0\). Rescaling this component gives the normalization used in the OPE.}
\GSfield{Reconstruction}{This calculation distinguishes contravariant from covariant components rather than silently using the same symbol for both.}
\GSfield{Fate}{The holomorphic stress-tensor construction is standard in \citet[Section~2, equations~(2.6)--(2.10)]{belavinpolyakovzamolodchikov1984infinite}. Direct tensor transformation and the two real conservation equations expose the notebook index inconsistency.}
\GSfield{Present status}{As of 11 November 2026: The standard chiral conclusion is correct after an explicit index/chirality correction. The separate equality \(T^{z\bar z}=T^\mu{}_\mu\) on P001433 is consistent with contravariant indices and \(ds^2=dz\,d\bar z\); it should not be independently changed by a factor of two.}
\GSfield{Assessment}{literal statement: false as written; publication match: partial; verification: read the result; development: annotated; confidence: high}
\end{GSitem}
\begin{GSitem}{1987-736.7}{E}{Composite fields and the limits of Grassmann nilpotence}
\GSfield{Kind; source}{remark; \GSpage{P001431}}
\GSfield{The notes}{The example lists electron and photon composites, says that powers \(\psi^5,\psi^6\) vanish because the spinor has four components, and then warns that a pointwise \(A^2\) need not exist.}
\GSfield{Statement}{For four classical Grassmann generators at one point every homogeneous product of degree above four vanishes. This elementary exterior-algebra statement does not bound all quantum composites: derivatives give additional generators, and operator-valued distributions require a specified renormalized product. Even the product of two bosonic fields at the same point is not automatically defined.}
\GSfield{Reconstruction}{The distinction between finite Grassmann generators and quantum local fields supplies the necessary scope.}
\GSfield{Fate}{The field-system/OPE formulation of \citet[Introduction, equation~(1.6)]{belavinpolyakovzamolodchikov1984infinite} includes derivatives and composite fields. The finite Grassmann calculation is an elementary reconstruction, not a general theorem about electron-field operators.}
\GSfield{Present status}{As of 11 November 2026: The warning about coincident products is sound. The nilpotence assertion needs the restricted finite Grassmann-algebra interpretation; it is not an identity for every quantum spinor composite.}
\GSfield{Assessment}{literal statement: correct with missing hypotheses; publication match: thematic only; verification: read the result; development: none; confidence: high}
\end{GSitem}
\begin{GSitem}{1987-736.8}{E}{Contour modes of the stress tensor}
\GSfield{Kind; source}{construction; \GSpage{P001431}}
\GSfield{The notes}{The modes are defined by \(\oint z^{n+1}T(z)A(w)\,dz/(2\pi i)\) on a contour around w; \(L_{-1}=\partial_w\) is then asserted.}
\GSfield{Statement}{In a local coordinate \(t=z-w\), define \(L_n^{(w)}A(w)=\operatorname{Res}_{z=w}(z-w)^{n+1}T(z)A(w)\). The simple-pole Ward identity gives \(L_{-1}^{(w)}A=\partial_wA\). The global modes about the origin use \(z^{n+1}\); they act on fields at w by coordinate-dependent combinations and must be distinguished from local descendant modes.}
\GSfield{Reconstruction}{The local/global coordinate distinction is made explicit.}
\GSfield{Fate}{\citet[Section~2, equations~(2.18)--(2.21); Section~3]{belavinpolyakovzamolodchikov1984infinite} gives the contour generators and Laurent modes. The earlier occurrence 1982-573.17 records the same stress-tensor construction.}
\GSfield{Present status}{As of 11 November 2026: Established. The source formula is valid as a global charge expression, but the local highest-weight formulas later on require a coordinate centred at the insertion or an insertion at zero.}
\GSfield{Assessment}{literal statement: correct with missing hypotheses; publication match: exact; verification: read the result; development: none; confidence: high}
\end{GSitem}
\begin{GSitem}{1987-736.9}{E}{The singular stress-tensor operator product}
\GSfield{Kind; source}{theorem; \GSpage{P001432}}
\GSfield{The notes}{The Laurent expansion is assumed to start at order \(-4\); the coefficients are stated to be \(T_{-4}=cI/2\), \(T_{-3}=0\), \(T_{-2}=2T\), \(T_{-1}=\partial T\).}
\GSfield{Statement}{For the normalized holomorphic stress tensor of an ordinary local CFT with a scalar central term and the conformal Ward identity, \[T(z)T(w)\sim\frac{c/2}{(z-w)^4}+\frac{2T(w)}{(z-w)^2}+\frac{\partial T(w)}{z-w}.\] The absence of a third-order pole and the scalar fourth-order coefficient use the stress-tensor transformation law and vacuum assumptions. Reflection positivity alone does not determine this OPE.}
\GSfield{Reconstruction}{The notebook attributes more to positivity than it alone supplies; the Ward/vacuum assumptions are restored.}
\GSfield{Fate}{The equivalent infinitesimal transformation with the Schwarzian central term is given in \citet[Section~2, equations~(2.12), (2.19)--(2.21)]{belavinpolyakovzamolodchikov1984infinite}, before this lecture.}
\GSfield{Present status}{As of 11 November 2026: Standard exact singular OPE under the stated CFT assumptions; regular terms contain further composite fields.}
\GSfield{Assessment}{literal statement: correct with missing hypotheses; publication match: exact; verification: read the result; development: none; confidence: high}
\end{GSitem}
\begin{GSitem}{1987-736.10}{E}{The two commuting Virasoro algebras}
\GSfield{Kind; source}{theorem; \GSpage{P001432}}
\GSfield{The notes}{The boxed formula is \([L_n,L_m]=(n-m)L_{n+m}+\frac c{12}(n^3-n)\delta_{n+m,0}\), and an antiholomorphic copy is described.}
\GSfield{Statement}{Taking residues in the stress OPE yields the displayed commutator, with the central term multiplying the identity. The antiholomorphic modes satisfy the analogous relation, and left and right modes commute in the local theory. Modulo its centre this is the Witt algebra; integration of the central term requires a central extension or projective representation of the diffeomorphism group, not simply the unextended group.}
\GSfield{Reconstruction}{The identity operator in the central term and the group-extension qualification are explicit.}
\GSfield{Fate}{\citet[Section~2, equation~(2.21)]{belavinpolyakovzamolodchikov1984infinite} states precisely these relations. Prior occurrences include 1982-573.1 and 1985-657.1--4. No reducibility or singular-vector problem is posed in this passage, so the distinct earlier problem GP-1982-03 is not relabelled as a new problem here.}
\GSfield{Present status}{As of 11 November 2026: Established algebraic relation. The restriction to \(n=-1,0,1\) has zero central term and recovers the projective conformal subalgebra.}
\GSfield{Assessment}{literal statement: correct with missing hypotheses; publication match: exact; verification: read the result; development: none; confidence: high}
\end{GSitem}
\begin{GSitem}{1987-736.11}{F}{An asserted universal smallest length}
\GSfield{Kind; source}{remark; \GSpage{P001432-P001433}}
\GSfield{The notes}{The speaker says that lengths below \(10^{-31}\) cm have no meaning, independently of theory, from relations among universal constants, and asks what space could mean with such a restriction.}
\GSfield{Statement}{No axioms defining meaningful distance, observable localization or an admissible geometry are stated. A dimensional combination of constants gives a candidate scale; it does not by itself prove a lower bound on all measured distances or a theorem about spaces. Consequently no precise universal theorem can be reconstructed from this assertion.}
\GSfield{Reconstruction}{No metric or observational definition is invented.}
\GSfield{Fate}{The conformal-field-theory source \citet{belavinpolyakovzamolodchikov1984infinite} identifies the surrounding lecture but provides no identified result establishing this assertion. The search and source check did not identify an exact mathematical theorem behind it.}
\GSfield{Present status}{As of 11 November 2026: The literal philosophical/physical claim is not a definite mathematical open problem. The scanned exponent is retained; no updated numerical constant is substituted for it.}
\GSfield{Assessment}{literal statement: not determinable; publication match: none found; verification: unresolved lead; development: none; confidence: high}
\end{GSitem}
\begin{GSitem}{1987-736.12}{F}{An unspecified adelic determinant algorithm}
\GSfield{Kind; source}{remark; \GSpage{P001433}}
\GSfield{The notes}{The interjection says that some determinants are best computed by an adelic method.}
\GSfield{Statement}{The determinant family, coefficient domain, input size, arithmetic-cost model and algorithm are all unspecified. Therefore the comparative claim best cannot be converted into a definite complexity theorem or numerical-analysis problem.}
\GSfield{Fate}{The surrounding references to p-adic physics and the prior adelic lecture do not identify a particular determinant algorithm. Searches of the stated phrase and determinant/adelic terminology did not produce an exact attributable match.}
\GSfield{Present status}{As of 11 November 2026: Unidentified as a precise mathematical claim; neither an open-status judgment nor a claim of computational superiority is warranted.}
\GSfield{Assessment}{literal statement: not determinable; publication match: none found; verification: unresolved lead; development: none; confidence: medium}
\end{GSitem}
\begin{GSitem}{1987-736.13}{F}{A broad comparison of polygon and polyhedron theory}
\GSfield{Kind; source}{remark; \GSpage{P001433}}
\GSfield{The notes}{The notes contrast complete knowledge of two-dimensional polyhedra with supposed little progress on three- and four-dimensional ones in two thousand years.}
\GSfield{Statement}{The objects and questions are unspecified: convexity, combinatorial type, realizability, face enumeration and metric rigidity lead to different problems. The passage states no quantified theorem that can be assigned a precise present status.}
\GSfield{Fate}{No exact problem is recoverable from the broad historical comparison. It is kept as a source remark rather than attached to an unrelated later polytope theorem.}
\GSfield{Present status}{As of 11 November 2026: Not identifiable as a specific research problem. This is not evidence that finite-dimensional polytope theory has had no development.}
\GSfield{Assessment}{literal statement: not determinable; publication match: none found; verification: unresolved lead; development: none; confidence: high}
\end{GSitem}
\begin{GSitem}{1987-736.14}{E}{Coordinate variation and the conformal Ward identity}
\GSfield{Kind; source}{method; \GSpage{P001433}}
\GSfield{The notes}{Varying \(x^\mu\) by \(\varepsilon^\mu(x)\) gives a sum of field variations plus \(\int\partial_\mu\varepsilon_\nu\langle T^{\mu\nu}A_1\cdots A_n\rangle\). Integration by parts turns the variation into contours around insertions.}
\GSfield{Statement}{For separated insertions, a stress tensor satisfying the distributional Ward identity gives the stated contour action. If \(\partial_\mu\varepsilon_\nu+\partial_\nu\varepsilon_\mu=g(x)\delta_{\mu\nu}\) and the stress tensor is symmetric and traceless, the bulk contraction vanishes away from insertions. Boundary terms, contact terms and any anomalies must be handled in the chosen theory; the unregularized path integral is only a motivation.}
\GSfield{Reconstruction}{The contact and boundary terms omitted by the formal integration by parts are stated.}
\GSfield{Fate}{\citet[Sections~2--3, especially equation~(3.4)]{belavinpolyakovzamolodchikov1984infinite} supplies the conformal Ward identities before the talk.}
\GSfield{Present status}{As of 11 November 2026: Established local derivation under the symmetry, regularity and boundary hypotheses. It does not establish conformal invariance for an arbitrary action.}
\GSfield{Assessment}{literal statement: correct with missing hypotheses; publication match: exact; verification: read the result; development: none; confidence: high}
\end{GSitem}
\begin{GSitem}{1987-736.15}{E}{Primary fields, descendants and conformal families}
\GSfield{Kind; source}{construction; \GSpage{P001434}}
\GSfield{The notes}{Primaries satisfy \(L_n\phi=\bar L_n\phi=0\) for \(n>0\), \(L_0\phi=h\phi,\bar L_0\phi=\bar h\phi\), with \(h,\bar h\ge0\); the algebra is written as a direct sum of their descendant families.}
\GSfield{Statement}{At the origin a highest-weight state has the displayed annihilation and eigenvalue conditions. Applying negative modes generates its highest-weight module, with null states quotiented out in an irreducible unitary family. In a unitary positive-energy theory \(h,\bar h\ge0\). A direct sum of such families requires an appropriate discrete semisimple spectrum; it is not automatic for every conformal theory or every noncompact model.}
\GSfield{Reconstruction}{Coordinate origin, null-state quotient and the conditions for a direct-sum decomposition are supplied.}
\GSfield{Fate}{\citet[Section~3, equations~(3.8)--(3.13)]{belavinpolyakovzamolodchikov1984infinite} presents primary fields and conformal families. The result predates this lecture.}
\GSfield{Present status}{As of 11 November 2026: Standard within the stated unitary highest-weight framework. The algebraic direct-sum sentence needs spectral/semisimplicity hypotheses, and the source does not pose a new singular-vector classification problem.}
\GSfield{Assessment}{literal statement: correct with missing hypotheses; publication match: exact; verification: read the result; development: none; confidence: high}
\end{GSitem}
\begin{GSitem}{1987-736.16}{E}{Scaling weights, spin and the identity field}
\GSfield{Kind; source}{claim; \GSpage{P001434}}
\GSfield{The notes}{The identity has weights zero. Under \(z\mapsto\lambda z,\bar z\mapsto\bar\lambda z\), the notes assign the primary the factor \(\lambda^h\bar\lambda^{\bar h}\) and call \(h-\bar h\) its spin.}
\GSfield{Statement}{The second coordinate must be \(\bar z\mapsto\bar\lambda\bar z\). For a primary, \(h+\bar h\) is its scaling dimension and \(h-\bar h\) its spin. In a passive coordinate-change convention \(\phi_{\rm new}(\lambda z,\bar\lambda\bar z)=\lambda^{-h}\bar\lambda^{-\bar h}\phi(z,\bar z)\); active conventions reverse this placement. The identity has \((h,\bar h)=(0,0)\).}
\GSfield{Reconstruction}{The missing coordinate conjugation is a literal error; the sign of the transformation factor is a convention that the source leaves unstated.}
\GSfield{Fate}{\citet[Introduction, equations~(1.16)--(1.17); Section~3]{belavinpolyakovzamolodchikov1984infinite} gives these transformation laws and weights.}
\GSfield{Present status}{As of 11 November 2026: Established after correcting the missing bar and specifying the active/passive convention. No sign correction of the weight factor is imposed without that convention.}
\GSfield{Assessment}{literal statement: false as written; publication match: exact; verification: read the result; development: none; confidence: high}
\end{GSitem}

\GStalk{737}{1988-02-22}{I. M. Gelfand; notes by A. K. Zvonkin}{P001435-P001437}{Complete invariants called angles: a problem sheet}
The handout defines an angle to mean a complete orbit invariant, including its attainable values. It passes from Euclidean pairs and configurations under \(GL\), through indefinite quadratic spaces, to a tensor construction and closed Hilbert subspaces. The finite Euclidean and Hilbert questions were already classical; finite indefinite pairs have a later precise canonical-form treatment. An indefinite infinite-dimensional version is too underspecified to classify as an open problem.

\begin{GSitem}{1987-737.1}{E}{An angle invariant for a line and a plane in three dimensions}
\GSfield{Kind; source}{problem; \GSpage{P001435}}
\GSfield{The notes}{The preamble calls an angle a complete orbit invariant and asks for its attainable values. Problem 1 asks for the angle between a line and a plane in three-dimensional space, explicitly separating intersecting and nonintersecting cases.}
\GSfield{Statement}{For an affine line \(\ell\) and an affine plane H under all Euclidean isometries, the direction angle \(\alpha\in[0,\pi/2]\), together with their distance \(\delta\ge0\), is complete, with the restriction \(\alpha>0\Rightarrow\delta=0\). When \(\alpha=0\), delta distinguishes a contained line from a disjoint parallel line. For linear subspaces only \(\delta=0\) occurs. Under origin-preserving orthogonal maps alone, additional positional data are needed for affine objects.}
\GSfield{Reconstruction}{Allowing translations is an explicit interpretation of the nonintersecting case, not an unannounced change from the preamble about orthogonal maps.}
\GSfield{Fate}{The directional angle is the elementary case of the singular-value characterization in \citet[Section~3, Theorem~1]{bjorckgolub1973angles}. The affine distance correction follows directly by moving the closest point or intersection to the origin.}
\GSfield{Present status}{As of 11 November 2026: Classical and solved after fixing the acting group. An angle alone is not a complete invariant of disjoint affine configurations.}
\GSfield{Assessment}{literal statement: correct with missing hypotheses; publication match: partial; verification: read the result; development: none; confidence: high}
\end{GSitem}
\begin{GSitem}{1987-737.2}{E}{An angle invariant for two planes in three dimensions}
\GSfield{Kind; source}{problem; \GSpage{P001435}}
\GSfield{The notes}{Problem 2 asks for the angle between two planes in three-dimensional space.}
\GSfield{Statement}{For linear two-planes U,W in \(\mathbb R^3\), their two principal angles are \(0,\theta\), with \(0\le\theta\le\pi/2\); theta is the acute angle between their normal lines and is a complete O(3) invariant. If affine planes are intended, distinct parallel planes additionally require their distance \(\delta>0\); nonparallel planes intersect, and their sole nonzero directional angle suffices under Euclidean isometries.}
\GSfield{Reconstruction}{The linear and affine interpretations are separated.}
\GSfield{Fate}{The linear statement is a dimension-three case of \citet[Section~3, Theorem~1]{bjorckgolub1973angles}. The affine supplement is elementary and is needed because the immediately preceding task includes disjoint affine objects.}
\GSfield{Present status}{As of 11 November 2026: Solved classical exercise. The source should not be read as claiming that one scalar classifies every origin-dependent affine placement.}
\GSfield{Assessment}{literal statement: correct with missing hypotheses; publication match: partial; verification: read the result; development: none; confidence: high}
\end{GSitem}
\begin{GSitem}{1987-737.3}{E}{Two principal angles classify two planes in four dimensions}
\GSfield{Kind; source}{problem; \GSpage{P001435}}
\GSfield{The notes}{Problem 3 repeats the question about two two-dimensional subspaces of a four-dimensional inner-product space, without choosing coordinate axes.}
\GSfield{Statement}{The complete O(4) invariant is \(0\le\theta_1\le\theta_2\le\pi/2\), with cosines the singular values of the cross-Gram matrix of orthonormal frames. Every such pair is realized by \(U=\operatorname{span}(e_1,e_2)\) and \(W=\operatorname{span}(\cos\theta_1e_1+\sin\theta_1e_3,\cos\theta_2e_2+\sin\theta_2e_4)\).}
\GSfield{Fate}{\citet[Section~3, Theorem~1]{bjorckgolub1973angles} gives the established singular-value formulation. This is the exact recurrence of 1987-735.2, with the range request supplied by the problem-sheet preamble.}
\GSfield{Present status}{As of 11 November 2026: Solved, including the full range of attainable invariants. The treatment for SO(4) requires the separate orientation discussion below.}
\GSfield{Assessment}{literal statement: correct; publication match: exact; verification: read the result; development: none; confidence: high; problem \hyperref[prob:GP-1987-28]{GP-1987-28}}
\end{GSitem}
\begin{GSitem}{1987-737.4}{E}{A vector and a subspace require magnitude as well as direction}
\GSfield{Kind; source}{problem; \GSpage{P001435}}
\GSfield{The notes}{Problem 4 gives \(v\in L\), a k-dimensional subspace V, and an inner product on the n-dimensional space L; it asks for their angle.}
\GSfield{Statement}{Under O(L), a pair (v,V) is classified by \(n,k,\|P_Vv\|,\|P_{V^\perp}v\|\). For \(v\ne0\), these are equivalently the norm of v and the angle \(\alpha\) defined by \(\cos\alpha=\|P_Vv\|/\|v\|\). Every pair of nonnegative component norms is possible when \(0<k<n\); if k=0 the first is zero, and if k=n the second is zero. The zero vector is a separate orbit with no direction angle.}
\GSfield{Reconstruction}{The source does not require a unit vector; its length is therefore retained.}
\GSfield{Fate}{The projection formula is the rank-one case of \citet[Section~3]{bjorckgolub1973angles}; completeness follows by independently rotating the two orthogonal components.}
\GSfield{Present status}{As of 11 November 2026: Solved. The problem-sheet definition of a complete invariant makes normalization material: angle alone does not distinguish vectors of different lengths.}
\GSfield{Assessment}{literal statement: correct with missing hypotheses; publication match: partial; verification: read the result; development: none; confidence: high}
\end{GSitem}
\begin{GSitem}{1987-737.5}{E}{Finite-dimensional pairs: complete invariants and their exact range}
\GSfield{Kind; source}{problem; \GSpage{P001436}}
\GSfield{The notes}{Problem 5 asks for the angle between subspaces V,W of dimensions k,m in an n-dimensional inner-product space.}
\GSfield{Statement}{Let \(r=\min(k,m)\) and \(q=\max(0,k+m-n)\). The ordered principal angles \(0\le\theta_1\le\cdots\le\theta_r\le\pi/2\), with at least q zero entries, classify ordered pairs under O(n), or under U(n) in the complex version. Every list satisfying exactly these restrictions is attainable. The dimensions of the ambient space and of both subspaces are also part of the classification when they are not fixed.}
\GSfield{Reconstruction}{The exact dimension-forced zero-angle condition supplies the range requested in the preamble.}
\GSfield{Fate}{The cross-Gram SVD in \citet[Section~3, Theorem~1]{bjorckgolub1973angles} gives the principal angles. Orthogonal principal two-planes plus intersection and unmatched directions give both completeness and the stated dimension restriction.}
\GSfield{Present status}{As of 11 November 2026: Solved classical finite-dimensional problem, including attainability. Replacing the list by its minimum, maximum or average angle discards information.}
\GSfield{Assessment}{literal statement: correct; publication match: exact; verification: read the result; development: none; confidence: high; problem \hyperref[prob:GP-1987-28]{GP-1987-28}}
\end{GSitem}
\begin{GSitem}{1987-737.6}{E}{One, two and three subspaces under the general linear group}
\GSfield{Kind; source}{problem; \GSpage{P001436}}
\GSfield{The notes}{Problem 6 removes the inner product and asks for complete invariants of one subspace, a pair and an ordered triple, with GL(n) as the acting group.}
\GSfield{Statement}{For one subspace the dimension suffices. For a pair, the dimensions of U,V and \(U\cap V\), together with n, suffice. Over \(\mathbb R\) or \(\mathbb C\), an ordered triple decomposes uniquely into eight one-dimensional membership types indexed by subsets of \(\{1,2,3\}\), and copies of the two-dimensional configuration of three distinct lines \(\operatorname{span}e_1,\operatorname{span}e_2,\operatorname{span}(e_1+e_2)\). The nine multiplicities are complete invariants. All nonnegative multiplicities occur, subject to the ambient dimension \(\sum_Sa_S+2b=n\) and any prescribed subspace dimensions.}
\GSfield{Reconstruction}{The source omits the base field. The stated elementary decomposition is for real or complex vector spaces; its precise multiplicities answer the completeness and range requests.}
\GSfield{Fate}{The introduction of \citet[Section~1]{gelfandponomarev1971quadruples} separates the easy r=2,3 cases from the four-subspace problem. The exact nine-type classification already appears in the prior index at 1972-277.5--6; it is established finite-type linear algebra, not a new 1988 conjecture.}
\GSfield{Present status}{As of 11 November 2026: Solved, with a discrete finite list of building-block types and arbitrary admissible multiplicities. Pairwise intersection dimensions alone do not capture the two-dimensional triple block.}
\GSfield{Assessment}{literal statement: correct; publication match: exact; verification: read the result; development: none; confidence: high; problem \hyperref[prob:GP-1987-35]{GP-1987-35}}
\end{GSitem}
\begin{GSitem}{1987-737.7}{E}{Two lines in the Minkowski plane}
\GSfield{Kind; source}{problem; \GSpage{P001436}}
\GSfield{The notes}{Problem 7 fixes the quadratic form \(-x_1^2+x_2^2\) and asks for complete invariants of two lines under its preserving group.}
\GSfield{Statement}{For coincident lines the causal type (negative, positive, null) is complete. For two distinct unoriented lines choose nonzero spanning vectors u,v and form \(G=\left(\begin{smallmatrix}\langle u,u\rangle&\langle u,v\rangle\\\langle u,v\rangle&\langle v,v\rangle\end{smallmatrix}\right)\). Its equivalence class under \(G\mapsto DGD\), with invertible real diagonal D, is a complete O(1,1) invariant. The attainable matrices are precisely the symmetric two-by-two matrices of signature (1,1), equivalently determinant less than zero. This formulation includes null lines without division by their zero norms.}
\GSfield{Reconstruction}{The diagonal rescaling identifies lines rather than chosen vectors; no orientation or time-orientation constraint is imposed.}
\GSfield{Fate}{This two-dimensional answer follows directly by sending the basis (u,v) to another basis with the same Gram matrix. It is a simple case of the broader form-preserving classification later treated by \citet[Introduction and Theorem~1]{sergeichuk1990pairs}; the elementary case was already available before the talk.}
\GSfield{Present status}{As of 11 November 2026: Solved. An ordinary cosine alone cannot cover spacelike, timelike and null cases.}
\GSfield{Assessment}{literal statement: correct; publication match: exact; verification: read the result; development: none; confidence: high; problem \hyperref[prob:GP-1987-31]{GP-1987-31}}
\end{GSitem}
\begin{GSitem}{1987-737.8}{B}{Pairs in finite-dimensional indefinite and degenerate forms}
\GSfield{Kind; source}{problem; \GSpage{P001436}}
\GSfield{The notes}{Problem 8 extends the pair problem to higher Minkowski dimensions, arbitrary indefinite signature, and degenerate forms such as \(x_1^2+x_2^2\) on \(\mathbb R^3\).}
\GSfield{Statement}{Fix a finite-dimensional vector space over a field of characteristic different from two with a symmetric, skew-symmetric or Hermitian form F, allowing degeneracy, and two labelled subspaces. Classify triples \((V_F;U_1,U_2)\) up to form-preserving linear isomorphisms carrying each labelled subspace to its counterpart. Sergeichuk gives a complete canonical decomposition for this problem, up to the classification of the Hermitian forms over the relevant fields; for the real symmetric setting these residual forms are determined by their inertia. Isometry blocks, their multiplicities and their allowed field parameters are the complete data, rather than only ordinary principal-angle numbers.}
\GSfield{Reconstruction}{The field, form type, labelled nature and isometry category are fixed explicitly. The notebook asks only the real symmetric cases; the cited theorem treats a broader family.}
\GSfield{Fate}{\citet[Theorem~1]{sergeichuk1990pairs} supplies a precise later classification covering the finite-dimensional indefinite and degenerate formulations. The paper was received on 13 September 1988 and published in 1990; its introduction mentions a partial earlier announcement, so this date is not asserted as the first discovery of every component.}
\GSfield{Present status}{As of 11 November 2026: The finite-dimensional problem has a canonical-form solution. The opened result does not settle the subsequent infinite-dimensional indefinite version, and no such extrapolation is made.}
\GSfield{Assessment}{literal statement: correct; publication match: exact; verification: read the result; development: none; confidence: high; problem \hyperref[prob:GP-1987-31]{GP-1987-31}}
\end{GSitem}
\begin{GSitem}{1987-737.9}{E}{Passing from the full orthogonal group to the special group}
\GSfield{Kind; source}{problem; \GSpage{P001436}}
\GSfield{The notes}{After the indefinite variants the sheet asks that all listed problems be considered for both O(n) and SO(n).}
\GSfield{Statement}{For a nondegenerate real quadratic space, an O-orbit splits into two SO-orbits exactly when the stabilizer has no determinant-minus-one element; this follows because SO has index two. In the Euclidean ordered-pair case, it splits exactly when all four intersections \(U\cap W,U\cap W^\perp,U^\perp\cap W,U^\perp\cap W^\perp\) vanish. Equivalently, \(n=2k=2m\) and every principal angle lies strictly between zero and \(\pi/2\). Otherwise a reflection in one of those intersections keeps the orbit unsplit. For a degenerate form the full preserving group can have arbitrary nonzero determinant on its radical, so this index-two argument cannot simply be transferred.}
\GSfield{Reconstruction}{The stabilizer criterion and the degeneracy caveat are supplied; the latter prevents an incorrect universal two-orbit claim.}
\GSfield{Fate}{The criterion is an elementary index-two orbit calculation, with the Euclidean stabilizer read from the principal-angle blocks of \citet[Section~3]{bjorckgolub1973angles}. The invariant is augmented by a binary orientation choice exactly in the split case.}
\GSfield{Present status}{As of 11 November 2026: Solved for the stated nondegenerate groups. The notation O(n),SO(n) in the degenerate extension needs a definition and cannot be used to infer an index-two subgroup there.}
\GSfield{Assessment}{literal statement: correct with missing hypotheses; publication match: partial; verification: read the result; development: annotated; confidence: high}
\end{GSitem}
\begin{GSitem}{1987-737.10}{E}{Four subspaces under GL: the Gelfand--Ponomarev classification}
\GSfield{Kind; source}{problem; \GSpage{P001436}}
\GSfield{The notes}{Problem 9 asks for complete invariants of a quadruple of subspaces and explicitly calls it the content of the work of Gelfand and Ponomarev.}
\GSfield{Statement}{For an ordered quadruple in a finite-dimensional vector space over an algebraically closed field of characteristic zero, decomposition into indecomposable quadruples reduces the classification to the discrete families and one-parameter families listed by Gelfand--Ponomarev. With defect \(d=\sum_i\dim U_i-2\dim V\), indecomposable defects are \(-2,-1,0,1,2\); the continuous regular families have defect zero. Their block multiplicities and parameters give complete invariants. Real versions require real canonical blocks rather than silently assuming every eigenvalue is real.}
\GSfield{Reconstruction}{The algebraically closed characteristic-zero hypothesis is the one actually read in the original paper. The labelled order of the four subspaces is retained.}
\GSfield{Fate}{\citet[Section~2, Assertions~1--2; Section~3, Theorem~1]{gelfandponomarev1971quadruples} gives the explicit list and decomposition. It was received on 21 October 1970 and published in 1971. Earlier occurrences are 1970-198.1--13 and 1972-277.7.}
\GSfield{Present status}{As of 11 November 2026: Solved before the talk. The source explicitly attributes existing work; it is not an open classification request in the historical sense.}
\GSfield{Assessment}{literal statement: correct with missing hypotheses; publication match: exact; verification: read the result; development: none; confidence: high; problem \hyperref[prob:GP-1987-36]{GP-1987-36}}
\end{GSitem}
\begin{GSitem}{1987-737.11}{E}{Principal angles of the two tensor subspaces}
\GSfield{Kind; source}{problem; \GSpage{P001436}}
\GSfield{The notes}{Problem 10, attributed to A. A. Kirillov junior, gives an inner-product space L, tensors \(a,b\in L\otimes L\), and asks for the angle between \(a\otimes L\) and \(L\otimes b\) in \(L^{\otimes3}\). The preamble explicitly asks which invariant values are attainable.}
\GSfield{Statement}{For a finite-dimensional real or complex Hilbert space of dimension d and nonzero tensors a,b, write their coefficient matrices as A,B in one orthonormal basis. The d principal cosines are the singular values of \((\overline A B)^T/(\|A\|_F\|B\|_F)\). If d is at least two, the complete range is all ordered \(s_1\ge\cdots\ge s_d\ge0\) with \(\sum_js_j\le1\). If d=1 the only cosine is 1. Zero tensors change the subspace dimensions and are excluded from this d-angle statement.}
\GSfield{Reconstruction}{Finite dimension, nonzero tensors and normalization are explicit additions. The real source case is covered; the complex Hilbert variant is also stated. The complete range proof is the chapter annotation, not a result claimed to be in the 1973 paper.}
\GSfield{Fate}{The singular-value interpretation is an application of the established principal-angle theorem of \citet[Section~3, Theorem~1]{bjorckgolub1973angles}. The chapter annotation at Section~\ref{dev:tensor-angles} proves the normalized cross-Gram formula and the exact range directly. It is an explicit solution of this finite-dimensional exercise, with no claim of novelty.}
\GSfield{Present status}{As of 11 November 2026: Solved in the stated finite-dimensional reconstruction, including attainable values. The source only says inner-product space; finite dimension is inferred from its placement before the separate Hilbert-space section, and is recorded rather than hidden.}
\GSfield{Assessment}{literal statement: correct with missing hypotheses; publication match: partial; verification: read the result; development: annotated; confidence: high; problem \hyperref[prob:GP-1987-29]{GP-1987-29}}
\end{GSitem}
\begin{GSitem}{1987-737.12}{E}{Hilbert pairs, discrete angles and continuous spectral data}
\GSfield{Kind; source}{problem; \GSpage{P001437}}
\GSfield{The notes}{Problem 11 asks for a complete invariant of two Hilbert subspaces and for examples with a countable set or a continuum of angle values.}
\GSfield{Statement}{For closed M,N, use the four intersection dimensions and the generic positive contraction C from 1987-735.3. On a separable Hilbert space its spectral measure class together with multiplicity gives the full continuous invariant. For discrete examples take \(C=\operatorname{diag}(\cos\theta_j)\) with \(0<\theta_j<\pi/2\); for a continuum take \(K=L^2(0,1)\), \(Cf(t)=tf(t)\), and \(N=\{(Cf,Sf):f\in K\}\) in \(K\oplus K\). The latter has continuous spectrum [0,1] for C but no eigenvectors at those spectral values.}
\GSfield{Reconstruction}{The projection formulation requires closed subspaces. Separability is imposed only for the familiar measure-and-multiplicity formulation and the displayed examples.}
\GSfield{Fate}{\citet[Theorems~1--3]{halmos1969subspaces} proves the canonical form and completeness. Its 1969 chronology makes this an established answer when the sheet was written.}
\GSfield{Present status}{As of 11 November 2026: Solved for closed Hilbert subspaces. A countable or continuous set of numbers without spectral type and multiplicity is not a complete orbit invariant.}
\GSfield{Assessment}{literal statement: correct with missing hypotheses; publication match: exact; verification: read the result; development: none; confidence: high; problem \hyperref[prob:GP-1987-32]{GP-1987-32}}
\end{GSitem}
\begin{GSitem}{1987-737.13}{F}{The indefinite Hilbert-space version is not specified}
\GSfield{Kind; source}{problem; \GSpage{P001437}}
\GSfield{The notes}{Problem 12 asks the same question for a Hilbert space with indefinite metric.}
\GSfield{Statement}{To formulate a definite classification one must specify, for example, a bounded nondegenerate fundamental symmetry J, the topology of the subspaces, whether they are regular for the indefinite form, and the equivalence group of bounded J-unitaries. A degenerate or unbounded form leads to different categories. These choices are absent, and the finite-dimensional canonical-block theorem does not answer all of them.}
\GSfield{Reconstruction}{The possible hypotheses illustrate why the source is insufficient; none is selected and attributed to the author.}
\GSfield{Fate}{\citet{halmos1969subspaces} treats positive Hilbert geometry; \citet[Introduction and Theorem~1]{sergeichuk1990pairs} treats finite-dimensional form spaces. Neither identifies a single exact infinite-dimensional assertion from the unspecified notebook wording.}
\GSfield{Present status}{As of 11 November 2026: No definite remaining open problem can be assigned from this statement alone. This entry is unidentifiable in category, not a declaration that all Krein-space pair classifications are unknown.}
\GSfield{Assessment}{literal statement: not determinable; publication match: thematic only; verification: unresolved lead; development: none; confidence: high}
\end{GSitem}
\begin{GSitem}{1987-737.14}{E}{Realize a prescribed invariant relative to a fixed subspace}
\GSfield{Kind; source}{problem; \GSpage{P001437}}
\GSfield{The notes}{The first comment emphasizes the inverse task: given a subspace and an angle invariant, construct another subspace having that invariant.}
\GSfield{Statement}{For a fixed k-plane U in \(\mathbb R^n\), any m-plane principal-angle list satisfying 1987-737.5 is realized by rotating orthogonal principal blocks, using the forced zero-angle intersections to fit the ambient dimension. For a fixed closed Hilbert subspace M, the Halmos data are realizable exactly when the required orthogonal summands fit: \(\dim M=d_{11}+d_{10}+\dim K\) and \(\dim M^\perp=d_{01}+d_{00}+\dim K\). On the generic summands set \(N=\{(Cx,(I-C^2)^{1/2}x):x\in K\}\), and add the chosen intersection summands.}
\GSfield{Reconstruction}{The finite and Hilbert versions are made explicit; no realization claim is made for the undefined indefinite category.}
\GSfield{Fate}{The construction is the inverse of the canonical forms in \citet[Section~3]{bjorckgolub1973angles} and \citet[Theorem~2]{halmos1969subspaces}. Both were established before the seminar.}
\GSfield{Present status}{As of 11 November 2026: Solved for the finite positive and closed Hilbert settings just specified. The cardinal-dimension restrictions are necessary, even when an abstract angle operator has already been chosen.}
\GSfield{Assessment}{literal statement: correct with missing hypotheses; publication match: exact; verification: read the result; development: none; confidence: high; problem \hyperref[prob:GP-1987-28]{GP-1987-28}}
\end{GSitem}
\begin{GSitem}{1987-737.15}{E}{Encode operator similarity by four subspaces}
\GSfield{Kind; source}{construction; \GSpage{P001437}}
\GSfield{The notes}{The final comment says that a linear operator can be specified by a quadruple of subspaces.}
\GSfield{Statement}{For \(T:E\to E\), the ordered quadruple \(E\oplus0,\ 0\oplus E,\ \{(x,x)\},\ \{(x,Tx)\}\) in \(E\oplus E\) determines T up to similarity. Any isomorphism preserving the first two subspaces is block diagonal; preservation of the diagonal forces the two blocks to coincide, and preservation of the last graph is exactly the similarity relation. Conversely any similarity gives such a configuration isomorphism.}
\GSfield{Fate}{\citet[Section~1, Example~1]{gelfandponomarev1971quadruples} explicitly gives the graph construction, with the final two labels interchanged. It is already indexed as 1970-198.3.}
\GSfield{Present status}{As of 11 November 2026: Established exact reduction. The source is an expository reminder rather than a new representation theorem.}
\GSfield{Assessment}{literal statement: correct; publication match: exact; verification: read the result; development: none; confidence: high; problem \hyperref[prob:GP-1987-36]{GP-1987-36}}
\end{GSitem}

\GStalk{738}{1988-03-21}{A. O. Radul}{P001438-P001443}{Quantum Gelfand--Dikii brackets after Lukyanov}
Radul reports quantum Miura constructions, preceded by local variational calculus, formal pseudodifferential symbols and the two Gelfand--Dikii brackets. Gelfand interrupts about intrinsic locality, products of local functionals, and a hierarchy based on differentials. Lukyanov's paper was submitted before this talk and published later in 1988; it proves closure of the quantum algebra. The notebook's informal hierarchy is not thereby solved. The next heading, Kapranov at P001444, belongs to record 739.

\begin{GSitem}{1987-738.1}{E}{Local polynomial functionals and the variational derivative}
\GSfield{Kind; source}{definition; \GSpage{P001438}}
\GSfield{The notes}{The operator is printed as \(L=\partial^N+\sum_{i=0}^{N-1}u_{N-1}\partial^i\); functionals are integrals of polynomials in coefficient jets. Equation (3) defines the variational derivative by first variation.}
\GSfield{Statement}{Fix a differential algebra \(\mathcal A\) of differential polynomials in \(a_0,\ldots,a_{N-1}\), with total derivative \(D\), and write \(L=D^N+\sum_{j=0}^{N-1}a_jD^j\). A formal local functional is a class \(\int P\in\mathcal A/D\mathcal A\). Its Euler derivative is \(\delta F/\delta a_j=\sum_{k\ge0}(-D)^k\,\partial P/\partial a_j^{(k)}\). For smooth periodic coefficients this agrees with integration by parts in \(\delta F=\int\sum_j(\delta F/\delta a_j)\delta a_j\).}
\GSfield{Reconstruction}{The independent coefficients are renamed \(a_j\) to avoid silently mixing the source conventions \(u_j\) and \(u_{N-j}\). Periodicity or the formal quotient removes boundary terms.}
\GSfield{Fate}{This formal calculus is established in \citet[Section~1]{gelfanddikii1978jets} and \citet[Section~2]{gelfanddorfman1979hamiltonian}; both precede the seminar. The repeated coefficient subscript in (1) is incompatible with the later independent variations and Miura formula.}
\GSfield{Present status}{As of 11 November 2026: Standard variational calculus. Local functionals form a vector-space quotient; multiplication of densities does not descend to multiplication of their quotient classes.}
\GSfield{Assessment}{literal statement: false as written; publication match: exact; verification: read the result; development: none; confidence: high}
\end{GSitem}
\begin{GSitem}{1987-738.2}{B}{Disjoint additivity does not imply a smooth local density from continuity alone}
\GSfield{Kind; source}{claim; \GSpage{P001438}; I. M. Gelfand}
\GSfield{The notes}{Gelfand says that increments supported in two disjoint places add independently, and that this property together with continuity forces a finite-jet integral of the displayed type.}
\GSfield{Statement}{On an open subset \(U\subset C^\infty(M)\), the additivity condition is \(F(\phi+\psi+\chi)=F(\phi+\psi)+F(\psi+\chi)-F(\psi)\) when \(\operatorname{supp}\phi\cap\operatorname{supp}\chi=\varnothing\) and all arguments belong to \(U\). Continuity alone is insufficient: \(F(u)=u(x_0)\) is continuous and additive, but its differential is the Dirac distribution, not integration against a smooth density. A precise characterization additionally assumes Bastiani smoothness, each \(DF_u\) a compactly supported smooth density, and smooth dependence of this density on \(u\) in the test-function topology. The conclusion is a finite-jet integral in a neighborhood of each configuration, with order allowed to depend on that neighborhood.}
\GSfield{Reconstruction}{The function space, topology, smooth density requirement and locally bounded jet order are reconstructed from the later theorem; they are not present in the notebook.}
\GSfield{Fate}{\citet[Theorems~I.2 and VI.3]{brouderetal2018functionals}, first circulated 3 May 2017 and published in 2018, proves this characterization with the stated regularity hypotheses. Its locally finite-order conclusion is weaker than a single global bound on jet order.}
\GSfield{Present status}{As of 11 November 2026: The literal continuity-only claim is false. A precise regular version is proved; no general continuity theorem or global finite-order theorem is inferred from it.}
\GSfield{Assessment}{literal statement: false as written; publication match: partial; verification: read the result; development: annotated; confidence: high; problem \hyperref[prob:GP-1987-33]{GP-1987-33}}
\end{GSitem}
\begin{GSitem}{1987-738.3}{E}{Length and area as local geometric functionals}
\GSfield{Kind; source}{construction; \GSpage{P001439}; I. M. Gelfand}
\GSfield{The notes}{Length of a curve and area of a surface are listed as local examples before eigenvalue functionals are called transcendental.}
\GSfield{Statement}{For a smooth immersed curve in a Riemannian manifold, length is \(\int |\dot\gamma(t)|\,dt\). For a smooth immersed surface, area is the integral of the square root of the determinant of the induced metric. These densities depend on first jets and are smooth on the open set of immersions. Disjointly supported perturbations therefore have additive increments, provided the perturbations remain in that open set.}
\GSfield{Reconstruction}{The target metric and nondegeneracy of the immersion are supplied.}
\GSfield{Fate}{These are classical examples of the local variational calculus formalized in \citet[Section~1]{gelfanddikii1978jets}; their smooth, rather than polynomial, jet dependence is the geometric extension explicitly allowed in \citet[Section~2, footnote]{gelfanddorfman1979hamiltonian}.}
\GSfield{Present status}{As of 11 November 2026: Established local examples. Smoothness at a degenerate map is not implied: the immersion hypothesis matters for the square-root densities.}
\GSfield{Assessment}{literal statement: correct with missing hypotheses; publication match: thematic only; verification: read the result; development: none; confidence: high}
\end{GSitem}
\begin{GSitem}{1987-738.4}{F}{A hierarchy of transcendence defined by variational differentiation}
\GSfield{Kind; source}{question; \GSpage{P001439}; I. M. Gelfand}
\GSfield{The notes}{Type zero means local; type one means that \(\delta F/\delta u\) is local; higher types are recursive. The first eigenvalue is a motivation, and the speaker predicts functionals of all types in characteristic-class theory.}
\GSfield{Statement}{The notes do not specify how a function-valued differential, a distributional one-form or a higher variation is to belong to the same class as a scalar functional. Nor do they specify the operator and boundary conditions defining \(\lambda_1\), the treatment of spectral multiplicities, equivalence modulo constants, or the global configuration space. The recursive words therefore do not determine a single classification problem or a theorem that the hierarchy is strict at every level.}
\GSfield{Reconstruction}{No modern differential-cohomology or spectral-invariant theory is silently substituted for the proposed hierarchy.}
\GSfield{Fate}{The named antecedent is identifiable as \citet[Introduction and Section~1]{gabrielovgelfandlosik1976atiyah}, which studies characteristic-class boundary functionals. That concrete paper does not in the sections read formulate the unrestricted recursive hierarchy asserted here; later local-functional characterizations likewise do not settle an unspecified notion of transcendence.}
\GSfield{Present status}{As of 11 November 2026: The general hierarchy remains unidentifiable as a definite theorem from this source. This F classification concerns the missing definition, not an assertion that the named characteristic-class work was never developed.}
\GSfield{Assessment}{literal statement: not determinable; publication match: thematic only; verification: unresolved lead; development: none; confidence: medium}
\end{GSitem}
\begin{GSitem}{1987-738.5}{E}{The identified Atiyah--Patodi--Singer functional antecedent}
\GSfield{Kind; source}{remark; \GSpage{P001439}; I. M. Gelfand}
\GSfield{The notes}{Gelfand recalls work with Gabrielov and an uncertainly written Losik concerning characteristic classes and functionals whose variations have a local character.}
\GSfield{Statement}{The identifiable setting is a compact Riemannian manifold with piecewise smooth boundary and a top-degree characteristic form of the tangent bundle. One seeks face terms such that the characteristic integral plus boundary-face integrals depends on the intrinsic boundary metric and topological data. The referenced paper constructs differential forms on faces and a generalized de Rham cocycle formulation. This is a specific geometric problem, distinct from an unrestricted classification of all scalar functionals by successive variations.}
\GSfield{Reconstruction}{The surname is identified bibliographically as M. V. Losik; the precise geometric setting is supplied by the opened primary paper.}
\GSfield{Fate}{\citet[Introduction and Section~1, Definitions~1.1--1.12]{gabrielovgelfandlosik1976atiyah}, received 22 December 1975 and published in 1976, is the exact author/topic match. Its introduction explicitly separates the proposed boundary-functional structure from what is known in general.}
\GSfield{Present status}{As of 11 November 2026: The historical reference is securely identified and predates the talk. Only the introduction and the foundational definitions were read, so no general boundary decomposition beyond their stated scope is certified.}
\GSfield{Assessment}{literal statement: correct with missing hypotheses; publication match: partial; verification: read the result; development: none; confidence: high}
\end{GSitem}
\begin{GSitem}{1987-738.6}{B}{Variational calculus for sums and products of local functionals}
\GSfield{Kind; source}{problem; \GSpage{P001439-P001440}; I. M. Gelfand}
\GSfield{The notes}{Gelfand observes that local functionals are not an algebra, asks for calculus on sums and products, and writes a double integral involving jets at two points. He rejects treating the issue as only a formal adjunction of products.}
\GSfield{Statement}{On a specified smooth field space, the algebra generated by local scalar functionals consists of finite sums of finite products. Its differential obeys \(D(FG)=G\,DF+F\,DG\); higher derivatives obey the iterated Leibniz rule and have kernels on products of the base. Products are generally nonlocal: for \(F(u)=\int u\), the mixed second derivative of \(F^2\) in disjointly supported directions \(h,k\) is \(2(\int h)(\int k)\), which need not vanish. A general smooth double-integral kernel need not be a finite sum of separated kernels, so the displayed two-point expression can describe a broader class.}
\GSfield{Reconstruction}{Finite products, arbitrary two-point kernels, and the unspecified philosophical demand for new geometry are distinguished explicitly.}
\GSfield{Fate}{The elementary product calculus predates the talk. A precise later framework for smooth local and multilocal functionals, including locality criteria on products of manifolds, is given in \citet[Sections~II and VII.B--C]{brouderetal2018functionals}, circulated in 2017 and published in 2018. It supplies a useful rigorous setting without identifying a unique geometric axiom that the notebook demands.}
\GSfield{Present status}{As of 11 November 2026: The stated finite-product differentiation problem is settled; a chosen topology and a distinction between finite products and general multilocal kernels remain essential. The broader request to change geometry is not treated as one unresolved conjecture.}
\GSfield{Assessment}{literal statement: correct with missing hypotheses; publication match: partial; verification: read the result; development: none; confidence: high; problem \hyperref[prob:GP-1987-37]{GP-1987-37}}
\end{GSitem}
\begin{GSitem}{1987-738.7}{E}{Hamiltonian formalism without multiplication of local functionals}
\GSfield{Kind; source}{method; \GSpage{P001440}; I. M. Gelfand}
\GSfield{The notes}{The proposed way out invokes structures introduced by Gelfand and Dorfman, followed by connections, de Rham cohomology and Yang--Baxter brackets.}
\GSfield{Statement}{One precise construction starts with a Lie algebra \(\mathfrak g\) acting on a module \(M\), with no multiplication assumed on \(M\). Its Chevalley--Eilenberg forms and differential give a Cartan calculus. A Hamiltonian operator \(H\) from suitable one-forms to \(\mathfrak g\), with Lie-closed image and closed induced skew form, defines \(\{f,g\}=(Hdf)g\), satisfying Jacobi. In the variational example, \(M=\mathcal A/D\mathcal A\) and \(\mathfrak g\) consists of evolutionary derivations.}
\GSfield{Reconstruction}{The exact Lie-algebra/module construction is identified; the later programme words are not converted into unsupported claims.}
\GSfield{Fate}{\citet[Section~1, Theorems~1.1 and 1.3; Section~2, Theorem~2.1]{gelfanddorfman1979hamiltonian}, received 16 April 1979, states and proves precisely this module-based construction. The subsequent bare references to connections and Yang--Baxter equations do not specify additional theorems to identify.}
\GSfield{Present status}{As of 11 November 2026: Established framework from 1979, directly relevant to the failure of a density quotient to be an algebra. No nonexistent multiplication on the quotient is needed.}
\GSfield{Assessment}{literal statement: correct; publication match: partial; verification: read the result; development: none; confidence: high}
\end{GSitem}
\begin{GSitem}{1987-738.8}{E}{Formal pseudodifferential symbols and composition}
\GSfield{Kind; source}{construction; \GSpage{P001440}}
\GSfield{The notes}{Equations (4)--(8) introduce symbol composition, a Laurent series bounded above in powers of \(\xi\), its differential part and its residue. The series is misleadingly labelled \(C(z)[[\xi]]\) and its summation indices disagree.}
\GSfield{Statement}{Over a characteristic-zero differential coefficient algebra, formal pseudodifferential symbols have the form \(A=\sum_{i\le N}a_i(z)\xi^i\), with arbitrarily negative but finitely many positive powers. Composition is \(A\circ B=\sum_{k\ge0}(\partial_\xi^k A)(\partial_z^k B)/k!\); every resulting coefficient is a finite sum. Put \(A_+=\sum_{i\ge0}a_i\xi^i\) and \(\operatorname{res}A=a_{-1}\). The inverse symbol \(D^{-1}\) is formal and is not a chosen analytic integration operator with boundary conditions.}
\GSfield{Reconstruction}{The coefficient algebra, upper boundedness, matching dummy index and distinction between formal and analytic inverse derivatives are supplied.}
\GSfield{Fate}{This symbolic calculus is already used in \citet[Sections~2--5]{gelfanddikii1976fractional}; the original 1976 article and its composition convention were opened. It is foundational background for both Hamiltonian brackets.}
\GSfield{Present status}{As of 11 November 2026: Established formal algebra. The appropriate Laurent notation is \(\mathcal A((\xi^{-1}))\), with the specified noncommutative composition, rather than the ordinary power-series ring in positive powers of \(\xi\).}
\GSfield{Assessment}{literal statement: false as written; publication match: exact; verification: read the result; development: none; confidence: high}
\end{GSitem}
\begin{GSitem}{1987-738.9}{E}{Residue trace and the pseudodifferential variational gradient}
\GSfield{Kind; source}{construction; \GSpage{P001440-P001441}}
\GSfield{The notes}{The gradient is \(\sum_{j=0}^{N-1}\xi^{-j-1}\circ\delta f/\delta u_j\), followed by \(\delta f=\oint\operatorname{res}(\delta L\circ\nabla f)\).}
\GSfield{Statement}{With \(a_j\) the coefficient of \(D^j\), set \(\operatorname{Tr}A=\int\operatorname{res}A\), modulo total derivatives or on the circle, and \(X_F=\sum_{j=0}^{N-1}D^{-j-1}\circ(\delta F/\delta a_j)\). Integration by parts gives \(\operatorname{Tr}[A,B]=0\) and \(\delta F=\operatorname{Tr}(\delta L X_F)\). The trace pairing determines the cotangent representative modulo the annihilator of variations of the monic differential operator.}
\GSfield{Reconstruction}{A single coefficient convention and the formal trace quotient are fixed; the residue trace is not confused with the operator trace on a Hilbert space.}
\GSfield{Fate}{The residue-symbol and variational calculus occurs before the talk in \citet{gelfanddikii1976fractional,gelfanddikii1978jets}. The formula is also a direct coefficient check using cyclicity of the integrated residue; no analytic trace-class assertion is involved.}
\GSfield{Present status}{As of 11 November 2026: Exact in formal or periodic variational calculus. The source gradient uses coefficient order \(u_jD^j\), whereas the later Miura line uses \(u_{N-j}D^j\), so the reindexing must be explicit.}
\GSfield{Assessment}{literal statement: correct with missing hypotheses; publication match: exact; verification: read the result; development: none; confidence: high}
\end{GSitem}
\begin{GSitem}{1987-738.10}{E}{The first Gelfand--Dikii bracket}
\GSfield{Kind; source}{theorem; \GSpage{P001441}}
\GSfield{The notes}{Equations (11) and (13) assert that \(\{F,G\}_1=\operatorname{Tr}([L,X_F]X_G)\) gives a Lie bracket on the stated functionals.}
\GSfield{Statement}{For monic scalar differential operators of fixed positive order over the differential-polynomial algebra, use the residue gradient of item 1987-738.9. The map \(J_{1,L}(X)=[L,X]_+\) is Hamiltonian, and the displayed pairing is a skew bracket satisfying Jacobi on local functionals. The negative part of \([L,X_F]\) does not contribute when paired with the negative-order representative \(X_G\). The leading coefficient is fixed; the coefficient of \(D^{N-1}\) is constant along these Hamiltonian vector fields.}
\GSfield{Reconstruction}{The fixed-order space, quotient by total derivatives and harmless projection under the trace are made explicit.}
\GSfield{Fate}{This is established before the talk. \citet[Section~5, Examples~2--3]{kupershmidtwilson1981modifying} states the exact formula and explains its restriction to zero subleading coefficient; \citet[Sections~10--11]{gelfanddikii1976fractional} supplies an earlier formal variational Lie-bracket construction.}
\GSfield{Present status}{As of 11 November 2026: The Jacobi assertion is proved in the specified formal setting; it is not a new open question in this report. Finite products of actual scalar functionals can be included by the Leibniz rule.}
\GSfield{Assessment}{literal statement: correct with missing hypotheses; publication match: exact; verification: read the result; development: none; confidence: high}
\end{GSitem}
\begin{GSitem}{1987-738.11}{E}{The second bracket, compatibility, and the need for reduction}
\GSfield{Kind; source}{theorem; \GSpage{P001441}}
\GSfield{The notes}{Equation (12) gives \(J_{2,L}(X)=L(XL)_+-(LX)_+L\), and equation (13) asserts the Lie property. The subsequent examples put the subleading coefficient equal to zero.}
\GSfield{Statement}{The pairing \(\{F,G\}_2=\operatorname{Tr}(J_{2,L}(X_F)X_G)\) is Poisson on monic scalar differential operators with all lower coefficients variable. It is compatible with the first bracket: every constant linear combination is Poisson. To pass to \(a_{N-1}=0\), the cotangent representative must be corrected by a \(D^{-N}\) term so that \(\operatorname{res}[X,L]=0\); the resulting Hamiltonian vector field is tangent to the constraint. Merely inserting \(a_{N-1}=0\) in the unreduced formula is insufficient.}
\GSfield{Reconstruction}{Compatibility and the tangency correction are supplied from the proved formulation. The source sign is retained, and its missing reduction is not concealed.}
\GSfield{Fate}{\citet[Section~5, Examples~4--6; Theorems~8.6 and 11.2]{kupershmidtwilson1981modifying} gives the bracket with the opposite overall sign, proves the Hamiltonian property and explains the reduction. Its bibliography identifies the 1978 Gelfand--Dikii IPM preprint136 mentioned in the notebook. \citet[Sections~1 and 6]{khovanova1986gelfand} uses precisely the notebook sign.}
\GSfield{Present status}{As of 11 November 2026: Established before March1988. Overall sign changes preserve the Poisson property, but an omitted reduction changes the actual bracket and the Virasoro normalization.}
\GSfield{Assessment}{literal statement: correct with missing hypotheses; publication match: exact; verification: read the result; development: none; confidence: high}
\end{GSitem}
\begin{GSitem}{1987-738.12}{E}{The order-two example is Virasoro, but the printed mode constants disagree}
\GSfield{Kind; source}{claim; \GSpage{P001441}}
\GSfield{The notes}{For \(L=D^2+u\), the notes set \(F_n=\frac12(2\pi i)^{-1}\oint u(z)z^{n+1}dz\) and assert \(\{F_n,F_m\}=(n-m)F_{n+m}+(n^3-n)\delta_{n+m,0}/12\).}
\GSfield{Statement}{Fix \(\operatorname{Tr}=(2\pi i)^{-1}\oint\operatorname{res}\) and exactly the notebook pairing \(\{F,G\}=\operatorname{Tr}((J_2X_F)X_G)\). After reduction, the coefficient operator is \(J_u(f)=-\frac12D^3f-2uDf-(Du)f\). For the displayed \(F_n\), integration by parts gives \[\{F_n,F_m\}=\frac{m-n}{2}F_{n+m}-\frac{n^3-n}{8}\delta_{n+m,0}.\] Reversing the pairing reverses both signs. Suitable rescalings recover a Virasoro bracket, but they cannot be omitted while retaining the stated modes and claiming the printed constants.}
\GSfield{Reconstruction}{The normalized residue, reduction and order of the pairing are all stated so that the correction is reproducible. The ordinary Fourier/Laurent identity is used, not a quantum normal-ordering correction.}
\GSfield{Fate}{\citet[Section~7, Lemma~6 and its preceding formulas]{khovanova1986gelfand} explicitly gives the reduced order-two operator and its coadjoint Virasoro interpretation, published in1986. The mode calculation above follows directly from that operator and was independently checked during this task.}
\GSfield{Present status}{As of 11 November 2026: The structural Virasoro identification is established. Under the explicit fixed convention the source mode formula is false; its numerical constants are not certified by an appeal to the structural identification alone.}
\GSfield{Assessment}{literal statement: false as written; publication match: partial; verification: read the result; development: annotated; confidence: high}
\end{GSitem}
\begin{GSitem}{1987-738.13}{E}{The order-three bracket and the primary spin-three field}
\GSfield{Kind; source}{claim; \GSpage{P001441}}
\GSfield{The notes}{For \(L=D^3+u_2D+u_3\), the notes give a Virasoro bracket for \(u_2\), a mixed coefficient \(2n-m\), and an incomplete \(u_3,u_3\) formula containing a convolution quadratic in \(u_2\).}
\GSfield{Statement}{The reduced order-three second Gelfand--Dikii bracket is the classical \(W_3\) Poisson algebra. For the raw operator coefficients, the spin-three primary is \(w_3=u_3-\frac12Du_2\). Indeed, infinitesimal coordinate covariance gives \(\delta_f u_2=fDu_2+2(Df)u_2+2D^3f\) and \(\delta_f u_3=fDu_3+3(Df)u_3+u_2D^2f+D^4f\), so \(\delta_f w_3=fDw_3+3(Df)w_3\). The pure spin-three mixed mode formula belongs to this corrected field. The self-bracket has quadratic stress-tensor terms, but the notebook ellipsis is not a complete set of defining relations.}
\GSfield{Reconstruction}{The reduction and primary field change are supplied explicitly. No exact coefficient for the unfinished quadratic convolution is invented.}
\GSfield{Fate}{The Virasoro inclusion is treated in \citet[Section~9]{khovanova1986gelfand}. \citet[Sections~1 and 3]{lukyanov1988quantization} relates these classical brackets to their quantum \(W\)-algebra construction. The displayed primary change follows directly from the two coordinate-variation equations and clarifies the raw-coefficient ambiguity.}
\GSfield{Present status}{As of 11 November 2026: The classical algebra was known before the talk. Its identification does not validate the omitted terms or the numerical central coefficients of the notebook. A pure primary mixed bracket for the unchanged raw \(u_3\) is incorrect.}
\GSfield{Assessment}{literal statement: false as written; publication match: partial; verification: read the result; development: none; confidence: high}
\end{GSitem}
\begin{GSitem}{1987-738.14}{E}{The classical Miura map is Poisson}
\GSfield{Kind; source}{theorem; \GSpage{P001442}}
\GSfield{The notes}{The product \(L=(D+v_1)\cdots(D+v_N)\) defines the coefficient fields, and the modes of the \(v_i\) have independent Heisenberg brackets \(h\,n\delta_{ij}\delta_{n+m,0}\).}
\GSfield{Statement}{Take the \(v_i\) as independent differential variables with the constant diagonal first-order Hamiltonian structure. Expanding the ordered product makes every coefficient of \(L\) a differential polynomial in them. The induced bracket on those coefficients is the second Gelfand--Dikii bracket, with the common overall scale and sign fixed by the chosen pairing. Reduction to zero subleading coefficient requires the corresponding constraint \(\sum_i v_i=0\) and the reduced current bracket; independent currents cannot simply retain their unreduced brackets after that constraint.}
\GSfield{Reconstruction}{Formal differential variables, factor order, normalization and the separate constrained-current reduction are specified.}
\GSfield{Fate}{\citet[Theorem~11.5, p.~436; Theorem~8.6 for the reduced case]{kupershmidtwilson1981modifying} proves the constant-current realization. The original cover gives1981, whereas current Springer metadata gives1980; receipt was February1980 and revision18August1980. \citet[equation~(1.5)]{lukyanov1988quantization} recalls the classical theorem before quantizing it.}
\GSfield{Present status}{As of 11 November 2026: Established Poisson realization. It is a nonlinear differential-polynomial map, not a globally invertible change of coefficient functions on every analytic function space.}
\GSfield{Assessment}{literal statement: correct with missing hypotheses; publication match: exact; verification: read the result; development: none; confidence: high}
\end{GSitem}
\begin{GSitem}{1987-738.15}{E}{Known first-order factors give a solution by quadratures}
\GSfield{Kind; source}{method; \GSpage{P001442}; I. M. Gelfand}
\GSfield{The notes}{Gelfand observes that \((D+v_1)(D+v_2)y=f\) can be solved by quadratures and compares differential Galois theory with solving algebraic equations by radicals.}
\GSfield{Statement}{If continuous \(v_1,v_2,f\) are known on an interval, put \(w=(D+v_2)y\) and \(V_i(x)=\int_{x_0}^xv_i(t)\,dt\). Then \(w=e^{-V_1}(C_1+\int_{x_0}^x e^{V_1(t)}f(t)\,dt)\) and \(y=e^{-V_2}(C_2+\int_{x_0}^xe^{V_2(t)}w(t)\,dt)\). These formulas solve the factored equation with two initial constants. The assertion presupposes known factors; it does not supply them over a prescribed differential field or impose global periodic boundary conditions.}
\GSfield{Reconstruction}{The two elementary integrating-factor formulas give the precise content of the remark without attributing a general differential-Galois theorem to it.}
\GSfield{Fate}{This is the classical integrating-factor calculation. The factorization and differential-field viewpoint are also explicit in \citet[Section~9]{kupershmidtwilson1981modifying}. The brief mention of Picard--Vessiot theory is an analogy, not a complete solvability criterion stated in the source.}
\GSfield{Present status}{As of 11 November 2026: Correct under the stated coefficient and factorization assumptions. Factoring a general differential operator can be the difficult step, and quadratures need not lie in the original coefficient field.}
\GSfield{Assessment}{literal statement: correct with missing hypotheses; publication match: thematic only; verification: read the result; development: none; confidence: high}
\end{GSitem}
\begin{GSitem}{1987-738.16}{E}{Nonlinear Poisson relations do not destroy the Lie property}
\GSfield{Kind; source}{claim; \GSpage{P001442}}
\GSfield{The notes}{The notes say that a Lie algebra has become isomorphic to a quadratic algebra that is not a Lie algebra; Khovanova is named, and an unspecific comparison with Dynkin follows.}
\GSfield{Statement}{A Poisson algebra is always a Lie algebra under its bracket, even when brackets of selected generators are quadratic polynomials rather than linear combinations of those generators. Thus the span of the chosen coefficient fields may fail to be Lie-closed, while the full Poisson algebra still satisfies Jacobi. The Miura map is a Poisson realization, and the existence of a nonlinear formula does not make it an isomorphism onto the entire current algebra.}
\GSfield{Reconstruction}{The meanings of generator span, Poisson algebra and realization are separated; no specific Dynkin theorem is guessed.}
\GSfield{Fate}{\citet[Section~1, Lemma~1; Section~7]{khovanova1986gelfand} explicitly constructs Gelfand--Dikii Lie algebras, directly contradicting the literal denial of a Lie structure. \citet[Theorem~11.5]{kupershmidtwilson1981modifying} states the Miura restriction theorem, not a universal isomorphism. The Dynkin remark names no definite algebra or theorem.}
\GSfield{Present status}{As of 11 November 2026: The literal formulation is false. The meaningful distinction is between a linearly closed generating space and nonlinear Poisson relations in the larger algebra.}
\GSfield{Assessment}{literal statement: false as written; publication match: partial; verification: read the result; development: none; confidence: high}
\end{GSitem}
\begin{GSitem}{1987-738.17}{B}{Quantum Miura coefficients from Heisenberg currents}
\GSfield{Kind; source}{construction; \GSpage{P001442-P001443}}
\GSfield{The notes}{The factors are normal ordered, \(L(z)=:(D+v_1(z))\cdots(D+v_N(z)):\), and the question is which relations the coefficient fields satisfy.}
\GSfield{Statement}{Work with formal radially ordered Heisenberg currents whose singular products are \(v_i(z)v_j(w)\sim h\delta_{ij}(z-w)^{-2}\). Choose the usual mode normal ordering and expand \(L=D^N+\sum_{j=1}^N W_jD^{N-j}\). This defines coefficient fields as normally ordered differential polynomials. The problem is to compute their singular operator products and prove closure, rather than to multiply arbitrary unregularized operator-valued distributions at a point.}
\GSfield{Reconstruction}{The formal mode algebra, radial order and normal-ordering convention supply the operator meaning suppressed by the notation.}
\GSfield{Fate}{\citet[Section~1, equations~(1.1)--(1.5)]{lukyanov1988quantization} makes precisely this construction. The paper was received25February1988, before the21March report, and published later in1988; it acknowledges Radul. The theorem belongs to Lukyanov, whom Radul explicitly reports.}
\GSfield{Present status}{As of 11 November 2026: The construction is established. The subsequent items identify the regularity equation and closure theorem answering the stated relation problem.}
\GSfield{Assessment}{literal statement: correct with missing hypotheses; publication match: exact; verification: read the result; development: none; confidence: high; problem \hyperref[prob:GP-1987-34]{GP-1987-34}}
\end{GSitem}
\begin{GSitem}{1987-738.18}{B}{Lukyanov conjugated regularity relation}
\GSfield{Kind; source}{theorem; \GSpage{P001443}}
\GSfield{The notes}{With \(\sigma f(z,w)=f(w,z)\), the notes set \(\tau=\frac{\nu}{2}\log(z-w)^2\,\sigma\), \(\nu(\nu-1)=h\), and assert \(e^{-\tau}L(z)L(w)e^\tau[f]=O(1)\) for every holomorphic test function. The final symbol is marked uncertain between O and zero.}
\GSfield{Statement}{In the formal differential-operator algebra, define \(\sigma f(z,w)=f(w,z)\) and \(\tau=\frac{\nu}{2}\log((z-w)^2)\sigma\). Let \(f\) be holomorphic on the product of punctured coordinate planes, including the diagonal. For the quantum Miura operator and \(\nu(\nu-1)=h\), the conjugated expression has no singular part as \(z\to w\): \(e^{-\tau}L(z)L(w)e^\tau[f]=O(1)\). The logarithm is used locally with a compatible branch in the formal calculation. The conclusion is regularity, not that the expression vanishes identically.}
\GSfield{Reconstruction}{The formal differential-operator action and regularity meaning are supplied from the original paper. The published formula identifies the intended O(1), without claiming the handwriting alone is unambiguous.}
\GSfield{Fate}{\citet[Section~2.1, Proposition~2 and equations~(2.11)--(2.17)]{lukyanov1988quantization} is an exact match, including the parameter relation and the conjugation formula (2.15). The published formula resolves the notebook ambiguity in favour of \(O(1)\).}
\GSfield{Present status}{As of 11 November 2026: Proved in the original1988 paper. This exact primary match removes the need to classify the formula as an unidentified exchange relation.}
\GSfield{Assessment}{literal statement: correct with missing hypotheses; publication match: exact; verification: read the result; development: none; confidence: high; problem \hyperref[prob:GP-1987-34]{GP-1987-34}}
\end{GSitem}
\begin{GSitem}{1987-738.19}{B}{Quadratic closure and explicit hypergeometric coefficients}
\GSfield{Kind; source}{theorem; \GSpage{P001443}}
\GSfield{The notes}{Equation (24) gives a sum of singular powers with coefficient products, but the bracket signs are ambiguous; the constants are said to involve hypergeometric functions.}
\GSfield{Statement}{The quantum Miura coefficient fields satisfy quadratic singular-product relations of the form \[W_i(z)W_j(w)+\sum_{k=2}^{i+j}\frac{\sum_{p+q=i+j-k}C_{ij}^{pq}(k)\,W_p(z)W_q(w)}{(z-w)^k}=O(1),\] with \(W_0=1\) and numerical coefficients depending on the quantum parameter. Their explicit formulas are obtained from the auxiliary coefficients expressed through generalized hypergeometric functions. This is a statement about ordered operator products and their singular parts; blindly interpreting every square bracket in the notebook as a commutator does not give the published theorem.}
\GSfield{Reconstruction}{The ordered-product interpretation and finite singular sum are supplied from the publication; the unverified handwritten sign choices are not silently certified.}
\GSfield{Fate}{\citet[Theorem~1, equation~(1.6), and equations~(2.18)--(2.21)]{lukyanov1988quantization} proves this closure and gives the coefficient prescription. Theorem 2 gives the associated reduced \(W_N\) construction after the current constraint is imposed with its appropriate covariance.}
\GSfield{Present status}{As of 11 November 2026: The definite relation problem is solved in the identified1988 work. The source formula records the correct quadratic structure, but its ambiguous bracket notation is not a fully reliable literal presentation.}
\GSfield{Assessment}{literal statement: correct with missing hypotheses; publication match: partial; verification: read the result; development: none; confidence: high; problem \hyperref[prob:GP-1987-34]{GP-1987-34}}
\end{GSitem}
\begin{GSitem}{1987-738.20}{F}{A schematic quantum inverse-scattering exchange analogy}
\GSfield{Kind; source}{remark; \GSpage{P001443}}
\GSfield{The notes}{The final sentence compares the regularity formula with the quantum inverse-scattering method and writes \(L(z)L(w)=R^{-1}L(w)L(z)R\), without defining \(R\).}
\GSfield{Statement}{No auxiliary space, operator \(R\), spectral-parameter dependence, tensor-leg convention or equality modulo regular terms is specified. The written identity therefore does not determine a separate Yang--Baxter or exchange theorem. In particular it cannot replace the already identified conjugated regularity relation by an exact unrestricted operator equality.}
\GSfield{Fate}{\citet[Section~2.1, after equation~(2.6)]{lukyanov1988quantization} explicitly remarks on an analogy with a classical r-matrix, while its proved statement is the regularity relation identified above. No precise additional R-matrix theorem matching the final bare formula was established by the source check.}
\GSfield{Present status}{As of 11 November 2026: Unidentifiable as a distinct exact theorem; retained as an analogy. The surrounding quantum closure result is known and is not assigned open status because this final shorthand lacks a definition.}
\GSfield{Assessment}{literal statement: not determinable; publication match: thematic only; verification: unresolved lead; development: none; confidence: medium}
\end{GSitem}

\GStalk{739}{1988-03-21}{M. M. Kapranov}{P001444-P001445}{Quivers, reflection functors, and hearts of derived categories}
Kapranov reports Parthasarathy on derived reflection functors and on abelian hearts with equivalent bounded derived categories. The notes contrast an ordinary cokernel with the complex retaining its kernel. The final three-versus-four distinction does not identify the exact heart or algebra. The catalogue supplies the date; the local heading is undated. P001442--P001443 continue Radul and do not belong to this talk.

\begin{GSitem}{1987-739.1}{E}{Representations of an acyclic quiver}
\GSfield{Kind; source}{definition; \GSpage{P001444}}
\GSfield{The notes}{A quiver is called a directed graph without cycles; \(M(\Lambda,\Omega)\) is its representation category.}
\GSfield{Statement}{For a finite quiver \(Q=(Q_0,Q_1)\) with no oriented cycles over a field \(k\), let \(\operatorname{rep}_kQ\) consist of finite-dimensional vertex spaces and linear arrow maps. It is an abelian hereditary category, equivalent to finite-dimensional modules over \(kQ\).}
\GSfield{Reconstruction}{Finiteness, the coefficient field, and oriented rather than undirected cycles are made explicit.}
\GSfield{Fate}{This is the framework of \citet[\S1]{bernsteingelfandponomarev1973coxeter}; the derived-category report uses it in \citet[\S1]{parthasarathy1988tstructures}.}
\GSfield{Present status}{As of 11 November 2026: established under the stated hypotheses.}
\GSfield{Assessment}{literal statement: correct with missing hypotheses; publication match: exact; verification: read the result; development: none; confidence: high}
\end{GSitem}
\begin{GSitem}{1987-739.2}{E}{Orientation independence and its graph hypothesis}
\GSfield{Kind; source}{claim; \GSpage{P001444}}
\GSfield{The notes}{The lemma says \(\mathcal D(\Lambda,\Omega_1)\simeq\mathcal D(\Lambda,\Omega_2)\) for different orientations, with reflections as proof.}
\GSfield{Statement}{A reflection at a source or sink induces an equivalence of bounded derived categories. Hence orientations joined by such reflections have equivalent categories; in particular this holds for any two orientations of a finite tree, including every Dynkin tree.}
\GSfield{Reconstruction}{The tree or reflection-equivalence hypothesis is essential; the drawing alone does not establish it.}
\GSfield{Fate}{\citet[Propositions~1.4 and~1.6]{parthasarathy1988tstructures} constructs the derived reflection equivalences from BGP functors. An unrestricted assertion for all acyclic orientations of an arbitrary graph is not supplied by this argument. Derived equivalence under tilting was already proved in \citet[\S1.7, Corollary, p.~344]{happel1987derived}, published in 1987 after an Ottawa 1984 announcement; applied to a source/sink reflection, it supplies the pre-talk derived-category formulation.}
\GSfield{Present status}{As of 11 November 2026: derived reflection invariance is established; the notebook' s unrestricted orientation wording should not be used as a theorem.}
\GSfield{Assessment}{literal statement: correct with missing hypotheses; publication match: partial; verification: read the result; development: annotated; confidence: high}
\end{GSitem}
\begin{GSitem}{1987-739.3}{B}{Hearts with the same derived category}
\GSfield{Kind; source}{problem; \GSpage{P001444-P001445}}
\GSfield{The notes}{Find abelian subcategories \(\mathcal A\subset\mathcal D\) with \(\mathcal D^\ell(\mathcal A)=\mathcal D\); the notes say all are found and characterized by Dynkin graphs.}
\GSfield{Statement}{Let k be an algebraically closed field and Q a finite acyclic quiver. Describe bounded t-structures on \(D^b(\operatorname{rep}_kQ)\) whose heart \(\mathcal A\) is a finite-length abelian k-category with a projective generator and finite-dimensional Hom spaces, and whose realization functor \(D^b(\mathcal A)\to D^b(\operatorname{rep}_kQ)\) is an equivalence.}
\GSfield{Reconstruction}{The superscript is read as bounded derived category, and abelian subcategory is interpreted as a heart. The stated finiteness conditions prevent a claim about arbitrary embedded abelian categories.}
\GSfield{Fate}{\citet[Theorems~3.3, 7.1 and~8.5]{parthasarathy1988tstructures} gives a criterion for the specified t-structure data, relates the constructed hearts to iterated tilted algebras, and proves completeness in the Dynkin case. It was received 15 February 1988, revised 13 May, and published in December 1988 by Parthasarathy, whose work Kapranov reports here; receipt alone does not establish public pre-talk circulation of the criterion. The Dynkin algebra classification itself was already \citet[Theorem~5.12, p.~358]{happel1987derived}, published in 1987 after an Ottawa 1984 announcement. Thus the later publication match is the explicit heart-realization approach, not a claim that the Dynkin classification first appeared in 1988. The notebook blanket sentence is broader than either precise formulation.}
\GSfield{Present status}{As of 11 November 2026: the identified finite-type/iterated-tilting classification is established; an unrestricted classification is not asserted.}
\GSfield{Assessment}{literal statement: correct with missing hypotheses; publication match: partial; verification: read the result; development: none; confidence: medium; problem \hyperref[prob:GP-1987-41]{GP-1987-41}}
\end{GSitem}
\begin{GSitem}{1987-739.4}{E}{A source reflection retains a two-term complex}
\GSfield{Kind; source}{construction; \GSpage{P001444-P001445}}
\GSfield{The notes}{At a source \(\alpha\), replace \(V_\alpha\) by the cokernel of \(V_\alpha\to\bigoplus_{\alpha\to\beta}V_\beta\); vectors killed by every arrow are lost, so retain the map as a complex.}
\GSfield{Statement}{The ordinary reflection uses that cokernel and reverses the incident arrows. Its total derived functor uses the two-term complex, with a degree convention fixed once, and preserves the kernel information. This gives the source reflection equivalence.}
\GSfield{Reconstruction}{The vertex is required to be a source and the cone shift is conventional.}
\GSfield{Fate}{BGP's source/sink functors appear in \citet[\S1]{bernsteingelfandponomarev1973coxeter}; the derived construction is explicit in \citet[\S1, Proposition~1.6]{parthasarathy1988tstructures}. Derived equivalence under tilting was already proved in \citet[\S1.7, Corollary, p.~344]{happel1987derived}, published in 1987 after an Ottawa 1984 announcement; applied to a source/sink reflection, it supplies the pre-talk derived-category formulation.}
\GSfield{Present status}{As of 11 November 2026: established under the stated hypotheses.}
\GSfield{Assessment}{literal statement: correct with missing hypotheses; publication match: exact; verification: read the result; development: none; confidence: high}
\end{GSitem}
\begin{GSitem}{1987-739.5}{E}{Derived categories as a language for projective bundles}
\GSfield{Kind; source}{remark; \GSpage{P001444}; I. M. Gelfand}
\GSfield{The notes}{Gelfand says classification of vector bundles on \(\mathbb{CP}^n\) can only be done in the language of derived categories.}
\GSfield{Statement}{An established precise tool is the exceptional collection
\[
 (\mathcal O,\mathcal O(1),\ldots,\mathcal O(n))
\]
on \(\mathbb P^n\) and the resulting resolution of coherent sheaves by explicit complexes. This does not prove the literal exclusivity of a mathematical language or classify all bundles.}
\GSfield{Reconstruction}{The resolution is an illustrative reconstruction of the method, not an assertion that Gelfand stated a classification theorem.}
\GSfield{Fate}{\citet{beilinson1978coherent} supplies the relevant pre-1988 construction. The word only is methodological emphasis rather than a theorem.}
\GSfield{Present status}{As of 11 November 2026: established under the stated hypotheses.}
\GSfield{Assessment}{literal statement: not determinable; publication match: thematic only; verification: indirect; development: none; confidence: medium}
\end{GSitem}
\begin{GSitem}{1987-739.6}{F}{The unexplained three-versus-four distinction}
\GSfield{Kind; source}{remark; \GSpage{P001445}}
\GSfield{The notes}{Parthasarathy is said to answer when the derived category coincides with this one: for three, but not for four.}
\GSfield{Statement}{No definite equivalence assertion is reconstructible: the category being compared and the meaning of three or four are not specified beyond the adjacent star diagram.}
\GSfield{Fate}{The explicit source-star construction and the Dynkin classification in \citet[\S\S1,8]{parthasarathy1988tstructures} explain possible background, but do not license an exact numerical theorem from this sentence.}
\GSfield{Present status}{As of 11 November 2026: unidentified at the level of an exact mathematical statement.}
\GSfield{Assessment}{literal statement: not determinable; publication match: thematic only; verification: read the result; development: none; confidence: low}
\end{GSitem}

\GStalk{740}{1988-03-23}{A. V. Zelevinsky}{P001446-P001447}{Matroidal transversals and Hall's theorem}
An independence example leads to Rado's rank formula for a matroid induced by a relation. Hall's marriage theorem is recovered from the free matroid, with an algorithmic remark and an application to simultaneous left and right coset representatives. Two type errors survive in the handwritten reduction.

\begin{GSitem}{1987-740.1}{E}{Linear independence as a matroid}
\GSfield{Kind; source}{definition; \GSpage{P001446}; I. M. Gelfand}
\GSfield{The notes}{Gelfand gives linearly independent systems of vectors as the example of a matroid.}
\GSfield{Statement}{For a finite indexed family \((v_e)_{e\in E}\) in a vector space, let \(\mathcal I=\{I\subset E:(v_e)_{e\in I}\text{ is linearly independent}\}\). Heredity and the exchange axiom make \((E,\mathcal I)\) a matroid.}
\GSfield{Fate}{This is the elementary representable-matroid example in \citet[opening definitions]{oxleybriefly}. Compare the matroid/lattice occurrence 1985-658.12; no new problem is posed here.}
\GSfield{Present status}{As of 11 November 2026: established under the stated hypotheses.}
\GSfield{Assessment}{literal statement: correct; publication match: exact; verification: read the result; development: none; confidence: high}
\end{GSitem}
\begin{GSitem}{1987-740.2}{E}{Rado transversals and the induced rank formula}
\GSfield{Kind; source}{theorem; \GSpage{P001446}}
\GSfield{The notes}{For a relation \(\varphi\subset M_1\times M_2\), the asserted rank is \(r_2(A)=\min_{B\subset A}(r_1(\varphi(B))+|A-B|)\).}
\GSfield{Statement}{Let \(M_1\) be a finite matroid and \(M_2\) finite. A subset \(I\subset M_2\) is independent when it admits distinct representatives in \(M_1\), related to its elements by \(\varphi\), that form an independent set of the first matroid. These sets form a matroid, and the displayed minimum is its rank on \(A\).}
\GSfield{Reconstruction}{Finite ground sets and the independent-transversal definition are supplied explicitly.}
\GSfield{Fate}{The independence-transversal theorem is \citet{rado1942independence}; Hall's special case predates it. The publisher-deposited Crossref metadata for Rado were opened and checked; the full original theorem was not successfully opened in this run.}
\GSfield{Present status}{As of 11 November 2026: established under the stated hypotheses.}
\GSfield{Assessment}{literal statement: correct with missing hypotheses; publication match: exact; verification: publisher metadata checked; development: none; confidence: medium}
\end{GSitem}
\begin{GSitem}{1987-740.3}{E}{Hall's marriage theorem and the free-matroid reduction}
\GSfield{Kind; source}{theorem; \GSpage{P001446-P001447}}
\GSfield{The notes}{The union condition for sets \(B(n)\subset M\) yields an injection \(f:N\to M\); the notes write \(f(n)\subset B(n)\) and declare all subsets of \(N\) independent.}
\GSfield{Statement}{For finite \(N\), an injection satisfying \(f(n)\in B(n)\) exists iff \(|\bigcup_{n\in K}B(n)|\geq|K|\) for every \(K\subset N\). To derive this from Rado with representatives in \(M\), use the free matroid on \(M\), whose rank is cardinality.}
\GSfield{Reconstruction}{Membership replaces subset inclusion, and the free matroid belongs on the representative ground set \(M\).}
\GSfield{Fate}{\citet[Theorem~1, pp.~27--29]{hall1935representatives} proves precisely this condition; the paper was received and read in April 1934 and published in 1935.}
\GSfield{Present status}{As of 11 November 2026: established under the stated hypotheses.}
\GSfield{Assessment}{literal statement: false as written; publication match: exact; verification: read the result; development: annotated; confidence: high}
\end{GSitem}
\begin{GSitem}{1987-740.4}{E}{An algorithmic proof of the marriage theorem}
\GSfield{Kind; source}{method; \GSpage{P001447}}
\GSfield{The notes}{There is a proof containing an algorithm.}
\GSfield{Statement}{For a finite bipartite graph, start with a matching and repeatedly augment along alternating paths. If no augmenting path exists while a left vertex is unmatched, its alternating reachable sets violate Hall's condition. Thus the procedure finds a complete left matching under that condition.}
\GSfield{Reconstruction}{The notes do not name a running time or a particular implementation; no algorithmic novelty is inferred.}
\GSfield{Fate}{\citet[proof of Theorem~1]{hall1935representatives} gives a constructive replacement argument. The augmenting-path formulation is an explicit modern reconstruction of the terse algorithmic remark.}
\GSfield{Present status}{As of 11 November 2026: established under the stated hypotheses.}
\GSfield{Assessment}{literal statement: correct; publication match: partial; verification: read the result; development: none; confidence: high}
\end{GSitem}
\begin{GSitem}{1987-740.5}{E}{A common system of left and right coset representatives}
\GSfield{Kind; source}{theorem; \GSpage{P001447}}
\GSfield{The notes}{A finite subgroup has representatives that work simultaneously for its left and right cosets; the statement is attributed to Birkhoff.}
\GSfield{Statement}{If \(H\leq G\) and \(G\) is finite, there is a set meeting each left coset and each right coset of \(H\) once. More generally finite \(H\) admits such a transversal in arbitrary \(G\), by applying the finite equal-block matching argument within each double coset and choosing these transversals.}
\GSfield{Reconstruction}{For infinite \(G\), each double coset has equally many left and right cosets because \(H\) is finite; a choice of the finite matchings is used.}
\GSfield{Fate}{\citet[Theorem~3 and its equal-size corollary, pp.~29--30]{hall1935representatives} supplies the partition theorem; its application to cosets is immediate. Hall credits the equal-block result to Konig, so the notebook's attribution is not independently confirmed.}
\GSfield{Present status}{As of 11 November 2026: established under the stated hypotheses.}
\GSfield{Assessment}{literal statement: correct with missing hypotheses; publication match: partial; verification: read the result; development: none; confidence: high}
\end{GSitem}

\GStalk{741}{1988-03-23}{I. M. Gelfand}{P001447}{Famous problems as sources of mathematical theories}
Gelfand lists the parallel postulate, classical construction problems, Fermat's problem, the four-colour problem, and the Riemann hypothesis. These are historical and methodological remarks, not claims of new proofs. Faltings appears in the margin; the finiteness of rational points must not be confused with Fermat's theorem.

\begin{GSitem}{1987-741.1}{E}{The parallel postulate and non-Euclidean models}
\GSfield{Kind; source}{remark; \GSpage{P001447}}
\GSfield{The notes}{The parallel axiom is listed as a seemingly fruitless problem that produced non-Euclidean geometry.}
\GSfield{Statement}{Models of hyperbolic geometry satisfy the other standard incidence/congruence axioms but not Euclid's parallel axiom; relative to the consistency of the background geometry, the axiom cannot be deduced from those remaining axioms.}
\GSfield{Reconstruction}{The formal independence statement is relative to an axiom system and consistency assumption, neither specified in the remark.}
\GSfield{Fate}{\citet{beltrami1869interpretation} is the opened 1869 French publication of Beltrami's 1868 interpretation. This supports the mathematical endpoint, not an exhaustive causal history.}
\GSfield{Present status}{As of 11 November 2026: established under the stated hypotheses.}
\GSfield{Assessment}{literal statement: correct; publication match: partial; verification: publisher metadata checked; development: none; confidence: high}
\end{GSitem}
\begin{GSitem}{1987-741.2}{E}{Classical construction impossibilities}
\GSfield{Kind; source}{remark; \GSpage{P001447}}
\GSfield{The notes}{Angle trisection and circle quadrature are linked with the development of Galois theory.}
\GSfield{Statement}{With an unmarked straightedge and compass, an arbitrary given angle cannot be trisected, and a square of the same area as a given circle cannot be constructed. The first follows from a cubic obstruction, for example trisection of \(60^\circ\); the second follows from the transcendence of \(\pi\).}
\GSfield{Reconstruction}{The allowed instruments and the general rather than every-angle quantifier are supplied.}
\GSfield{Fate}{\citet[\S1 and the applications]{wantzel1837constructions} treats algebraic constructibility and trisection. The two impossibilities have different proofs; quadrature is not a consequence solely of finite Galois theory. The opened publisher record verifies \citet{lindemann1882pi}; its original proof was not read in this run.}
\GSfield{Present status}{As of 11 November 2026: established under the stated hypotheses.}
\GSfield{Assessment}{literal statement: correct with missing hypotheses; publication match: partial; verification: read the result; development: none; confidence: medium}
\end{GSitem}
\begin{GSitem}{1987-741.3}{B}{Fermat's equation and the later modularity proof}
\GSfield{Kind; source}{remark; \GSpage{P001447}}
\GSfield{The notes}{Fermat's problem is linked to Kummer's ideals and algebraic geometry; Faltings appears in a marginal annotation.}
\GSfield{Statement}{For every integer \(n>2\), there are no positive integers \(a,b,c\) with \(a^n+b^n=c^n\). Finiteness of rational points on a fixed genus-greater-than-one curve does not by itself prove this nonexistence statement.}
\GSfield{Reconstruction}{The exact Fermat statement is the standard meaning of the named problem; the marginal Faltings attribution is not turned into a claim of a solution.}
\GSfield{Fate}{\citet{wiles1995modular} establishes the semistable modularity theorem underlying the proof; the final argument was published in 1995 after the 1993 announcement and the repair with Taylor. The opened Annals record confirms publication, but the full proof was not read in this run.}
\GSfield{Present status}{As of 11 November 2026: Fermat' s last theorem is proved; the historical chain in the notes remains a compressed motivation.}
\GSfield{Assessment}{literal statement: correct; publication match: partial; verification: publisher metadata checked; development: none; confidence: high; problem \hyperref[prob:GP-1987-42]{GP-1987-42}}
\end{GSitem}
\begin{GSitem}{1987-741.4}{E}{Four colours and computer-assisted proof}
\GSfield{Kind; source}{remark; \GSpage{P001447}}
\GSfield{The notes}{The four-colour problem is connected with computer mathematics.}
\GSfield{Statement}{Every finite loopless planar graph admits a proper vertex colouring with at most four colours. The 1976 Appel--Haken proof combined unavoidable configurations with certified reducibility computations.}
\GSfield{Fate}{The reducibility component was published as \citet{appelhakenkoch1977reducibility}; \citet[\S\S2--3]{gonthier2005fourcolour} later formalized the complete theorem. Fuchs returns to the proof in record 745.}
\GSfield{Present status}{As of 11 November 2026: established under the stated hypotheses.}
\GSfield{Assessment}{literal statement: correct; publication match: partial; verification: read the result; development: none; confidence: high; problem \hyperref[prob:GP-1987-43]{GP-1987-43}}
\end{GSitem}
\begin{GSitem}{1987-741.5}{C}{The Riemann hypothesis}
\GSfield{Kind; source}{conjecture; \GSpage{P001447}}
\GSfield{The notes}{The Riemann hypothesis is the fifth named example; no proof or strengthened formulation follows.}
\GSfield{Statement}{Every zero \(\rho\) of the analytically continued Riemann zeta function with \(0<\Re\rho<1\) satisfies \(\Re\rho=1/2\).}
\GSfield{Fate}{The Clay Mathematics Institute's current official problem page, \citet{clay2026riemann}, still labels the problem unsolved. No narrower new question or mechanism is present in these notes.}
\GSfield{Present status}{As of 11 November 2026: open; no seminar-specific partial result is claimed.}
\GSfield{Assessment}{literal statement: correct; publication match: exact; verification: read the result; development: none; confidence: high; problem \hyperref[prob:GP-1987-44]{GP-1987-44}}
\GSfield{Worth developing}{The famous conjecture is worth research in general, but this occurrence offers no bounded route for a new contribution in the present chapter.}
\end{GSitem}

\GStalk{742}{1988-03-23}{A. N. Rudakov}{P001448}{Exceptional bundles, Markov triples, and braid actions}
Rudakov describes exceptional triples on the projective plane, their Markov ranks, and mutation. A quadric analogue involves exceptional quadruples and an unnamed Diophantine system; the missing system prevents an exact fate for the asserted converse. The higher-projective-space remark concerns mutation actions, without claiming transitivity in every dimension.

\begin{GSitem}{1987-742.1}{E}{Exceptional bundles and ordered exceptional collections}
\GSfield{Kind; source}{definition; \GSpage{P001448}}
\GSfield{The notes}{For bundles \(E_1,E_2,E_3\) on \(\mathbb P^2\), the notes impose scalar endomorphisms, vanishing positive self-Ext, and backward Ext vanishing.}
\GSfield{Statement}{An object \(E\in D^b(\operatorname{Coh}X)\) is exceptional when \(\operatorname{Hom}(E,E)=k\) and \(\operatorname{Ext}^p(E,E)=0\) for \(p\ne0\). An ordered exceptional collection satisfies \(\operatorname{Ext}^p(E_j,E_i)=0\) for every \(p\) and \(j>i\). Here the objects are algebraic vector bundles on \(\mathbb P^2\).}
\GSfield{Reconstruction}{The general derived definition includes all integer shifts; for vector bundles on the plane only Ext degrees 0,1,2 can occur.}
\GSfield{Fate}{\citet[\S1]{rudakov1988markov} states the bundle framework, already developed by the time of the talk. The corresponding 1987 projective-space paper is \citet{gorodentsevrudakov1987exceptional}.}
\GSfield{Present status}{As of 11 November 2026: established under the stated hypotheses.}
\GSfield{Assessment}{literal statement: correct; publication match: exact; verification: read the result; development: none; confidence: high}
\end{GSitem}
\begin{GSitem}{1987-742.2}{E}{Exceptional ranks satisfy the Markov equation}
\GSfield{Kind; source}{theorem; \GSpage{P001448}}
\GSfield{The notes}{The ranks \(x_1,x_2,x_3\) of the exceptional triple satisfy \(x_1^2+x_2^2+x_3^2=3x_1x_2x_3\).}
\GSfield{Statement}{For a full exceptional triple of vector bundles on \(\mathbb P^2\), its positive ranks form a positive integer solution of the Markov equation.}
\GSfield{Reconstruction}{Fullness and positive ranks are stated explicitly; this is not a rank equation for arbitrary triples of exceptional bundles.}
\GSfield{Fate}{\citet[Theorem~3.2]{rudakov1988markov} gives the rank identity. The Russian original appeared in 1988, volume 52, issue 1; the English translation in 1989. Its submission was in April 1986.}
\GSfield{Present status}{As of 11 November 2026: established under the stated hypotheses.}
\GSfield{Assessment}{literal statement: correct with missing hypotheses; publication match: exact; verification: read the result; development: none; confidence: high}
\end{GSitem}
\begin{GSitem}{1987-742.3}{E}{Realization of positive Markov triples by bundles}
\GSfield{Kind; source}{theorem; \GSpage{P001448}}
\GSfield{The notes}{Conversely, the notes say every integer solution is obtained.}
\GSfield{Statement}{Every positive integer Markov triple is realized by the ranks of an exceptional triple of bundles on \(\mathbb P^2\), obtained from the standard triple by the appropriate mutations and ordering conventions. Zero or negative integral solutions are not positive bundle ranks.}
\GSfield{Reconstruction}{The word integer is restricted to positive integer, as required by ranks.}
\GSfield{Fate}{\citet[Introduction and Theorem~4.2]{rudakov1988markov} constructs the Markov tree of adjacent exceptional bundles. This does not assert the separate Markov uniqueness conjecture that a largest entry determines the triple.}
\GSfield{Present status}{As of 11 November 2026: established under the stated hypotheses.}
\GSfield{Assessment}{literal statement: correct with missing hypotheses; publication match: exact; verification: read the result; development: annotated; confidence: high}
\end{GSitem}
\begin{GSitem}{1987-742.4}{E}{Mutation and the braid action on the plane}
\GSfield{Kind; source}{construction; \GSpage{P001448}}
\GSfield{The notes}{The braid group acts on exceptional triples, transitively in the plane case.}
\GSfield{Statement}{Left and right mutations of adjacent exceptional objects obey the braid relation and act on full exceptional collections. For the plane, the bundle/helix formulation produces the exceptional triples from a standard helix, with the usual shifts and simultaneous line-bundle twists accounted for.}
\GSfield{Reconstruction}{The action is naturally on derived exceptional collections; a statement restricted to bundles requires the helix conventions, which the notes suppress.}
\GSfield{Fate}{\citet{gorodentsevrudakov1987exceptional} is the 1987 projective-space source; \citet[\S\S3--4]{rudakov1988markov} gives the explicit plane reconstruction corresponding to Markov moves.}
\GSfield{Present status}{As of 11 November 2026: established under the stated hypotheses.}
\GSfield{Assessment}{literal statement: correct with missing hypotheses; publication match: partial; verification: read the result; development: none; confidence: medium}
\end{GSitem}
\begin{GSitem}{1987-742.5}{A}{Exceptional quadruples on a quadric}
\GSfield{Kind; source}{claim; \GSpage{P001448}}
\GSfield{The notes}{For \(\mathbb P^1\times\mathbb P^1\), exceptional quadruples also have a braid action and satisfy a more complicated Diophantine system.}
\GSfield{Statement}{On the smooth quadric surface, the four line bundles
\[
 \mathcal O,\quad\mathcal O(1,0),\quad\mathcal O(0,1),\quad\mathcal O(1,1)
\]
generate a standard helix. Mutations construct other helices; every helix and every exceptional bundle is constructible in Rudakov's precise sense.}
\GSfield{Reconstruction}{The standard helix and the exact later theorems are supplied as a bounded publication match; they do not reconstruct all numerical equations in the notebook.}
\GSfield{Fate}{\citet[Theorems~7.1 and~10.1]{rudakov1988quadric} proves these construction results in the Russian paper published in issue 4 of 1988, translated in 1989. The English original received-date line reads 6 July 1986. This is a precise later result, not identification of an omitted numerical system.}
\GSfield{Present status}{As of 11 November 2026: established under the stated hypotheses.}
\GSfield{Assessment}{literal statement: correct with missing hypotheses; publication match: partial; verification: read the result; development: none; confidence: high}
\end{GSitem}
\begin{GSitem}{1987-742.6}{F}{The unspecified converse for quadric numerical data}
\GSfield{Kind; source}{question; \GSpage{P001448}}
\GSfield{The notes}{The notes say it has not been proved that all solutions of the quadric system arise from exceptional quadruples.}
\GSfield{Statement}{An exact arithmetic realization question cannot be recovered, because neither the system nor its positivity, integrality and Chern-class restrictions is written.}
\GSfield{Fate}{\citet[\S\S5--7 and closing remark]{rudakov1988quadric} distinguishes helix construction, symmetric numerical equations, and a question about extending exceptional pairs. None can be silently substituted for the unnamed system here.}
\GSfield{Present status}{As of 11 November 2026: the exact notebook question is unidentified; it is not indexed as an open numerical problem.}
\GSfield{Assessment}{literal statement: not determinable; publication match: thematic only; verification: read the result; development: none; confidence: low}
\end{GSitem}
\begin{GSitem}{1987-742.7}{E}{Mutation of collections on higher projective spaces}
\GSfield{Kind; source}{construction; \GSpage{P001448}}
\GSfield{The notes}{For \(\mathbb P^n\), take \(n+1\) bundles with the same exceptional conditions; a braid group again acts.}
\GSfield{Statement}{Mutations act on ordered full exceptional collections of length \(n+1\) in \(D^b(\operatorname{Coh}\mathbb P^n)\). The standard line-bundle collection is an example. This statement does not claim transitivity on all collections or classification of all vector bundles when \(n>2\).}
\GSfield{Fate}{The pre-talk source is \citet{gorodentsevrudakov1987exceptional}; its bibliographic record was opened. The construction is consistent with the explicit mutation formalism in \citet[\S2]{rudakov1988quadric}.}
\GSfield{Present status}{As of 11 November 2026: established under the stated hypotheses.}
\GSfield{Assessment}{literal statement: correct with missing hypotheses; publication match: partial; verification: read the result; development: none; confidence: high}
\end{GSitem}

\GStalk{743}{1988-03-23}{M. M. Kapranov}{P001448-P001449}{Mutation of exceptional pairs}
Kapranov explains exceptional pairs and their mutations. Gelfand relates the discussion to representation type and quiver reflections. The kernel sequence in the notes has its resulting exceptional pair in the wrong order. The remark that an unsolvable classification always has an open orbit is also too broad.

\begin{GSitem}{1987-743.1}{E}{Finite, tame, and wild representation problems}
\GSfield{Kind; source}{remark; \GSpage{P001448}; I. M. Gelfand}
\GSfield{The notes}{Gelfand distinguishes no parameters, one parameter, and at least two parameters, comparing the last with classifying pairs of matrices.}
\GSfield{Statement}{For a fixed finite-dimensional algebra over an algebraically closed field, finite representation type means finitely many indecomposable isomorphism classes; tame type permits finitely many one-parameter families in each dimension, with finitely many exceptions; wildness means that the problem of modules over the free algebra in two generators embeds in the representation problem. A raw count of parameters does not define these alternatives.}
\GSfield{Reconstruction}{The standard representation-theoretic meanings replace the informal parameter count.}
\GSfield{Fate}{For quivers, the finite/Dynkin and extended-Dynkin examples already occur in 1972-277.8 and 1972-278.9. \citet{bernsteingelfandponomarev1973coxeter} supplies the finite-type setting. The remark is not a proof of a general dichotomy theorem.}
\GSfield{Present status}{As of 11 November 2026: established under the stated hypotheses.}
\GSfield{Assessment}{literal statement: correct with missing hypotheses; publication match: thematic only; verification: read the result; development: none; confidence: medium}
\end{GSitem}
\begin{GSitem}{1987-743.2}{E}{An open orbit is not automatic in a wild problem}
\GSfield{Kind; source}{claim; \GSpage{P001448}; I. M. Gelfand}
\GSfield{The notes}{Gelfand says an unsolvable classification always contains a solvable part, adding an open orbit.}
\GSfield{Statement}{Existence of a dense open orbit is a separate geometric condition. In simultaneous conjugation of pairs of \(n\times n\) complex matrices, the invariant \(\operatorname{tr}(A)\) varies, so no orbit can be dense in the full affine space of pairs. This already refutes an unconditional open-orbit assertion.}
\GSfield{Reconstruction}{\ours{} The text is read literally as claiming an open orbit in the original representation space; a more restricted stratum could behave differently.}
\GSfield{Fate}{This is an elementary editorial check: a polynomial invariant is constant on an orbit and its closure. No external classification theorem is needed. The BGP finite-type examples give context only.}
\GSfield{Present status}{As of 11 November 2026: false without a specified stratum or additional hypotheses.}
\GSfield{Assessment}{literal statement: false as written; publication match: thematic only; verification: indirect; development: annotated; confidence: high}
\end{GSitem}
\begin{GSitem}{1987-743.3}{E}{Self-Ext vanishing and infinitesimal rigidity}
\GSfield{Kind; source}{claim; \GSpage{P001449}}
\GSfield{The notes}{\(\operatorname{Ext}^1(E,E)=0\) is interpreted as rigidity of a vector bundle.}
\GSfield{Statement}{For a vector bundle on a proper scheme over a field, first-order deformations over the dual numbers are parametrized by \(\operatorname{Ext}^1(E,E)\). Its vanishing implies infinitesimal rigidity; the stronger assertion of a unique global isomorphism class does not follow.}
\GSfield{Reconstruction}{The meaning of rigidity is restricted to infinitesimal deformations, with the usual fixed base and deformation functor.}
\GSfield{Fate}{The deformation-free exceptional-bundle setting is explicit in \citet[Introduction and \S1]{rudakov1988markov}.}
\GSfield{Present status}{As of 11 November 2026: established under the stated hypotheses.}
\GSfield{Assessment}{literal statement: correct with missing hypotheses; publication match: exact; verification: read the result; development: none; confidence: high}
\end{GSitem}
\begin{GSitem}{1987-743.4}{E}{Exceptional pairs require exceptional individual objects}
\GSfield{Kind; source}{definition; \GSpage{P001449}}
\GSfield{The notes}{A pair is called exceptional if \(\operatorname{Ext}^j(F,E)=0\); the individual conditions are not repeated.}
\GSfield{Statement}{An exceptional pair \((E,F)\) consists of two exceptional objects and has \(\operatorname{RHom}(F,E)=0\). Backward orthogonality alone is insufficient to imply either object is exceptional.}
\GSfield{Reconstruction}{The individual scalar-endomorphism and self-Ext conditions are explicitly reinstated.}
\GSfield{Fate}{\citet[\S2]{rudakov1988quadric} gives the admissible/exceptional pair formalism; the earlier individual conditions on P001448 provide likely context, but must be restored when this item is read alone.}
\GSfield{Present status}{As of 11 November 2026: established under the stated hypotheses.}
\GSfield{Assessment}{literal statement: correct with missing hypotheses; publication match: exact; verification: read the result; development: none; confidence: high}
\end{GSitem}
\begin{GSitem}{1987-743.5}{E}{The kernel mutation has the wrong pair order in the notes}
\GSfield{Kind; source}{construction; \GSpage{P001449}}
\GSfield{The notes}{The sequence \(0\to R\to E\otimes\operatorname{Hom}(E,F)\to F\) is followed by the assertion that \((F,R)\) is exceptional.}
\GSfield{Statement}{For an exceptional pair with \(\operatorname{RHom}(E,F)\) concentrated in degree zero and surjective evaluation, the kernel \(R\) gives the left mutation, with the appropriate shift convention. The resulting exceptional pair is \((R,E)\), not \((F,R)\). In general use the distinguished triangle \(L_EF\to\operatorname{RHom}(E,F)\otimes E\to F\to L_EF[1]\).}
\GSfield{Reconstruction}{\ours{} Surjectivity and the strong-pair hypothesis are needed for a kernel bundle; the triangle supplies the general derived construction.}
\GSfield{Fate}{\citet[\S2]{rudakov1988quadric} gives left/right mutation. A direct check on \(\mathbb P^1\), \(E=\mathcal O\), \(F=\mathcal O(1)\), has \(R=\mathcal O(-1)\) and \(\operatorname{Hom}(R,F)\ne0\), contradicting the displayed pair in the notebook.}
\GSfield{Present status}{As of 11 November 2026: the standard mutation is established; the literal order in the source is false.}
\GSfield{Assessment}{literal statement: false as written; publication match: exact; verification: read the result; development: annotated; confidence: high}
\end{GSitem}
\begin{GSitem}{1987-743.6}{E}{Braid relations and the quiver-reflection analogy}
\GSfield{Kind; source}{construction; \GSpage{P001449}}
\GSfield{The notes}{Mutation gives the braid-group action; Gelfand recalls BGP and recommends studying quiver representations and derived categories for spectral theory.}
\GSfield{Statement}{Adjacent mutations on exceptional collections satisfy \(\sigma_i\sigma_{i+1}\sigma_i=\sigma_{i+1}\sigma_i\sigma_{i+1}\), with distant generators commuting. BGP reflections provide a related, but differently defined, transformation of quiver representations.}
\GSfield{Reconstruction}{The braid relation is stated for the derived mutation operations; the analogy is not an identification of the two actions.}
\GSfield{Fate}{\citet[\S2, Proposition~2.1 and the reconstruction formalism]{rudakov1988quadric} treats mutations; \citet[\S1]{bernsteingelfandponomarev1973coxeter} is the earlier reflection construction. The recommendation for spectral theory states no further definite theorem.}
\GSfield{Present status}{As of 11 November 2026: established under the stated hypotheses.}
\GSfield{Assessment}{literal statement: correct; publication match: partial; verification: read the result; development: none; confidence: high}
\end{GSitem}

\GStalk{744}{1988-03-23}{A. A. Beilinson}{P001449-P001452}{Constructible sheaves, singular support, and microlocal Morse groups}
The report, attributed to work of S. Gelfand, MacPherson and Vilonen, proceeds from constructibility to singular support, Lagrangian geometry and a Morse characterization of perversity. Several schematic assertions need the distinction between one test function and a conic neighbourhood, local and global finiteness, and Maslov-normalized Morse data. The zero section cannot simply be added over points outside the support.

\begin{GSitem}{1987-744.1}{F}{Gelfand's geometric configuration question}
\GSfield{Kind; source}{question; \GSpage{P001449}; I. M. Gelfand}
\GSfield{The notes}{Before the perverse-sheaf exposition, Gelfand asks about regions/\allowbreak configurations and Siegel domains, with a drawing in the real projective plane.}
\GSfield{Statement}{The invariant, parameter space, equivalence relation and intended assertion are not specified well enough to formulate a definite classification problem.}
\GSfield{Fate}{The neighbouring microlocal discussion does not supply the missing data. No exact publication is identified, and no open problem is inferred from a topic-level resemblance.}
\GSfield{Present status}{As of 11 November 2026: unidentified as a mathematical question.}
\GSfield{Assessment}{literal statement: not determinable; publication match: none found; verification: unresolved lead; development: none; confidence: low}
\end{GSitem}
\begin{GSitem}{1987-744.2}{E}{Constructible sheaves and complexes}
\GSfield{Kind; source}{definition; \GSpage{P001450}}
\GSfield{The notes}{A sheaf is locally constant of finite rank on each stratum; for a complex the condition is on its cohomology sheaves, not on its terms.}
\GSfield{Statement}{For a complex algebraic variety with a finite algebraic Whitney stratification \(\mathcal S\), \(D^b_{\mathcal S}(X;k)\) consists of bounded sheaf complexes whose cohomology restricts to finite-rank local systems on every stratum. Analytic versions use locally finite stratifications.}
\GSfield{Reconstruction}{Boundedness, coefficients, Whitney regularity and finiteness conventions are supplied.}
\GSfield{Fate}{\citet[Chapter~8]{kashiwaraschapira1985microlocal} treats the constructibility conditions. The notebook correctly distinguishes cohomological constructibility from a particular complex of representatives.}
\GSfield{Present status}{As of 11 November 2026: established under the stated hypotheses.}
\GSfield{Assessment}{literal statement: correct with missing hypotheses; publication match: exact; verification: read the result; development: none; confidence: high}
\end{GSitem}
\begin{GSitem}{1987-744.3}{E}{Singular support requires uniform local tests}
\GSfield{Kind; source}{definition; \GSpage{P001450}}
\GSfield{The notes}{A nonzero covector is excluded from \(SS(\mathcal F)\) when a relative cohomology group for a function with that differential vanishes; the notes say the choice of function does not matter.}
\GSfield{Statement}{For a bounded constructible complex on a real manifold, \((x,\xi)\notin SS(F)\) means there is an open conic neighbourhood on which all local cohomology tests \((R\Gamma_{\{\phi\geq\phi(y)\}}F)_y\) vanish whenever \((y,d\phi_y)\) lies in that neighbourhood. Holomorphic Morse tests describe the corresponding complex-constructible microsupport, with their normalizations.}
\GSfield{Reconstruction}{The precise real-test definition is supplied to avoid interpreting a schematic single-function test as an invariant definition.}
\GSfield{Fate}{The uniform-neighbourhood condition is in \citet[Definition~0.1 and \S3.1]{kashiwaraschapira1985microlocal}. Vanishing for a single arbitrary chosen function is not the complete definition.}
\GSfield{Present status}{As of 11 November 2026: established under the stated hypotheses.}
\GSfield{Assessment}{literal statement: correct with missing hypotheses; publication match: exact; verification: read the result; development: annotated; confidence: high}
\end{GSitem}
\begin{GSitem}{1987-744.4}{E}{The zero-section part of singular support}
\GSfield{Kind; source}{claim; \GSpage{P001450}}
\GSfield{The notes}{If the complex is nonzero, the entire zero section is added; an acyclic complex has empty singular support.}
\GSfield{Statement}{One has \(SS(F)\cap T_X^*X=\operatorname{supp}(F)\), where the right side is embedded as zero covectors. Thus \(SS(0)=\varnothing\), but a nonzero skyscraper sheaf does not contain zero covectors over other points.}
\GSfield{Reconstruction}{\ours{} }
\GSfield{Fate}{\citet[\S3.1, basic properties of microsupport]{kashiwaraschapira1985microlocal} gives the support relation. The skyscraper example on P001452 itself detects the error.}
\GSfield{Present status}{As of 11 November 2026: the support-restricted statement is standard; the unconditional addition of the zero section is false.}
\GSfield{Assessment}{literal statement: false as written; publication match: exact; verification: read the result; development: annotated; confidence: high}
\end{GSitem}
\begin{GSitem}{1987-744.5}{E}{Complex constructibility gives a conic Lagrangian support}
\GSfield{Kind; source}{theorem; \GSpage{P001451}}
\GSfield{The notes}{\(SS(F)\) is an analytic conic Lagrangian set with finitely many irreducible components; ideal involutivity is presented as equivalent.}
\GSfield{Statement}{For complex-constructible \(F\), its microsupport, identified with the complex cotangent bundle, is a closed complex-conic Lagrangian analytic subset. Components are locally finite; global finiteness follows in the finite algebraic setting. Involutivity of the ideal gives coisotropicity, and is equivalent to the Lagrangian condition only with the half-dimensional hypothesis.}
\GSfield{Reconstruction}{Complex constructibility and the dimension/finiteness qualifications are made explicit.}
\GSfield{Fate}{\citet[Theorem~6.4.1 and Chapter~8]{kashiwaraschapira1985microlocal} supplies involutivity and the constructible Lagrangian characterization.}
\GSfield{Present status}{As of 11 November 2026: established under the stated hypotheses.}
\GSfield{Assessment}{literal statement: correct with missing hypotheses; publication match: exact; verification: read the result; development: annotated; confidence: high}
\end{GSitem}
\begin{GSitem}{1987-744.6}{E}{Irreducible conic Lagrangians are conormal closures}
\GSfield{Kind; source}{theorem; \GSpage{P001451}}
\GSfield{The notes}{An irreducible conic Lagrangian \(Y\subset T^*X\) is \(N^*Z\) for \(Z=\pi(Y)\), with closure over the smooth locus understood.}
\GSfield{Statement}{In the algebraic setting, a closed irreducible conic Lagrangian subset is the closure of \(T^*_{Z_{\mathrm{sm}}}X\) for the irreducible subvariety \(Z=\overline{\pi(Y)}\). The analytic assertion is local. This follows by annihilating tangent vectors to a generic smooth image under the canonical symplectic form and comparing dimensions.}
\GSfield{Reconstruction}{Closedness, algebraic/local analytic context and closure of the image are specified.}
\GSfield{Fate}{This is the conormal geometry underlying \citet[Chapter~8]{kashiwaraschapira1985microlocal}. The argument and the closure conventions are elementary reconstruction of the lemma as stated.}
\GSfield{Present status}{As of 11 November 2026: established under the stated hypotheses.}
\GSfield{Assessment}{literal statement: correct with missing hypotheses; publication match: partial; verification: read the result; development: none; confidence: high}
\end{GSitem}
\begin{GSitem}{1987-744.7}{E}{Morse transversality to a singular support}
\GSfield{Kind; source}{definition; \GSpage{P001451}}
\GSfield{The notes}{The graph of \(df\) must meet \(Y\) only at smooth points and transversely; examples are the zero section and a smooth conormal.}
\GSfield{Statement}{For a holomorphic function \(f\) and conic Lagrangian \(Y\), a microlocal Morse point is a transverse intersection of \(\Gamma_{df}\) with \(Y_{\mathrm{sm}}\), away from the singular part of \(Y\). For \(Y=T_X^*X\) this is an ordinary nondegenerate critical point; for \(Y=T_Z^*X\) with smooth \(Z\), it is a nondegenerate critical point of \(f|_Z\).}
\GSfield{Reconstruction}{The criterion is local at an isolated transverse intersection, not a demand that a global function exist with every required property.}
\GSfield{Fate}{This is the test-function construction used in \citet[Introduction and \S1]{gelfandmacphersonvilonen2005microlocal}, building on the earlier microlocal theory. The crossed-out extra demand \(df\ne0\) is not part of the general conormal criterion.}
\GSfield{Present status}{As of 11 November 2026: established under the stated hypotheses.}
\GSfield{Assessment}{literal statement: correct; publication match: exact; verification: read the result; development: none; confidence: high}
\end{GSitem}
\begin{GSitem}{1987-744.8}{E}{Perversity detected by Morse groups}
\GSfield{Kind; source}{definition; \GSpage{P001452}}
\GSfield{The notes}{All Morse relative hypercohomology groups must occur in degree zero; a parenthetical remark instead says first degree.}
\GSfield{Statement}{With the convention that \(k_X[\dim_{\mathbb C}X]\) is perverse, a constructible bounded complex is perverse exactly when its normalized local Morse groups are concentrated in degree zero. For an isolated Morse test, these can be represented by \(\mathbb H^j(B,B\cap f^{-1}(\varepsilon);F)\), after fixing the vanishing-cycle shift convention.}
\GSfield{Reconstruction}{The normalized Morse functor rather than an unlabelled cone fixes the degree convention.}
\GSfield{Fate}{\citet[Theorem~9.5.2 and Introduction]{kashiwaraschapira1985microlocal} gives the microlocal characterization. The degree-one parenthesis can reflect the unshifted vanishing-cycle cone, and cannot coexist literally with degree zero without specifying the shift.}
\GSfield{Present status}{As of 11 November 2026: established under the stated hypotheses.}
\GSfield{Assessment}{literal statement: correct with missing hypotheses; publication match: exact; verification: read the result; development: annotated; confidence: high}
\end{GSitem}
\begin{GSitem}{1987-744.9}{E}{Basic perverse sheaves and the closed-embedding hypothesis}
\GSfield{Kind; source}{construction; \GSpage{P001452}}
\GSfield{The notes}{Examples are a skyscraper, \(\mathbb C_X[n]\), and \(\mathbb C_Z[\dim Z]\) extended by zero from a smooth subset, or an intersection-cohomology complex.}
\GSfield{Statement}{A finite-dimensional skyscraper is perverse. On a smooth complex \(n\)-fold a local system shifted by \(n\) is perverse. For a closed smooth embedding \(i:Z\hookrightarrow X\), \(i_*L[\dim Z]\) is perverse. For an arbitrary open or locally closed embedding, ordinary extension by zero need not be perverse; intermediate extension is the appropriate canonical construction.}
\GSfield{Reconstruction}{Closedness is restored to the zero-extension example; the singular case uses intersection cohomology rather than a constant sheaf.}
\GSfield{Fate}{\citet[Theorem~9.5.2]{kashiwaraschapira1985microlocal} gives the perverse framework, and \citet{goreskymacpherson1983intersection} is the intersection-cohomology source.}
\GSfield{Present status}{As of 11 November 2026: established under the stated hypotheses.}
\GSfield{Assessment}{literal statement: correct with missing hypotheses; publication match: partial; verification: read the result; development: none; confidence: high}
\end{GSitem}
\begin{GSitem}{1987-744.10}{E}{Singular support does not canonically recover every stratum}
\GSfield{Kind; source}{remark; \GSpage{P001452}}
\GSfield{The notes}{An explicitly approximate remark identifies strata with projections of irreducible components of \(SS(F)\).}
\GSfield{Statement}{The projections recover certain closed singular-support loci, not a unique stratification. A compatible Whitney stratification may be refined without changing \(F\) or \(SS(F)\); even the constant sheaf admits many such refinements.}
\GSfield{Reconstruction}{\ours{} The notebook already flags the assertion as approximate; its strongest literal interpretation is rejected.}
\GSfield{Fate}{\citet[Chapter~8]{kashiwaraschapira1985microlocal} relates constructibility to unions of conormals, which is weaker than a unique reconstruction of the chosen strata.}
\GSfield{Present status}{As of 11 November 2026: only the conormal-support relation is valid; no canonical full stratification follows.}
\GSfield{Assessment}{literal statement: false as written; publication match: partial; verification: read the result; development: annotated; confidence: high}
\end{GSitem}
\begin{GSitem}{1987-744.11}{B}{Morse groups form a microlocal local system with a twist}
\GSfield{Kind; source}{construction; \GSpage{P001452}}
\GSfield{The notes}{At a smooth covector \(z\), define \(\Phi(F)_z=H^0(U,f^{-1}(\varepsilon);F)\) for \(df=z\), claiming no dependence on \(f\).}
\GSfield{Statement}{On the smooth locus of a fixed singular support, normalized Morse data vary locally constantly after accounting for the Maslov/orientation data of the test-function Hessian. The natural object is a twisted local system (or a local system on the appropriate bundle of Morse data), not canonically an untwisted local system from the covector alone.}
\GSfield{Reconstruction}{The twisting is a modern clarification, not wording supplied by the notebook.}
\GSfield{Fate}{\citet[\S1, the bundle of test Lagrangians and Maslov data]{gelfandmacphersonvilonen2005microlocal} makes the suppressed structure explicit. Its preprint appeared in September 2005; the 1988 talk already attributes the programme to the same three authors.}
\GSfield{Present status}{As of 11 November 2026: established under the stated hypotheses.}
\GSfield{Assessment}{literal statement: correct with missing hypotheses; publication match: partial; verification: read the result; development: annotated; confidence: high}
\end{GSitem}

\GStalk{745}{1988-04-11}{D. B. Fuchs}{P001453-P001457}{Reducibility and discharging in the four-colour theorem}
Fuchs explains the Appel--Haken strategy through minimal counterexamples, reducible configurations, Kempe changes and Euler charge. Gelfand asks for a parallel program for the reducibility test and digresses on spherical functions. The notebook gives examples and incomplete discharging rules, not the full unavoidable set or a reproducible implementation.

\begin{GSitem}{1987-745.1}{E}{Configurations, reducibility, and unavoidability}
\GSfield{Kind; source}{definition; \GSpage{P001453}}
\GSfield{The notes}{A configuration is a triangulated disc; an unavoidable collection occurs in every triangulation, and a reducible member cannot occur in a smallest counterexample.}
\GSfield{Statement}{In a finite simple plane triangulation, a configuration specifies a finite near-triangulation together with the required degrees of its internal vertices and a boundary ring. Reducibility means that an occurrence can be replaced by a smaller graph whose four-colouring can be lifted, possibly using Kempe changes. An unavoidable set is one occurring in every triangulation of the relevant reduced class.}
\GSfield{Reconstruction}{Degree and boundary data and the relevant class of minimal counterexamples are restored.}
\GSfield{Fate}{\citet[\S3]{gonthier2005fourcolour} states the established proof structure, already present in the 1977 Appel--Haken proof. An arbitrary disc triangulation without embedding and degree data is not the complete configuration specification.}
\GSfield{Present status}{As of 11 November 2026: established under the stated hypotheses.}
\GSfield{Assessment}{literal statement: correct with missing hypotheses; publication match: exact; verification: read the result; development: none; confidence: high}
\end{GSitem}
\begin{GSitem}{1987-745.2}{E}{An unavoidable reducible family proves four-colourability}
\GSfield{Kind; source}{method; \GSpage{P001453}}
\GSfield{The notes}{The main idea is to combine the two definitions and induct on the number of vertices.}
\GSfield{Statement}{If every minimal non-four-colourable plane triangulation contains a member of a fixed family, and every member is reducible, no minimal counterexample exists. After deleting duplicate edges, a finite loopless planar graph with at least three vertices extends to a simple plane triangulation; smaller graphs are immediate. Restricting a four-colouring yields the theorem.}
\GSfield{Fate}{The proof scheme was realized in \citet{appelhakenkoch1977reducibility}; \citet[\S3, parts I--II]{gonthier2005fourcolour} states its formalized version.}
\GSfield{Present status}{As of 11 November 2026: established under the stated hypotheses.}
\GSfield{Assessment}{literal statement: correct; publication match: exact; verification: read the result; development: none; confidence: high; problem \hyperref[prob:GP-1987-43]{GP-1987-43}}
\end{GSitem}
\begin{GSitem}{1987-745.3}{E}{The cosine addition identity as a spherical-function example}
\GSfield{Kind; source}{remark; \GSpage{P001454}; I. M. Gelfand}
\GSfield{The notes}{Gelfand says his spherical-function theory arose from \(\cos(\alpha+\beta)+\cos(\alpha-\beta)=2\cos\alpha\cos\beta\), using translations and reflections of the line.}
\GSfield{Statement}{For \(G=\mathbb R\rtimes\{\pm1\}\), \(K=\{\pm1\}\), the functions \(\varphi_t(x)=\cos(tx)\) are normalized \(K\)-bi-invariant spherical functions and satisfy \(\frac12[\varphi_t(x+y)+\varphi_t(x-y)]=\varphi_t(x)\varphi_t(y)\). This is a concrete product-formula model.}
\GSfield{Reconstruction}{The parameter t, subgroup K and normalized averaging are explicit reconstruction of the stated example.}
\GSfield{Fate}{The general spherical-function setting is represented by \citet{berezingelfand1956spherical}; the displayed semidirect-product calculation directly verifies the specific example. The personal account of invention is not independently checked.}
\GSfield{Present status}{As of 11 November 2026: established under the stated hypotheses.}
\GSfield{Assessment}{literal statement: correct; publication match: thematic only; verification: indirect; development: none; confidence: high}
\end{GSitem}
\begin{GSitem}{1987-745.4}{E}{Degree-four vertices are reducible by a Kempe argument}
\GSfield{Kind; source}{theorem; \GSpage{P001454}}
\GSfield{The notes}{For four neighbours with all four colours, either the 1--3 Kempe path is absent, or its planar separation prevents a 2--4 path; a colour swap frees a colour.}
\GSfield{Statement}{In a plane triangulation, a vertex of degree four cannot belong to a minimal non-four-colourable graph. Remove it and colour the smaller graph. If the four cyclic neighbours use four different colours, disjoint alternating 1--3 and 2--4 connecting paths cannot both exist in the punctured plane; swap on a component of one colour pair and extend the colouring.}
\GSfield{Reconstruction}{The cyclic order and planarity are the essential topological hypotheses.}
\GSfield{Fate}{\citet[\S3, the square case]{gonthier2005fourcolour} records this classical Kempe reduction. The valid degree-four argument does not validate Kempe's failed degree-five argument.}
\GSfield{Present status}{As of 11 November 2026: established under the stated hypotheses.}
\GSfield{Assessment}{literal statement: correct with missing hypotheses; publication match: exact; verification: read the result; development: none; confidence: high}
\end{GSitem}
\begin{GSitem}{1987-745.5}{E}{Boundary colourings and Kempe closure}
\GSfield{Kind; source}{method; \GSpage{P001455}}
\GSfield{The notes}{Some boundary colourings extend into a configuration, others do not; a computer searches for recolourings as in the previous example.}
\GSfield{Statement}{For a fixed finite ring configuration, enumerate proper four-colourings of the ring and the subset extendible through the interior. Reducibility tests combine this finite data with the permitted noncrossing Kempe-chain pairing information from the exterior. Arbitrary swaps of boundary colours without such compatibility do not establish reducibility.}
\GSfield{Reconstruction}{The exterior noncrossing compatibility, not spelled out in the sketch, is essential.}
\GSfield{Fate}{\citet[\S5.2]{gonthier2005fourcolour} formalizes the Kempe-chain and reducibility calculation. The 1977 proof already used finite reducibility algorithms; the later proof verifies a specific implementation.}
\GSfield{Present status}{As of 11 November 2026: established under the stated hypotheses.}
\GSfield{Assessment}{literal statement: correct with missing hypotheses; publication match: partial; verification: read the result; development: none; confidence: high}
\end{GSitem}
\begin{GSitem}{1987-745.6}{E}{The drawn configuration is the Birkhoff diamond}
\GSfield{Kind; source}{theorem; \GSpage{P001455}}
\GSfield{The notes}{The scan draws a hexagonal ring around four interior vertices and calls it reducible, adding that Fuchs could not prove it.}
\GSfield{Statement}{The four interior vertices span two triangles sharing an edge, and each has degree five in the completed patch. Together with the six-vertex boundary this is the free completion of the Birkhoff diamond, a reducible configuration.}
\GSfield{Reconstruction}{Identification uses the graph incidence and degree-five internal vertices visible in the scan; no missing proof is attributed to Fuchs.}
\GSfield{Fate}{The scan was compared with the structural description in \citet[\S3, part g, pp.~7--8]{gonthier2005fourcolour}. Its reducibility is a classical pre-1988 result, not a new unresolved problem; the attribution to the name Birkhoff is editorial.}
\GSfield{Present status}{As of 11 November 2026: established under the stated hypotheses.}
\GSfield{Assessment}{literal statement: correct; publication match: exact; verification: read the result; development: none; confidence: medium}
\end{GSitem}
\begin{GSitem}{1987-745.7}{B}{A parallel program for a fixed reducibility test}
\GSfield{Kind; source}{problem; \GSpage{P001455}; I. M. Gelfand}
\GSfield{The notes}{Gelfand's numbered problem asks for a program determining reducibility on a parallel computer.}
\GSfield{Statement}{Given a finite ring configuration, implement a specified sound finite reducibility test based on boundary colourings and compatible Kempe data, allowing independent colourings/configurations to be processed in parallel. No machine model, complexity bound or completeness criterion for all notions of reducibility is stated.}
\GSfield{Reconstruction}{A specified finite sufficient test is distinguished from universal semantic reducibility; without that distinction the requested input/output problem is ambiguous.}
\GSfield{Fate}{Computer reducibility tests already existed in 1977, and \citet[\S5.2]{gonthier2005fourcolour} later gives a formally checked implementation. That is only a partial match: this search did not identify a publication answering Gelfand's exact parallel-program request. Finite enumeration itself can be parallelized, so the text does not define a new unsolved mathematical decision problem.}
\GSfield{Present status}{As of 11 November 2026: the finite testing method and verified implementations are established; no exact historical parallel implementation has been identified.}
\GSfield{Assessment}{literal statement: correct with missing hypotheses; publication match: partial; verification: read the result; development: none; confidence: medium; problem \hyperref[prob:GP-1987-45]{GP-1987-45}}
\end{GSitem}
\begin{GSitem}{1987-745.8}{E}{Euler curvature and a first unavoidable set}
\GSfield{Kind; source}{theorem; \GSpage{P001455-P001456}}
\GSfield{The notes}{Using \(V-E+F=2\), \(2E=\sum i\nu_i\), \(3F=2E\), the notes derive \(\sum(6-i)\nu_i=12\), so a degree-three, four or five vertex exists.}
\GSfield{Statement}{For a finite simple triangulation of the sphere with at least four vertices, \(\sum_v(6-\deg v)=12\). Its minimum degree is at most five. Degree-three and degree-four reduction therefore leaves minimum degree at least five in any minimal counterexample.}
\GSfield{Reconstruction}{Sphere topology, simplicity and the small-size exclusion justify the degree range.}
\GSfield{Fate}{This is the Euler-counting step recalled in \citet[\S3]{gonthier2005fourcolour}. The numerical identity follows directly from the three displayed counts and does not depend on computation.}
\GSfield{Present status}{As of 11 November 2026: established under the stated hypotheses.}
\GSfield{Assessment}{literal statement: correct with missing hypotheses; publication match: exact; verification: read the result; development: none; confidence: high}
\end{GSitem}
\begin{GSitem}{1987-745.9}{E}{Discharging preserves the total charge}
\GSfield{Kind; source}{method; \GSpage{P001456-P001457}}
\GSfield{The notes}{Assign charge \(60(6-i)\) to degree i, redistribute 30 along specified edges, and seek a contradiction by making every final charge nonpositive.}
\GSfield{Statement}{The initial charge sum is \(720\). Any pairwise transfer conserves it. Therefore, if absence of a finite configuration family forces every final charge to be nonpositive under a fully specified set of transfer rules, that family is unavoidable. The single sample rule in the notes does not by itself prove the required universal implication.}
\GSfield{Reconstruction}{Conservation is exact; completeness of the discharge rules is not reconstructed from the incomplete examples.}
\GSfield{Fate}{\citet[\S3, discharging steps 1--2]{gonthier2005fourcolour} gives the precise proof pattern. Appel--Haken's completed rule system is a much larger object than this board sketch. For the manuscript computer-time question, the contemporary digest \citet[abstract]{bernhart1977digest} reports 1200 hours for nearly 1900 configurations.}
\GSfield{Present status}{As of 11 November 2026: established under the stated hypotheses.}
\GSfield{Assessment}{literal statement: correct with missing hypotheses; publication match: partial; verification: read the result; development: none; confidence: high}
\end{GSitem}

\GStalk{746}{1988-04-11}{A. A. Beilinson}{P001457-P001459}{Perverse sheaves by vanishing cycles}
The continuation defines vanishing cycles and repeats the microlocal Morse description. It poses the definite reconstruction problem for perverse sheaves from such data. The solution-complex and intersection-cohomology examples connect the discussion to the Riemann--Hilbert correspondence and intermediate extension. Later literature gives a fully faithful embedding and codimension-three extension, with explicit essential-image data still a separate problem.

\begin{GSitem}{1987-746.1}{E}{Zero microsupport detects local constancy}
\GSfield{Kind; source}{theorem; \GSpage{P001457}}
\GSfield{The notes}{\(SS(F)\) equals the zero section iff all cohomology sheaves are locally constant; the whole zero section is again inserted for any nonzero complex.}
\GSfield{Statement}{For a constructible complex, \(SS(F)\subset T_X^*X\) iff all its cohomology sheaves are locally constant. Equality holds on connected components where the complex is nonzero. The zero complex has empty support, so equality cannot be used without the component/support qualification.}
\GSfield{Fate}{\citet[\S3.1]{kashiwaraschapira1985microlocal} gives the local-constancy criterion. This is a second occurrence of the support error corrected in 1987-744.3.}
\GSfield{Present status}{As of 11 November 2026: established under the stated hypotheses.}
\GSfield{Assessment}{literal statement: false as written; publication match: exact; verification: read the result; development: annotated; confidence: high}
\end{GSitem}
\begin{GSitem}{1987-746.2}{E}{Vanishing cycles as relative hypercohomology}
\GSfield{Kind; source}{construction; \GSpage{P001457-P001458}}
\GSfield{The notes}{Take \(f(x)=0\), \(df_x=\nu\), and compare cohomology on a small neighbourhood with a nearby nonzero fibre; relative cohomology is defined by a cone.}
\GSfield{Statement}{For a sufficiently small Milnor neighbourhood and a constructible complex adapted to \(f\), the relative derived sections \(R\Gamma(B,B\cap f^{-1}(\varepsilon);F)\) are the fibre of the restriction map. Equivalently they are its cone shifted by \(-1\). They describe local vanishing-cycle data, with the chosen convention recorded.}
\GSfield{Reconstruction}{The Milnor neighbourhood, nearby regular value and shifted-cone convention are supplied.}
\GSfield{Fate}{The microlocal construction in \citet[\S1]{gelfandmacphersonvilonen2005microlocal} uses these local Morse data. The cone shift matters for the degree-zero definition on the next line.}
\GSfield{Present status}{As of 11 November 2026: established under the stated hypotheses.}
\GSfield{Assessment}{literal statement: correct with missing hypotheses; publication match: exact; verification: read the result; development: none; confidence: high}
\end{GSitem}
\begin{GSitem}{1987-746.3}{E}{Repeated microlocal characterization of perversity}
\GSfield{Kind; source}{definition; \GSpage{P001458}}
\GSfield{The notes}{A holomorphic Morse function has graph transverse to singular support; vanishing cycles occur in one degree and form a local system over the smooth support.}
\GSfield{Statement}{Use isolated transverse intersections with the smooth support, degree-zero normalized Morse groups, and the Maslov-twisted local-system construction. These are the same conditions as in 1987-744.7--744.11, not a new criterion.}
\GSfield{Reconstruction}{The normalization and twist from the earlier occurrence remain necessary.}
\GSfield{Fate}{\citet[Theorem~9.5.2]{kashiwaraschapira1985microlocal} and \citet[\S1]{gelfandmacphersonvilonen2005microlocal} supply the precise formulations.}
\GSfield{Present status}{As of 11 November 2026: established under the stated hypotheses.}
\GSfield{Assessment}{literal statement: correct with missing hypotheses; publication match: partial; verification: read the result; development: none; confidence: high}
\end{GSitem}
\begin{GSitem}{1987-746.4}{C}{Reconstruct perverse sheaves from microlocal Morse data}
\GSfield{Kind; source}{problem; \GSpage{P001458-P001459}}
\GSfield{The notes}{The proposed third method is a microlocal description by vanishing cycles; the report says the work solves half of the problem.}
\GSfield{Statement}{Let \(X\) be a complex manifold and \(\Lambda\subset T^*X\) a specified closed complex analytic conic Lagrangian. With coefficients in \(\mathbb C\), describe the category of perverse sheaves with singular support in \(\Lambda\) by local Morse data on the smooth pieces, gluing at codimension-one loci, and explicit compatibility conditions at singular loci. A complete answer must characterize the essential image, not only embed the category faithfully.}
\GSfield{Reconstruction}{The fixed-support categorical formulation and essential-image distinction make the informal reconstruction question precise; no arbitrary missing codimension-two formula is invented.}
\GSfield{Fate}{\citet[Theorem~5.1 and Theorem~6.4]{gelfandmacphersonvilonen2005microlocal}, first circulated on arXiv in September 2005, proves an equivalence away from the codimension-two locus and a global fully faithful embedding in melded Morse data. \citet[Theorem~1.1]{kashiwaravilonen2014codimension}, announced in 2010 with a full preprint in 2012 and publication in 2014, proves unique extension across analytic subsets of codimension at least three in the Lagrangian support. The exact remaining task is formulated as Problem~9.3 of that paper (published p.~617): describe the conditions imposed by the codimension-two locus on the essential image. The extension theorem does not erase these conditions. \citet[Introduction]{guminov2025perverse} confirms that explicit category descriptions remain available only for restricted families.}
\GSfield{Present status}{As of 11 November 2026: partially developed over complex coefficients in the explicit microlocal-data sense stated here. Kashiwara--Vilonen Problem 9.3 explicitly isolates the codimension-two essential-image question. The consulted later primary literature does not settle that general geometric description; this is a bounded search conclusion, not a claim that no categorical reconstruction theorem of any kind exists.}
\GSfield{Assessment}{literal statement: correct with missing hypotheses; publication match: partial; verification: read the result; development: none; confidence: medium; problem \hyperref[prob:GP-1987-46]{GP-1987-46}}
\GSfield{Worth developing}{A fixed singular support with concrete codimension-two links would make a useful bounded project, but the notebook supplies no such geometry to calculate here.}
\end{GSitem}
\begin{GSitem}{1987-746.5}{E}{Holonomic differential equations and perverse solution complexes}
\GSfield{Kind; source}{construction; \GSpage{P001458}}
\GSfield{The notes}{One construction uses overdetermined differential equations, a solution complex, and a Dolbeault resolution.}
\GSfield{Statement}{For a regular holonomic left \(\mathcal D_X\)-module on a smooth complex n-fold, the shifted derived solution complex \(R\mathcal Hom_{\mathcal D_X}(M,\mathcal O_X)[n]\) is perverse. Regular holonomic modules and complex-constructible perverse sheaves correspond contravariantly. A merely overdetermined system need not be holonomic, and an ordinary solution sheaf need not equal its derived solution complex.}
\GSfield{Reconstruction}{Regular holonomicity is used for this safe exact equivalence; the notebook does not specify the module class or solution shift.}
\GSfield{Fate}{\citet[main theorem, p.~1]{kashiwara1980constructible} announces the derived equivalence in the lecture of 22 April 1980; \citet{kashiwara1984riemann} is the full later publication. This predates the seminar.}
\GSfield{Present status}{As of 11 November 2026: established under the stated hypotheses.}
\GSfield{Assessment}{literal statement: correct with missing hypotheses; publication match: exact; verification: read the result; development: none; confidence: high}
\end{GSitem}
\begin{GSitem}{1987-746.6}{E}{Intersection cohomology and intermediate extension}
\GSfield{Kind; source}{construction; \GSpage{P001459}}
\GSfield{The notes}{The Goresky--MacPherson combinatorial method extends sheaves to singular points and yields perverse sheaves.}
\GSfield{Statement}{If \(j:U\hookrightarrow X\) is a smooth dense stratum of complex dimension d and \(L\) a finite-rank local system, its perverse intermediate extension \(j_{!*}L[d]\) is the intersection-cohomology complex. Arbitrary extensions across singular strata need not be perverse or canonical.}
\GSfield{Reconstruction}{Intermediate extension is the precise reconstruction of the broad extension statement.}
\GSfield{Fate}{The intersection-cohomology construction is \citet{goreskymacpherson1983intersection}; the extension viewpoint is also explicit in \citet{beilinson1987glue}.}
\GSfield{Present status}{As of 11 November 2026: established under the stated hypotheses.}
\GSfield{Assessment}{literal statement: correct with missing hypotheses; publication match: exact; verification: indirect; development: none; confidence: high}
\end{GSitem}

\GStalk{747}{1988-04-18}{A. A. Beilinson}{P001461}{Semi-infinite Morse theory and relative index}
Beilinson contrasts the finite-index Morse theory of geodesics with strongly indefinite functionals, relating the latter to Atiyah's discussion of gauge theory. P001460 contains a Bourbaki reading list, not a sequence of asserted theorems. The detailed gauge-theoretic continuation occurs in record 751.

\begin{GSitem}{1987-747.1}{E}{The finite Morse index of a geodesic}
\GSfield{Kind; source}{theorem; \GSpage{P001461}}
\GSfield{The notes}{The energy functional on the path space has finite negative index, equal to the number of conjugate points along the geodesic.}
\GSfield{Statement}{For a geodesic segment in a smooth Riemannian manifold with fixed endpoints, the Hessian of energy has finite Morse index. It equals the sum of the multiplicities of conjugate points to the initial endpoint in the open segment; conjugacy at the final endpoint contributes to nullity, not to the negative index.}
\GSfield{Reconstruction}{Fixed endpoints and the endpoint/nullity convention are supplied.}
\GSfield{Fate}{\citet[\S\S1--3]{duistermaat1974morse} formulates the index in variational and Lagrangian terms. Boundary conditions and multiplicities, omitted in the board sentence, determine the exact result.}
\GSfield{Present status}{As of 11 November 2026: established under the stated hypotheses.}
\GSfield{Assessment}{literal statement: correct with missing hypotheses; publication match: exact; verification: read the result; development: none; confidence: high}
\end{GSitem}
\begin{GSitem}{1987-747.2}{B}{Relative grading for strongly indefinite functionals}
\GSfield{Kind; source}{method; \GSpage{P001461}}
\GSfield{The notes}{Critical points may have infinitely many positive and negative directions, but the numbers of negative directions are said to differ by a finite number; homology is counted from infinity over two.}
\GSfield{Statement}{For a continuous path of self-adjoint Fredholm Hessians with invertible ends and appropriate operator-domain control, spectral flow supplies a finite relative index. Infinite positive and negative indices alone do not imply that an index difference exists. A Floer complex additionally requires compactness, transversality, coherent orientation and control of broken trajectories.}
\GSfield{Reconstruction}{The Fredholm and analytic hypotheses replace subtraction of two infinite cardinalities.}
\GSfield{Fate}{\citet[pp.~286--294]{atiyah1988invariants} explains the gauge-theoretic Morse picture, circulated as a 1987 lecture/preprint. \citet{floer1988instanton}, received 3 May and published June 1988, establishes the associated homology-sphere invariant. This is exposition of a programme under development, not a general theorem for arbitrary functionals.}
\GSfield{Present status}{As of 11 November 2026: relative-index Floer theories are established under specific analytic hypotheses; the unrestricted formulation is invalid.}
\GSfield{Assessment}{literal statement: correct with missing hypotheses; publication match: partial; verification: read the result; development: none; confidence: high}
\end{GSitem}
\begin{GSitem}{1987-747.3}{E}{Relation with Donaldson invariants}
\GSfield{Kind; source}{remark; \GSpage{P001461}}
\GSfield{The notes}{The striking claim is that these invariants are related to Donaldson invariants of four-manifolds; gauge fields and flat connections are drawn.}
\GSfield{Statement}{The relevant functional on three-dimensional connections is Chern--Simons: its critical points are flat connections and its cylinder gradient equation is the anti-self-dual equation in four dimensions. Donaldson theory can therefore interact with the boundary Floer groups through cutting and gluing.}
\GSfield{Reconstruction}{The functional is reconstructed as Chern--Simons from the named source and later continuation; Yang--Mills criticality alone would be insufficient.}
\GSfield{Fate}{\citet[\S\S2--5]{atiyah1988invariants} is the named source; record 751 supplies the detailed seminar continuation. The ordinary Yang--Mills energy on a three-manifold can have nonflat critical points, so the functional in the sketch must not be identified with it without correction.}
\GSfield{Present status}{As of 11 November 2026: established under the stated hypotheses.}
\GSfield{Assessment}{literal statement: correct with missing hypotheses; publication match: partial; verification: read the result; development: none; confidence: high}
\end{GSitem}

\GStalk{748}{1988-04-18}{I. M. Gelfand}{P001461}{Polarizations beyond the loop space}
Gelfand connects positive and negative Fourier modes with tau-functions and semi-infinite cohomology, then suggests replacing the circle by another manifold. No elliptic operator, polarization class, or target invariant is specified for that extension.

\begin{GSitem}{1987-748.1}{E}{Positive and negative Fourier polarizations on the circle}
\GSfield{Kind; source}{construction; \GSpage{P001461}}
\GSfield{The notes}{Gelfand recalls decomposing \(C^\infty(S^1)\) into boundary values extending inside and outside, with applications to tau-functions and semi-infinite cohomology.}
\GSfield{Statement}{A smooth complex function on the circle splits into its Fourier modes \(\sum_{n\geq0}a_nz^n\) and \(\sum_{n<0}a_nz^n\). With constants included in both pieces the intersection is one-dimensional; with the negative part vanishing at infinity the direct-sum convention removes that overlap. Appropriate Sobolev completions give Hardy polarizations.}
\GSfield{Reconstruction}{Boundary regularity, the constant mode and the outside normalization are made explicit.}
\GSfield{Fate}{The elementary Fourier construction verifies the concrete circle example. \citet[the loop-space discussion]{atiyah1988invariants} supplies the context of relative polarizations; no new tau-function identity is asserted.}
\GSfield{Present status}{As of 11 November 2026: established under the stated hypotheses.}
\GSfield{Assessment}{literal statement: correct with missing hypotheses; publication match: thematic only; verification: read the result; development: none; confidence: high}
\end{GSitem}
\begin{GSitem}{1987-748.2}{F}{Replacing the circle by another manifold}
\GSfield{Kind; source}{remark; \GSpage{P001461}}
\GSfield{The notes}{Gelfand says that remarkably the circle can be replaced by other manifolds, and postpones the explanation.}
\GSfield{Statement}{No definite generalization is determined without a chosen operator or polarization, coefficient bundle, function space and intended semi-infinite invariant.}
\GSfield{Fate}{The adjacent Atiyah discussion motivates spectral polarizations, but this sentence alone does not identify one. No missing theorem or open problem has been manufactured from it.}
\GSfield{Present status}{As of 11 November 2026: unidentified beyond a research direction.}
\GSfield{Assessment}{literal statement: not determinable; publication match: thematic only; verification: unresolved lead; development: none; confidence: low}
\end{GSitem}

\GStalk{749}{1988-04-18}{N. Yu. Reshetikhin}{P001461-P001467}{Quantum groups, braiding, and coloured link invariants}
Reshetikhin gives the rank-one quantum enveloping algebra, coproduct and representations, then the universal R-matrix, braid actions and diagrammatic link evaluation. The core Hopf-algebra and braiding framework was public by Drinfeld's 1986 congress report. Reshetikhin's LOMI preprints E-4-87 and E-17-87 closely match the talk, but the scanned editions bear a 1988 imprint; a 1987 series number and recommendation are not an independently established circulation date. Antipode and normalization formulas in the notes need correction; generic deformation parameters and framing conventions are indispensable. The assertion that the rank-one invariants suffice has no stated separation criterion.

\begin{GSitem}{1987-749.1}{E}{The rank-one quantum enveloping algebra and its completion}
\GSfield{Kind; source}{definition; \GSpage{P001461-P001462}}
\GSfield{The notes}{The generators \(H,X^\pm\) satisfy
\[
 [H,X^\pm]=\pm2X^\pm,\qquad
 [X^+,X^-]=\frac{\sinh(hH/2)}{\sinh(h/2)},
\]
with \(q=e^h\); powers of both \(H\) and \(q^{H/4}\) are called a vector-space basis.}
\GSfield{Statement}{One coherent formulation is the h-adic algebra \(U_h(\mathfrak{sl}_2)\), topologically free over \(\mathbb C[[h]]\), with the displayed relations and PBW monomials \((X^-)^aH^b(X^+)^c\). Alternatively, for a numerical generic parameter use the algebraic Cartan generator \(K^{\pm1}\) and a PBW basis in \(X^-,K,X^+\). Exponentials are interpreted in a completion, not added as unspecified independent formal generators.}
\GSfield{Reconstruction}{The completion and base ring are fixed explicitly; the notebook does not specify which algebra is meant by its basis.}
\GSfield{Fate}{\citet[I, Definition~1.1 and Theorem~1.1]{reshetikhin1987quantizedi} gives the close publication match. Its series number is E-4-87 but the scanned imprint is 1988; the June 1987 recommendation does not by itself establish first public circulation. Independently, \citet[Example~6.2, p.~807]{drinfeld1987quantum}, a 1986 congress report published in 1987, gives the completed quantum enveloping algebra and symmetric coproduct.}
\GSfield{Present status}{As of 11 November 2026: established under the stated hypotheses.}
\GSfield{Assessment}{literal statement: correct with missing hypotheses; publication match: partial; verification: read the result; development: none; confidence: high}
\end{GSitem}
\begin{GSitem}{1987-749.2}{E}{The symmetric coproduct and counit}
\GSfield{Kind; source}{construction; \GSpage{P001462}}
\GSfield{The notes}{\(\Delta H=H\otimes1+1\otimes H\), \(\Delta X^\pm=X^\pm\otimes q^{H/4}+q^{-H/4}\otimes X^\pm\); this is an algebra homomorphism and coassociative.}
\GSfield{Statement}{Put \(L=e^{hH/4}\), so \(\Delta L=L\otimes L\). The displayed coproduct extends to the completed algebra, is coassociative, and has counit \(\epsilon(H)=\epsilon(X^\pm)=0\), \(\epsilon(1)=1\).}
\GSfield{Reconstruction}{All tensor products use the same h-adic completion.}
\GSfield{Fate}{These formulas are \citet[I, Theorem~1.1, equations~(1.4)--(1.5)]{reshetikhin1987quantizedi}. Coassociativity on \(X\) follows by expanding both sides as \(X\otimes L\otimes L+L^{-1}\otimes X\otimes L+L^{-1}\otimes L^{-1}\otimes X\). The corresponding pre-seminar structure is also explicit in \citet[Example~6.2, p.~807]{drinfeld1987quantum} (1986 lecture; 1987 publication).}
\GSfield{Present status}{As of 11 November 2026: established under the stated hypotheses.}
\GSfield{Assessment}{literal statement: correct with missing hypotheses; publication match: exact; verification: read the result; development: none; confidence: high}
\end{GSitem}
\begin{GSitem}{1987-749.3}{E}{The antipode needs a minus sign and half the exponent}
\GSfield{Kind; source}{claim; \GSpage{P001462}}
\GSfield{The notes}{The scan states \(S(X^\pm)=q^{H/2}X^\pm q^{-H/2}\), with \(S(H)=-H\).}
\GSfield{Statement}{For the displayed coproduct, \(S(X^\pm)=-L X^\pm L^{-1}=-q^{\pm1/2}X^\pm\), where \(L=q^{H/4}\). Indeed \(m(S\otimes\mathrm{id})\Delta X=S(X)L+LX=0\). The source expression fails this identity and has the wrong classical limit.}
\GSfield{Reconstruction}{\oursComputed{} }
\GSfield{Fate}{\citet[I, Theorem~1.1, equation~(1.6), with \(\rho=H/2\)]{reshetikhin1987quantizedi} confirms the normalization. The correction is to a faithfully transcribed formula.}
\GSfield{Present status}{As of 11 November 2026: the corrected antipode is established; the printed one is incompatible with the printed coproduct.}
\GSfield{Assessment}{literal statement: false as written; publication match: exact; verification: read the result; development: annotated; confidence: high}
\end{GSitem}
\begin{GSitem}{1987-749.4}{E}{The Cartan-generator substitution changes the denominator}
\GSfield{Kind; source}{definition; \GSpage{P001462}}
\GSfield{The notes}{After putting \(K=q^{H/2}\), the notes keep \(KX^\pm=q^{\pm1}X^\pm K\), but write the commutator denominator \(q-q^{-1}\).}
\GSfield{Statement}{With the unchanged parameter \(q=e^h\), the initial relation becomes \([X^+,X^-]=(K-K^{-1})/(q^{1/2}-q^{-1/2})\). If instead the standard parameter is \(t=q^{1/2}\), then \(KX^\pm=t^{\pm2}X^\pm K\) and the denominator is \(t-t^{-1}\). A rescaling of generators is another possible convention, but must be stated.}
\GSfield{Reconstruction}{\oursComputed{} The correction keeps the original generators and parameter; no silent rescaling is made.}
\GSfield{Fate}{Compare the fixed relation in \citet[I, equation~(1.2)]{reshetikhin1987quantizedi}. Direct substitution detects the factor mismatch.}
\GSfield{Present status}{As of 11 November 2026: established under the stated hypotheses.}
\GSfield{Assessment}{literal statement: false as written; publication match: exact; verification: read the result; development: annotated; confidence: high}
\end{GSitem}
\begin{GSitem}{1987-749.5}{E}{Finite-dimensional spin representations and the classical limit}
\GSfield{Kind; source}{construction; \GSpage{P001463}}
\GSfield{The notes}{\((X^+)^*=X^-\), \(H^*=H\), and a square-root formula is given on \(e_m^j\), with dimension \(2j+1\); Gelfand asks for the \(q=1\) limit.}
\GSfield{Statement}{For real \(q>0\), define \([a]=(q^{a/2}-q^{-a/2})/(q^{1/2}-q^{-1/2})\), with its limiting value at 1. For \(j\in\frac12\mathbb Z_{\geq0}\), \(m=-j,-j+1,\ldots,j\), let \(He_m=2me_m\) and \(X^\pm e_m=\sqrt{[j\mp m][j\pm m+1]}e_{m\pm1}\). These give the compact real-form spin modules and tend to the usual \(\mathfrak{su}_2\) formulas as \(q\to1\).}
\GSfield{Reconstruction}{Nonnegative half-integer j, unit steps for m, and a positive real parameter for the unitary interpretation are restored.}
\GSfield{Fate}{The highest-weight framework is \citet[I, Propositions~1.1--1.2]{reshetikhin1987quantizedi}; the rank-one coefficients can be checked by substitution into the relations.}
\GSfield{Present status}{As of 11 November 2026: established under the stated hypotheses.}
\GSfield{Assessment}{literal statement: correct with missing hypotheses; publication match: exact; verification: read the result; development: none; confidence: high}
\end{GSitem}
\begin{GSitem}{1987-749.6}{E}{Expanding vector fields on the sphere}
\GSfield{Kind; source}{question; \GSpage{P001463}; I. M. Gelfand}
\GSfield{The notes}{Gelfand asks how to expand vector fields on a sphere in spherical harmonics, beginning with \(SU(2)\). The last word is partly unclear.}
\GSfield{Statement}{For the standard oriented round \(S^2\), a smooth tangent field is expanded in \(\nabla_{S^2}Y_{\ell m}\) and \(n\times\nabla_{S^2}Y_{\ell m}\), \(\ell\geq1\). There is no harmonic one-form term because \(H^1(S^2)=0\). For an ambient \(\mathbb R^3\)-valued field also include the radial terms \(nY_{\ell m}\). Smooth expansions converge in the smooth topology with the usual coefficient decay.}
\GSfield{Reconstruction}{The standard round sphere and tangent/ambient distinction are supplied; they are not explicit in the abbreviated question.}
\GSfield{Fate}{The scalar normalization is given in \citet[\S14.30]{nist2026dlmf}; gradient and rotated-gradient components follow from the Hodge decomposition on the sphere. This answers a classical exercise, not a new quantum-group problem.}
\GSfield{Present status}{As of 11 November 2026: established under the stated hypotheses.}
\GSfield{Assessment}{literal statement: correct with missing hypotheses; publication match: partial; verification: indirect; development: none; confidence: medium}
\end{GSitem}
\begin{GSitem}{1987-749.7}{E}{The universal R series needs deformation factors and a completion}
\GSfield{Kind; source}{claim; \GSpage{P001463}}
\GSfield{The notes}{The formula has \(R=e^{hH\otimes H/4}\sum_{n\geq0}q^n e^n\otimes f^n/[n]_q!\), with \(e=q^{H/4}X^+\), \(f=q^{-H/4}X^-\).}
\GSfield{Statement}{The formal quasitriangular structure has \(R\in U_h\widehat\otimes U_h\), \(R\equiv1\pmod h\), and intertwines the two coproducts. Its nth nilpotent term contains a factor vanishing to order n in h, conventionally a power of \(q^{1/2}-q^{-1/2}\). The printed series with the printed e,f cannot have the required classical limit.}
\GSfield{Reconstruction}{\oursComputed{} The exact missing coefficient is not guessed; only the required deformation order and completion are asserted.}
\GSfield{Fate}{\citet[I, Theorems~1.2--1.3]{reshetikhin1987quantizedi} and \citet[\S3]{reshetikhinturaev1990ribbon} give the quasitriangular construction. The q-exponent depends on generator conventions; this chapter does not invent an exact corrected coefficient from the abbreviated series. The corresponding pre-seminar structure is also explicit in \citet[\S13, pp.~816--817]{drinfeld1987quantum} (1986 lecture; 1987 publication).}
\GSfield{Present status}{As of 11 November 2026: established under the stated hypotheses.}
\GSfield{Assessment}{literal statement: false as written; publication match: partial; verification: read the result; development: annotated; confidence: high}
\end{GSitem}
\begin{GSitem}{1987-749.8}{E}{Quasitriangular identities and the Yang--Baxter equation}
\GSfield{Kind; source}{theorem; \GSpage{P001463}}
\GSfield{The notes}{The intertwining equation and Yang--Baxter equation are stated, but both coproduct-factor identities have their order reversed.}
\GSfield{Statement}{With \(\Delta^{\mathrm{op}}(a)=R\Delta(a)R^{-1}\), use \((\Delta\otimes\mathrm{id})R=R_{13}R_{23}\), \((\mathrm{id}\otimes\Delta)R=R_{13}R_{12}\). These imply \(R_{12}R_{13}R_{23}=R_{23}R_{13}R_{12}\). Opposite conventions require changing the accompanying definitions consistently.}
\GSfield{Reconstruction}{The correction fixes the intertwining convention already printed in the notes.}
\GSfield{Fate}{The stated order is exactly \citet[I, Theorem~1.3, equations~(1.8a)--(1.8c)]{reshetikhin1987quantizedi}, also \citet[equations~(3.1.1)--(3.1.3)]{reshetikhinturaev1990ribbon}. The corresponding pre-seminar structure is also explicit in \citet[\S10, equations~(20)--(21)]{drinfeld1987quantum} (1986 lecture; 1987 publication).}
\GSfield{Present status}{As of 11 November 2026: established under the stated hypotheses.}
\GSfield{Assessment}{literal statement: false as written; publication match: exact; verification: read the result; development: annotated; confidence: high}
\end{GSitem}
\begin{GSitem}{1987-749.9}{E}{The q-gamma function}
\GSfield{Kind; source}{question; \GSpage{P001463}; I. M. Gelfand}
\GSfield{The notes}{After the q-factorial definition Gelfand asks what \(\Gamma_q\) is.}
\GSfield{Statement}{For \(0<q<1\), \(\Gamma_q(z)=(1-q)^{1-z}(q;q)_\infty/(q^z;q)_\infty\), with \((a;q)_\infty=\prod_{r\geq0}(1-aq^r)\). It satisfies \(\Gamma_q(z+1)=(1-q^z)\Gamma_q(z)/(1-q)\), \(\Gamma_q(1)=1\), and tends to \(\Gamma(z)\) as \(q\uparrow1\) away from poles.}
\GSfield{Reconstruction}{The choice \(0<q<1\) and the nonsymmetric q-factorial convention make the exercise precise.}
\GSfield{Fate}{The standard q-gamma definition and recurrence are tabulated in \citet[\S5.18, equations~(5.18.4)--(5.18.7)]{nist2026dlmf}. The symmetric q-number in the preceding quantum algebra differs by an elementary q-power normalization.}
\GSfield{Present status}{As of 11 November 2026: established under the stated hypotheses.}
\GSfield{Assessment}{literal statement: correct with missing hypotheses; publication match: exact; verification: read the result; development: none; confidence: high}
\end{GSitem}
\begin{GSitem}{1987-749.10}{E}{Braiding and braid-group representations}
\GSfield{Kind; source}{construction; \GSpage{P001464}}
\GSfield{The notes}{\(R^{j_1j_2}=P(\pi_{j_1}\otimes\pi_{j_2})(R)\), mapping the tensor factors in the opposite order, satisfies the braid relation.}
\GSfield{Statement}{For finite-dimensional modules in the chosen quasitriangular category, \(c_{V,W}=P\circ(\pi_V\otimes\pi_W)(R)\) is a natural module isomorphism \(V\otimes W\to W\otimes V\). Adjacent c operators on a repeated tensor power obey the braid relations and define a braid-group representation.}
\GSfield{Reconstruction}{On differently coloured strands one has a coloured braid-groupoid action; an ordinary braid group acts on a fixed repeated tensor power.}
\GSfield{Fate}{\citet[I, \S1]{reshetikhin1987quantizedi} supplies the R-matrix framework; \citet[\S\S3--5]{reshetikhinturaev1990ribbon} formalizes the graphical functor. The corresponding pre-seminar structure is also explicit in \citet[\S10, pp.~811--812]{drinfeld1987quantum} (1986 lecture; 1987 publication).}
\GSfield{Present status}{As of 11 November 2026: established under the stated hypotheses.}
\GSfield{Assessment}{literal statement: correct with missing hypotheses; publication match: exact; verification: read the result; development: none; confidence: high}
\end{GSitem}
\begin{GSitem}{1987-749.11}{E}{Generic Clebsch--Gordan decomposition}
\GSfield{Kind; source}{theorem; \GSpage{P001464}}
\GSfield{The notes}{All finite-dimensional representations are called completely reducible, and the tensor product is the classical spin sum.}
\GSfield{Statement}{For a generic parameter, restricted to finite-dimensional type-1 modules, \(V^{j_1}\otimes V^{j_2}\simeq\bigoplus_{r=0}^{2\min(j_1,j_2)}V^{j_1+j_2-r}\). This category is semisimple. Root-of-unity representation categories need not be semisimple, so the omitted parameter restriction changes the theorem.}
\GSfield{Reconstruction}{Generic q and type-1 are specified, and the summation proceeds in steps of one.}
\GSfield{Fate}{\citet[I, Propositions~1.1--1.2 and Theorem~1.5]{reshetikhin1987quantizedi} gives the highest-weight and tensor-product setting. The seminar's Clebsch--Gordan coefficients are change-of-basis matrices for this decomposition.}
\GSfield{Present status}{As of 11 November 2026: established under the stated hypotheses.}
\GSfield{Assessment}{literal statement: correct with missing hypotheses; publication match: exact; verification: read the result; development: none; confidence: high}
\end{GSitem}
\begin{GSitem}{1987-749.12}{E}{Graphical tensor calculus, duality, and ribbon structure}
\GSfield{Kind; source}{construction; \GSpage{P001464-P001466}}
\GSfield{The notes}{Composition joins diagrams; straight strands denote identities, crossings R and its inverse, and cups/caps have signed q-weight formulas.}
\GSfield{Statement}{In a rigid braided tensor category, composition and tensor product are represented by stacking and juxtaposition. Evaluation and coevaluation must satisfy the two zigzag identities; a ribbon twist supplies compatibility with framed isotopy. Explicit cup/cap coefficients depend on chosen identifications with duals and bases, particularly for half-integer spins.}
\GSfield{Reconstruction}{The categorical identities are identified precisely; no unique phase normalization is inferred from the fragmentary coefficients.}
\GSfield{Fate}{\citet[\S\S2--5]{reshetikhinturaev1990ribbon} proves the precise ribbon-graph functor. Reshetikhin's 1987-numbered LOMI preprints (1988 imprint) present the quantum-group graphical programme. The two inconsistent-looking q-powers in the notes cannot be repaired uniquely without those basis choices. The scanned LOMI editions carry a 1988 imprint despite their 1987 series identifiers; an exact public-circulation date before this talk is not established by that cover evidence.}
\GSfield{Present status}{As of 11 November 2026: established under the stated hypotheses.}
\GSfield{Assessment}{literal statement: correct with missing hypotheses; publication match: partial; verification: read the result; development: none; confidence: medium}
\end{GSitem}
\begin{GSitem}{1987-749.13}{E}{The coloured link state sum and framing correction}
\GSfield{Kind; source}{construction; \GSpage{P001467}}
\GSfield{The notes}{Sum products of matrix entries over spin labels on elementary diagram pieces; multiply by a q-power involving under/over counts to obtain a link invariant containing Jones's invariant.}
\GSfield{Statement}{The ribbon-category evaluation gives an invariant of coloured framed oriented links. For a blackboard-framed diagram, pass to the zero-framed version by multiplying the contribution of each component coloured by j by \(\theta_j^{-w_i}\), where \(\theta_j\) is its twist scalar and \(w_i\) its signed self-writhe, in a fixed convention. An unweighted under/over count is not the required spin-dependent twist correction. The fundamental \(\mathfrak{sl}_2\) colour yields the Jones invariant up to standard variable and normalization choices.}
\GSfield{Reconstruction}{Framed versus ambient isotopy, signed self-writhe and colour-dependent twist are explicit repairs.}
\GSfield{Fate}{\citet[\S\S5--6]{reshetikhinturaev1990ribbon} gives the functor and link evaluation. The 1987-numbered preprint (1988 imprint) \citet{reshetikhin1987quantizedii} establishes the earlier link-invariant programme. The notebook's displayed normalization is not copied as a valid formula. The scanned LOMI editions carry a 1988 imprint despite their 1987 series identifiers; an exact public-circulation date before this talk is not established by that cover evidence.}
\GSfield{Present status}{As of 11 November 2026: established under the stated hypotheses.}
\GSfield{Assessment}{literal statement: false as written; publication match: partial; verification: read the result; development: annotated; confidence: medium}
\end{GSitem}
\begin{GSitem}{1987-749.14}{E}{Other simple Lie algebras and 6j recoupling}
\GSfield{Kind; source}{construction; \GSpage{P001467}}
\GSfield{The notes}{Every simple Lie algebra supplies link invariants; 6j-symbols are also mentioned.}
\GSfield{Statement}{For each complex simple Lie algebra at generic deformation parameter, finite-dimensional highest-weight modules and the universal R-matrix yield a family of coloured ribbon-link invariants. Associativity and chosen decompositions of triple tensor products produce quantum 6j change-of-basis matrices, with multiplicity labels when tensor products are not multiplicity-free.}
\GSfield{Reconstruction}{The category and multiplicity conventions are supplied; the brief 6j remark does not state a numerical formula.}
\GSfield{Fate}{\citet[I, Introduction and \S1]{reshetikhin1987quantizedi} states the all-simple-Lie-algebra construction. \citet[\S\S5--6]{reshetikhinturaev1990ribbon} supplies the general ribbon formulation. The scanned LOMI editions carry a 1988 imprint despite their 1987 series identifiers; an exact public-circulation date before this talk is not established by that cover evidence.}
\GSfield{Present status}{As of 11 November 2026: established under the stated hypotheses.}
\GSfield{Assessment}{literal statement: correct with missing hypotheses; publication match: exact; verification: read the result; development: none; confidence: medium}
\end{GSitem}
\begin{GSitem}{1987-749.15}{F}{What does sufficiency of the SL2 invariants mean?}
\GSfield{Kind; source}{conjecture; \GSpage{P001467}}
\GSfield{The notes}{There is a conjecture that the \(SL_2\)-invariants suffice; no object or equivalence relation follows the verb.}
\GSfield{Statement}{If interpreted as separation of all oriented knot types by all colours of the Jones polynomial, the assertion is false. But the notes could instead mean recovery of another family of quantum invariants, or another weaker property; the exact conjecture is not recoverable.}
\GSfield{Reconstruction}{The separation interpretation is deliberately conditional and is not entered as the original conjecture in the problem index.}
\GSfield{Fate}{\citet[Theorem~1.1]{stoimenowtanaka2008mutation} constructs infinitely many pairs of hyperbolic knots with equal coloured Jones polynomials that are not mutants. The paper records a first manuscript in December 2005 and arXiv circulation in 2006; the opened version is from 2008. This refutes the explicit separation interpretation, not every possible reading of ``suffice.}
\GSfield{Present status}{As of 11 November 2026: the exact statement is unidentified; the most natural complete-knot-invariant interpretation has published counterexamples.}
\GSfield{Assessment}{literal statement: not determinable; publication match: partial; verification: read the result; development: annotated; confidence: medium}
\end{GSitem}

\GStalk{750}{1988-04-25}{A. G. Ushveridze}{P001468-P001472}{Riccati equations and quasi-exact solvability}
The talk constructs finite families of Riccati solutions from poles satisfying Bethe-type equations, with rational, trigonometric and elliptic addition laws. It discusses electrostatics, common eigenspaces and conversion to Schrodinger equations. The author's 1989 review is a close primary source. Two addition-law constants and the root equations in the manuscript require explicit repair.

\begin{GSitem}{1987-750.1}{A}{A rational pole ansatz and its Bethe equations}
\GSfield{Kind; source}{construction; \GSpage{P001468}}
\GSfield{The notes}{The Riccati equation is \(y^{\prime}+y^2+2By+C=0\), with \(B=\sum b_\alpha\Delta(\lambda-a_\alpha)\), \(C=\sum c_\alpha\Delta(\lambda-a_\alpha)\), and \(y=\sum_i\Delta(\lambda-\lambda_i)\). The first root equation has inconsistent indices and omits a \(\Delta\).}
\GSfield{Statement}{For \(\Delta(z)=1/z\), distinct fixed points and distinct roots avoiding them, the ansatz solves the equation exactly iff \(\sum_{j\ne i}(\lambda_i-\lambda_j)^{-1}+\sum_\alpha b_\alpha(\lambda_i-a_\alpha)^{-1}=0\) for each i and \(c_\alpha=2b_\alpha\sum_i(\lambda_i-a_\alpha)^{-1}\). Double poles cancel automatically; these conditions cancel all simple poles, and the remaining proper rational function vanishes.}
\GSfield{Reconstruction}{\oursComputed{} The exact equivalence is deliberately restricted to the rational kernel. Other kernels require additional constant/elliptic compatibility conditions absent from the notebook.}
\GSfield{Fate}{\citet[\S3, Theorem~4, pp.~1197--1198, and \S4]{ushveridze1989quasiexact} develops the analytic pole method, referring back to the author's 1988 FIAN announcements. The residue calculation gives an exact match in the rational specialization, not an unqualified elliptic theorem.}
\GSfield{Present status}{As of 11 November 2026: established under the stated hypotheses.}
\GSfield{Assessment}{literal statement: false as written; publication match: partial; verification: read the result; development: annotated; confidence: high}
\end{GSitem}
\begin{GSitem}{1987-750.2}{E}{The addition identity and its corrected special functions}
\GSfield{Kind; source}{claim; \GSpage{P001468}}
\GSfield{The notes}{For \(x+y+z=0\), \(\Delta(x)\Delta(y)+\Delta(y)\Delta(z)+\Delta(z)\Delta(x)+\Omega(x)+\Omega(y)+\Omega(z)=0\), with odd \(\Delta\) and even \(\Omega\). The trigonometric constant is printed positive and the elliptic expression lacks a half.}
\GSfield{Statement}{Correct examples are \((\Delta,\Omega)=(1/z,0)\), \((\omega\cot(\omega z),-\omega^2/3)\), and \((\zeta(z),(\zeta(z)^2-\wp(z))/2)\). In the last case use \((\zeta(x)+\zeta(y)+\zeta(z))^2=\wp(x)+\wp(y)+\wp(z)\) when \(x+y+z=0\).}
\GSfield{Reconstruction}{\oursComputed{} }
\GSfield{Fate}{The addition identity is classical; \citet[\S23.10]{nist2026dlmf} gives the Weierstrass identity. \citet[\S3 and \S11]{ushveridze1989quasiexact} places such scalar triangle identities in the quasi-exactly solvable construction. Direct substitution checks both emendations.}
\GSfield{Present status}{As of 11 November 2026: the corrected examples satisfy the stated equation; the two literal examples do not.}
\GSfield{Assessment}{literal statement: false as written; publication match: partial; verification: read the result; development: annotated; confidence: high}
\end{GSitem}
\begin{GSitem}{1987-750.3}{E}{A local Riccati identity from the addition law}
\GSfield{Kind; source}{theorem; \GSpage{P001469}}
\GSfield{The notes}{A solution with a singularity at zero satisfies \(\Delta^{\prime}+\Delta^2=2\Omega+\Omega(0)\).}
\GSfield{Statement}{Assume \(\Delta\) is odd meromorphic with \(\Delta(t)=1/t+O(t)\), and \(\Omega\) even and holomorphic at zero. Put \(y=\varepsilon\), \(z=-x-\varepsilon\) in the addition law and let \(\varepsilon\to0\). The finite part gives exactly \(\Delta^{\prime}(x)+\Delta(x)^2=2\Omega(x)+\Omega(0)\). A different pole residue changes the derivative coefficient.}
\GSfield{Reconstruction}{\oursComputed{} A simple pole with residue one and regularity of \(\Omega\) at zero are indispensable extra hypotheses.}
\GSfield{Fate}{The normalization argument for the constant-\(\Omega\) scalar triangle equation is given in \citet[\S3, equation~(27), p.~1197]{ushveridze1989quasiexact}. The limiting expansion supplies the displayed general form.}
\GSfield{Present status}{As of 11 November 2026: established under the stated hypotheses.}
\GSfield{Assessment}{literal statement: correct with missing hypotheses; publication match: partial; verification: read the result; development: none; confidence: high}
\end{GSitem}
\begin{GSitem}{1987-750.4}{F}{The claimed general kernel family is underspecified}
\GSfield{Kind; source}{claim; \GSpage{P001469}}
\GSfield{The notes}{The general solution is written \(\Delta_0(b\lambda)+a\lambda\), with \(\Delta_0\) one of the rational, trigonometric or elliptic examples.}
\GSfield{Statement}{A valid normalized family is obtained from a residue-one kernel by the rescaling \(b\Delta_0(b\lambda)\) and a linear addition, together with the correspondingly transformed \(\Omega\). A classification of every solution additionally requires regularity, pole and degeneracy hypotheses. The notebook neither gives these nor states the transformed \(\Omega\).}
\GSfield{Reconstruction}{The prefactor b is needed only when maintaining residue one; the original unnormalized family can still satisfy an appropriately rescaled functional equation.}
\GSfield{Fate}{\citet[\S3 and the final elliptic discussion in \S11]{ushveridze1989quasiexact} confirms the functional-equation method, but the retrieved review does not supply an exact theorem with the notebook's unrestricted quantifiers.}
\GSfield{Present status}{As of 11 November 2026: the displayed families are meaningful, but the claimed exhaustive classification is not identified precisely enough to adopt.}
\GSfield{Assessment}{literal statement: not determinable; publication match: partial; verification: read the result; development: none; confidence: medium}
\end{GSitem}
\begin{GSitem}{1987-750.5}{E}{Logarithmic derivatives linearize the Riccati equation}
\GSfield{Kind; source}{method; \GSpage{P001469-P001470}}
\GSfield{The notes}{Putting \(y=f^{\prime}/f\) is said to give a Schrodinger equation; the later page writes \(f^{\prime\prime}+2Bf^{\prime}+Cf=0\).}
\GSfield{Statement}{On an interval where \(f\ne0\), the substitution is equivalent to \(f^{\prime\prime}+2Bf^{\prime}+Cf=0\). Conversely, any Riccati solution gives a local nonzero f by exponentiating its integral. A gauge and coordinate change are still needed to obtain Schrodinger normal form; the first substitution alone leaves a first derivative.}
\GSfield{Reconstruction}{\oursComputed{} The local nonvanishing condition and distinction between a linear ODE and Schrodinger normal form are explicit.}
\GSfield{Fate}{\citet[\S\S2--3]{ushveridze1989quasiexact} gives the analytic normal-form method. The equivalence follows directly from \(y^{\prime}+y^2=f^{\prime\prime}/f\).}
\GSfield{Present status}{As of 11 November 2026: established under the stated hypotheses.}
\GSfield{Assessment}{literal statement: correct with missing hypotheses; publication match: exact; verification: read the result; development: none; confidence: high}
\end{GSitem}
\begin{GSitem}{1987-750.6}{A}{Electrostatic interpretation of the rational roots}
\GSfield{Kind; source}{construction; \GSpage{P001469}}
\GSfield{The notes}{The fixed charges are \(b_\alpha\) at \(a_\alpha\), moving charges are one at \(\lambda_i\), and the potential is written as a positive logarithm. Trigonometric and theta analogues are mentioned.}
\GSfield{Statement}{For real distinct ordered roots and real fixed points, the rational Bethe equations are the stationarity equations of \(U=-\sum_{i<j}\log|\lambda_i-\lambda_j|-\sum_{i,\alpha}b_\alpha\log|\lambda_i-a_\alpha|\). Positive b produces repulsion with this sign. The logarithm of sine gives the trigonometric force; an elliptic force must use a sigma/theta logarithmic derivative with its possible linear term specified.}
\GSfield{Reconstruction}{\oursComputed{} The real domain and sign convention are fixed. The elliptic theta kernel is not silently equated with Weierstrass zeta without its linear correction.}
\GSfield{Fate}{\citet[\S5, equations~(59)--(63), pp.~1210--1212]{ushveridze1989quasiexact} gives the electrostatic interpretation and real-interval stability assumptions. A positive logarithmic potential with minimization reverses the repulsive-energy sign.}
\GSfield{Present status}{As of 11 November 2026: established under the stated hypotheses.}
\GSfield{Assessment}{literal statement: correct with missing hypotheses; publication match: partial; verification: read the result; development: annotated; confidence: high}
\end{GSitem}
\begin{GSitem}{1987-750.7}{A}{Quantum and classical stability are not an unrestricted equivalence}
\GSfield{Kind; source}{claim; \GSpage{P001469}}
\GSfield{The notes}{The same \(\lambda_i\) are called wavefunction nodes, charge positions and Bethe rapidities; quantum stability is asserted equivalent to Coulomb stability.}
\GSfield{Statement}{The algebraic root equations admit these interpretations in the matched model. Spectral admissibility also depends on the Hilbert-space measure, self-adjoint domain, and endpoint behaviour. It is not equivalent merely to existence of a stationary electrostatic configuration.}
\GSfield{Reconstruction}{Only the shared algebraic Bethe equations are identified; physical stability notions require separate definitions.}
\GSfield{Fate}{\citet[\S5, p.~1212, comparing equations~(45)--(46) with (62)--(63)]{ushveridze1989quasiexact} explicitly says the classical and quantum interval conditions almost coincide, with shifted numerical thresholds in that paper's normalization. Thus the author's later primary account itself prevents the literal biconditional.}
\GSfield{Present status}{As of 11 November 2026: the model-dependent correspondence is established; the unqualified stability equivalence is false.}
\GSfield{Assessment}{literal statement: false as written; publication match: partial; verification: read the result; development: annotated; confidence: high}
\end{GSitem}
\begin{GSitem}{1987-750.8}{A}{A common eigenspace yields partial exact solvability}
\GSfield{Kind; source}{method; \GSpage{P001469-P001470}}
\GSfield{The notes}{For commuting \(H_0,H_1,H_2\), form \(H=H_0+L_1(H_1-c_1)+L_2(H_2-c_2)\). Common eigenvectors with the chosen eigenvalues remain known.}
\GSfield{Statement}{On a common invariant domain, let \(\mathcal V=\ker(H_1-c_1)\cap\ker(H_2-c_2)\) and suppose \(H_0\mathcal V\subset\mathcal V\). Then \(H|_{\mathcal V}=H_0|_{\mathcal V}\). In particular every known \(H_0\)-eigenvector in \(\mathcal V\) remains an H-eigenvector; no commutation of the \(L_j\) with \(H_i\) is needed. A finite-dimensional nonzero \(\mathcal V\) provides an algebraically solvable sector, not the full spectrum.}
\GSfield{Reconstruction}{Domain and invariance hypotheses are supplied; self-adjointness of H is an additional analytic requirement, not automatic from the formula.}
\GSfield{Fate}{\citet[\S9, equation~(113), p.~1232]{ushveridze1989quasiexact} states exactly this operator construction, in a general finite sum. The review appeared in 1989 after the FIAN work cited in the notes.}
\GSfield{Present status}{As of 11 November 2026: established under the stated hypotheses.}
\GSfield{Assessment}{literal statement: correct with missing hypotheses; publication match: exact; verification: read the result; development: none; confidence: high}
\end{GSitem}
\begin{GSitem}{1987-750.9}{F}{The nonlinear-Schrodinger and oscillator analogy}
\GSfield{Kind; source}{remark; \GSpage{P001470}}
\GSfield{The notes}{Two lines compare quantum nonlinear Schrodinger theory with the harmonic oscillator, and energy with a particle-number operator; spectra are said to come from Bethe equations.}
\GSfield{Statement}{The Hamiltonians, representation spaces, boundary conditions and precise correspondence are absent, so no exact spectral equivalence can be stated.}
\GSfield{Fate}{\citet[\S\S7--9]{ushveridze1989quasiexact} supplies related integrable-spin and quasi-exact constructions, but does not identify these two bare correspondences as written.}
\GSfield{Present status}{As of 11 November 2026: unidentified beyond the integrable-system analogy.}
\GSfield{Assessment}{literal statement: not determinable; publication match: thematic only; verification: read the result; development: none; confidence: low}
\end{GSitem}
\begin{GSitem}{1987-750.10}{E}{The Liouville transform to Schrodinger form}
\GSfield{Kind; source}{construction; \GSpage{P001470}}
\GSfield{The notes}{Multiply by an arbitrary A, set \(x=\int d\lambda/\sqrt A\), eliminate the first derivative, and obtain a potential depending on A.}
\GSfield{Statement}{On a real interval with \(A(\lambda)>0\), set \(dx/d\lambda=A^{-1/2}\), \(P=2B\sqrt A-A^{\prime}/(2\sqrt A)\), and \(\Psi=A^{-1/4}\exp(\int B\,d\lambda)f\). Then \(-\Psi_{xx}+V\Psi=0\), where \(V=\tfrac12 P_x+\tfrac14P^2-AC\). Complex versions need a local branch of \(\sqrt A\).}
\GSfield{Reconstruction}{\oursComputed{} The gauge factor, positivity/branch condition and derivatives are made explicit; a coordinate change alone does not remove the first derivative.}
\GSfield{Fate}{\citet[\S2, transformation formulas leading to equation~(24)]{ushveridze1989quasiexact} develops the coordinate/gauge procedure. The displayed formula is a direct reconstruction from the notebook's equation (12).}
\GSfield{Present status}{As of 11 November 2026: established under the stated hypotheses.}
\GSfield{Assessment}{literal statement: correct with missing hypotheses; publication match: exact; verification: read the result; development: none; confidence: high}
\end{GSitem}
\begin{GSitem}{1987-750.11}{A}{Separating a spectral constant from the potential}
\GSfield{Kind; source}{claim; \GSpage{P001470-P001471}}
\GSfield{The notes}{The product \(A=\prod_\alpha\Xi(\lambda-a_\alpha)\), \(\Xi^{\prime}/\Xi=\Delta\), is said to put all root information into an x-independent energy.}
\GSfield{Statement}{For the rational kernel take A proportional to \(\prod_\alpha(\lambda-a_\alpha)\). Then \(AC\) is polynomial of degree at most \(N-2\), because the Bethe equations imply \(\sum c_\alpha=0\). For \(N=2\) it is a constant energy. For \(N>2\), additional root-dependent coefficients generally remain in the potential and must be fixed to obtain several eigenstates of one Hamiltonian.}
\GSfield{Reconstruction}{\oursComputed{} A nonzero scalar multiple can be chosen for positivity on the desired real interval; the source' s bare product may have negative sign there.}
\GSfield{Fate}{\citet[\S4, equations~(38)--(43), pp.~1199--1201]{ushveridze1989quasiexact} distinguishes the infinitely degenerate spectral parameter from the others and makes the N-dependent distinction. The product choice alone does not prove an all-N single-energy separation.}
\GSfield{Present status}{As of 11 November 2026: established under the stated hypotheses.}
\GSfield{Assessment}{literal statement: correct with missing hypotheses; publication match: partial; verification: read the result; development: annotated; confidence: high}
\end{GSitem}
\begin{GSitem}{1987-750.12}{A}{Two fixed positive charges and Jacobi polynomials}
\GSfield{Kind; source}{theorem; \GSpage{P001471}}
\GSfield{The notes}{For \(N=2\), \(b_1,b_2>0\), repulsion traps M particles between \(a_1<a_2\), giving a solution with M nodes for every M and an exactly solvable model.}
\GSfield{Statement}{Let \(a_1<a_2\) be real, \(b_1,b_2>0\), and \(M\geq0\) an integer. On the chamber \(a_1<\lambda_1<\cdots<\lambda_M<a_2\), the negative-log energy is strictly convex and tends to infinity at its boundary, hence has a unique minimizer. The monic polynomial \(f_M=\prod_i(\lambda-\lambda_i)\) satisfies \(Af_M^{\prime\prime}+2BAf_M^{\prime}+E_Mf_M=0\), \(A=(\lambda-a_1)(a_2-\lambda)\), \(E_M=M(M-1+2b_1+2b_2)\). Here \(B(\lambda)=b_1/(\lambda-a_1)+b_2/(\lambda-a_2)\); for M=0 the configuration is empty and \(f_0=1\). These are rescaled Jacobi polynomials, orthogonal with weight \((\lambda-a_1)^{2b_1-1}(a_2-\lambda)^{2b_2-1}\).}
\GSfield{Reconstruction}{\ours{} The chosen chamber, energy sign, self-adjoint weighted realization and positive A fix the implicit analytic setup.}
\GSfield{Fate}{\citet[\S4, N=2 case, and \S5]{ushveridze1989quasiexact} gives the exactly-solvable and electrostatic interpretations. The convexity and coefficient comparison justify this precise reconstruction; physical boundary conditions after the Liouville transform must still be selected.}
\GSfield{Present status}{As of 11 November 2026: established under the stated hypotheses.}
\GSfield{Assessment}{literal statement: correct with missing hypotheses; publication match: partial; verification: read the result; development: none; confidence: high}
\end{GSitem}
\begin{GSitem}{1987-750.13}{A}{Exterior intervals and finitely many admissible levels}
\GSfield{Kind; source}{claim; \GSpage{P001471-P001472}}
\GSfield{The notes}{For M particles to the right of the second fixed charge, the stability condition is written \(b_1+b_2+M<0\); failure for large M yields only finitely many levels.}
\GSfield{Statement}{With \(b_2>0\), the displayed inequality is a sufficient confining condition for the negative-log energy on the exterior chamber, but it is not generally the sharp threshold. Already for M=1, \(a_1=0,a_2=1,b_1=-1,b_2=1/2\), \(U=\log\lambda-\tfrac12\log(\lambda-1)\) has a strict minimum at \(\lambda=2\), although \(b_1+b_2+M=1/2\). Finiteness of normalizable quantum states requires its own endpoint and measure analysis.}
\GSfield{Reconstruction}{\oursComputed{} The counterexample only rejects necessity of the displayed condition; it does not reject that condition as a sufficient bound.}
\GSfield{Fate}{\citet[\S5, equations~(62)--(63) and the following comparison, pp.~1211--1212]{ushveridze1989quasiexact} treats stable intervals and distinguishes quantum conditions. The notebook's conclusion is qualitative; no exact level count or necessary-and-sufficient bound is inferred.}
\GSfield{Present status}{As of 11 November 2026: finite admissible spectra occur in the relevant models; the bare inequality is not a universal sharp quantum bound.}
\GSfield{Assessment}{literal statement: correct with missing hypotheses; publication match: partial; verification: read the result; development: annotated; confidence: high}
\end{GSitem}

\GStalk{751}{1988-04-25}{A. A. Beilinson}{P001472-P001477}{Morse theory, Chern--Simons theory, and three-manifold invariants}
Beilinson reports Atiyah's picture linking the Morse complex, flat connections, spectral flow, Lagrangian intersections and four-dimensional instantons. The false equivalence between ordinary and twisted first-cohomology vanishing is corrected separately. The scan at P001478 starts a new session dated 12 September 1988, and P001479 continues that next academic year; neither belongs to this record.

\begin{GSitem}{1987-751.1}{E}{The finite-dimensional Morse complex}
\GSfield{Kind; source}{construction; \GSpage{P001472-P001473}}
\GSfield{The notes}{For a compact manifold and generic metric, trajectories from \(x_1\) to \(x_2\) have dimension \(\operatorname{ind}(x_1)-\operatorname{ind}(x_2)-1\), and signed index-one trajectories define the differential. The vector field is written \(+\operatorname{grad}f\).}
\GSfield{Statement}{For a Morse function on a closed manifold and a Morse--Smale metric, use the negative gradient. The unparametrized trajectory space from p to q has dimension \(\operatorname{ind}p-\operatorname{ind}q-1\). Choose orientations of unstable manifolds; signed counts for index difference one define \(d:C_i\to C_{i-1}\), with \(d^2=0\), and its homology is the ordinary homology. Mod-two coefficients avoid the orientation choices.}
\GSfield{Reconstruction}{Closedness, Morse--Smale transversality, quotient by time translation and coherent orientation choices are stated explicitly.}
\GSfield{Fate}{\citet[\S3, pp.~287--288]{atiyah1988invariants} recalls this construction after \citet[\S2]{witten1982supersymmetry}. The notebook combines positive-gradient notation with the negative-gradient index convention.}
\GSfield{Present status}{As of 11 November 2026: established under the stated hypotheses.}
\GSfield{Assessment}{literal statement: false as written; publication match: exact; verification: read the result; development: annotated; confidence: high}
\end{GSitem}
\begin{GSitem}{1987-751.2}{E}{The irreducible gauge quotient}
\GSfield{Kind; source}{construction; \GSpage{P001473}}
\GSfield{The notes}{The main manifold is \(X=C/G\), irreducible Hermitian connections on the trivial rank-two bundle, with \(G=SU(2)^M\).}
\GSfield{Statement}{Let M be a closed connected oriented smooth three-manifold. On the trivial \(SU(2)\)-bundle, connections form \(\mathcal A=d+\Omega^1(M;\mathfrak{su}_2)\). An irreducible connection has stabilizer the centre \(\{\pm1\}\); after Sobolev completion, the effective gauge group \(\mathcal G/\{\pm1\}\) acts freely and the quotient has local Coulomb slices. Reducible points are excluded from this manifold assertion.}
\GSfield{Reconstruction}{Orientation, completion and the ineffective central stabilizer are restored.}
\GSfield{Fate}{\citet[\S2, pp.~286--287]{atiyah1988invariants} gives precisely the irreducible quotient used in the Casson/Floer picture. The trivial-bundle condition is automatic for \(SU(2)\)-bundles over a three-dimensional CW complex, but is retained from the notes.}
\GSfield{Present status}{As of 11 November 2026: established under the stated hypotheses.}
\GSfield{Assessment}{literal statement: correct with missing hypotheses; publication match: exact; verification: read the result; development: none; confidence: high}
\end{GSitem}
\begin{GSitem}{1987-751.3}{E}{The Chern--Simons one-form and its periods}
\GSfield{Kind; source}{construction; \GSpage{P001473-P001474}}
\GSfield{The notes}{\(\nu_A(B)=\int_M\operatorname{tr}(F_A\wedge B)\) descends to the quotient and is closed. A primitive is identified with a transgression integral on \(M\times[0,1]\); Gelfand refers to Gabrielov--Gelfand--Losik.}
\GSfield{Statement}{For closed oriented M, gauge invariance and the Bianchi identity make \(\nu\) basic. On \(\mathcal A\), a primitive is \(\tfrac12\int_M\operatorname{tr}(A\wedge dA+\tfrac23A\wedge A\wedge A)\). Without the half, its differential is \(2\nu\). With a conventional integral normalization, large gauge transformations change the primitive by periods, so on the quotient the functional is circle-valued or defined on a covering.}
\GSfield{Reconstruction}{\oursComputed{} The factor two and integral normalization are distinguished; the notebook suppresses both.}
\GSfield{Fate}{\citet[\S3, pp.~287--288]{atiyah1988invariants} describes the closed one-form, the Chern--Simons functional and its integral periods. The historical Gelfand attribution is recorded as an intervention, not treated as an independent new result.}
\GSfield{Present status}{As of 11 November 2026: established under the stated hypotheses.}
\GSfield{Assessment}{literal statement: correct with missing hypotheses; publication match: exact; verification: read the result; development: annotated; confidence: high}
\end{GSitem}
\begin{GSitem}{1987-751.4}{E}{Flat critical points and the tangent quotient}
\GSfield{Kind; source}{construction; \GSpage{P001474}}
\GSfield{The notes}{Critical points are flat connections, hence representations of \(\pi_1(M)\). The quotient tangent space is written with differential forms on X instead of M.}
\GSfield{Statement}{The zero set of \(\nu\) consists of flat connections. Gauge-equivalence classes correspond to conjugacy classes of \(\rho:\pi_1(M)\to SU(2)\). At an irreducible connection, the tangent quotient is \(\Omega^1(M;\operatorname{ad}P)/d_A\Omega^0(M;\operatorname{ad}P)\), with a Coulomb-slice realization after completion. The base of these differential forms is M, not the infinite-dimensional quotient X.}
\GSfield{Reconstruction}{The notation error X is replaced by M; the tangent quotient is for the irreducible stratum.}
\GSfield{Fate}{\citet[\S2, pp.~286--287]{atiyah1988invariants} states the flat-connection/representation correspondence; the quotient formula follows by differentiating the gauge action.}
\GSfield{Present status}{As of 11 November 2026: established under the stated hypotheses.}
\GSfield{Assessment}{literal statement: false as written; publication match: exact; verification: read the result; development: annotated; confidence: high}
\end{GSitem}
\begin{GSitem}{1987-751.5}{E}{Twisted first cohomology is not ordinary homology}
\GSfield{Kind; source}{claim; \GSpage{P001475}}
\GSfield{The notes}{The Hessian is \(\int_M\operatorname{tr}(d_Aa\wedge a)\). Nondegeneracy is identified with \(H^1(M;E)=0\) for every local system, and then with \(H_1(M;\mathbb Z)=0\). The Russian says \ru{Это эквивалентно}.}
\GSfield{Statement}{At a specified irreducible flat connection with holonomy \(\rho\), the Hessian kernel modulo gauge is \(H^1(M;\operatorname{ad}\rho)\). An integral homology sphere need not have this group zero. Conversely, vanishing of \(H^1(M;E)\) for every finite-dimensional characteristic-zero local system does not detect integral torsion.}
\GSfield{Reconstruction}{The local system governing this Hessian is the adjoint system of the chosen representation, not all local systems at once.}
\GSfield{Fate}{\citet[p.~287 and p.~289]{atiyah1988invariants} explicitly requires perturbations when flat critical points are degenerate. Section~\ref{dev:twisted-h1} gives independently checked counterexamples: \(\mathbb{RP}^3\) for the reverse implication, and the connected sum of two Poincare spheres with an irreducible \(SU(2)\) representation for the forward implication.}
\GSfield{Present status}{As of 11 November 2026: the pointwise twisted-cohomology criterion is established; both directions of the notebook' s ordinary/twisted equivalence are false.}
\GSfield{Assessment}{literal statement: false as written; publication match: exact; verification: read the result; development: annotated; confidence: high}
\end{GSitem}
\begin{GSitem}{1987-751.6}{B}{Spectral flow and the relative grading modulo eight}
\GSfield{Kind; source}{construction; \GSpage{P001475}}
\GSfield{The notes}{Index difference is called the number of points where a zero mode appears, then spectral asymmetry; it is said to be well defined modulo eight.}
\GSfield{Statement}{For a path of self-adjoint Fredholm gauge-fixed Hessians with nondegenerate ends, spectral flow is the signed crossing count with multiplicity (or its perturbation-invariant definition), not the unsigned number of zero events. In the \(SU(2)\) instanton theory of an integral homology sphere, changing a path by a generating gauge loop changes this index by eight; hence the relative grading descends modulo eight.}
\GSfield{Reconstruction}{Signs, multiplicities, gauge fixing and the homology-sphere gauge-theory setup are required for the stated eight-periodicity.}
\GSfield{Fate}{\citet[\S4, p.~289]{atiyah1988invariants} gives the spectral-flow calculation and its index-theorem origin. \citet{floer1988instanton} establishes the resulting mod-eight graded homology, received 3 May and published June 1988; only its publisher abstract, not its full proof, was read here.}
\GSfield{Present status}{As of 11 November 2026: established under the stated hypotheses.}
\GSfield{Assessment}{literal statement: correct with missing hypotheses; publication match: partial; verification: read the result; development: annotated; confidence: high}
\end{GSitem}
\begin{GSitem}{1987-751.7}{E}{A cotangent intersection recovers the Euler characteristic}
\GSfield{Kind; source}{construction; \GSpage{P001476}}
\GSfield{The notes}{Take the zero section and the graph of df; intersections are critical points, with signed total written \(\chi(T^*L)\).}
\GSfield{Statement}{For a closed smooth manifold L and Morse f, the zero section and \(\Gamma_{df}\) in \(T^*L\) meet transversely at the critical points. With compatible orientations the local sign is \((-1)^{\operatorname{ind}f}\), and their sum is \(\chi(L)=\chi(T^*L)\), the last equality following from deformation retraction.}
\GSfield{Reconstruction}{Compactness of L, Morse nondegeneracy and the orientation convention are supplied.}
\GSfield{Fate}{\citet[\S5]{atiyah1988invariants} uses Lagrangian intersections as the finite-dimensional model. The equality with the total-space Euler characteristic is harmless here, although the total space is noncompact.}
\GSfield{Present status}{As of 11 November 2026: established under the stated hypotheses.}
\GSfield{Assessment}{literal statement: correct with missing hypotheses; publication match: exact; verification: read the result; development: none; confidence: high}
\end{GSitem}
\begin{GSitem}{1987-751.8}{E}{The closed action one-form and Maslov relative index}
\GSfield{Kind; source}{construction; \GSpage{P001476-P001477}}
\GSfield{The notes}{On paths from \(L_1\) to \(L_2\), define \(\nu_\gamma(\xi)=\int_0^1\omega(\dot\gamma,\xi)dt\); zeros are intersections, and a relative index is obtained from Maslov theory.}
\GSfield{Statement}{Let \((N,\omega)\) be symplectic and \(L_1,L_2\) Lagrangian. The action one-form on a suitable Sobolev path space is closed because its exterior derivative is the endpoint difference of \(\omega\), which vanishes on the Lagrangians. Its zeros are constant paths in \(L_1\cap L_2\). Transverse intersections give nondegenerate Fredholm Hessians; relative grading uses the Maslov index, with ambiguity determined by the relevant loop classes. Exactness or a covering is needed for a single-valued action.}
\GSfield{Reconstruction}{Lagrangian boundary conditions, Sobolev/Fredholm setup and the action-period ambiguity are explicit; existence of a full Floer homology needs further compactness hypotheses.}
\GSfield{Fate}{\citet[\S5, pp.~292--293]{atiyah1988invariants} describes the symplectic Floer construction and cites Floer's 1987 preprints. The index formalism has earlier variational background in \citet[\S2]{duistermaat1974morse}.}
\GSfield{Present status}{As of 11 November 2026: established under the stated hypotheses.}
\GSfield{Assessment}{literal statement: correct with missing hypotheses; publication match: exact; verification: read the result; development: none; confidence: high}
\end{GSitem}
\begin{GSitem}{1987-751.9}{E}{Heegaard splitting as a Lagrangian representation-space intersection}
\GSfield{Kind; source}{construction; \GSpage{P001477}}
\GSfield{The notes}{A three-manifold is cut into two pieces along a surface; fundamental groups and their representations give the finite-dimensional counterpart of the connection construction.}
\GSfield{Statement}{For a genus-g Heegaard splitting \(M=W_1\cup_\Sigma W_2\), each handlebody group is free of rank g and van Kampen presents \(\pi_1M\) as their amalgam over \(\pi_1\Sigma\). On the smooth irreducible locus, the \(SU(2)\) character variety of \(\Sigma\) has its cup-product symplectic form. The handlebody representation loci are Lagrangian where smooth, and their intersection represents flat connections on M. Reducibles, singularities and transversality cannot be ignored.}
\GSfield{Reconstruction}{The character-variety smooth locus and Lagrangian hypotheses are supplied rather than claiming that every full representation quotient is a smooth symplectic manifold.}
\GSfield{Fate}{\citet[\S5, pp.~290--293]{atiyah1988invariants} is the direct publication match. It discusses the conjectural relation of the two Floer constructions, but the notebook does not explicitly formulate that homology-isomorphism conjecture, so none is added to this year's problem index.}
\GSfield{Present status}{As of 11 November 2026: established under the stated hypotheses.}
\GSfield{Assessment}{literal statement: correct with missing hypotheses; publication match: exact; verification: read the result; development: none; confidence: high}
\end{GSitem}
\begin{GSitem}{1987-751.10}{B}{Instantons and cutting a four-manifold along a three-manifold}
\GSfield{Kind; source}{method; \GSpage{P001477}}
\GSfield{The notes}{Four-dimensional Yang--Mills equations and cutting a four-manifold into two three-dimensional ones are invoked as the connection to Donaldson theory.}
\GSfield{Statement}{On \(M\times\mathbb R\), an anti-self-dual connection in temporal gauge satisfies \(\dot A=-*_MF_A\) for the product orientation, a Chern--Simons gradient equation after choosing the overall sign convention. A decomposition \(Z=Z_1\cup_MZ_2\) has two four-dimensional pieces sharing a three-dimensional boundary. Finite-energy instanton moduli spaces and their boundary data are the objects used in Floer/Donaldson gluing; arbitrary Yang--Mills critical points do not suffice.}
\GSfield{Reconstruction}{Anti-self-duality, finite energy, temporal gauge and the dimensions of the pieces are corrected explicitly.}
\GSfield{Fate}{\citet[\S4, equation~(4.1), and \S7]{atiyah1988invariants} explains the cylinder equation and the gluing programme. \citet{floer1988instanton} is the subsequent rigorous homology construction. The notebook's dimensions are a literal slip, and its analogy alone is not a complete gluing theorem.}
\GSfield{Present status}{As of 11 November 2026: established under the stated hypotheses.}
\GSfield{Assessment}{literal statement: false as written; publication match: partial; verification: read the result; development: annotated; confidence: high}
\end{GSitem}

\section{Developments}\label{dev:1987}
\subsection{Joint spectra of stochastic matrices: a bounded development of GP-1987-21}
\label{dev:stochastic}

\paragraph{The source question and its scope.}
In item \hyperref[item:1987-728.4]{1987-728.4}, Gelfand asks for the
\ru{набор соб. значений в \(\mathbb C^n\)} and emphasizes that projection onto
one coordinate may conceal its geometry. The scan of P001418 supports a
question about all eigenvalues of a real nonnegative row-stochastic matrix.
It does not impose symmetry, double stochasticity, irreducibility, a specified
zero pattern, or a prescribed Jordan form. Those are different problems.
The strongest precise formulation supported by the page is to describe,
for each positive integer \(n\), the permutation-invariant set of lists
for which there exists a matrix \(P\) satisfying
\[
 \Sigma_n=\left\{(\lambda_1,\ldots,\lambda_n)\in\mathbb C^n:
 \begin{array}{l}
 \det(tI-P)=\prod_{j=1}^n(t-\lambda_j),\\
 P\in\mathbb R^{n\times n},\quad
 P_{ij}\geq0,\quad P\mathbf1=\mathbf1
 \end{array}\right\}.
\]
Multiplicities are retained. This is a reconstruction of a shorthand
question, not a conjectured theorem attributed to the speaker.

\paragraph{State of the art and the bounded contribution.}
The order-three characterization predates the seminar:
\citet{loewylondon1978inverse} appeared in 1978; its attribution and
explicit alternatives are recorded in \citet[p.~506]{ranteng2024patterns}.
The latter paper additionally classifies prescribed order-three zero patterns
and appeared in July 2024 (received May 2023).
\citet[Theorems~3.9--3.12 and Sections~4--5, 8]{johnsonpaparella2025perron}
establish geometric properties of joint stochastic spectra and polyhedral
families from fixed diagonalizing matrices; their preprint first circulated
on 12 September 2024 and the journal paper appeared in 2025.
Their nearly complete geometric picture in order four does not solve
all orders. As of 11 November 2026, the sources checked provide no general
structural solution in orders at least five. The distinction between
diagonalizable and unrestricted realization is substantive;
\citet[Section~2 and Proposition~1]{croninlaffey2018diagonalizable}
first circulated in January 2017 and appeared in July 2018.

For a fixed \(n\), polynomial feasibility already makes membership
decidable for algebraic input. Thus \emph{open} does not mean that there is
no finite decision procedure. What remains is an explicit, usable
description of the joint region, its boundary and realizing matrices.
The contribution here is a self-contained, checked low-dimensional atlas
and precise obstructions to extrapolating necessary tests.
None of the results below is claimed to be new.

\begin{proposition}[Geometry and a finite formulation \ours]
\label{prop:stochastic-geometry}
For every fixed \(n\), \(\Sigma_n\) is compact and semialgebraic in
\(\mathbb R^{2n}\). It is invariant under coordinate permutations and
complex conjugation, and star-shaped at \((1,\ldots,1)\).
The set
\[
 \widehat\Sigma_n=\{(\lambda_2,\ldots,\lambda_n):
                         (1,\lambda_2,\ldots,\lambda_n)\in\Sigma_n\}
\]
is star-shaped at the origin.
\end{proposition}
\begin{proof}
The stochastic matrices form a compact product of \(n\) probability
simplices. Every eigenvalue has modulus at most one: if \(Pv=\lambda v\)
and \(|v_i|=\max_j|v_j|>0\), then
\(|\lambda v_i|\leq\sum_jP_{ij}|v_j|\leq|v_i|\).
For a convergent sequence of lists in \(\Sigma_n\), choose corresponding
matrices and a convergent subsequence of these matrices. Continuity of
characteristic-polynomial coefficients shows that the limiting matrix
realizes the limiting list. Thus \(\Sigma_n\) is closed and bounded,
without requiring a continuously chosen ordering of eigenvalues.

Writing \(\lambda_j=x_j+iy_j\), equality of coefficients in
\(\det(tI-P)=\prod_j(t-x_j-iy_j)\) is a finite system of real
polynomial equalities. Append \(P_{ij}\geq0\) and
\(\sum_jP_{ij}=1\). Projecting onto the \(x,y\) coordinates and applying
real quantifier elimination gives semialgebraicity and a decision
algorithm for effectively encoded algebraic input.
The background theorem is the standard one treated in
\citet[Chapters~2, 13--14]{basupollackroy2006algorithms};
we do not reprove that theorem or claim an efficient complexity bound.

Permutation invariance follows from the characteristic polynomial;
conjugation invariance follows from its real coefficients.
For \(0\leq s\leq1\), \((1-s)I+sP\) is stochastic with eigenvalues
\(1-s+s\lambda_j\), proving the first star-shapedness assertion.
Put \(J=\mathbf1\mathbf1^{\mathsf T}/n\) and \(Q=sP+(1-s)J\).
On the invariant line spanned by \(\mathbf1\), \(Q\) is the identity.
On the quotient by that line, it induces \(s\) times the map induced
by \(P\), because \(J\) has image in the line.
In a basis beginning with \(\mathbf1\), its characteristic polynomial
is therefore \((t-1)\prod_{j=2}^n(t-s\lambda_j)\).
This also covers repeated eigenvalues and nondiagonalizable matrices.
\end{proof}

\begin{lemma}[Power-sum inequalities \ours]
\label{lem:jll}
For a list in \(\Sigma_n\), put \(s_k=\sum_j\lambda_j^k\).
Then
\[
 s_k\geq0,\qquad s_k^m\leq n^{m-1}s_{km}
 \quad(k,m\in\mathbb N).
\]
The list contains \(1\) and is closed under conjugation.
\end{lemma}
\begin{proof}
The last assertions follow from reality and \(P\mathbf1=\mathbf1\).
Let \(B=P^k\geq0\). Then \(s_k=\operatorname{tr}B\geq0\).
The term in \((B^m)_{ii}\) that remains at \(i\) at every step is
\(B_{ii}^m\); every other term is nonnegative. Consequently
\(\sum_iB_{ii}^m\leq\operatorname{tr}(B^m)=s_{km}\).
The power-mean inequality gives
\((\sum_iB_{ii})^m\leq n^{m-1}\sum_iB_{ii}^m\).
Combining the inequalities proves the claim, including \(m=1\).
\end{proof}
These are the classical Johnson--Loewy--London inequalities, independently
derived here; the historical references are discussed in
\citet[Section~1.1]{croninlaffey2018diagonalizable}.

\begin{theorem}[The complete two- and three-state regions \ours]
\label{thm:lowdim-stochastic}
A two-element multiset is the spectrum of a real row-stochastic
\(2\times2\) matrix exactly when it is \(\{1,r\}\), \(-1\leq r\leq1\).
A three-element multiset is the spectrum of a real row-stochastic
\(3\times3\) matrix exactly when it has one of the forms
\[
\begin{array}{ll}
\{1,r,s\},&r,s\in[-1,1],\quad 1+r+s\geq0;\\[2pt]
\{1,a+ib,a-ib\},& b\ne0,\quad -\tfrac12\leq a\leq1,\quad
 |b|\leq(1-a)/\sqrt3 .
\end{array}
\]
Weak inequalities retain every boundary case. The second line describes
the triangle with vertices \(1,-1/2\pm i\sqrt3/2\), with its intersection
with the real axis omitted because \(b\ne0\).
\end{theorem}
\begin{proof}
In dimension two, reality and the eigenvalue \(1\) force the other
eigenvalue to be real, and the modulus bound gives necessity.
The matrix
\[
 C(r)=\frac12\begin{pmatrix}1+r&1-r\\1-r&1+r\end{pmatrix}
\]
is nonnegative and row-stochastic for \(-1\leq r\leq1\).
Its eigenvectors \((1,1)\), \((1,-1)\) have eigenvalues \(1,r\).

In dimension three, the remaining pair is either real or conjugate.
For a real pair, the modulus and trace conditions are necessary.
For a conjugate pair, the trace gives \(a\geq-1/2\), and the
modulus bound gives \(a\leq1\).
Lemma~\ref{lem:jll} with \(k=1,m=2,n=3\) gives
\[
 (1+2a)^2\leq3\{1+2(a^2-b^2)\}
 \quad\Longleftrightarrow\quad 3b^2\leq(1-a)^2 .
\]
These are the claimed necessary inequalities.

For real-pair sufficiency, if one member is nonnegative, rename it
\(r\in[0,1]\) and write the other as \(s\in[-1,1]\). Then
\[
 P=\begin{pmatrix}
 (1+s)/2&(1-s)/2&0\\
 (1-s)/2&(1+s)/2&0\\
 1-r&0&r
 \end{pmatrix}
\]
is nonnegative and stochastic; block triangularity gives eigenvalues
\(1,s,r\). If both are negative, write \(r=-u,s=-v\), where
\(u,v>0\) and \(u+v\leq1\). The matrix
\[
 P=\begin{pmatrix}
 0&1&0\\0&0&1\\uv&u+v-uv&1-u-v
 \end{pmatrix}
\]
is stochastic and nonnegative, since
\(u+v-uv=u(1-v)+v\geq0\). Its characteristic polynomial is
\[
 t^3-(1-u-v)t^2-(u+v-uv)t-uv=(t-1)(t+u)(t+v).
\]
Zero eigenvalues were covered by the first construction.

For a conjugate pair, set
\[
 p_0=\frac{1+2a}{3},\qquad
 p_1=\frac{1-a}{3}+\frac{b}{\sqrt3},\qquad
 p_2=\frac{1-a}{3}-\frac{b}{\sqrt3}.
\]
The inequalities make these nonnegative with sum one.
For \(\omega=e^{2\pi i/3}\), the circulant
\[
 P=\begin{pmatrix}p_0&p_1&p_2\\p_2&p_0&p_1\\p_1&p_2&p_0\end{pmatrix}
\]
has eigenvectors \((1,1,1)\), \((1,\omega,\omega^2)\),
\((1,\omega^2,\omega)\), with eigenvalues
\(1,p_0+p_1\omega+p_2\omega^2,p_0+p_1\omega^2+p_2\omega\).
These are \(1,a+ib,a-ib\), completing the sufficiency proof.
\end{proof}

\begin{example}[Two insufficient tests \ours]
Each of \(-3/4\) and \(1\) occurs in a stochastic three-state
spectrum: use \(C(-3/4)\oplus(1)\).
Yet \(\{1,-3/4,-3/4\}\) cannot be a joint spectrum because its
trace is \(-1/2\). A product of marginal regions is therefore too large.
Even the modulus, conjugation and all nonnegative power-sum tests
can be insufficient. The list \(\{1,i/\sqrt2,-i/\sqrt2\}\) has
\(s_{2m+1}=1\), \(s_{2m}=1+2(-1/2)^m\geq0\), but
\(s_1^2=1>3s_2=0\). Both exclusions are exact.
\end{example}

\begin{proposition}[An order-five obstruction surviving all JLL tests \ours]
\label{prop:five-obstruction}
The list \(\Lambda=\{1,1,-2/3,-2/3,-2/3\}\) satisfies every
power-sum and JLL inequality for \(n=5\), together with the Perron and
conjugation conditions. It is nevertheless not the spectrum of any
nonnegative \(5\times5\) matrix, and hence is not in \(\Sigma_5\).
\end{proposition}
\begin{proof}
Put \(q=2/3\), so \(s_k=2+3(-q)^k\).
For even \(k\), apply the power-mean inequality to
\(1,1,q^k,q^k,q^k\), obtaining
\(s_k^m\leq5^{m-1}s_{km}\).
For odd \(k\), \(0\leq s_k\leq2\), since \(q^k\leq2/3\).
If \(m\) is odd, \(s_{km}\geq s_k\), so
\(s_k^m\leq2^{m-1}s_k\leq5^{m-1}s_{km}\).
If \(m\) is even, \(s_{km}\geq2\), so
\(s_k^m\leq2^m\leq5^{m-1}s_{km}\).
This proves all inequalities for every positive \(k,m\).

Suppose a nonnegative matrix realizes the list. Put it in Frobenius
block-triangular form, with irreducible nonnegative diagonal blocks,
including one-by-one blocks. Each block has its nonnegative spectral
radius as an eigenvalue. The only nonnegative member of \(\Lambda\)
is \(1\), so each block must contain a copy of \(1\).
The Perron eigenvalue of an irreducible block is simple; there must
therefore be exactly two blocks, each using one copy.
One block receives at least two of the three negative eigenvalues.
Its trace is at most \(1-4/3=-1/3\), impossible for a nonnegative block.
\end{proof}
This is the classical list recorded in
\citet[Section~2]{croninlaffey2018diagonalizable}.
The all-\(k,m\) check and block proof are supplied to make the
obstruction independently verifiable, without claiming novelty.

\paragraph{Independent verification \oursComputed.}
The script \path{code/verify_development.py} checks the characteristic
polynomials symbolically and then checks the constructions using a
separate exact three-by-three determinant formula on rational grids.
It also tests independently generated stochastic matrices against the
inequalities and checks the order-five example with exact fractions
over a finite exponent range. The finite tests are diagnostic;
the universally quantified claims are proved above.
The output is \path{logs/development_checks.json}.
An independent derivation of the triangular bound uses
\(\operatorname{tr}P^2=\sum_iP_{ii}^2+
2\sum_{i<j}P_{ij}P_{ji}\), followed by Cauchy--Schwarz.
The Fourier eigenvector proof checks the converse separately from
the characteristic-polynomial calculation. Both checks agree,
including boundary cases and reducible real-pair realizations.

\paragraph{What remains and what another run could accomplish.}
The bottleneck is joint compatibility: one nonnegative matrix must
realize all roots. Reducibility and repeated Perron roots constrain
how roots can be grouped. Projection to one eigenvalue, all power-sum
inequalities, or one fixed diagonalizing basis loses information.
The calculation does not resolve the general \(n\geq5\) region.
There is no evidence that this question was abandoned: the checked
2024--2025 papers show continued work.
A bounded next run could construct an exact atlas for one specified
family of four-state zero patterns, compare it with the known
order-four characterization and the constructions in
\citet[Section~8]{johnsonpaparella2025perron}, and certify every
boundary piece by a realizing matrix or polynomial obstruction.
That would provide a reusable companion/website contribution.
Its novelty would have to be established by that comparison.

\subsection{A Poisson-pencil repair for the Neumann calculation}
\label{dev:poisson-repair}

This annotation of item \hyperref[item:1987-716.6]{1987-716.6}
records a failure found by the independent check. With
\(\varepsilon_{123}=1\), the second printed table on P001391 has
\(\{M_i,M_j\}=\varepsilon_{ijk}a_kM_k\),
\(\{M_i,p_j\}=\varepsilon_{ijk}a_kp_k\), and
\(\{p_i,p_j\}=\varepsilon_{ijk}M_k\).
Its cyclic Jacobi expression at \((M_1,p_1,p_2)\) is
\((a_3-a_2)M_2\), so it is not a Poisson bracket for general
diagonal \(A\). Reversing time cannot repair this defect.
The scan confirms the printed coefficients. The following is a
consistent algebraic emendation, not a claim to recover a uniquely
intended historical formula.

\begin{proposition}[A compatible repaired bracket \ours]
For arbitrary real \(a_1,a_2,a_3\), the table
\[
 \{M_i,M_j\}_2=\varepsilon_{ijk}a_kM_k,\qquad
 \{M_i,p_j\}_2=\varepsilon_{ijk}a_jp_k,\qquad
 \{p_i,p_j\}_2=-\varepsilon_{ijk}M_k
\]
defines a linear Poisson bracket. It is compatible with the standard
\(\mathfrak e(3)^*\) bracket
\(\{M_i,M_j\}_1=\varepsilon_{ijk}M_k\),
\(\{M_i,p_j\}_1=\varepsilon_{ijk}p_k\),
\(\{p_i,p_j\}_1=0\).
With \(F=p^2/2\), the repaired bracket gives
\(\dot p=M\times p\), \(\dot M=(Ap)\times p\), exactly the
first bracket's vector field for
\(H=(M^2+p^{\mathsf T}Ap)/2\).
\end{proposition}
\begin{proof}
Let \(E_{ij}\) denote the skew matrix with entry \(1\) in position
\((i,j)\), entry \(-1\) in position \((j,i)\), and all other entries
zero. For any diagonal \(4\times4\) matrix \(D\), the product
\(X*Y=XDY\) is associative. Its commutator
\([X,Y]_D=XDY-YDX\) is consequently a Lie bracket.
Skew matrices are closed under this bracket since \(D^{\mathsf T}=D\).
Choose the basis
\[
 m_1=-E_{23},\quad m_2=-E_{31},\quad m_3=-E_{12},\qquad
 q_i=E_{i4}\quad(1\leq i\leq3).
\]
For \(D=\operatorname{diag}(a_1,a_2,a_3,-1)\), direct matrix
multiplication gives precisely the three displayed structure-constant
formulas. The associated Lie--Poisson bracket on the dual therefore
satisfies Jacobi, including when \(D\) is singular.
For \(D_0=\operatorname{diag}(1,1,1,0)\), the same basis gives the
first bracket. Every linear combination is again the commutator for
the corresponding linear combination of \(D,D_0\), proving compatibility.
Finally, contracting the repaired table with
\(\partial F/\partial p_j=p_j\) gives the two stated vector fields.
Contracting the first table with the derivatives of \(H\) gives the
same fields, independently checking their signs.
\end{proof}

\paragraph{Second check \oursComputed.}
The coordinate computation in the early-record verification script
tests every coordinate Jacobi identity for the repaired bracket and
the full pencil, and checks equality of the Hamiltonian fields.
The associative-matrix construction above is an independent proof
of those identities. Its agreement with the symbolic output confirms
the emendation; the failed original table is retained in that output.

\subsection{The complete attainable range in the tensor-angle exercise}
\label{dev:tensor-angles}

This is an elementary annotation of item
\hyperref[item:1987-737.11]{1987-737.11} and GP-1987-29, classified E;
no priority claim is made. The preamble on P001435 explicitly asks
which values the invariants in all the exercises can take. On P001436,
the tensor exercise specifies a Euclidean space \(L\) and compares
\(a\otimes L\) with \(L\otimes b\) in \(L^{\otimes3}\).
We reconstruct the finite-dimensional setting of that part of the
sheet, require \(a,b\ne0\), and normalize them to unit length;
scaling does not change either subspace. The same proof applies to
complex Hilbert spaces, although the source says Euclidean.

\begin{theorem}[Principal cosines for two tensor subspaces \ours]
Let \(L\) be a real Euclidean or complex Hilbert space of dimension
\(d\geq2\). As unit tensors \(a,b\in L\otimes L\) vary, the ordered
principal cosines \(s_1\geq\cdots\geq s_d\geq0\) between
\(a\otimes L\) and \(L\otimes b\) are exactly the lists satisfying
\[
                  \sum_{j=1}^d s_j\leq1.
\]
For \(d=1\), the sole possible list is \((1)\).
\end{theorem}
\begin{proof}
In an orthonormal basis write
\(a=\sum_{i,j}A_{ij}e_i\otimes e_j\) and
\(b=\sum_{j,k}B_{jk}e_j\otimes e_k\), so
\(\|A\|_F=\|B\|_F=1\).
The maps \(Uu=a\otimes u\), \(Vv=v\otimes b\) are isometries.
With the complex inner product conjugate-linear in its first entry,
\[
 (U^*V)_{ki}=\langle a\otimes e_k,e_i\otimes b\rangle
            =\sum_j\overline{A_{ij}}B_{jk}.
\]
Thus the cross-Gram matrix is \((\overline A B)^{\mathsf T}\),
and its singular values are the principal cosines. The trace norm
inequality
\(\|\overline A B\|_*\leq\|A\|_F\|B\|_F=1\)
gives necessity. To check this inequality directly, choose a unitary
matrix from a singular value decomposition so that its trace pairing
with \(\overline A B\) is the sum of the singular values, and apply
Cauchy--Schwarz to the matrix entries. Left or right multiplication
by a unitary matrix preserves the Frobenius norm.

Conversely, put \(T=\sum_js_j\leq1\), and let \(r\) be the number
of positive entries. All the following matrices are real, so their
constructions work over either field. If \(T=0\), take
\(A_{11}=B_{22}=1\), with all other entries zero;
both matrices have norm one and \(AB=0\).
If \(0<r<d\), set
\[
 A_{ii}=\sqrt{s_i/T},\qquad B_{ii}=\sqrt{Ts_i}
 \quad(1\leq i\leq r),\qquad
 B_{r+1,1}=\sqrt{1-T^2},
\]
and set all other entries to zero. The additional entry lies in a
row annihilated by \(A\), so
\(AB=\operatorname{diag}(s_1,\ldots,s_d)\).
Moreover, \(\|A\|_F^2=1\) and
\(\|B\|_F^2=T^2+(1-T^2)=1\).

It remains to consider \(r=d\). On the open probability simplex
\(x_j>0\), \(\sum_jx_j=1\), define
\(f(x)=\sum_js_j^2/x_j\). Cauchy--Schwarz gives \(f(x)\geq T^2\),
with equality at \(x_j=s_j/T\).
If \(T=1\), use this minimizing point. If \(T<1\), the minimum
is below one, whereas \(f(x)\) tends to infinity along a segment
approaching a face, since every \(s_j>0\).
Continuity supplies a point with \(f(x)=1\).
Now take \(A_{jj}=\sqrt{x_j}\) and
\(B_{jj}=s_j/\sqrt{x_j}\). Both Frobenius norms equal one, and
\(AB=\operatorname{diag}(s_j)\), as required.
When \(d=1\), both tensor subspaces are the entire one-dimensional
ambient space; their principal cosine is one.
\end{proof}

\paragraph{Independent verification and scope \oursComputed.}
The cross-Gram formula was re-derived entry by entry and checked
against explicit tensor embeddings, including complex coefficients.
An independent mathematical review checked all rank cases and the
simplex-continuity argument. The script
\path{code/verify_development.py} tests both real and complex
random tensors, rank-deficient constructions, and the exceptional
case \(d=1\); its output records the observed residuals.
This resolves the stated finite-dimensional attainable-range exercise.
It does not classify simultaneous invariants of several unrelated
subspaces or the Hilbert-space cases on the next part of the sheet.

\subsection{Homology spheres and twisted cohomology: checking a literal error}
\label{dev:twisted-h1}

This is an annotation of P001475, not an extra class-C or class-D
research problem. The source passes from the Hessian at a flat
\(\mathrm{SU}(2)\) connection to all local systems and asserts an
equivalence with integral first-homology vanishing.
At a specified irreducible flat connection with holonomy \(\rho\),
the relevant infinitesimal space is \(H^1(M;\operatorname{ad}\rho)\).
Ordinary and twisted homology cannot be interchanged.
Already \citet[p.~287]{atiyah1988invariants} requires perturbations
to obtain simple zeros in the homology-sphere construction.
The next two checks disprove both directions of the literal
equivalence over characteristic zero.

\begin{lemma}[Averaging one-cocycles \ours]
For a finite group \(G\) acting linearly on a real or complex vector
space \(V\), every one-cocycle \(c:G\to V\) is principal.
\end{lemma}
\begin{proof}
The identity is \(c(gh)=c(g)+g c(h)\).
Put \(S=\sum_{h\in G}c(h)\); summing gives
\(S=|G|c(g)+gS\). Thus \(c(g)=w-gw\) for \(w=S/|G|\).
The characteristic-zero hypothesis permits division by \(|G|\).
\end{proof}

\begin{example}[Twisted vanishing misses integral torsion \ours]
The closed connected orientable manifold \(M=\mathbb{RP}^3=L(2,1)\)
has \(\pi_1M=\mathbb Z/2\) and \(H_1(M;\mathbb Z)=\mathbb Z/2\).
It is not an integral homology sphere. Nevertheless every
finite-dimensional real or complex local system \(E\) satisfies
\[
 H^1(M;E)=H^1(\pi_1M;E)=0
\]
by averaging. The equality in degree one follows because the simply
connected universal cover gives an exact lifted cellular resolution
through the degrees needed for \(H^1\).
The ordinary homology calculation is also the cellular one in
\citet[Examples~2.42--2.43]{hatcher2002algebraic}.
\end{example}

\begin{proposition}[A homology sphere with a degenerate irreducible connection \ours]
There is a closed connected orientable integral homology
three-sphere \(M\) and an irreducible representation
\(\rho:\pi_1M\to\mathrm{SU}(2)\) such that
\(\dim_{\mathbb R}H^1(M;\operatorname{ad}\rho)=3\).
\end{proposition}
\begin{proof}
Let \(G=2I\subset\mathrm{SU}(2)\) be the binary icosahedral group
and \(Y=S^3/G\), with the free orientation-preserving action by
left multiplication on unit quaternions.
The order-\(120\) group has presentation
\(\langle a,b\mid a^5=b^3=(ab)^2\rangle\).
Its abelianization vanishes: taking the two relators as rows gives the integer
matrix \(\left(\begin{smallmatrix}5&-3\\-2&1\end{smallmatrix}\right)\),
of determinant \(-1\).
This presentation and calculation are recorded in
\citet[Example~2.38]{hatcher2002algebraic}.
Hence \(H_1(Y;\mathbb Z)=0\); orientability, connectedness and
Poincare duality give the integral homology of \(S^3\).

Set \(M=Y\#Y\). Van Kampen gives \(\pi_1M=G*G\);
abelianization and duality again give the integral homology of \(S^3\).
Define \(\rho\) to be the inclusion into \(\mathrm{SU}(2)\) on each
free factor. Its image is nonabelian and acts irreducibly on
\(\mathbb C^2\). For \(V=\mathfrak{su}_2\) with the adjoint action,
\(V^G=0\): the icosahedral rotation group has no common fixed axis.

A one-cocycle on \(G*G\) is uniquely determined by its restrictions
to the two free factors; extend by the cocycle identity along reduced
words. Averaging makes these restrictions
\(w_1-gw_1\) and \(w_2-gw_2\).
Since \(V^G=0\), the \(w_i\) are unique. Thus
\(Z^1(G*G;V)\cong V\oplus V\); principal cocycles form its diagonal
copy of \(V\). Consequently
\[
 H^1(M;\operatorname{ad}\rho)\cong
 (V\oplus V)/\operatorname{diag}V\cong V.
\]
The associated \(\mathrm{SU}(2)\) principal bundle is trivial, as
required in the source: over a three-dimensional CW complex a section
exists because \(\pi_0,\pi_1,\pi_2\) of \(\mathrm{SU}(2)\cong S^3\)
all vanish. This verifies the bundle and irreducibility conditions.
\end{proof}

\paragraph{Independent second check.}
Decompose \(M\) into two punctured copies of \(Y\) intersecting
in a collar of \(S^2\). With these coefficients each piece has
\(H^0=0\) and \(H^1=0\), while the intersection has \(H^0=V\)
and \(H^1=0\). Twisted Mayer--Vietoris gives
\(H^1(M;V)\cong V\) again. Puncturing does not change the
fundamental group or degree-one local-system cohomology.
This agrees with the cocycle quotient and computes the exact dimension;
it is stronger than merely exhibiting a family of representations.

\section{Problem index}\label{problems:1987}
Earlier-year occurrence identifiers are retained as text; current-year identifiers link to their entries.

\subsection{GP-1987-01: Character realizations of Schubert polynomials}\label{prob:GP-1987-01}
For each finite permutation \(w\), find a finite-dimensional polynomial representation, in a specified group category, whose torus character equals \(\mathfrak S_w\).
\par\noindent\textit{Occurrences:} \hyperref[item:1987-720.21]{1987-720.21}.
\par\noindent\textit{Status (2026-11):} As of November 2026, the Borel-module formulation is solved by Kra\'skiewicz--Pragacz, in a 1986 preprint announced in 1987 and fully published in 2004. These modules need not be irreducible; the seminar does not specify a category for its stronger irreducibility wording. \citep[Theorem 4.1 and Introduction]{kraskiewiczpragacz2004schubert}

\subsection{GP-1987-02: Real and complex representability of a finite matroid}\label{prob:GP-1987-02}
Given a finite matroid by its bases, determine whether it is representable over \(\mathbb R\) or over \(\mathbb C\).
\par\noindent\textit{Occurrences:} 1985-658.13, 1985-658.16, \hyperref[item:1987-725.3]{1987-725.3}.
\par\noindent\textit{Status (2026-11):} As of November 2026, the finite decision problem is settled by polynomial encoding and quantifier elimination, available before 1987. The broader phrase classification is not made into a separate unresolved claim without a specified structural target. This is distinct from earlier problems on realization-space topology and determinant-sign isotopy. \citep[\S\S 2--3]{tarski1951decision}

\subsection{GP-1987-03: Axiomatizing the nonzero-minor data of a bimatroid}\label{prob:GP-1987-03}
Characterize abstract collections \(\mathcal R\) of equally sized row-column subset pairs by axioms generalizing nonzero square minors.
\par\noindent\textit{Occurrences:} \hyperref[item:1987-725.6]{1987-725.6}.
\par\noindent\textit{Status (2026-11):} As of November 2026, the collection is a bimatroid exactly when \(\{(E\setminus I)\sqcup J:(I,J)\in\mathcal R\}\) is a matroid basis family on \(E\sqcup F\) containing \(E\). The checked modern exposition attributes the theory to Kung 1978. \citep[Proposition 2.2]{rohrleulirsch2026bimatroids}

\subsection{GP-1987-04: Information in a bimatroid beyond its row and column matroids}\label{prob:GP-1987-04}
Determine what data of a matrix bimatroid are lost when only its row matroid and column matroid are retained.
\par\noindent\textit{Occurrences:} \hyperref[item:1987-725.12]{1987-725.12}.
\par\noindent\textit{Status (2026-11):} As of November 2026, the separate matroids do not determine the bimatroid; a two-by-two counterexample is given in the occurrence. Joint submatrix ranks, equivalently the nonzero square-minor pairs subject to bimatroid axioms, supply the missing data. \citep[Proposition 2.6]{rohrleulirsch2026bimatroids}

\subsection{GP-1987-05: Generic irreducible Hecke characters in symmetric functions}\label{prob:GP-1987-05}
For the semisimple type-A Hecke algebra \(H_N(q)\), give a symmetric-function formula for irreducible character values on standard cycle elements \(T_{w_\mu}\).
\par\noindent\textit{Occurrences:} \hyperref[item:1987-726.8]{1987-726.8}.
\par\noindent\textit{Status (2026-11):} As of November 2026, a Frobenius formula is available; Ram published one in 1991, and Halverson--Ram give the exact normalization and credit several independent discoveries. This is a generic-parameter answer and does not settle the different root-of-unity decomposition problem. \citep[equation (4.10), Theorem 4.11]{halversonram1999bitraces}

\subsection{GP-1987-21: Joint spectra of row-stochastic matrices}\label{prob:GP-1987-21}
For each \(n\), give necessary and sufficient conditions on a complex multiset \((1,\lambda_2,\ldots,\lambda_n)\) to be the spectrum of a real nonnegative \(n\times n\) matrix with every row sum equal to one.
\par\noindent\textit{Occurrences:} \hyperref[item:1987-728.4]{1987-728.4}.
\par\noindent\textit{Status (2026-11):} As of 11 November 2026: partly characterized. The order-three region was known before the talk; recent geometric and prescribed-pattern results do not give a general explicit solution in orders at least five \citep[p.~506]{ranteng2024patterns}\citep[Sections~4--5, 8]{johnsonpaparella2025perron}. Section~\ref{dev:stochastic} supplies complete two- and three-state realizations and proves that even all JLL inequalities are insufficient in order five. Fixed-order algebraic feasibility is decidable, so the unresolved goal is a useful structural characterization with realizing matrices, rather than bare decidability.

\subsection{GP-1987-22: Geometric invariants of equivariant toric bundles from filtrations}\label{prob:GP-1987-22}
Compute cohomology, weight characters and Chern classes of an equivariant vector bundle on a smooth complete toric variety directly from its compatible ray filtrations.
\par\noindent\textit{Occurrences:} \hyperref[item:1987-730.3]{1987-730.3}.
\par\noindent\textit{Status (2026-11):} 2026-11: solved in Klyachko's 1989 paper (English translation 1990), Theorems 3.2.1,4.1.1,4.2.1.

\subsection{GP-1987-23: Moduli of compatible filtrations and toric bundles}\label{prob:GP-1987-23}
For fixed fan, rank and equivariant Chern data, construct and analyze the moduli of compatible ray filtrations modulo change of basis; distinguish this from a finite indecomposable list for an arbitrary simplicial complex.
\par\noindent\textit{Occurrences:} \hyperref[item:1987-730.4]{1987-730.4}.
\par\noindent\textit{Status (2026-11):} 2026-11: toric fixed-data moduli spaces are constructed; rank-three examples satisfy Murphy's law (Payne2008). This supplies a moduli formulation, with no assertion of finite-type representation classification for every complex.

\subsection{GP-1987-24: Cauchy double-alternant determinant}\label{prob:GP-1987-24}
Evaluate \(\det(1/(x_i-y_j))\) under the nonvanishing denominator hypotheses.
\par\noindent\textit{Occurrences:} 1981-541.3, \hyperref[item:1987-733.1]{1987-733.1}.
\par\noindent\textit{Status (2026-11):} 2026-11: classical solved identity; the factorization and sign convention are given in the occurrence entry.

\subsection{GP-1987-25: Real zeros in the Legendre Rodrigues formula}\label{prob:GP-1987-25}
Prove that the nth derivative of \((x^2-1)^n\) has n real simple zeros in \((-1,1)\).
\par\noindent\textit{Occurrences:} 1968-160.15, \hyperref[item:1987-733.2]{1987-733.2}.
\par\noindent\textit{Status (2026-11):} 2026-11: classical, by Rolle's theorem or orthogonality.

\subsection{GP-1987-26: Possible discontinuities of Baire-class-one functions}\label{prob:GP-1987-26}
For a pointwise limit of real continuous functions on an interval, determine whether discontinuities can occur, be dense, or fill the interval.
\par\noindent\textit{Occurrences:} \hyperref[item:1987-733.3]{1987-733.3}.
\par\noindent\textit{Status (2026-11):} 2026-11: yes to the first two; no to the last, by Baire's theorem.

\subsection{GP-1987-27: Normalized conformal maps of the slit and dumbbell domains}\label{prob:GP-1987-27}
Construct a disk map for a radially slit disk, and establish a normalized disk map for a simply connected dumbbell domain with its marked interior point sent to0.
\par\noindent\textit{Occurrences:} \hyperref[item:1987-733.4]{1987-733.4}, \hyperref[item:1987-733.5]{1987-733.5}.
\par\noindent\textit{Status (2026-11):} 2026-11: classical mapping theory settles existence; the slit case has an elementary composition. The picture alone does not specify a numerical dumbbell map.

\subsection{GP-1987-28: Principal-angle classification and realization in finite dimension}\label{prob:GP-1987-28}
For two subspaces of fixed dimensions k,m in an n-dimensional positive inner-product space, determine a complete orthogonal or unitary orbit invariant and all its attainable values, including realization relative to a fixed first subspace.
\par\noindent\textit{Occurrences:} 1964-46.2, \hyperref[item:1987-735.2]{1987-735.2}, \hyperref[item:1987-737.3]{1987-737.3}, \hyperref[item:1987-737.5]{1987-737.5}, \hyperref[item:1987-737.14]{1987-737.14}.
\par\noindent\textit{Status (2026-11):} 2026-11: solved classical problem. The ordered principal angles lie in [0,pi/2], with at least max(0,k+m-n) zeros; these are exactly the restrictions.

\subsection{GP-1987-29: Angles and attainable invariants of tensor subspaces}\label{prob:GP-1987-29}
For nonzero \(a,b\in L\otimes L\) with finite-dimensional positive inner-product L, classify the pair \(a\otimes L,L\otimes b\subset L^{\otimes3}\) by principal angles and determine their full possible range.
\par\noindent\textit{Occurrences:} \hyperref[item:1987-737.11]{1987-737.11}.
\par\noindent\textit{Status (2026-11):} 2026-11: the chapter annotation at Section~\ref{dev:tensor-angles} proves that for \(\dim L=d\ge2\), precisely the nonnegative ordered principal cosines of sum at most one occur; for d=1 only cosine one occurs. The result is presented as a worked solution of an established linear-algebra exercise, without a novelty claim.

\subsection{GP-1987-31: Pairs of subspaces in finite-dimensional indefinite or degenerate forms}\label{prob:GP-1987-31}
Fix a finite-dimensional real vector space with a symmetric bilinear form, possibly indefinite or degenerate. Classify ordered pairs of subspaces under the full group preserving that form, including null subspaces.
\par\noindent\textit{Occurrences:} \hyperref[item:1987-737.7]{1987-737.7}, \hyperref[item:1987-737.8]{1987-737.8}.
\par\noindent\textit{Status (2026-11):} 2026-11: finite-dimensional canonical forms are available in Sergeichuk (1990). This does not transfer to the unspecified infinite-dimensional indefinite variant.

\subsection{GP-1987-32: Complete invariants and realization for closed Hilbert-subspace pairs}\label{prob:GP-1987-32}
Classify an ordered pair of closed Hilbert subspaces under ambient isometries, retaining both discrete intersections and the generic spectral invariant; determine the admissible data and give discrete and continuous examples.
\par\noindent\textit{Occurrences:} \hyperref[item:1987-735.3]{1987-735.3}, \hyperref[item:1987-737.12]{1987-737.12}, \hyperref[item:1987-737.14]{1987-737.14}.
\par\noindent\textit{Status (2026-11):} 2026-11: solved by the two-projection canonical form (Halmos, 1969). The invariant includes spectral measure class and multiplicity, not just the spectrum as a set; dimension constraints govern realization.

\subsection{GP-1987-35: Configurations of at most three linear subspaces}\label{prob:GP-1987-35}
Classify one subspace, ordered pairs and ordered triples of subspaces of a finite-dimensional real or complex vector space under GL, and characterize the attainable invariant data.
\par\noindent\textit{Occurrences:} 1972-277.5, 1972-277.6, \hyperref[item:1987-737.6]{1987-737.6}.
\par\noindent\textit{Status (2026-11):} 2026-11: solved finite-type classification; triples have eight one-dimensional membership blocks and one two-dimensional three-line block, with nine nonnegative multiplicities.

\subsection{GP-1987-36: Four-subspace classification and operator encoding}\label{prob:GP-1987-36}
Classify ordered quadruples in a finite-dimensional vector space up to simultaneous linear isomorphism, and include the similarity problem for one operator through its graph configuration.
\par\noindent\textit{Occurrences:} 1970-198.1, 1970-198.3, 1970-198.6, 1970-198.13, 1972-277.7, \hyperref[item:1987-737.10]{1987-737.10}, \hyperref[item:1987-737.15]{1987-737.15}.
\par\noindent\textit{Status (2026-11):} 2026-11: the algebraically closed characteristic-zero problem was classified by Gelfand--Ponomarev in 1971 using explicit indecomposable families. It is an established result explicitly recalled in this problem sheet.

\subsection{GP-1987-33: An intrinsic regularity criterion for local functionals}\label{prob:GP-1987-33}
Characterize when a scalar functional on an open subset of \(C^\infty(M)\) is, locally in configuration space, an integral of a smooth finite-jet density; determine what must supplement disjoint-support additivity.
\par\noindent\textit{Occurrences:} \hyperref[item:1987-738.2]{1987-738.2}.
\par\noindent\textit{Status (2026-11):} 2026-11: continuity and additivity alone are insufficient, as point evaluation shows. Under Bastiani smoothness, smooth-density differential and smooth dependence of the gradient, the exact local finite-jet characterization is proved by Brouder et al. (preprint2017; publication2018).

\subsection{GP-1987-34: Relations among quantum Miura coefficients}\label{prob:GP-1987-34}
For independent Heisenberg currents and \(L=:(D+v_1)\cdots(D+v_N):\), determine the singular operator products of the coefficient fields and prove closure of the resulting quantum Gelfand--Dikii algebra.
\par\noindent\textit{Occurrences:} \hyperref[item:1987-738.17]{1987-738.17}, \hyperref[item:1987-738.18]{1987-738.18}, \hyperref[item:1987-738.19]{1987-738.19}.
\par\noindent\textit{Status (2026-11):} 2026-11: solved in Lukyanov (1988), received25February1988 before the Radul report. Theorem 1 gives quadratic closure, Section 2 its regularity relation and explicit coefficients, and Theorem 2 the reduced \(W_N\) construction.

\subsection{GP-1987-37: Variational calculus for products and multilocal functionals}\label{prob:GP-1987-37}
Extend variational differentiation from local scalar functionals to finite sums and products, while distinguishing that generated algebra from arbitrary smooth kernels depending on jets at several base points.
\par\noindent\textit{Occurrences:} \hyperref[item:1987-738.6]{1987-738.6}.
\par\noindent\textit{Status (2026-11):} 2026-11: finite products satisfy the ordinary iterated Leibniz calculus. Rigorous smooth-functional and multilocal settings are developed, including Brouder et al. (2017/2018). A topology and a specified class of kernels are needed for questions about completion; the notebook states no unique such completion problem.

\subsection{GP-1987-41: Hearts realizing the derived category of a quiver}\label{prob:GP-1987-41}
Let k be an algebraically closed field and Q a finite acyclic quiver. Describe bounded t-structures on \(D^b(\operatorname{rep}_kQ)\) whose heart \(\mathcal A\) is a finite-length abelian k-category with a projective generator and finite-dimensional Hom spaces, and whose realization functor \(D^b(\mathcal A)\to D^b(\operatorname{rep}_kQ)\) is an equivalence.
\par\noindent\textit{Occurrences:} \hyperref[item:1987-739.3]{1987-739.3}.
\par\noindent\textit{Status (2026-11):} The Dynkin algebra endpoint was established by \citet[Theorem~5.12]{happel1987derived} in 1987, after an Ottawa 1984 announcement. Parthasarathy's December 1988 publication supplies the precise t-structure-data realization criterion and construction reported here; this later refinement is the class-B publication match. Receipt on 15 February 1988 is not itself proof of public circulation before the seminar. Arbitrary embedded abelian subcategories are outside the identified theorem.

\subsection{GP-1987-42: Fermat's last theorem}\label{prob:GP-1987-42}
For every integer \(n>2\), there are no positive integers \(a,b,c\) with \(a^n+b^n=c^n\). Finiteness of rational points on a fixed genus-greater-than-one curve does not by itself prove this nonexistence statement.
\par\noindent\textit{Occurrences:} \hyperref[item:1987-741.3]{1987-741.3}.
\par\noindent\textit{Status (2026-11):} Settled by the modularity method published in 1995; the 1988 marginal mention of Faltings did not settle the equation.

\subsection{GP-1987-43: Four-colour theorem and reducibility}\label{prob:GP-1987-43}
Does every finite loopless planar graph admit a proper four-colouring?
\par\noindent\textit{Occurrences:} \hyperref[item:1987-741.4]{1987-741.4}, \hyperref[item:1987-745.2]{1987-745.2}.
\par\noindent\textit{Status (2026-11):} Settled by Appel and Haken in 1976, published in 1977; a complete formal proof was subsequently constructed by Gonthier.

\subsection{GP-1987-44: Riemann hypothesis}\label{prob:GP-1987-44}
Every zero \(\rho\) of the analytically continued Riemann zeta function with \(0<\Re\rho<1\) satisfies \(\Re\rho=1/2\).
\par\noindent\textit{Occurrences:} \hyperref[item:1987-741.5]{1987-741.5}.
\par\noindent\textit{Status (2026-11):} Open as of 11 November 2026, as confirmed by the official Clay problem page.

\subsection{GP-1987-45: Parallel implementation of configuration reducibility tests}\label{prob:GP-1987-45}
Given a finite ring configuration, implement a specified sound finite reducibility test based on boundary colourings and compatible Kempe data, allowing independent colourings/configurations to be processed in parallel. No machine model, complexity bound or completeness criterion for all notions of reducibility is stated.
\par\noindent\textit{Occurrences:} \hyperref[item:1987-745.7]{1987-745.7}.
\par\noindent\textit{Status (2026-11):} Established finite reducibility algorithms and later verified implementations give a partial match; the exact parallel architecture and requested performance were never specified. No new complexity claim is made.

\subsection{GP-1987-46: Microlocal reconstruction of perverse sheaves}\label{prob:GP-1987-46}
Let \(X\) be a complex manifold and \(\Lambda\subset T^*X\) a specified closed complex analytic conic Lagrangian. With coefficients in \(\mathbb C\), describe the category of perverse sheaves with singular support in \(\Lambda\) by local Morse data on the smooth pieces, gluing at codimension-one loci, and explicit compatibility conditions at singular loci. A complete answer must characterize the essential image, not only embed the category faithfully.
\par\noindent\textit{Occurrences:} \hyperref[item:1987-746.4]{1987-746.4}.
\par\noindent\textit{Status (2026-11):} Partial over complex coefficients: GMV 2005 proves the fully faithful embedding and the codimension-one equivalence; Kashiwara--Vilonen 2014 (announced 2010, full preprint 2012) proves extension across codimension at least three. Their Problem~9.3, p.~617, explicitly asks for the geometric codimension-two conditions on the essential image. No general solution to that remaining task was identified through 11 November 2026.

\section{Corrections to the transcription and translation}\label{corr:1987}
The evidence distinguishes a transcription correction from a mathematical emendation of wording faithfully preserved in the notebook.
\subsection*{P001386 (both; faithful error): sentence following flabby-sheaf definition}
\noindent\textit{Transcription reads.} Locally constant sheaves are not flabby\par
\noindent\textit{Critical reading.} The displayed example obstructs flabbiness only with additional hypotheses; zero sheaves and sheaves on discrete spaces are flabby.\par
\noindent\textit{Evidence.} The Russian text and scan make the unrestricted assertion. The counterexamples follow immediately from the definition.\par
\subsection*{P001389 (both; faithful error): interpolation substitution}
\noindent\textit{Transcription reads.} \(f(x)\equiv0\)\par
\noindent\textit{Critical reading.} \(f(x)\equiv1\)\par
\noindent\textit{Evidence.} The scan supports the recorded zero; Good 1970, p.~1884, explicitly interpolates the function one, and only that gives the displayed identity.\par
\subsection*{P001389 (both; provenance): 1970 proof attribution}
\noindent\textit{Transcription reads.} Vuda / \ru{Вуда}\par
\noindent\textit{Critical reading.} I.~J. Good (\ru{Гуд}); retain the handwritten form in the diplomatic source\par
\noindent\textit{Evidence.} Good 1970 reproduces the exact interpolation and recurrence. This is a bibliographic emendation, not an asserted palaeographic rereading.\par
\subsection*{P001391 (both; faithful error): definition of angular momentum}
\noindent\textit{Transcription reads.} \(M=[\dot x,x]\), with the displayed standard brackets and \(\dot f=\{f,H\}\)\par
\noindent\textit{Critical reading.} \(M=[x,\dot x]\), or explicitly reverse the cross-product convention\par
\noindent\textit{Evidence.} Direct epsilon contraction gives \(\dot p=M\times p\); on \(p^2=1\) and \(M\cdot p=0\), this forces \(M=p\times\dot p\). The scan and Russian transcription agree.\par
\subsection*{P001391 (both; faithful error): second bracket and Hamiltonian}
\noindent\textit{Transcription reads.} Mixed bracket \(\varepsilon_{ijk}a_kp_k\) and momentum bracket \(\varepsilon_{ijk}M_k\); \(F=p^2/2\)\par
\noindent\textit{Critical reading.} The table fails Jacobi. One verified repair uses \(\varepsilon_{ijk}a_jp_k\) and \(-\varepsilon_{ijk}M_k\), keeping \(F=p^2/2\)\par
\noindent\textit{Evidence.} The scan confirms both literal entries. \(\operatorname{Jac}(M_1,p_1,p_2)=(a_3-a_2)M_2\). Exact symbolic checks show the repaired table is Poisson and produces the first Neumann vector field; the repair is a reconstruction, not a claimed unique intention.\par
\subsection*{P001395 (both; faithful error): last derivative telescoping row}
\noindent\textit{Transcription reads.} The endpoint difference of \(f^{(p)}\) begins with \(S_p\)\par
\noindent\textit{Critical reading.} It begins with \(S_{p+1}\)\par
\noindent\textit{Evidence.} Differentiating to order \(p\) and Taylor expanding \(f^{(p)}(j+1)-f^{(p)}(j)\) makes the first term \(f^{(p+1)}(j)\); the scan confirms the printed index.\par
\subsection*{P001395 (both; faithful error): exponential generating function}
\noindent\textit{Transcription reads.} \(\big((n+1)^{-1}\exp(t(n+1))-e\big)/(\exp(t)-1)\)\par
\noindent\textit{Critical reading.} \(\sum_{j=1}^{n}\exp(tj)=\frac{\exp(t(n+1))-\exp(t)}{\exp(t)-1}\)\par
\noindent\textit{Evidence.} Direct geometric-series identity; the notebook fraction is inconsistent even at \(n=1\). Treat this as a mathematical emendation, not a silent transcription correction.\par
\subsection*{P001395 (both; faithful error): scope of summation formula}
\noindent\textit{Transcription reads.} For any functions\par
\noindent\textit{Critical reading.} For polynomials the finite formula is exact; other functions require convergence hypotheses or a finite Euler--Maclaurin remainder\par
\noindent\textit{Evidence.} NIST DLMF 2.10.1 supplies explicit analytic assumptions and remainder.\par
\subsection*{P001397 (both; faithful error): mapping-cone differential}
\noindent\textit{Transcription reads.} \(C_f^i=Y^i\oplus X^{i+1}\), with \(d(x,y)=(-dx,dy+fx)\)\par
\noindent\textit{Critical reading.} With the declared ordering, \(d(y,x)=(d_Yy+fx,-d_Xx)\)\par
\noindent\textit{Evidence.} The scan and Russian text combine opposite orderings; direct verification of \(d^2=0\) fixes the consistent form.\par
\subsection*{P001399 (both; faithful error): matrix terminology}
\noindent\textit{Transcription reads.} Hankel matrix\par
\noindent\textit{Critical reading.} Upper triangular Toeplitz matrix\par
\noindent\textit{Evidence.} The displayed entry is \(h_{j-i}\), constant along diagonals, whereas a Hankel matrix depends on \(i+j\). The scan confirms the heading.\par
\subsection*{P001399 (both; faithful error): alphabet with three variables}
\noindent\textit{Transcription reads.} \(P=\alpha u+\beta u+\gamma w\)\par
\noindent\textit{Critical reading.} \(P=\alpha u+\beta v+\gamma w\)\par
\noindent\textit{Evidence.} The generating product immediately below uses distinct \(u,v,w\). The scan repeats \(u\) in the first line; this is a mathematical emendation.\par
\subsection*{P001404 (both; faithful error): dimension following the six-element example}
\noindent\textit{Transcription reads.} \(n!\)\par
\noindent\textit{Critical reading.} \((n+1)!\) for the displayed \(n+1\)-variable convention\par
\noindent\textit{Evidence.} The explicit code bounds give \((n+1)n\cdots1\) choices. The \(n=2\) example lists the six \(S_3\) Schubert classes.\par
\subsection*{P001405 (both; ambiguity): character request}
\noindent\textit{Transcription reads.} Irreducible representations with Schubert-function characters\par
\noindent\textit{Critical reading.} Cyclic Borel modules give the established realization; the intended irreducible category is unspecified\par
\noindent\textit{Evidence.} Kraśkiewicz--Pragacz, Theorem~4.1, matches the character request but not irreducibility. Its introduction records the 1986 preprint and 1987 announcement.\par
\subsection*{P001405-P001406 (both; ambiguity): heading and four kinds of data}
\noindent\textit{Transcription reads.} 2-category\par
\noindent\textit{Critical reading.} Double category for the displayed data; the later parallel-arrow definition is a strict 2-category\par
\noindent\textit{Evidence.} Brown--Spencer 1976, \S~1, defines precisely the objects, two arrow directions and squares shown here.\par
\subsection*{P001407 (both; faithful error): opening Ore diagram}
\noindent\textit{Transcription reads.} The outgoing arrow \(X\to Z^{\prime}\) is labelled quasi-isomorphism\par
\noindent\textit{Critical reading.} The outgoing arrow \(Y\to Z^{\prime}\) is the quasi-isomorphism\par
\noindent\textit{Evidence.} The scan labels the left outgoing leg. Homotopy pushout along the incoming quasi-isomorphism \(Z\to X\) makes the opposite leg \(Y\to P\) a quasi-isomorphism.\par
\subsection*{P001409 (both; faithful error): affine map with unit multiplier}
\noindent\textit{Transcription reads.} If \(|a|=1\), all points lie on invariant circles carrying rotations\par
\noindent\textit{Critical reading.} Require \(a\ne1\), or \(b=0\); \(a=1\), \(b\ne0\) is a translation\par
\noindent\textit{Evidence.} The scan explicitly writes \(|a|=1\). The affine fixed-point equation exposes the missing exceptional case.\par
\subsection*{P001410 (both; faithful error): normal-family definition}
\noindent\textit{Transcription reads.} The set of iterates on \(U\) is compact\par
\noindent\textit{Critical reading.} The iterates form a normal family, equivalently are relatively compact for locally uniform spherical convergence\par
\noindent\textit{Evidence.} The handwritten word is compact; the standard normal-family definition concerns subsequential limits, which need not be iterates themselves.\par
\subsection*{P001410 (both; faithful error): Sullivan statement, part (a)}
\noindent\textit{Transcription reads.} Every component orbit stabilizes\par
\noindent\textit{Critical reading.} Every component orbit is eventually periodic\par
\noindent\textit{Evidence.} Sullivan 1985, Theorem~1, p.~404, uses eventual periodicity. Attracting cycles of period greater than one prevent eventual constancy.\par
\subsection*{P001410 (both; ambiguity): disk-packing questions}
\noindent\textit{Transcription reads.} Inscribe disks; residual area \(0\) and dimension \(<2\)\par
\noindent\textit{Critical reading.} Retain both questions with the insertion rule explicitly marked missing\par
\noindent\textit{Evidence.} Visual inspection finds no selection or iteration rule besides the drawing and the word inscribe. Arbitrary disjoint disks do not imply either conclusion.\par
\subsection*{P001411 (both; ambiguity): general representability problem}
\noindent\textit{Transcription reads.} Which matroids are real or complex representable is unsolved\par
\noindent\textit{Critical reading.} The finite decision problem is decidable; a stronger intended structural classification is unspecified\par
\noindent\textit{Evidence.} Tarski 1951, \S\S~2--3, states real quantifier elimination and its extension to complex algebra; encode the finitely many required vanishing and nonvanishing maximal minors.\par
\subsection*{P001412 (both; provenance): name in bimatroid comparison}
\noindent\textit{Transcription reads.} \ru{Кунтом}, tentatively transcribed as Knuth\par
\noindent\textit{Critical reading.} Retain the uncertain handwritten name; Kung 1978 is the relevant bimatroid lead\par
\noindent\textit{Evidence.} Visual inspection does not settle the surname. Kung authored \emph{Bimatroids and invariants} in 1978, whereas the Gelfand--Serganova passage remains unidentified.\par
\subsection*{P001412 (both; faithful error): equation(6) determinant exchange}
\noindent\textit{Transcription reads.} \(\sum_i(-1)^{d-i}[y_1,x_2,\ldots,x_d][y_1,\ldots,y_{i-1},x_i,y_{i+1},\ldots,y_d]\)\par
\noindent\textit{Critical reading.} \(\sum_i[y_i,x_2,\ldots,x_d][y_1,\ldots,y_{i-1},x_1,y_{i+1},\ldots,y_d]\)\par
\noindent\textit{Evidence.} The scan and authoritative transcription agree on the erroneous signs and indices. Expanding \(x_1\) in the \(y\) basis proves the emended identity; a two-dimensional exact counterexample disproves the literal one.\par
\subsection*{P001412 (both; faithful error): basis-exchange cross-reference}
\noindent\textit{Transcription reads.} This gives exchange axiom (3)\par
\noindent\textit{Critical reading.} This gives exchange axiom (2)\par
\noindent\textit{Evidence.} The preceding page numbers the basis-exchange axiom (2) and the bimatroid discussion (3); the scan preserves the wrong reference.\par
\subsection*{P001413 (both; provenance): recommended author name}
\noindent\textit{Transcription reads.} Tipper / \ru{Типпера}\par
\noindent\textit{Critical reading.} Richard Dipper / \ru{Диппера}\par
\noindent\textit{Evidence.} The original handwritten name is preserved in both transcriptions; the opened Wiley publisher record identifies Richard Dipper and Gordon James, PLMS 52(1), 20--52, January 1986.\par
\subsection*{P001414 (both; faithful error): equation(9) maximum range}
\noindent\textit{Transcription reads.} \(1\le k\le N\)\par
\noindent\textit{Critical reading.} \(1\le k<N\)\par
\noindent\textit{Evidence.} \(r_k\) uses cells \(k\) and \(k+1\) of a tableau with \(N\) cells; Kerov 1987, \S~2, p.~157, has \(k=1,\ldots,N-1\). The scan confirms the notebook upper limit \(N\).\par
\subsection*{P001415 (both; faithful error): equations(16)--(17) recurrence range}
\noindent\textit{Transcription reads.} \(r\ge0\)\par
\noindent\textit{Critical reading.} \(r\ge1\), with \(W_{-1}=1\) and \(W_1=0\)\par
\noindent\textit{Evidence.} The scan explicitly includes zero. The rational formula at \(r=0\) has denominator \(q^0-1=0\); the original Kerov appendix starts the recurrence at positive \(r\).\par
\subsection*{P001415 (both; faithful error): equation(19)}
\noindent\textit{Transcription reads.} \(W_r(\tau)W_{1-r}(\tau)=1\)\par
\noindent\textit{Critical reading.} \(W_r(\tau)W_{1-r}(\tau)=\tau\), for \(r\ge2\)\par
\noindent\textit{Evidence.} The scan and Russian editorial note confirm the right side \(1\). Kerov 1987, Appendix (P5), and independent simplification of notebook (20)--(21) give \(\tau=q/(q+1)^2\).\par
\subsection*{P001416 (both; faithful error): weighted zeta definition, last display}
\noindent\textit{Transcription reads.} \(\prod_{k=0}^{n-1}\varphi(T^k x)\) with sum over \(\xi\)\par
\noindent\textit{Critical reading.} \(\prod_{k=0}^{n-1}\varphi(T^k\xi)\)\par
\noindent\textit{Evidence.} Scan inspected: the notebook itself has x while the outer summation variable is xi. Dummy-variable emendation, not a mistranscription.\par
\subsection*{P001416 (both; provenance): Mayer attribution}
\noindent\textit{Transcription reads.} Lindon Mayer (?)\par
\noindent\textit{Critical reading.} D. H. Mayer (bibliographic identification; retain uncertainty about the reporter)\par
\noindent\textit{Evidence.} The scan has a doubtful name; the exact title, journal volume 14 and year match the opened Springer article by Dieter H. Mayer.\par
\subsection*{P001418 (both; faithful error): weighted zeta, equation (5)}
\noindent\textit{Transcription reads.} \(\prod_{k=0}^{n}\varphi(T^k x)\)\par
\noindent\textit{Critical reading.} \(\prod_{k=0}^{n-1}\varphi(T^k x)\)\par
\noindent\textit{Evidence.} Scan inspected: n is in the notebook. A one-symbol shift with weight c falsifies the determinant identity with n+1 factors; this is an emendation of the mathematical formula.\par
\subsection*{P001419 (both; faithful error): product formula domain}
\noindent\textit{Transcription reads.} \(a\in\mathbb Q\)\par
\noindent\textit{Critical reading.} \(a\in\mathbb Q^\times\)\par
\noindent\textit{Evidence.} Scan and both transcriptions omit the nonzero condition; the formula fails for zero.\par
\subsection*{P001419 (both; faithful error): gamma ratio, equation (3)}
\noindent\textit{Transcription reads.} \(\pi^{-s+1}\)\par
\noindent\textit{Critical reading.} \(\pi^{1/2-s}\)\par
\noindent\textit{Evidence.} The scan visibly has -s+1. Dividing the two sides of equation (2), also confirmed against DLMF 25.4, gives 1/2-s. Faithful notebook error.\par
\subsection*{P001420 (both; ambiguity): adelic product notation}
\noindent\textit{Transcription reads.} \(\mathrm{SL}(n,\mathbb A)\sim\mathrm{SL}(n,\mathbb R)\times\mathrm{SL}(n,\mathbb Q_2)\times\cdots\)\par
\noindent\textit{Critical reading.} Use the restricted product with respect to \(\mathrm{SL}(n,\mathbb Z_p)\).\par
\noindent\textit{Evidence.} The displayed ordinary product is shorthand; the restricted-product condition is required for adeles. Mathematical completion, not alteration of handwriting.\par
\subsection*{P001422 (both; faithful error): Deligne formula}
\noindent\textit{Transcription reads.} \(\det H^\bullet=\det F_K\)\par
\noindent\textit{Critical reading.} An inverse is needed for determinant weights under the stated convention; a full corrected epsilon-line formula needs additional hypotheses\par
\noindent\textit{Evidence.} Scan inspected: no inverse is written. A constant local system has weight 2-2g on the left and 2g-2 on the right. Do not present this sign correction as a proved universal replacement.\par
\subsection*{P001423 (both; ambiguity): filtration equivalence}
\noindent\textit{Transcription reads.} Distributivity on every cone, with no smoothness condition\par
\noindent\textit{Critical reading.} For smooth toric varieties this is sufficient; for a singular cone impose an integral character decomposition\par
\noindent\textit{Evidence.} Klyachko Theorems 0.1.1 and2.2.1, full English paper read. The scan states the short condition without this distinction.\par
\subsection*{P001424 (both; ambiguity): tetrahedral decomposition claim}
\noindent\textit{Transcription reads.} An arbitrary representation decomposes into one-dimensional objects and the indicated exterior powers\par
\noindent\textit{Critical reading.} State the short-filtration restriction of Klyachko Theorem 6.4.1\par
\noindent\textit{Evidence.} The source scan has an unrestricted assertion; the matching published theorem explicitly restricts filtration length. A mathematical qualification, not a transcription correction.\par
\subsection*{P001426 (both; faithful error): principal reduction, transitivity}
\noindent\textit{Transcription reads.} \(G\) acts transitively on \(P\).\par
\noindent\textit{Critical reading.} \(G\) acts freely and transitively on every fibre \(P_x\).\par
\noindent\textit{Evidence.} Scan inspected: global transitivity is written. The principal-bundle definition and dimension of the base require a fibrewise assertion.\par
\subsection*{P001427 (both; faithful error): quotient below the coframe map}
\noindent\textit{Transcription reads.} \(T_pM/T_p^vM\)\par
\noindent\textit{Critical reading.} \(T_pP/T_p^vP\)\par
\noindent\textit{Evidence.} Scan inspected: \(M\) appears in the quotient; \(p\) is a point of \(P\), and its vertical tangent is in \(T_pP\). Both transcriptions follow the manuscript.\par
\subsection*{P001428 (both; ambiguity): torsion equivariance}
\noindent\textit{Transcription reads.} Hom(V,g)-equivariance with unspecified target action\par
\noindent\textit{Critical reading.} The change-of-complement action on torsion is affine via the Spencer alternation\par
\noindent\textit{Evidence.} Scan inspected; Sections 6.1--6.2 of Alekseevsky--Michor give the precise prolongation construction.\par
\subsection*{P001430 (both; faithful error): repeated translation Ward identity}
\noindent\textit{Transcription reads.} \(\varepsilon_{\mu\nu}\,dx^\lambda\)\par
\noindent\textit{Critical reading.} \(\varepsilon_{\mu\lambda}\,dx^\lambda\)\par
\noindent\textit{Evidence.} Scan inspected. The preceding displayed contour has the lambda contraction; the repetition leaves lambda uncontracted and repeats the free index nu.\par
\subsection*{P001431 (both; faithful error): complex stress components and equation (11)}
\noindent\textit{Transcription reads.} \(T=T^{zz},\quad\partial_{\bar z}T=0\)\par
\noindent\textit{Critical reading.} Use the covariant component \(T_{zz}\) for the holomorphic field, or exchange the asserted chiralities.\par
\noindent\textit{Evidence.} Scan inspected: the indices are written upstairs. Direct tensor transformation for \(z=x^1+ix^2\) gives \(\partial_zT^{zz}=0\). The trace identity on P001433 is consistent with upstairs indices and is not separately erroneous.\par
\subsection*{P001431 (both; ambiguity): contour definition of descendant modes}
\noindent\textit{Transcription reads.} \(z^{n+1}\) around insertion \(z_1\)\par
\noindent\textit{Critical reading.} \((z-z_1)^{n+1}\) for local descendant modes; \(z^{n+1}\) for global modes about the origin.\par
\noindent\textit{Evidence.} The scan faithfully uses \(z^{n+1}\); the distinction becomes material for the later highest-weight formulas.\par
\subsection*{P001434 (both; faithful error): barred scaling coordinate}
\noindent\textit{Transcription reads.} \(\bar z\mapsto\bar\lambda z\)\par
\noindent\textit{Critical reading.} \(\bar z\mapsto\bar\lambda\bar z\)\par
\noindent\textit{Evidence.} Scan inspected: the final z has no bar; both transcriptions preserve the missing conjugation.\par
\subsection*{P001435 (both; ambiguity): school problems and acting group}
\noindent\textit{Transcription reads.} Nonintersecting line/plane pairs after a definition using orthogonal maps\par
\noindent\textit{Critical reading.} Specify affine Euclidean isometries for the disjoint cases, or retain additional origin-dependent position data\par
\noindent\textit{Evidence.} Scans P001435--P001436 inspected. Linear subspaces always meet at zero; disjoint affine pairs need at least a distance invariant.\par
\subsection*{P001436 (both; ambiguity): O(n),SO(n) after the degenerate-form variants}
\noindent\textit{Transcription reads.} Use the two group names for all variants without definition\par
\noindent\textit{Critical reading.} Specify the preserving group for the degenerate form; determinant-one need not have index two\par
\noindent\textit{Evidence.} Scan inspected. Arbitrary invertible maps on the radical preserve a degenerate form and can have any nonzero determinant.\par
\subsection*{P001438 (both; faithful error): operator (1), repeated coefficient subscript}
\noindent\textit{Transcription reads.} \(\sum_{i=0}^{N-1}u_{N-1}(z)\partial_z^i\)\par
\noindent\textit{Critical reading.} \(\sum_{i=0}^{N-1}u_{N-i}(z)\partial_z^i\), or rename consistently as \(\sum a_i(z)\partial_z^i\).\par
\noindent\textit{Evidence.} Scan inspected. The repeated N-1 is in the source; the later Miura line uses N-i and the gradient uses the opposite coefficient ordering. This is a mathematical reindexing, not a claim of mistranscription.\par
\subsection*{P001438 (both; faithful error): continuity and locality assertion}
\noindent\textit{Transcription reads.} Continuity plus disjoint additivity forces a finite-jet smooth integral\par
\noindent\textit{Critical reading.} Add smooth-density and smooth-gradient hypotheses; the valid theorem is local in configuration space\par
\noindent\textit{Evidence.} Point evaluation is a continuous additive counterexample. Brouder et al. Theorem VI.3 was read and supplies the precise additional hypotheses.\par
\subsection*{P001439 (both; provenance): coauthor surname}
\noindent\textit{Transcription reads.} Gelfand, Gabrielov, Lo[shi?]k\par
\noindent\textit{Critical reading.} Identify the reference as Gabrielov, Gelfand and M. V. Losik, 1976\par
\noindent\textit{Evidence.} Scan inspected; the opened MathNet faa2146 record and original introduction give the exact author/topic match.\par
\subsection*{P001440 (both; faithful error): Laurent-symbol ring (6)}
\noindent\textit{Transcription reads.} \(\mathbb C(z)[[\xi]]\) with \(\sum_{j=-\infty}^{N}c_i(z)\xi^i\)\par
\noindent\textit{Critical reading.} \(\mathcal A((\xi^{-1}))\) with \(\sum_{i=-\infty}^{N}c_i(z)\xi^i\)\par
\noindent\textit{Evidence.} Scan inspected. The source bounds powers above while conventional positive-power-series notation means the opposite; dummy indices also disagree.\par
\subsection*{P001441 (both; faithful error): order-two reduction and mode formula (14)--(16)}
\noindent\textit{Transcription reads.} Set u1=0 and use the stated half-normalized modes with central term 1/12\par
\noindent\textit{Critical reading.} Specify the reduction and pairing; under normalized residue the printed modes have coefficient (m-n)/2 and central coefficient -1/8\par
\noindent\textit{Evidence.} Khovanova 1986 Section 7 gives the reduced third-order Hamiltonian operator. Direct integration with the source mode factor1/2 gives the corrected constants; independently checked.\par
\subsection*{P001441 (both; ambiguity): order-three mixed bracket (18)b}
\noindent\textit{Transcription reads.} The raw coefficient u3 transforms as a primary field of weight3\par
\noindent\textit{Critical reading.} Use \(w_3=u_3-\tfrac12\partial u_2\) for the primary mixed bracket.\par
\noindent\textit{Evidence.} The scan identifies \(u_3\) as the raw constant coefficient of \(D^3+u_2D+u_3\). Coordinate covariance produces second- and fourth-derivative terms until the field change.\par
\subsection*{P001442 (both; faithful error): Miura moral}
\noindent\textit{Transcription reads.} A Lie algebra is isomorphic to a quadratic algebra that is not a Lie algebra\par
\noindent\textit{Critical reading.} The coefficient generators need not be linearly closed; the full nonlinear Poisson algebra is still a Lie algebra, and Miura is a Poisson realization\par
\noindent\textit{Evidence.} Khovanova 1986 Lemma 1 and Kupershmidt--Wilson Theorem 11.5 were read. The source sentence confuses generator span with the full Poisson algebra.\par
\subsection*{P001443 (both; ambiguity): boundedness symbol in (23)}
\noindent\textit{Transcription reads.} O(1) or possibly 0(1)\par
\noindent\textit{Critical reading.} Intended O(1): no singular part at z=w\par
\noindent\textit{Evidence.} The handwriting is ambiguous, but Lukyanov 1988 equation 2.15 exactly matches the displayed relation and has O(1).\par
\subsection*{P001443 (both; ambiguity): brackets and signs in (24)}
\noindent\textit{Transcription reads.} Uncertain square brackets with annotation plus or minus\par
\noindent\textit{Critical reading.} Use the ordered-product singular relation of Lukyanov Theorem 1; do not read all brackets as commutators\par
\noindent\textit{Evidence.} Scanned ambiguity retained; original equations 1.6 and2.18--2.21 identify the intended quadratic closure, but do not establish each handwritten sign.\par
\subsection*{P001447 (both; faithful error): Hall conclusion}
\noindent\textit{Transcription reads.} \(f(n)\subset B(n)\)\par
\noindent\textit{Critical reading.} \(f(n)\in B(n)\)\par
\noindent\textit{Evidence.} The manuscript and both transcriptions use a subset sign although f(n) is a single representative; Hall 1935 Theorem 1 supplies the mathematical emendation.\par
\subsection*{P001447 (both; faithful error): Reduction to Rado}
\noindent\textit{Transcription reads.} all subsets of \(N\) are independent\par
\noindent\textit{Critical reading.} use the free matroid on \(M\)\par
\noindent\textit{Evidence.} Representatives belong to \(M\), whereas \(N\) indexes the sets \(B(n)\); this is a faithful error in the notebook, not a translation discrepancy.\par
\subsection*{P001448 (both; faithful error): Markov converse}
\noindent\textit{Transcription reads.} all integer solutions\par
\noindent\textit{Critical reading.} all positive integer solutions\par
\noindent\textit{Evidence.} Ranks of nonzero vector bundles are positive; Rudakov 1988, Introduction and Theorem 4.2, treats the positive Markov tree.\par
\subsection*{P001448 (both; faithful error): Open-orbit remark}
\noindent\textit{Transcription reads.} an unsolvable classification always has a solvable part (open orbit)\par
\noindent\textit{Critical reading.} an open orbit requires additional hypotheses or restriction to a specified stratum\par
\noindent\textit{Evidence.} Simultaneous conjugation of matrix pairs has the varying polynomial invariant \(\operatorname{tr}(A)\), so no orbit is dense in the full representation space.\par
\subsection*{P001449 (both; faithful error): Kernel mutation result}
\noindent\textit{Transcription reads.} \((F,R)\) is exceptional\par
\noindent\textit{Critical reading.} \((R,E)\) is exceptional, under the kernel hypotheses\par
\noindent\textit{Evidence.} The scan clearly writes \((F,R)\). On \(\mathbb P^1\), the evaluation \(\mathcal O^2\to\mathcal O(1)\) has kernel \(\mathcal O(-1)\); \(\operatorname{Hom}(\mathcal O(-1),\mathcal O(1))\ne0\), disproving the printed order.\par
\subsection*{P001450 (both; faithful error): Zero section in SS}
\noindent\textit{Transcription reads.} add the zero section whenever the complex is nonzero\par
\noindent\textit{Critical reading.} add zero covectors only over \(\operatorname{supp}(F)\)\par
\noindent\textit{Evidence.} Scan P001450 explicitly says to add the whole zero section. A skyscraper at \(x\) has singular support only over \(x\); the corrected basic identity is \(SS(F)\cap T_X^*X=\operatorname{supp}(F)\), identifying the zero section with \(X\).\par
\subsection*{P001451 (both; faithful error): Lagrangian versus involutive ideal}
\noindent\textit{Transcription reads.} an equivalent condition is involutivity of the ideal\par
\noindent\textit{Critical reading.} involutivity plus pure dimension \(\dim X\)\par
\noindent\textit{Evidence.} The full cotangent bundle has zero, hence involutive, defining ideal but is not Lagrangian when \(\dim X>0\). The notebook omits the dimension hypothesis.\par
\subsection*{P001452 (both; ambiguity): Perversity degree}
\noindent\textit{Transcription reads.} degree 0; parenthesis: first degree\par
\noindent\textit{Critical reading.} state the cone/vanishing-cycle shift convention\par
\noindent\textit{Evidence.} Both readings occur in the scan and transcriptions. The normalized Morse relative complex is concentrated in degree zero for \(\mathbb C_X[\dim X]\); raw vanishing cycles have a shifted convention.\par
\subsection*{P001452 (both; faithful error): Independence of test function}
\noindent\textit{Transcription reads.} the choice of f makes no difference\par
\noindent\textit{Critical reading.} Morse groups are identified after normalization and Maslov data\par
\noindent\textit{Evidence.} The chosen covector does not supply a canonical untwisted identification across loops in the space of Hessians; GMV 2005 section 1 states the extra structure.\par
\subsection*{P001457 (both; ambiguity): Computer time}
\noindent\textit{Transcription reads.} 12,000 hours, with (1200?) below\par
\noindent\textit{Critical reading.} retain the manuscript doubt; the contemporary published digest reports 1200 hours\par
\noindent\textit{Evidence.} The scan itself queries the larger figure. Bernhart 1977, publisher abstract, explicitly gives 1200 hours for nearly 1900 configurations. This is a factual emendation, not a transcription error.\par
\subsection*{P001457 (both; faithful error): Zero-section equivalence}
\noindent\textit{Transcription reads.} \(SS(F)\) equals the zero section iff every \(H^i(F)\) is locally constant\par
\noindent\textit{Critical reading.} \(SS(F)\) is contained in the zero section iff every \(H^i(F)\) is locally constant; equality only over the support\par
\noindent\textit{Evidence.} The zero object is an immediate counterexample to equality. On disconnected \(X\), a nonzero locally constant complex may vanish on a component.\par
\subsection*{P001462 (both; faithful error): Equation (4), antipode}
\noindent\textit{Transcription reads.} \(S(X^\pm)=q^{H/2}X^\pm q^{-H/2}\)\par
\noindent\textit{Critical reading.} \(S(X^\pm)=-q^{H/4}X^\pm q^{-H/4}\)\par
\noindent\textit{Evidence.} Scan verified. The displayed coproduct forces \(S(X)L+LX=0\) for \(L=q^{H/4}\); Reshetikhin I, equation (1.6), agrees after \(\rho=H/2\).\par
\subsection*{P001462 (both; faithful error): Cartan-generator commutator}
\noindent\textit{Transcription reads.} \((K-K^{-1})/(q-q^{-1})\), with \(K=q^{H/2}\)\par
\noindent\textit{Critical reading.} \((K-K^{-1})/(q^{1/2}-q^{-1/2})\) with the original generators\par
\noindent\textit{Evidence.} Direct substitution into P001461 gives this denominator. A different standard parameter is possible only if all \(q\)-relations are changed consistently.\par
\subsection*{P001463 (both; faithful error): Equation (8), universal R series}
\noindent\textit{Transcription reads.} coefficients \(q^n/[n]_q!\), with \(e=q^{H/4}X^+\) and \(f=q^{-H/4}X^-\)\par
\noindent\textit{Critical reading.} include the missing deformation factors and specify a completed tensor product\par
\noindent\textit{Evidence.} The scan agrees with the transcription. As \(h\to0\), the displayed nonconstant terms do not vanish, contradicting \(R\equiv1\pmod h\); no unique full coefficient is guessed.\par
\subsection*{P001463 (both; faithful error): Equation (10), coproduct identities}
\noindent\textit{Transcription reads.} \((\Delta\otimes\mathrm{id})R=R_{23}R_{13}\); \((\mathrm{id}\otimes\Delta)R=R_{12}R_{13}\)\par
\noindent\textit{Critical reading.} \((\Delta\otimes\mathrm{id})R=R_{13}R_{23}\); \((\mathrm{id}\otimes\Delta)R=R_{13}R_{12}\)\par
\noindent\textit{Evidence.} With \(\Delta^{\mathrm{op}}(a)=R\Delta(a)R^{-1}\), Reshetikhin I, Theorem 1.3, and Reshetikhin--Turaev 1990, equations (3.1.2)--(3.1.3), fix this order; scan verified.\par
\subsection*{P001465 (both; faithful error): Equation (20), crossing domain/codomain}
\noindent\textit{Transcription reads.} \(R^{j_1j_2}:V^{j_1}\otimes V^{j_1}\to V^{j_1}\otimes V^{j_1}\)\par
\noindent\textit{Critical reading.} \(R^{j_1j_2}:V^{j_1}\otimes V^{j_2}\to V^{j_2}\otimes V^{j_1}\)\par
\noindent\textit{Evidence.} The previous definition on P001464 and the differently coloured crossing require both labels; the repeated \(j_1\) is a faithful formula slip.\par
\subsection*{P001467 (both; faithful error): Equation (31), projection correction}
\noindent\textit{Transcription reads.} \(q\) raised to an under/over count without a colour-dependent twist\par
\noindent\textit{Critical reading.} use the ribbon twist scalar for each colour and the signed self-writhe\par
\noindent\textit{Evidence.} Scan verified; a Reidemeister-I move changes ribbon evaluation by the twist of the colour, not a universal \(q\) factor. Reshetikhin--Turaev 1990, Sections 5--6, gives the categorical normalization.\par
\subsection*{P001468 (both; faithful error): Equation (5), root condition}
\noindent\textit{Transcription reads.} the summation index \(i\) is also free; \(b_\alpha(\lambda_i-a_\alpha)\) lacks \(\Delta\)\par
\noindent\textit{Critical reading.} for each \(i\): \(\sum_{j\ne i}\Delta(\lambda_i-\lambda_j)+\sum_\alpha b_\alpha\Delta(\lambda_i-a_\alpha)=0\) in the rational specialization\par
\noindent\textit{Evidence.} Scan verified. Cancelling the residue at \(\lambda_i\) in \(y^{\prime}+y^2+2By+C\) forces this equation; the 1989 review, Section 3, gives the matching pole method.\par
\subsection*{P001468 (both; faithful error): Example (7)2}
\noindent\textit{Transcription reads.} \(\Omega=\omega^2/3\)\par
\noindent\textit{Critical reading.} \(\Omega=-\omega^2/3\)\par
\noindent\textit{Evidence.} Scan verified. For \(x+y+z=0\), \(\cot(\omega x)\cot(\omega y)+\cot(\omega y)\cot(\omega z)+\cot(\omega z)\cot(\omega x)=1\); the written plus-\(\Omega\) convention requires the negative constant.\par
\subsection*{P001468 (both; faithful error): Example (7)3}
\noindent\textit{Transcription reads.} \(\Omega=\zeta^2-\wp\)\par
\noindent\textit{Critical reading.} \(\Omega=(\zeta^2-\wp)/2\)\par
\noindent\textit{Evidence.} Scan verified. The identity \((\zeta(x)+\zeta(y)+\zeta(z))^2=\wp(x)+\wp(y)+\wp(z)\), for \(x+y+z=0\), gives the factor \(1/2\); NIST DLMF 23.10 supplies the addition identity.\par
\subsection*{P001469 (both; faithful error): Repulsive electrostatic potential and minimum}
\noindent\textit{Transcription reads.} \(q_iq_j\log|x_i-x_j|\), followed by repulsion and a minimum\par
\noindent\textit{Critical reading.} use \(-q_iq_j\log|x_i-x_j|\) for the repulsive energy convention\par
\noindent\textit{Evidence.} Differentiation shows the Bethe equations are critical points with either overall sign, but the trapping/minimum argument uses negative logarithms.\par
\subsection*{P001469 (both; faithful error): Stability assertion}
\noindent\textit{Transcription reads.} quantum stability iff Coulomb stability\par
\noindent\textit{Critical reading.} model-specific algebraic correspondence; check quantum domains and endpoint thresholds separately\par
\noindent\textit{Evidence.} Ushveridze 1989 section5 p1212 explicitly distinguishes the classical and quantum stability bounds, saying they almost coincide.\par
\subsection*{P001469 (both; faithful error): Equation (8), pole normalization}
\noindent\textit{Transcription reads.} a solution with a singularity at zero\par
\noindent\textit{Critical reading.} assume a simple pole of residue one and \(\Omega\) regular at zero\par
\noindent\textit{Evidence.} The epsilon-to-zero derivation has the derivative multiplied by the pole residue; unit residue is needed for the stated Riccati identity.\par
\subsection*{P001471 (both; faithful error): Exterior stability bound}
\noindent\textit{Transcription reads.} \(b_1+b_2+M<0\) as the stability condition\par
\noindent\textit{Critical reading.} a sufficient, not universal necessary-and-sufficient, bound; specify the energy and endpoints\par
\noindent\textit{Evidence.} For \(M=1\), \(a_1=0\), \(a_2=1\), \(b_1=-1\), \(b_2=1/2\), the function \(U=\log\lambda-\tfrac12\log(\lambda-1)\) has a strict minimum at \(\lambda=2\), although the displayed sum is \(1/2\). The condition is stronger than required here.\par
\subsection*{P001472 (both; faithful error): Gradient convention}
\noindent\textit{Transcription reads.} \(\tau=\operatorname{grad}f\), with \(\dim T(x_1,x_2)=\operatorname{ind}(x_1)-\operatorname{ind}(x_2)-1\)\par
\noindent\textit{Critical reading.} use \(\tau=-\operatorname{grad}f\) for the displayed source-to-target index difference\par
\noindent\textit{Evidence.} Along the positive gradient, \(f\) increases and the stable/unstable dimension formula reverses. The displayed chain differential lowers Morse index.\par
\subsection*{P001474 (both; faithful error): Primitive of nu}
\noindent\textit{Transcription reads.} \(\int_{M\times I}\operatorname{tr}(F\wedge F)\) as the unnormalized primitive\par
\noindent\textit{Critical reading.} insert the factor \(1/2\) relative to \(\nu(B)=\int_M\operatorname{tr}(F\wedge B)\), then choose integral normalization\par
\noindent\textit{Evidence.} Differentiating \(\int_M\operatorname{tr}(A\wedge dA+\tfrac23 A^3)\) gives \(2\int_M\operatorname{tr}(F\wedge B)\). Overall characteristic-class normalization is also suppressed in the notebook.\par
\subsection*{P001474-P001475 (both; faithful error): Forms and ordinary homology base}
\noindent\textit{Transcription reads.} \(\Omega^i(X)\) and \(H_1(X;\mathbb Z)\)\par
\noindent\textit{Critical reading.} \(\Omega^i(M)\) and \(H_1(M;\mathbb Z)\)\par
\noindent\textit{Evidence.} \(X\) denotes the infinite-dimensional gauge quotient; the covariant differential and homology-sphere condition are on the original three-manifold \(M\).\par
\subsection*{P001475 (both; faithful error): Equivalence (16)-(17)}
\noindent\textit{Transcription reads.} \(H^1\) with every local system vanishes iff \(M\) is a homology sphere\par
\noindent\textit{Critical reading.} Hessian nondegeneracy at \(\rho\) iff \(H^1(M;\operatorname{ad}\rho)=0\); ordinary homology is not equivalent\par
\noindent\textit{Evidence.} Scan independently checked. \(\mathbb{RP}^3\) has vanishing twisted \(H^1\) over characteristic zero for every finite-dimensional local system but \(H_1(-;\mathbb Z)=\mathbb Z/2\). The connected sum of two Poincare spheres has an irreducible \(SU(2)\) adjoint system with first cohomology of dimension three; see Section~\ref{dev:twisted-h1}.\par
\subsection*{P001475 (both; faithful error): Relative-index crossing count}
\noindent\textit{Transcription reads.} number of points where a zero mode appears = spectral asymmetry\par
\noindent\textit{Critical reading.} spectral flow: signed crossings counted with multiplicity\par
\noindent\textit{Evidence.} An eigenvalue crossing up and later down contributes zero net spectral flow but two unsigned events. Atiyah 1988, p.~289, gives the spectral-flow interpretation and modulo-eight ambiguity.\par
\subsection*{P001477 (both; faithful error): Dimensions in cutting}
\noindent\textit{Transcription reads.} cut a four-manifold into two three-dimensional ones\par
\noindent\textit{Critical reading.} two four-manifolds with common three-dimensional boundary\par
\noindent\textit{Evidence.} A codimension-one cut preserves the dimension of each piece; Atiyah 1988, Section 7, gives the intended decomposition.\par
\subsection*{P001478-P001479 (both; provenance): Academic-year boundary}
\noindent\textit{Transcription reads.} appended scans included in record751 page-range\par
\noindent\textit{Critical reading.} exclude these pages from the 1987-88 mathematical occurrences\par
\noindent\textit{Evidence.} P001478 explicitly begins 12 September 1988 with a new Gelfand determinants session; P001479 continues with G. Moore on conformal field theory. Their inclusion is archival overlap, not continuation of 25 April.\par
